% Options for packages loaded elsewhere
\PassOptionsToPackage{unicode}{hyperref}
\PassOptionsToPackage{hyphens}{url}
%
\documentclass[
]{book}
\usepackage{amsmath,amssymb}
\usepackage{iftex}
\ifPDFTeX
  \usepackage[T1]{fontenc}
  \usepackage[utf8]{inputenc}
  \usepackage{textcomp} % provide euro and other symbols
\else % if luatex or xetex
  \usepackage{unicode-math} % this also loads fontspec
  \defaultfontfeatures{Scale=MatchLowercase}
  \defaultfontfeatures[\rmfamily]{Ligatures=TeX,Scale=1}
\fi
\usepackage{lmodern}
\ifPDFTeX\else
  % xetex/luatex font selection
\fi
% Use upquote if available, for straight quotes in verbatim environments
\IfFileExists{upquote.sty}{\usepackage{upquote}}{}
\IfFileExists{microtype.sty}{% use microtype if available
  \usepackage[]{microtype}
  \UseMicrotypeSet[protrusion]{basicmath} % disable protrusion for tt fonts
}{}
\makeatletter
\@ifundefined{KOMAClassName}{% if non-KOMA class
  \IfFileExists{parskip.sty}{%
    \usepackage{parskip}
  }{% else
    \setlength{\parindent}{0pt}
    \setlength{\parskip}{6pt plus 2pt minus 1pt}}
}{% if KOMA class
  \KOMAoptions{parskip=half}}
\makeatother
\usepackage{xcolor}
\usepackage{color}
\usepackage{fancyvrb}
\newcommand{\VerbBar}{|}
\newcommand{\VERB}{\Verb[commandchars=\\\{\}]}
\DefineVerbatimEnvironment{Highlighting}{Verbatim}{commandchars=\\\{\}}
% Add ',fontsize=\small' for more characters per line
\usepackage{framed}
\definecolor{shadecolor}{RGB}{248,248,248}
\newenvironment{Shaded}{\begin{snugshade}}{\end{snugshade}}
\newcommand{\AlertTok}[1]{\textcolor[rgb]{0.94,0.16,0.16}{#1}}
\newcommand{\AnnotationTok}[1]{\textcolor[rgb]{0.56,0.35,0.01}{\textbf{\textit{#1}}}}
\newcommand{\AttributeTok}[1]{\textcolor[rgb]{0.13,0.29,0.53}{#1}}
\newcommand{\BaseNTok}[1]{\textcolor[rgb]{0.00,0.00,0.81}{#1}}
\newcommand{\BuiltInTok}[1]{#1}
\newcommand{\CharTok}[1]{\textcolor[rgb]{0.31,0.60,0.02}{#1}}
\newcommand{\CommentTok}[1]{\textcolor[rgb]{0.56,0.35,0.01}{\textit{#1}}}
\newcommand{\CommentVarTok}[1]{\textcolor[rgb]{0.56,0.35,0.01}{\textbf{\textit{#1}}}}
\newcommand{\ConstantTok}[1]{\textcolor[rgb]{0.56,0.35,0.01}{#1}}
\newcommand{\ControlFlowTok}[1]{\textcolor[rgb]{0.13,0.29,0.53}{\textbf{#1}}}
\newcommand{\DataTypeTok}[1]{\textcolor[rgb]{0.13,0.29,0.53}{#1}}
\newcommand{\DecValTok}[1]{\textcolor[rgb]{0.00,0.00,0.81}{#1}}
\newcommand{\DocumentationTok}[1]{\textcolor[rgb]{0.56,0.35,0.01}{\textbf{\textit{#1}}}}
\newcommand{\ErrorTok}[1]{\textcolor[rgb]{0.64,0.00,0.00}{\textbf{#1}}}
\newcommand{\ExtensionTok}[1]{#1}
\newcommand{\FloatTok}[1]{\textcolor[rgb]{0.00,0.00,0.81}{#1}}
\newcommand{\FunctionTok}[1]{\textcolor[rgb]{0.13,0.29,0.53}{\textbf{#1}}}
\newcommand{\ImportTok}[1]{#1}
\newcommand{\InformationTok}[1]{\textcolor[rgb]{0.56,0.35,0.01}{\textbf{\textit{#1}}}}
\newcommand{\KeywordTok}[1]{\textcolor[rgb]{0.13,0.29,0.53}{\textbf{#1}}}
\newcommand{\NormalTok}[1]{#1}
\newcommand{\OperatorTok}[1]{\textcolor[rgb]{0.81,0.36,0.00}{\textbf{#1}}}
\newcommand{\OtherTok}[1]{\textcolor[rgb]{0.56,0.35,0.01}{#1}}
\newcommand{\PreprocessorTok}[1]{\textcolor[rgb]{0.56,0.35,0.01}{\textit{#1}}}
\newcommand{\RegionMarkerTok}[1]{#1}
\newcommand{\SpecialCharTok}[1]{\textcolor[rgb]{0.81,0.36,0.00}{\textbf{#1}}}
\newcommand{\SpecialStringTok}[1]{\textcolor[rgb]{0.31,0.60,0.02}{#1}}
\newcommand{\StringTok}[1]{\textcolor[rgb]{0.31,0.60,0.02}{#1}}
\newcommand{\VariableTok}[1]{\textcolor[rgb]{0.00,0.00,0.00}{#1}}
\newcommand{\VerbatimStringTok}[1]{\textcolor[rgb]{0.31,0.60,0.02}{#1}}
\newcommand{\WarningTok}[1]{\textcolor[rgb]{0.56,0.35,0.01}{\textbf{\textit{#1}}}}
\usepackage{longtable,booktabs,array}
\usepackage{calc} % for calculating minipage widths
% Correct order of tables after \paragraph or \subparagraph
\usepackage{etoolbox}
\makeatletter
\patchcmd\longtable{\par}{\if@noskipsec\mbox{}\fi\par}{}{}
\makeatother
% Allow footnotes in longtable head/foot
\IfFileExists{footnotehyper.sty}{\usepackage{footnotehyper}}{\usepackage{footnote}}
\makesavenoteenv{longtable}
\usepackage{graphicx}
\makeatletter
\def\maxwidth{\ifdim\Gin@nat@width>\linewidth\linewidth\else\Gin@nat@width\fi}
\def\maxheight{\ifdim\Gin@nat@height>\textheight\textheight\else\Gin@nat@height\fi}
\makeatother
% Scale images if necessary, so that they will not overflow the page
% margins by default, and it is still possible to overwrite the defaults
% using explicit options in \includegraphics[width, height, ...]{}
\setkeys{Gin}{width=\maxwidth,height=\maxheight,keepaspectratio}
% Set default figure placement to htbp
\makeatletter
\def\fps@figure{htbp}
\makeatother
\setlength{\emergencystretch}{3em} % prevent overfull lines
\providecommand{\tightlist}{%
  \setlength{\itemsep}{0pt}\setlength{\parskip}{0pt}}
\setcounter{secnumdepth}{5}
% definitions for citeproc citations
\NewDocumentCommand\citeproctext{}{}
\NewDocumentCommand\citeproc{mm}{%
  \begingroup\def\citeproctext{#2}\cite{#1}\endgroup}
\makeatletter
 % allow citations to break across lines
 \let\@cite@ofmt\@firstofone
 % avoid brackets around text for \cite:
 \def\@biblabel#1{}
 \def\@cite#1#2{{#1\if@tempswa , #2\fi}}
\makeatother
\newlength{\cslhangindent}
\setlength{\cslhangindent}{1.5em}
\newlength{\csllabelwidth}
\setlength{\csllabelwidth}{3em}
\newenvironment{CSLReferences}[2] % #1 hanging-indent, #2 entry-spacing
 {\begin{list}{}{%
  \setlength{\itemindent}{0pt}
  \setlength{\leftmargin}{0pt}
  \setlength{\parsep}{0pt}
  % turn on hanging indent if param 1 is 1
  \ifodd #1
   \setlength{\leftmargin}{\cslhangindent}
   \setlength{\itemindent}{-1\cslhangindent}
  \fi
  % set entry spacing
  \setlength{\itemsep}{#2\baselineskip}}}
 {\end{list}}
\usepackage{calc}
\newcommand{\CSLBlock}[1]{\hfill\break\parbox[t]{\linewidth}{\strut\ignorespaces#1\strut}}
\newcommand{\CSLLeftMargin}[1]{\parbox[t]{\csllabelwidth}{\strut#1\strut}}
\newcommand{\CSLRightInline}[1]{\parbox[t]{\linewidth - \csllabelwidth}{\strut#1\strut}}
\newcommand{\CSLIndent}[1]{\hspace{\cslhangindent}#1}
\ifLuaTeX
  \usepackage{selnolig}  % disable illegal ligatures
\fi
\usepackage{bookmark}
\IfFileExists{xurl.sty}{\usepackage{xurl}}{} % add URL line breaks if available
\urlstyle{same}
\hypersetup{
  pdftitle={Manual de Ecologia Evolutiva, Biometria e Biogeografia},
  pdfauthor={Renan Maestri},
  hidelinks,
  pdfcreator={LaTeX via pandoc}}

\title{Manual de Ecologia Evolutiva, Biometria e Biogeografia}
\author{Renan Maestri}
\date{2024-07-11}

\begin{document}
\maketitle

\chapter*{Prefácio}\label{prefuxe1cio}
\addcontentsline{toc}{chapter}{Prefácio}

\begin{center}\includegraphics[width=1\linewidth]{images/Capa do Livro} \end{center}

Esse manual tem o objetivo de servir de apoio às disciplinas que leciono na UFRGS. O manual contém material de apoio das disciplinas de Introdução à Ecologia, Ecologia para Biotecnologia, Ecologia I (evolutiva) e Ecologia II (macroecologia e biogeografia), Ecologia e Evolução Morfológica, Fundamentos de Morfometria Geométrica, Macroecologia, Macroevolução e Métodos Filogenéticos Comparativos e Oficina de Ecologia Evolutiva. Ministro essas disciplinas na Universidade Federal do Rio Grande do Sul (UFRGS) desde 2018, algumas para a graduação e algumas para a pós-graduação, para diferentes cursos de graduação e Programas de Pós-Graduação do Instituto de Biociências da UFRGS. Algumas dessas disciplinas são ou foram desenvolvidas em colaboração com os professores Gonçalo Ferraz e Heinrich Hasenack (Ecologia I), Leandro Duarte e Luciane Crossetti (Ecologia II), Rodrigo Fornel (Fundamentos de Morfometria Geométrica) e Leandro Duarte (Oficina de Ecologia Evolutiva), a quem agradeço pela parceria. O livro contém uma parte do material visto nessas disciplinas, com foco nas atividades práticas.

O intuito de disponibilizar esse manual de \textbf{Ecologia Evolutiva, Biometria e Biogeografia} é que ele sirva de material de apoio às disciplinas que ministro. Tentei organizá-lo de modo unificado, para que ele seja útil simultaneamente às diferentes disciplinas e níveis de ensino.

Todos os exercícios práticos usam o programa R (\textbf{R-base?}). Há um capítulo que introduz o R.

Agradeço aos alunos com quem tive o prazer de interagir e que me incentivam a ser um professor um pouco melhor a cada dia, espero que esse manual seja útil a vocês durante e após as disciplinas. Tenho uma dívida impagável com as pessoas que fazem parte e que passaram pelo \href{https://www.ufrgs.br/lema/}{Laboratório de Ecomorfologia e Macroevolução} pela troca de experiências e ensinamentos que culminaram na elaboração deste manual.

Se você tiver qualquer comentário ou contribuição ao livro, por favor me escreva (\href{mailto:renanmaestri@gmail.com}{\nolinkurl{renanmaestri@gmail.com}}).

\part{PARTE I - Conceitos básicos}\label{part-parte-i---conceitos-buxe1sicos}

\chapter{Conceitos básicos}\label{conceitos-buxe1sicos}

Em construção.

\section{Processos evolutivos}\label{processos-evolutivos}

Em construção.

\section{Seleção natural}\label{seleuxe7uxe3o-natural}

Em construção.

\subsection{Seleção direcional}\label{seleuxe7uxe3o-direcional}

Em construção.

\subsection{Seleção disruptiva}\label{seleuxe7uxe3o-disruptiva}

Em construção.

\subsection{Seleção estabilizadora}\label{seleuxe7uxe3o-estabilizadora}

Em construção.

\section{Deriva genética}\label{deriva-genuxe9tica}

Em construção.

\subsection{Simulação de deriva e seleção}\label{simulauxe7uxe3o-de-deriva-e-seleuxe7uxe3o}

Em construção.

\section{Ecologia evolutiva do desenvolvimento}\label{ecologia-evolutiva-do-desenvolvimento}

Em construção.

\subsection{Interações genótipo-fenótipo-ambiente}\label{interauxe7uxf5es-genuxf3tipo-fenuxf3tipo-ambiente}

Em construção.

\subsection{Plasticidade fenotípica}\label{plasticidade-fenotuxedpica}

Em construção.

\subsection{Epigenética}\label{epigenuxe9tica}

Em construção.

\part{PARTE II - Introdução ao R}\label{part-parte-ii---introduuxe7uxe3o-ao-r}

\chapter{Uma breve introdução ao R}\label{cap2}

O R pode ser instalado a partir do website: \url{https://www.r-project.org/} Nós vamos usar também o RStudio: \url{https://www.rstudio.com/}

O R é um ambiente para análise computacional e produção de gráficos baseado em linguagem de programação.
Hoje, o R é o programa predominante em ecologia e evolução, em larga medida porque é e adaptável para várias necessidades, e existe uma comunidade grande de indivíduos dispostos a fornecer ajuda.
No contexto da biometria e de análises filogenéticas comparativas, o R tem sido a área primária de desenvolvimento de novos métodos.

Aqui, você vai encontrar um breve tutorial sobre o R para conseguir realizar as atividades práticas da disciplina.
No entanto, há material muito mais completo disponível.
Existem vários bons livros e cursos gratuitos de R e de estatística disponibilizados na web.
Tire proveito deles!
Por exemplo: - \href{https://www.youtube.com/playlist?list=PL3U7PUKvYnzAZAEoFcjGLVkIML2ZGGl6t}{Curso de Programação em R, prof. Bruno Vilela} - \href{https://analises-ecologicas.com/}{Livro Análises Ecológicas no R} - \href{http://estat.bio.br/}{Curso de Estatística para Pesquisa Científica, prof. Thiago Rangel}

\section{Definindo um diretório de trabalho}\label{definindo-um-diretuxf3rio-de-trabalho}

É conveniente criar uma pasta única com todos os arquivos das aulas práticas.
Nós vamos informar ao R para carregar os dados a partir desta pasta.
Isso pode ser feito no RStudio com o menu: Session\textgreater Set Working Directory\textgreater Choose Directory.
Ou com o código setwd().

\section{Roteiros (scripts)}\label{roteiros-scripts}

Um script é um bloco de texto que pode ser editado em qualquer editor de texto comum (ex: bloco de notas), ou com o editor de texto do próprio R ou RStudio.
Você vai escrever os códigos dentro do script, e depois enviar para o console do R - selecione a linha ou parte do código a ser enviada ao console do R e pressione `Run' ou use crtl+r ou ctrl+Enter.
Ao salvar o script você mantém um registro do código utilizado.
Se você quiser adicionar comentários ao código basta adicionar o sinal \# na frente do comentário.
Uma ótima alternativa ao script comum é utilizar o R Markdown, onde é possível diferenciar facilmente códigos e textos explicativos.

\section{Funcionamento básico do R}\label{funcionamento-buxe1sico-do-r}

Objetos e funções são os principais componentes do R.
Objetos são variáveis como dados ou resultados.
Eles podem estar formatados de diferentes maneiras, formando estruturas diferentes.
Por exemplo, vetores, matrizes, `dataframes', listas, etc. são todos exemplos de estruturas de objetos.
Funções são tipos especiais de objetos que usam um ou mais argumentos para fazer alguma coisa dentro do R.

Objetos são como proteínas e funções são como enzimas.
Os aminoácidos são números ou letras.

Compilados de funções são organizados dentro de pacotes criados por diversos autores.

O R funciona através de comandos, que são normalmente acompanhados de funções ou operações (+,-,*,etc).
Um comando retorna um objeto na tela, ou armazena o objeto na memória quando um nome é atribuído com o operador \textless- ou =.

\begin{Shaded}
\begin{Highlighting}[]
\DecValTok{2}\SpecialCharTok{+}\DecValTok{5}
\CommentTok{\#\textgreater{} [1] 7}
\NormalTok{x}\OtherTok{\textless{}{-}}\DecValTok{2}\SpecialCharTok{+}\DecValTok{5}
\NormalTok{x}
\CommentTok{\#\textgreater{} [1] 7}
\end{Highlighting}
\end{Shaded}

Cuidado, os objetos distinguem letras maiúsculas de minúsculas.

\begin{Shaded}
\begin{Highlighting}[]
\NormalTok{x}\OtherTok{\textless{}{-}}\DecValTok{1}
\NormalTok{X}\OtherTok{\textless{}{-}}\DecValTok{10}
\NormalTok{x}
\CommentTok{\#\textgreater{} [1] 1}
\NormalTok{X}
\CommentTok{\#\textgreater{} [1] 10}
\end{Highlighting}
\end{Shaded}

\section{Tipos de objeto no R}\label{tipos-de-objeto-no-r}

O R classifica objetos em tipos/classes dependendo de sua estrutura.
Vamos encontrar frequentemente alguns tipos de objetos mais comuns, como vector, factor, matrix, data.frame, array e list.

\begin{Shaded}
\begin{Highlighting}[]
\CommentTok{\# Vetores:}
\NormalTok{x}\OtherTok{\textless{}{-}}\FunctionTok{c}\NormalTok{(}\DecValTok{1}\NormalTok{,}\DecValTok{2}\NormalTok{,}\DecValTok{3}\NormalTok{,}\DecValTok{4}\NormalTok{,}\DecValTok{5}\NormalTok{)}
\NormalTok{x}\OtherTok{\textless{}{-}}\DecValTok{1}\SpecialCharTok{:}\DecValTok{5}
\NormalTok{x}
\CommentTok{\#\textgreater{} [1] 1 2 3 4 5}
\FunctionTok{class}\NormalTok{(x)}
\CommentTok{\#\textgreater{} [1] "integer"}
\FunctionTok{length}\NormalTok{(x)}
\CommentTok{\#\textgreater{} [1] 5}
\CommentTok{\# A função which é útil para selecionar alguns dados de um vetor}
\FunctionTok{which}\NormalTok{(x}\SpecialCharTok{\textgreater{}}\DecValTok{3}\NormalTok{)}
\CommentTok{\#\textgreater{} [1] 4 5}
\FunctionTok{which}\NormalTok{(x}\SpecialCharTok{\textgreater{}=}\DecValTok{3}\NormalTok{)}
\CommentTok{\#\textgreater{} [1] 3 4 5}
\end{Highlighting}
\end{Shaded}

\begin{Shaded}
\begin{Highlighting}[]
\CommentTok{\# Fatores possuem níveis}
\NormalTok{f}\OtherTok{\textless{}{-}}\FunctionTok{c}\NormalTok{(}\StringTok{"Sp1"}\NormalTok{,}\StringTok{"Sp1"}\NormalTok{,}\StringTok{"Sp2"}\NormalTok{,}\StringTok{"Sp3"}\NormalTok{,}\StringTok{"Sp3"}\NormalTok{)}
\NormalTok{f}
\CommentTok{\#\textgreater{} [1] "Sp1" "Sp1" "Sp2" "Sp3" "Sp3"}
\NormalTok{f}\OtherTok{\textless{}{-}}\FunctionTok{factor}\NormalTok{(f)}
\NormalTok{f}
\CommentTok{\#\textgreater{} [1] Sp1 Sp1 Sp2 Sp3 Sp3}
\CommentTok{\#\textgreater{} Levels: Sp1 Sp2 Sp3}
\FunctionTok{table}\NormalTok{(f)}
\CommentTok{\#\textgreater{} f}
\CommentTok{\#\textgreater{} Sp1 Sp2 Sp3 }
\CommentTok{\#\textgreater{}   2   1   2}
\FunctionTok{class}\NormalTok{(f)}
\CommentTok{\#\textgreater{} [1] "factor"}
\end{Highlighting}
\end{Shaded}

Os objetos podem conter nomes internos, o que é bastante útil com dados biológicos já que eles podem fornecer rótulos para os dados, ex: espécies ou localidades.

\begin{Shaded}
\begin{Highlighting}[]
\FunctionTok{names}\NormalTok{(x)}\OtherTok{=}\NormalTok{f}
\NormalTok{x}
\CommentTok{\#\textgreater{} Sp1 Sp1 Sp2 Sp3 Sp3 }
\CommentTok{\#\textgreater{}   1   2   3   4   5}
\end{Highlighting}
\end{Shaded}

\begin{Shaded}
\begin{Highlighting}[]
\CommentTok{\# Matrizes}
\NormalTok{m}\OtherTok{\textless{}{-}}\FunctionTok{matrix}\NormalTok{(}\DecValTok{1}\SpecialCharTok{:}\DecValTok{9}\NormalTok{,}\DecValTok{3}\NormalTok{,}\DecValTok{3}\NormalTok{) }\CommentTok{\# dados, n{-}linhas, n{-}colunas}
\NormalTok{m}
\CommentTok{\#\textgreater{}      [,1] [,2] [,3]}
\CommentTok{\#\textgreater{} [1,]    1    4    7}
\CommentTok{\#\textgreater{} [2,]    2    5    8}
\CommentTok{\#\textgreater{} [3,]    3    6    9}
\FunctionTok{class}\NormalTok{(m)}
\CommentTok{\#\textgreater{} [1] "matrix" "array"}
\NormalTok{m[,}\DecValTok{2}\NormalTok{] }\CommentTok{\# a segunda coluna}
\CommentTok{\#\textgreater{} [1] 4 5 6}
\NormalTok{m[}\DecValTok{3}\NormalTok{,] }\CommentTok{\# a terceira linha}
\CommentTok{\#\textgreater{} [1] 3 6 9}
\end{Highlighting}
\end{Shaded}

\begin{Shaded}
\begin{Highlighting}[]
\CommentTok{\# Data frames suportam vários tipos de dados distintos}
\NormalTok{d}\OtherTok{\textless{}{-}}\FunctionTok{data.frame}\NormalTok{(f,x)}
\NormalTok{d}
\CommentTok{\#\textgreater{}     f x}
\CommentTok{\#\textgreater{} 1 Sp1 1}
\CommentTok{\#\textgreater{} 2 Sp1 2}
\CommentTok{\#\textgreater{} 3 Sp2 3}
\CommentTok{\#\textgreater{} 4 Sp3 4}
\CommentTok{\#\textgreater{} 5 Sp3 5}
\FunctionTok{class}\NormalTok{(d)}
\CommentTok{\#\textgreater{} [1] "data.frame"}
\end{Highlighting}
\end{Shaded}

\begin{Shaded}
\begin{Highlighting}[]
\CommentTok{\# Lista é uma estrutura geral que pode carregar vários tipos de dados, incluindo, por ex. árvores filogenéticas}
\NormalTok{l}\OtherTok{\textless{}{-}}\FunctionTok{list}\NormalTok{(f,x,m,d)}
\NormalTok{l}
\CommentTok{\#\textgreater{} [[1]]}
\CommentTok{\#\textgreater{} [1] Sp1 Sp1 Sp2 Sp3 Sp3}
\CommentTok{\#\textgreater{} Levels: Sp1 Sp2 Sp3}
\CommentTok{\#\textgreater{} }
\CommentTok{\#\textgreater{} [[2]]}
\CommentTok{\#\textgreater{} Sp1 Sp1 Sp2 Sp3 Sp3 }
\CommentTok{\#\textgreater{}   1   2   3   4   5 }
\CommentTok{\#\textgreater{} }
\CommentTok{\#\textgreater{} [[3]]}
\CommentTok{\#\textgreater{}      [,1] [,2] [,3]}
\CommentTok{\#\textgreater{} [1,]    1    4    7}
\CommentTok{\#\textgreater{} [2,]    2    5    8}
\CommentTok{\#\textgreater{} [3,]    3    6    9}
\CommentTok{\#\textgreater{} }
\CommentTok{\#\textgreater{} [[4]]}
\CommentTok{\#\textgreater{}     f x}
\CommentTok{\#\textgreater{} 1 Sp1 1}
\CommentTok{\#\textgreater{} 2 Sp1 2}
\CommentTok{\#\textgreater{} 3 Sp2 3}
\CommentTok{\#\textgreater{} 4 Sp3 4}
\CommentTok{\#\textgreater{} 5 Sp3 5}
\NormalTok{l[[}\DecValTok{4}\NormalTok{]] }\CommentTok{\# quarto elemento da lista}
\CommentTok{\#\textgreater{}     f x}
\CommentTok{\#\textgreater{} 1 Sp1 1}
\CommentTok{\#\textgreater{} 2 Sp1 2}
\CommentTok{\#\textgreater{} 3 Sp2 3}
\CommentTok{\#\textgreater{} 4 Sp3 4}
\CommentTok{\#\textgreater{} 5 Sp3 5}
\end{Highlighting}
\end{Shaded}

A função ls() lista todos os objetos criados na seção.
Você também pode conferir os objetos na aba Environment, ao lado.

\begin{Shaded}
\begin{Highlighting}[]
\FunctionTok{ls}\NormalTok{()}
\end{Highlighting}
\end{Shaded}

Todas as funções no R possuem uma página de ajuda descrevendo a função.
A ajuda pode ser acessada com ?
seguido do nome da função.

\begin{Shaded}
\begin{Highlighting}[]
\NormalTok{?mean}
\CommentTok{\#\textgreater{} starting httpd help server ... done}
\NormalTok{?matrix}
\end{Highlighting}
\end{Shaded}

Um dos aspectos importantes negligenciados por iniciantes no R são as páginas de ajuda!
Sempre confira a ajuda da função que está tentando aplicar usando o operador ?
antes do nome da função.
Por exemplo, ?m
ean.
Preste atenção aos argumentos que a função usa, e quais tipos/estrutura de objetos são aceitos como argumentos.
Estude a página de ajuda e entenda como ela funciona.

Algumas operações são úteis para reproduzir códigos/funções várias vezes.
O \emph{for} loop repete uma determinada operação por um número fixo de vezes.
Por exemplo, para repetir a mesma operação em todas as filogenias de um conjunto de dados.

\begin{Shaded}
\begin{Highlighting}[]
\CommentTok{\# For loop para calcular a média de atributos, ex: colunas de uma matriz}
\NormalTok{X}\OtherTok{\textless{}{-}}\FunctionTok{matrix}\NormalTok{(}\DecValTok{1}\SpecialCharTok{:}\DecValTok{16}\NormalTok{,}\DecValTok{4}\NormalTok{,}\DecValTok{4}\NormalTok{)}
\NormalTok{X}
\CommentTok{\#\textgreater{}      [,1] [,2] [,3] [,4]}
\CommentTok{\#\textgreater{} [1,]    1    5    9   13}
\CommentTok{\#\textgreater{} [2,]    2    6   10   14}
\CommentTok{\#\textgreater{} [3,]    3    7   11   15}
\CommentTok{\#\textgreater{} [4,]    4    8   12   16}
\CommentTok{\# vetor vazio para salvar os resultados}
\NormalTok{medias}\OtherTok{\textless{}{-}}\FunctionTok{vector}\NormalTok{(}\AttributeTok{mode=}\StringTok{"numeric"}\NormalTok{,}\FunctionTok{length}\NormalTok{(}\FunctionTok{ncol}\NormalTok{(X)))}
\NormalTok{medias}
\CommentTok{\#\textgreater{} [1] 0}
\CommentTok{\# loop}
\ControlFlowTok{for}\NormalTok{(i }\ControlFlowTok{in} \DecValTok{1}\SpecialCharTok{:}\FunctionTok{ncol}\NormalTok{(X))\{ }
\NormalTok{ medias[i]}\OtherTok{\textless{}{-}}\FunctionTok{mean}\NormalTok{(X[,i])}
\NormalTok{\}}
\CommentTok{\# resultado}
\NormalTok{medias}
\CommentTok{\#\textgreater{} [1]  2.5  6.5 10.5 14.5}
\end{Highlighting}
\end{Shaded}

\section{Pacotes no R}\label{pacotes-no-r}

Os pacotes contêm compilados de funções que realizam tarefas sobre um determinado tema.
Esses pacotes, criados por autores diversos, estão disponíveis para serem baixados e carregados no R.
Os pacotes são baixados uma única vez no computador, porém, precisam ser carregados no console do R toda vez que se iniciar o R.

\section{Pacotes para morfometria geométrica}\label{pacotes-para-morfometria-geomuxe9trica}

Um dos pacotes mais completos para morfometria geométrica é o \emph{geomorph}.
Vamos instalar esse pacote e carregar no R.
Os pacotes podem ser instalados através de uma linha de comando ou clicando no botão Install da aba Packages.

\begin{Shaded}
\begin{Highlighting}[]
\FunctionTok{install.packages}\NormalTok{(}\StringTok{"geomorph"}\NormalTok{)}
\FunctionTok{library}\NormalTok{(geomorph)}
\end{Highlighting}
\end{Shaded}

O \emph{geomorph} é constantemente atualizado.
Por isso, é recomendável checar frequentemente por atualizações nas funções.
Com o comando ??
seguido do nome do pacote (??geomorph) é possível acessar um resumo do pacote e de todas as suas funções.

\section{Objetos importantes para morfometria geométrica}\label{objetos-importantes-para-morfometria-geomuxe9trica}

Via de regra, dados obtidos via morfometria geométrica vão estar em um de 2 formatos no R: como \textbf{3D array} ou como \textbf{matrix}.
No formato array, análogo ao formato .tps, cada indivíduo aparece de maneira individualizada e sequencial dentro do objeto; para cada indivíduo as linhas correspondem aos landmarks e as colunas às coordenadas x e y (e z se dados em 3D) de cada landmark.
No formato matrix todos os indivíduos aparecem de uma única vez na matriz; as linhas correspondem aos indivíduos e as colunas correspondem às coordenadas dos landmarks de maneira sequencial: a coluna 1 corresponde ao landmark x1 de cada indivíduo (coordenada x do landmark 1), a coluna 2 = y1, coluna 3 = x2, coluna 4 = y2, coluna 5 = x3 etc.

\begin{Shaded}
\begin{Highlighting}[]
\CommentTok{\# Matrix }
\NormalTok{m}\OtherTok{\textless{}{-}}\FunctionTok{matrix}\NormalTok{(}\DecValTok{1}\SpecialCharTok{:}\DecValTok{12}\NormalTok{,}\DecValTok{3}\NormalTok{,}\DecValTok{4}\NormalTok{) }\CommentTok{\# dados, n{-}linhas, n{-}colunas }
\NormalTok{m }\CommentTok{\# 3 indivíduos, 2 landmarks com coordenadas x,y }
\CommentTok{\#\textgreater{}      [,1] [,2] [,3] [,4]}
\CommentTok{\#\textgreater{} [1,]    1    4    7   10}
\CommentTok{\#\textgreater{} [2,]    2    5    8   11}
\CommentTok{\#\textgreater{} [3,]    3    6    9   12}
\CommentTok{\# Array }
\NormalTok{n}\OtherTok{\textless{}{-}}\FunctionTok{array}\NormalTok{(m,}\FunctionTok{c}\NormalTok{(}\DecValTok{4}\NormalTok{,}\DecValTok{2}\NormalTok{,}\DecValTok{3}\NormalTok{)) }\CommentTok{\#dados, dimensões(p,k,n) (p=pontos,k=dimensões,n=amostra) }
\NormalTok{n }\CommentTok{\# 3 indivíduos, 2 landmarks com coordenadas x,y }
\CommentTok{\#\textgreater{} , , 1}
\CommentTok{\#\textgreater{} }
\CommentTok{\#\textgreater{}      [,1] [,2]}
\CommentTok{\#\textgreater{} [1,]    1    5}
\CommentTok{\#\textgreater{} [2,]    2    6}
\CommentTok{\#\textgreater{} [3,]    3    7}
\CommentTok{\#\textgreater{} [4,]    4    8}
\CommentTok{\#\textgreater{} }
\CommentTok{\#\textgreater{} , , 2}
\CommentTok{\#\textgreater{} }
\CommentTok{\#\textgreater{}      [,1] [,2]}
\CommentTok{\#\textgreater{} [1,]    9    1}
\CommentTok{\#\textgreater{} [2,]   10    2}
\CommentTok{\#\textgreater{} [3,]   11    3}
\CommentTok{\#\textgreater{} [4,]   12    4}
\CommentTok{\#\textgreater{} }
\CommentTok{\#\textgreater{} , , 3}
\CommentTok{\#\textgreater{} }
\CommentTok{\#\textgreater{}      [,1] [,2]}
\CommentTok{\#\textgreater{} [1,]    5    9}
\CommentTok{\#\textgreater{} [2,]    6   10}
\CommentTok{\#\textgreater{} [3,]    7   11}
\CommentTok{\#\textgreater{} [4,]    8   12}
\NormalTok{n[,,}\DecValTok{1}\NormalTok{] }\CommentTok{\# indivíduo número 1}
\CommentTok{\#\textgreater{}      [,1] [,2]}
\CommentTok{\#\textgreater{} [1,]    1    5}
\CommentTok{\#\textgreater{} [2,]    2    6}
\CommentTok{\#\textgreater{} [3,]    3    7}
\CommentTok{\#\textgreater{} [4,]    4    8}
\end{Highlighting}
\end{Shaded}

\section{Pacotes filogenéticos no R}\label{pacotes-filogenuxe9ticos-no-r}

Existem muitos pacote lidando com métodos filogenéticos no R.
Alguns que serão usados para manipular e representar graficamente árvores filogenéticas serão instalados abaixo.
O pacote \emph{ape} é o pacote central mais importante, já que propõe a classe \emph{phylo}, que é utilizada pela grande maioria dos demais pacotes.

\begin{Shaded}
\begin{Highlighting}[]
\CommentTok{\# Instalando pacotes: }
\FunctionTok{install.packages}\NormalTok{(}\StringTok{"ape"}\NormalTok{) }
\FunctionTok{install.packages}\NormalTok{(}\StringTok{"phytools"}\NormalTok{) }
\FunctionTok{install.packages}\NormalTok{(}\StringTok{"geiger"}\NormalTok{)}
\end{Highlighting}
\end{Shaded}

\section{\texorpdfstring{O objeto \emph{phylo}}{O objeto phylo}}\label{o-objeto-phylo}

Árvores filogenéticas são lidas no R como um objeto \emph{list} de classe \emph{phylo}.

\begin{Shaded}
\begin{Highlighting}[]
\CommentTok{\# Carregar ape }
\FunctionTok{require}\NormalTok{(ape) }
\CommentTok{\#\textgreater{} Carregando pacotes exigidos: ape}
\CommentTok{\# Construindo uma árvore newick via texto }
\NormalTok{texto}\OtherTok{\textless{}{-}}\StringTok{"(((((((vaca, cavalo),baleia),(morcego,(camundongo,humano))),(pardal,tartaruga)),celacanto),nemo),tubarão);"}
\CommentTok{\# Função read.tree lê o arquivo contendo a árvore }
\NormalTok{arvore}\OtherTok{\textless{}{-}}\FunctionTok{read.tree}\NormalTok{(}\AttributeTok{text=}\NormalTok{texto)}
\end{Highlighting}
\end{Shaded}

Podemos representar a árvore graficamente, bem como estudar sua estrutura.

\begin{Shaded}
\begin{Highlighting}[]
\CommentTok{\# Representando a árvore graficamente }
\FunctionTok{plot}\NormalTok{(arvore,}\AttributeTok{edge.width=}\DecValTok{2}\NormalTok{) }
\end{Highlighting}
\end{Shaded}

\begin{center}\includegraphics{02-intro-ao-R_files/figure-latex/unnamed-chunk-16-1} \end{center}

\begin{Shaded}
\begin{Highlighting}[]

\CommentTok{\# Estrutura da árvore }
\NormalTok{arvore }
\CommentTok{\#\textgreater{} }
\CommentTok{\#\textgreater{} Phylogenetic tree with 11 tips and 10 internal nodes.}
\CommentTok{\#\textgreater{} }
\CommentTok{\#\textgreater{} Tip labels:}
\CommentTok{\#\textgreater{}   vaca, cavalo, baleia, morcego, camundongo, humano, ...}
\CommentTok{\#\textgreater{} }
\CommentTok{\#\textgreater{} Rooted; no branch lengths.}
\FunctionTok{str}\NormalTok{(arvore) }
\CommentTok{\#\textgreater{} List of 3}
\CommentTok{\#\textgreater{}  $ edge     : int [1:20, 1:2] 12 13 14 15 16 17 18 18 17 16 ...}
\CommentTok{\#\textgreater{}  $ Nnode    : int 10}
\CommentTok{\#\textgreater{}  $ tip.label: chr [1:11] "vaca" "cavalo" "baleia" "morcego" ...}
\CommentTok{\#\textgreater{}  {-} attr(*, "class")= chr "phylo"}
\CommentTok{\#\textgreater{}  {-} attr(*, "order")= chr "cladewise"}

\CommentTok{\# Decompondo a árvore }
\FunctionTok{class}\NormalTok{(arvore) }
\CommentTok{\#\textgreater{} [1] "phylo"}
\FunctionTok{plot}\NormalTok{(arvore) }
\FunctionTok{tiplabels}\NormalTok{() }
\FunctionTok{nodelabels}\NormalTok{() }
\end{Highlighting}
\end{Shaded}

\begin{center}\includegraphics{02-intro-ao-R_files/figure-latex/unnamed-chunk-16-2} \end{center}

\begin{Shaded}
\begin{Highlighting}[]

\CommentTok{\# ramos }
\NormalTok{arvore}\SpecialCharTok{$}\NormalTok{edge }\CommentTok{\# matriz contendo os índices do nó inicial e final para cada ramo }
\CommentTok{\#\textgreater{}       [,1] [,2]}
\CommentTok{\#\textgreater{}  [1,]   12   13}
\CommentTok{\#\textgreater{}  [2,]   13   14}
\CommentTok{\#\textgreater{}  [3,]   14   15}
\CommentTok{\#\textgreater{}  [4,]   15   16}
\CommentTok{\#\textgreater{}  [5,]   16   17}
\CommentTok{\#\textgreater{}  [6,]   17   18}
\CommentTok{\#\textgreater{}  [7,]   18    1}
\CommentTok{\#\textgreater{}  [8,]   18    2}
\CommentTok{\#\textgreater{}  [9,]   17    3}
\CommentTok{\#\textgreater{} [10,]   16   19}
\CommentTok{\#\textgreater{} [11,]   19    4}
\CommentTok{\#\textgreater{} [12,]   19   20}
\CommentTok{\#\textgreater{} [13,]   20    5}
\CommentTok{\#\textgreater{} [14,]   20    6}
\CommentTok{\#\textgreater{} [15,]   15   21}
\CommentTok{\#\textgreater{} [16,]   21    7}
\CommentTok{\#\textgreater{} [17,]   21    8}
\CommentTok{\#\textgreater{} [18,]   14    9}
\CommentTok{\#\textgreater{} [19,]   13   10}
\CommentTok{\#\textgreater{} [20,]   12   11}
\CommentTok{\# O número de linhas em tree$edge é igual o número de ramos; cada ramo começa e termina com um par de índices (nó{-}nó ou nó{-}terminal) }

\CommentTok{\# terminais }
\NormalTok{arvore}\SpecialCharTok{$}\NormalTok{tip.label }
\CommentTok{\#\textgreater{}  [1] "vaca"       "cavalo"     "baleia"     "morcego"   }
\CommentTok{\#\textgreater{}  [5] "camundongo" "humano"     "pardal"     "tartaruga" }
\CommentTok{\#\textgreater{}  [9] "celacanto"  "nemo"       "tubarão"}

\CommentTok{\# nós }
\NormalTok{arvore}\SpecialCharTok{$}\NormalTok{Nnode }
\CommentTok{\#\textgreater{} [1] 10}

\CommentTok{\# Um objeto "phylo" também pode conter o comprimento dos ramos: edge.length}
\end{Highlighting}
\end{Shaded}

\section{É eRrando que se aprende}\label{uxe9-errando-que-se-aprende}

A curva de aprendizado do R é bastante íngreme no início.
Quando as mensagens de erro começarem a aparecer, é importante não se desesperar.
Tente ler e interpretar as mensagens de erro.
Volte ao código e o leia com calma etapa por etapa.
Procure por erros comuns, como um nome de objeto errado, um espaço extra, uma vírgula mal colocada, etc.
Lembre-se que não estás aprendendo a ``mexer num programa'', estás aprendendo uma língua nova.

\part{PARTE III - Morfometria Geométrica}\label{part-parte-iii---morfometria-geomuxe9trica}

\chapter{Morfometria Geométrica}\label{cap3}

``Organismos são entidades integradas, e não coleções de objetos discretos'' (\textbf{gould1979?}).

A visão ``atomista'' dos organismos, que trata medidas lineares únicas como merecedoras de atenção particular, está errada do ponto de vista da genética e do desenvolvimento. A evolução morfológica é inerentemente multivarida. Por isso, precisamos usar ferramentas multivariadas para entender os processos biológicos que produzem morfologias distintas. A morfometria fornece maneiras quantitativas de comparar a forma dos organismos, obtidas a partir do uso de ferramentas matemáticas sobre matrizes de variáveis contínuas. Assim, a morfometria multivariada é a melhor maneira de entender morfogênese e evolução morfológica.

Contudo, a morfometria tradicional não oferece uma maneira precisa para acessar a geometria dos organismos, já que carece de uma definição precisa de forma e, assim, não distingue claramente forma e tamanho. Além disso, a morfometria tradicional não oferece maneiras de visualizar e interpretar com facilidade as diferenças morfológicas que estamos tentando entender.

Ferramentas de \textbf{morfometria geométrica} oferecem uma \emph{descrição precisa e acurada} da forma (geometria) dos organismos, podem ser combinadas com \emph{análises estatísticas} rigorosas, e nos permitem \emph{visualizar} as diferenças de \textbf{forma} entre organismos com extrema facilidade, o que não é alcançado facilmente por outros métodos (\textbf{zelditch2012?}).

\textbf{Forma}, em morfometria geométrica, é formamente definida como \emph{todas as características geométricas de uma configuração de pontos exceto por seu tamanho, posição e orientação} (\textbf{kendall?}; \textbf{bookstein1989?}; \textbf{dryden1998?}; \textbf{monteiro1999?}; \textbf{klingenberg2016?}). Essa formalidade é capturada matematicamente pela técnica conhecida como \emph{Análise Generalizada de Procrustes}, que veremos na sequência. Assim, \emph{morfometria geométrica faz uma distinção entre tamanho e forma}. \emph{Tamanho} pode ser resumido em um número que captura a dimensão ou escala geral de um objeto (\textbf{klingenberg2016?}) - veja a seção do manual sobre Tamanho do Centróide. \emph{Forma} captura as diferentes proporções de um objeto: as formas de dois objetos podem ser iguais se forem geométricamente idênticas, mesmo que seu tamanho seja diferente (\textbf{zelditch2012?}; \textbf{klingenberg2016?}).

Os \textbf{landmarks}, ou marcos anatômicos, são pontos anatômicos discretos, fudamentais em morfometria, que precisam ser reconhecíveis em todos os indivíduos/objetos da amostra. A definição dos landmarks é central, já que eles permitirão capturar as estruturas biológicas de interesse (entre os landmarks) e comparar essas estruturas entre organismos. Daí a importância de os landmarks serem homólogos entre os indivíduos no estudo. Em morfometria geométrica, as \textbf{coordenadas cartesianas} desses landmarks é que serão usadas nas análises. Há uma outra vantagem da morfometria geométrica que resulta do uso dessas coordenadas cartesianas ao invés do uso de medidas lineares entre landmarks. A análise das variáveis de forma, obtidas por morfometria geométrica, sobre determinado número N de landmarks, contém tanta informação de forma quanto teriam todas as medidas lineares entre os N pares de landmarks, e com um número muito menor de variáveis de forma, removendo a redundância que existiria na matriz de medidas lineares entre os N landmarks.

Além dos landmarks, também podem ser utilizados semilandmarks ou análises de contorno quando a estrutura biológica for muito complexa para ser inteiramente capturada por landmarks e/ou quando landmarks não puderem ser localizados com confiabilidade. Veja a seção sobre semilandmarks.

Para livros abrangentes sobre morfometria geométrica, veja: (\textbf{monteiro1999?}) e (\textbf{zelditch2012?}).

\chapter{Digitalização de landmarks}\label{digitalizauxe7uxe3o-de-landmarks}

\section{Criação de arquivo .tps com as imagens bi-dimensionais}\label{criauxe7uxe3o-de-arquivo-.tps-com-as-imagens-bi-dimensionais}

Esse roteiro é para a manipulação de imagens digitais (2D) e a digitalização de marcos anatômicos (landmarks) nas imagens usando \href{https://www.sbmorphometrics.org/}{softwares TPS}. As imagens, previamente capturadas de maneira padronizada, devem estar individualizadas, todas em uma mesma pasta no computador.

No TPSUtil, selecione a operação `Build tps file from images'. Clique em `Input' e selecione a pasta (selecione qualquer imagem dentro da pasta) contendo as fotos. No nosso caso, a pasta se chama `Atividade 1 -- Imagens'. Clique em `Output' e atribua um nome para o arquivo (use mandibula.tps). Clique em `Setup'. Todas as imagens da pasta estarão listadas. É possível conferir os nomes dos arquivos e modificar a ordem ou deletar imagens se necessário. Clique em `Create' para criar o arquivo .tps com a lista de imagens. Vá até o mesmo diretório (pasta) com as imagens e verifique o arquivo mandibula.tps (abra com o bloco de notas).

\section{Digitalização de landmarks no TPSDig2}\label{digitalizauxe7uxe3o-de-landmarks-no-tpsdig2}

No TPSDig2, para carregar as imagens clique em `File\textgreater Input source\textgreater File'. Selecione o arquivo .tps contendo a lista de imagens, no nosso caso `mandibula.tps'. Após clicar em `Abrir', a primeira foto do arquivo será exibida. As setas vermelhas na barra de ferramentas permitem navegar pelas fotos, que serão exibidas uma por vez. A barra de ferramentas também permite ajustar o zoom das imagens, e inúmeras outras configurações de nitidez da imagem e cores/tamanho dos landmarks.

Um passo importante é definir a escala da imagem. Para isso, clique em `Options\textgreater Image tools'. Na aba `Measure' defina o comprimento de referência e a unidade que desejar para o comprimento. Clique em `Set scale', minimize a aba de ferramentas, vá até a escala de referência na imagem (no nosso caso, o papel milimetrado) e defina, com dois pontos (cliques), o mesmo comprimento atribuído na barra de ferramentas. Volte para a aba `Measure' e clique em `Ok'. A calibração estará definida para esta imagem (o fator de escala será o valor de referência dividido pelo número de pixels necessários para obter a referência). Futuros valores de tamanho do centroide para esta imagem estarão na unidade definida como referência (e.g., cm, mm, etc.). O procedimento precisa ser repetido para cada imagem.

Para digitalizar os landmarks, i.e., atribuir uma coodenada x,y para cada um, selecione a ferramenta `Digitize landmarks' no menu `Modes' ou clicando na figura com o formato de alvo na barra de ferramentas. Atribua os landmarks clicando na imagem, na posição desejada. Os landmarks vão aparecer em sequência numérica a partir do 1, e em todas as imagens a sequência de digitalização dos landmarks (a ordem) deve ser a mesma (e.g., o landmark 8 estará na mesma posição anatômica em todas as imagens).

Após o término da digitalização clique em `File\textgreater Save data' e salve os dados. É extremamente recomendável repetir o salvamento após a digitalização de cada imagem; use o mesmo nome e clique em `Overwrite' para que os dados atualizados sejam gravados sobre os antigos, gerando um único arquivo ao final. Vá até a pasta de trabalho e visualize o arquivo abrindo-o com o bloco de notas.

Para a atividade, atribua os landmarks conforme a imagem abaixo:

\begin{figure}

{\centering \includegraphics[width=0.8\linewidth]{images/mandibula} 

}

\caption{Landmarks na mandibula de um roedor sigmodontíneo.}\label{fig:unnamed-chunk-1}
\end{figure}

Descrição dos landmarks:
L1: Borda anterior do alvéolo incisivo; L2: Extremo anterior superior da borda do alvéolo incisive; L3: Posição de maior inflexão do diastema; L4: Borda anterior do álveolo no primeiro molar; L5: Interseção entre a linha de molars e o processo coronóide em vista lateral; L6: Ponta do processo coronóide; L7: Ponto de máxima curvatura entre o processo coronóide e o condilar; L8: Margem dorsal da borda anterior da supercífie articular do processo condilar; L9: Borda ventral da superfície articular do processo condilar; L10: Ponto de máxima curvatura entre o processo condilar e o angular; L11: Ponta do processo angular; L12: Interseção entre o corpo da mandíbula e a crista massetérica; L13: Limite ventral da sínfise mandibular.

\section{Digitalização de landmarks em 3D}\label{digitalizauxe7uxe3o-de-landmarks-em-3d}

O software MeshLab permite a digitalização de landmarks em superfícies, além de outras operações. Após abrir o software, selecione `File\textgreater Import Mesh' e carregue o arquivo chamado `exemplo3d.ply' de dentro da pasta `R Dados'.

Para digitalizar landmarks na superfície clique em `Edit\textgreater PickPoints' ou no ícone PP. Ao clicar com o botão esquerdo do mouse (pode demorar um pouco), as coordenadas são capturadas. Após o término da digitalização é possível salvar um arquivo em formato xml (.pp) com as coordenadas em 3D, que poderá ser aberto no R para análises posteriores.

\chapter{Análise de Procrustes}\label{anuxe1lise-de-procrustes}

O passo primário fundamental em morfometria geométrica é a Análise Generalizada de Procrustes (GPA - \emph{Generalized Procrustes Analysis}), que remove os efeitos de rotação/orientação, translação/posição e tamanho dos objetos. O resultado é uma matriz com as variáveis de forma. Outros métodos de superposição de marcos anatômicos foram propostos, para revisões, veja (\textbf{monteiro1999?}); (\textbf{adams2004?}); (\textbf{zelditch2012?}), mas a GPA é o mais popular e mais adequado na maioria dos casos.

Na GPA, os efeitos de translação/posição são removidos pois (i) o centróide de cada configuração é transladado para a origem: as coordenadas do centróide de cada configuração são subtraídas da coordenada correspondente de cada landmark, gerando a centralização ou superposição das configurações. Os efeitos de tamanho são removidos pois (ii) as configurações são escaladas/proporcionalizadas para um tamanho de centróide único após a divisão das coordenadas dos landmarks pelo seu tamanho do centróide. E os efeitos de rotação/orientação são removidos pois (iii) a raiz quadrada da soma das distâncias ao quadrado entre os landmarks correspondentes é minimizada usando rotação otimizada por quadrados mínimos; o processo inicia rotacionando as configurações em relação à uma forma de referência, em seguida, uma nova forma de referência (a média da configuração de landmarks) é calculada, e o processo se repete iterativamente até que as diferenças de orientação sejam minimizadas.

Em duas dimensões, dois graus de liberdade são perdidos com a translação, um grau de liberdade é perdido com o escalonamento e um grau de liberdade é perdido com a rotação, totalizando quatro graus de liberdade perdidos. Em três dimensões, três graus de liberdade são perdidos com a translação, um grau de liberdade é perdido com o escalonamento, e três graus de liberdade são perdidos com a rotação, totalizando sete graus de liberdade perdidos.

\section{O objeto .tps no R}\label{o-objeto-.tps-no-r}

O arquivo .tps criado no TPSDig2 pode ser lido no R com a função readland.tps do pacote geomorph. Antes, lembre-se de colocar o arquivo .tps na pasta dadosmg e definir essa pasta como diretório de trabalho no R.

\begin{Shaded}
\begin{Highlighting}[]
\CommentTok{\# Carregar geomorph }
\FunctionTok{require}\NormalTok{(geomorph) }
\CommentTok{\#\textgreater{} Carregando pacotes exigidos: geomorph}
\CommentTok{\#\textgreater{} Warning: package \textquotesingle{}geomorph\textquotesingle{} was built under R version 4.3.2}
\CommentTok{\#\textgreater{} Carregando pacotes exigidos: RRPP}
\CommentTok{\#\textgreater{} Warning: package \textquotesingle{}RRPP\textquotesingle{} was built under R version 4.3.2}
\CommentTok{\#\textgreater{} Carregando pacotes exigidos: rgl}
\CommentTok{\#\textgreater{} Carregando pacotes exigidos: Matrix}
\CommentTok{\# Importar arquivo .tps }
\CommentTok{\# lembre{-}se de substituir por "nomequevocêdeu.tps"}
\NormalTok{tps}\OtherTok{\textless{}{-}}\FunctionTok{readland.tps}\NormalTok{(}\StringTok{"dadosmg/mandibula.dig\_curso.tps"}\NormalTok{,}\AttributeTok{specID =} \StringTok{"ID"}\NormalTok{, }\AttributeTok{readcurves =} \ConstantTok{FALSE}\NormalTok{) }
\FunctionTok{dim}\NormalTok{(tps) }
\NormalTok{tps}
\end{Highlighting}
\end{Shaded}

Podemos visualizar os dados brutos com a função plotAllSpecimens.

\begin{Shaded}
\begin{Highlighting}[]
\FunctionTok{plotAllSpecimens}\NormalTok{(tps)}
\end{Highlighting}
\end{Shaded}

\includegraphics{03.2-mggpa_files/figure-latex/unnamed-chunk-2-1.pdf}

Os círculos maiores representam a posição média do landmark, e os círculos menores representam a posição de cada landmark.

Use indexação com {[}{]} para localizar certos landmarks ou indivíduos. Por exemplo.

\begin{Shaded}
\begin{Highlighting}[]
\CommentTok{\# Deletar espécimes}
\NormalTok{tps.new}\OtherTok{\textless{}{-}}\NormalTok{tps[,,}\SpecialCharTok{{-}}\DecValTok{7}\NormalTok{] }\CommentTok{\# deleta espécime n 7}
\CommentTok{\# Deletar landmarks}
\NormalTok{tps.new}\OtherTok{\textless{}{-}}\NormalTok{tps[}\SpecialCharTok{{-}}\DecValTok{5}\NormalTok{,,] }\CommentTok{\# deleta landmark n 5}
\end{Highlighting}
\end{Shaded}

\section{Estimar landmarks faltantes}\label{estimar-landmarks-faltantes}

A função estimate.missing do pacote \emph{geomorph} implementa dois métodos para estimar landmarks faltantes.

\begin{Shaded}
\begin{Highlighting}[]
\NormalTok{tps.missing}\OtherTok{\textless{}{-}}\FunctionTok{readland.tps}\NormalTok{(}\StringTok{"dadosmg/mandibula.dig\_curso\_missing.tps"}\NormalTok{,}\AttributeTok{specID=}\StringTok{"ID"}\NormalTok{,}\AttributeTok{readcurves=}\ConstantTok{FALSE}\NormalTok{)}
\NormalTok{tps.missing}
\NormalTok{tps.estimated}\OtherTok{\textless{}{-}}\FunctionTok{estimate.missing}\NormalTok{(tps.missing, }\AttributeTok{method =} \StringTok{"TPS"}\NormalTok{)}
\NormalTok{tps.estimated}

\FunctionTok{writeland.tps}\NormalTok{(tps.estimated, }\StringTok{"dadosmg/tps.estimated.tps"}\NormalTok{) }\CommentTok{\# exportar arquivo tps com landmarks estimados}
\end{Highlighting}
\end{Shaded}

\section{Análise Generalizada de Procrustes}\label{anuxe1lise-generalizada-de-procrustes}

Aqui, faremos a análise de Procrustes (GPA) em dados sem simetria de objeto e atribuiremos nomes para os outputs de forma e tamanho do centróide.

\begin{Shaded}
\begin{Highlighting}[]
\NormalTok{gpa.object}\OtherTok{\textless{}{-}}\FunctionTok{gpagen}\NormalTok{(tps)}
\CommentTok{\#\textgreater{} }
\CommentTok{\#\textgreater{} Performing GPA}
\CommentTok{\#\textgreater{} }
  \SpecialCharTok{|}                                                        
  \ErrorTok{|}                                                  \ErrorTok{|}   \DecValTok{0}\NormalTok{\%}
  \SpecialCharTok{|}                                                        
  \ErrorTok{|============}                                      \ErrorTok{|}  \DecValTok{25}\NormalTok{\%}
  \SpecialCharTok{|}                                                        
  \ErrorTok{|=========================}                         \ErrorTok{|}  \DecValTok{50}\NormalTok{\%}
  \SpecialCharTok{|}                                                        
  \ErrorTok{|==================================================|} \DecValTok{100}\NormalTok{\%}
\CommentTok{\#\textgreater{} }
\CommentTok{\#\textgreater{} Making projections... Finished!}
\NormalTok{gpa.object}
\CommentTok{\#\textgreater{} }
\CommentTok{\#\textgreater{} Call:}
\CommentTok{\#\textgreater{} gpagen(A = tps) }
\CommentTok{\#\textgreater{} }
\CommentTok{\#\textgreater{} }
\CommentTok{\#\textgreater{} }
\CommentTok{\#\textgreater{} Generalized Procrustes Analysis}
\CommentTok{\#\textgreater{} with Partial Procrustes Superimposition}
\CommentTok{\#\textgreater{} }
\CommentTok{\#\textgreater{} 13 fixed landmarks}
\CommentTok{\#\textgreater{} 0 semilandmarks (sliders)}
\CommentTok{\#\textgreater{} 2{-}dimensional landmarks}
\CommentTok{\#\textgreater{} 2 GPA iterations to converge}
\CommentTok{\#\textgreater{} }
\CommentTok{\#\textgreater{} }
\CommentTok{\#\textgreater{} Consensus (mean) Configuration}
\CommentTok{\#\textgreater{} }
\CommentTok{\#\textgreater{}                X           Y}
\CommentTok{\#\textgreater{} 1  {-}0.3600904084 {-}0.03209176}
\CommentTok{\#\textgreater{} 2  {-}0.3838537500  0.08920190}
\CommentTok{\#\textgreater{} 3  {-}0.2678074781  0.01674324}
\CommentTok{\#\textgreater{} 4  {-}0.2101661423  0.05684845}
\CommentTok{\#\textgreater{} 5  {-}0.0001796958  0.03565707}
\CommentTok{\#\textgreater{} 6   0.1873051760  0.14767796}
\CommentTok{\#\textgreater{} 7   0.1660719417  0.08697077}
\CommentTok{\#\textgreater{} 8   0.2984433258  0.09285954}
\CommentTok{\#\textgreater{} 9   0.3486394512  0.01069463}
\CommentTok{\#\textgreater{} 10  0.2335846883 {-}0.07886796}
\CommentTok{\#\textgreater{} 11  0.2889938410 {-}0.17663723}
\CommentTok{\#\textgreater{} 12 {-}0.0281091968 {-}0.13445727}
\CommentTok{\#\textgreater{} 13 {-}0.2728317526 {-}0.11459935}
\NormalTok{shape}\OtherTok{\textless{}{-}}\NormalTok{gpa.object}\SpecialCharTok{$}\NormalTok{coords}
\NormalTok{size}\OtherTok{\textless{}{-}}\NormalTok{gpa.object}\SpecialCharTok{$}\NormalTok{Csize}
\end{Highlighting}
\end{Shaded}

Podemos visualizar a distribuição dos dados de tamanho e forma. Note a diferença no gráfico da forma antes e depois da análise de Procrustes.

\begin{Shaded}
\begin{Highlighting}[]
\FunctionTok{hist}\NormalTok{(size)}
\end{Highlighting}
\end{Shaded}

\begin{Shaded}
\begin{Highlighting}[]
\FunctionTok{plotAllSpecimens}\NormalTok{(shape)}
\end{Highlighting}
\end{Shaded}

\includegraphics{03.2-mggpa_files/figure-latex/unnamed-chunk-7-1.pdf}

\section{Espaço tangente}\label{espauxe7o-tangente}

A correlação do espaço de forma com o espaço tangente quase sempre será muito alta em dados de morfometria geométrica, justificando o uso de estatística Euclidiana.

\begin{Shaded}
\begin{Highlighting}[]
\FunctionTok{require}\NormalTok{(Morpho)}
\CommentTok{\#\textgreater{} Warning: package \textquotesingle{}Morpho\textquotesingle{} was built under R version 4.3.2}
\FunctionTok{regdist}\NormalTok{(shape)}
\end{Highlighting}
\end{Shaded}

\includegraphics{03.2-mggpa_files/figure-latex/unnamed-chunk-8-1.pdf}

\begin{verbatim}
#> $cor
#> [1] 0.9999991
#> 
#> $procSS
#> [1] 0.1789463
#> 
#> $tanSS
#> [1] 0.1788539
#> 
#> $rhoSS
#> [1] 0.1791391
\end{verbatim}

\section{Formato array e matrix}\label{formato-array-e-matrix}

Os dados de forma estão em formato array. Podemos nos mover facilmente entre os formatos array e matrix. Ambos carregam a mesma informação, mas algumas análises/operações só aceitam um dos dois formatos.

\section{GPA com simetria de objeto}\label{gpa-com-simetria-de-objeto}

Quando os dados tiverem simetria e objeto, vamos informar isso ao R com uma matriz contendo os pares de landmarks simétricos e realizar uma GPA considerando a simetria.

\begin{Shaded}
\begin{Highlighting}[]
\CommentTok{\# Carregar dados}
\NormalTok{tps.sim}\OtherTok{\textless{}{-}}\FunctionTok{readland.tps}\NormalTok{(}\StringTok{"dadosmg/Lista tuco dig 2.tps"}\NormalTok{,}\AttributeTok{specID =} \StringTok{"ID"}\NormalTok{, }\AttributeTok{readcurves =} \ConstantTok{FALSE}\NormalTok{)}
\FunctionTok{dim}\NormalTok{(tps.sim)}
\CommentTok{\# Matriz com landmarks simétricos}
\NormalTok{pairs.matrix}\OtherTok{\textless{}{-}}\FunctionTok{matrix}\NormalTok{(}\FunctionTok{c}\NormalTok{(}\DecValTok{2}\NormalTok{,}\DecValTok{3}\NormalTok{,}\DecValTok{5}\NormalTok{,}\DecValTok{6}\NormalTok{,}\DecValTok{7}\NormalTok{,}\DecValTok{8}\NormalTok{,}\DecValTok{10}\NormalTok{,}\DecValTok{11}\NormalTok{,}\DecValTok{12}\NormalTok{,}\DecValTok{13}\NormalTok{,}\DecValTok{14}\NormalTok{,}\DecValTok{15}\NormalTok{,}\DecValTok{16}\NormalTok{,}\DecValTok{17}\NormalTok{,}\DecValTok{18}\NormalTok{,}\DecValTok{19}\NormalTok{,}\DecValTok{20}\NormalTok{,}\DecValTok{21}\NormalTok{,}\DecValTok{23}\NormalTok{,}\DecValTok{24}\NormalTok{,}\DecValTok{25}\NormalTok{,}\DecValTok{26}\NormalTok{,}\DecValTok{27}\NormalTok{,}\DecValTok{28}\NormalTok{),}\AttributeTok{nrow=}\DecValTok{12}\NormalTok{,}\AttributeTok{ncol=}\DecValTok{2}\NormalTok{,}\AttributeTok{byrow=}\NormalTok{T)}
\NormalTok{pairs.matrix}
\CommentTok{\# Construir vetor com rótulo de indivíduos (neste caso, uma marcação por indivíduo)}
\NormalTok{ind}\OtherTok{\textless{}{-}}\FunctionTok{c}\NormalTok{(}\DecValTok{1}\SpecialCharTok{:}\FunctionTok{dim}\NormalTok{(tps.sim)[}\DecValTok{3}\NormalTok{])}
\CommentTok{\# GPA com simetria bilateral}
\NormalTok{b.s}\OtherTok{\textless{}{-}}\FunctionTok{bilat.symmetry}\NormalTok{(tps.sim,}\AttributeTok{ind=}\NormalTok{ind,}\AttributeTok{object.sym=}\ConstantTok{TRUE}\NormalTok{,}\AttributeTok{land.pairs=}\NormalTok{pairs.matrix)}

\NormalTok{shape.sym}\OtherTok{\textless{}{-}}\NormalTok{b.s}\SpecialCharTok{$}\NormalTok{symm.shape }\CommentTok{\# componente simétrico da forma}
\FunctionTok{plotAllSpecimens}\NormalTok{(shape.sym)}
\end{Highlighting}
\end{Shaded}

\includegraphics{03.2-mggpa_files/figure-latex/unnamed-chunk-10-1.pdf}

\section{Encontrando outliers}\label{encontrando-outliers}

Podemos encontrar indivíduos muito diferentes dos demais observando sua distância em relação à forma média da amostra. Em morfometria geométrica, é comum a marcação de landmarks trocados em alguns indivíduos, o que vai gerar óbvios outliers.

\begin{Shaded}
\begin{Highlighting}[]
\FunctionTok{plotOutliers}\NormalTok{(shape)}
\end{Highlighting}
\end{Shaded}

\includegraphics{03.2-mggpa_files/figure-latex/unnamed-chunk-11-1.pdf}

\chapter{Visualização da forma}\label{visualizauxe7uxe3o-da-forma}

Existem várias maneiras de visualizar mudanças de forma com dados de Morfometria Geométrica (ver (\textbf{klingenberg2013?})). Aliás, essa é uma das vantagens da morfometria geométrica em comparação com a morfometria tradicional. As formas mais comuns de visualização são por grids de deformação, vetores (ou lollipops), e através da construção de links entre os landmarks e extrapolação da forma via \emph{outlines} e \emph{surfaces}. Em qualquer caso, é importante perceber que a diferença de forma está contida apenas entre os landmarks/semilandmarks; as extrapolações de forma além dos landmarks (via grids e outlines, por exemplo) podem ser úteis para visualização, mas devem ser interpretadas com bastante cautela. É também importante perceber que a interpretação das mudanças de forma não deve considerar os landmarks individualmente, mas sim a sua mudança relativa a todos os outros landmarks.

Construção de forma média da amostra para comparação.

\begin{Shaded}
\begin{Highlighting}[]
\NormalTok{ref}\OtherTok{\textless{}{-}}\FunctionTok{mshape}\NormalTok{(shape)}
\end{Highlighting}
\end{Shaded}

Construção de links entre os landmarks.

\begin{Shaded}
\begin{Highlighting}[]
\NormalTok{links}\OtherTok{\textless{}{-}}\FunctionTok{matrix}\NormalTok{(}\FunctionTok{c}\NormalTok{(}\DecValTok{1}\NormalTok{,}\DecValTok{2}\NormalTok{,}\DecValTok{2}\NormalTok{,}\DecValTok{3}\NormalTok{,}\DecValTok{3}\NormalTok{,}\DecValTok{4}\NormalTok{,}\DecValTok{4}\NormalTok{,}\DecValTok{5}\NormalTok{,}\DecValTok{5}\NormalTok{,}\DecValTok{6}\NormalTok{,}\DecValTok{6}\NormalTok{,}\DecValTok{7}\NormalTok{,}\DecValTok{7}\NormalTok{,}\DecValTok{8}\NormalTok{,}\DecValTok{8}\NormalTok{,}\DecValTok{9}\NormalTok{,}\DecValTok{9}\NormalTok{,}\DecValTok{10}\NormalTok{,}\DecValTok{10}\NormalTok{,}\DecValTok{11}\NormalTok{,}\DecValTok{11}\NormalTok{,}\DecValTok{12}\NormalTok{,}\DecValTok{12}\NormalTok{,}\DecValTok{13}\NormalTok{,}\DecValTok{13}\NormalTok{,}\DecValTok{1}\NormalTok{),}\AttributeTok{nrow=}\DecValTok{13}\NormalTok{,}\AttributeTok{ncol=}\DecValTok{2}\NormalTok{,}\AttributeTok{byrow=}\NormalTok{T)}
\FunctionTok{plotAllSpecimens}\NormalTok{(shape,}\AttributeTok{mean=}\ConstantTok{TRUE}\NormalTok{,}\AttributeTok{links=}\NormalTok{links)}
\end{Highlighting}
\end{Shaded}

\includegraphics{03.2-mggpa_files/figure-latex/unnamed-chunk-13-1.pdf}

Os links entre os landmarks também podem ser definidos interativamente com a função define.links do pacote \emph{geomorph}.

\section{Método de grid de deformação.}\label{muxe9todo-de-grid-de-deformauxe7uxe3o.}

\begin{Shaded}
\begin{Highlighting}[]
\CommentTok{\# Definir cor dos landmarks e links}
\NormalTok{GP1}\OtherTok{\textless{}{-}}\FunctionTok{gridPar}\NormalTok{(}\AttributeTok{pt.bg=}\StringTok{"gray"}\NormalTok{,}\AttributeTok{link.col=}\StringTok{"gray"}\NormalTok{,}\AttributeTok{link.lty=}\DecValTok{1}\NormalTok{) }\CommentTok{\# cor dos landmarks e links}
\CommentTok{\# Gráfico de mudança de forma}
\FunctionTok{plotRefToTarget}\NormalTok{(ref,shape[,,}\DecValTok{8}\NormalTok{],}\AttributeTok{links=}\NormalTok{links,}\AttributeTok{method=}\StringTok{"TPS"}\NormalTok{)}
\end{Highlighting}
\end{Shaded}

\includegraphics{03.2-mggpa_files/figure-latex/unnamed-chunk-14-1.pdf}

\section{Método de vetores.}\label{muxe9todo-de-vetores.}

\begin{Shaded}
\begin{Highlighting}[]
\FunctionTok{plotRefToTarget}\NormalTok{(ref,shape[,,}\DecValTok{8}\NormalTok{],}\AttributeTok{links=}\NormalTok{links,}\AttributeTok{method=}\StringTok{"vector"}\NormalTok{)}
\end{Highlighting}
\end{Shaded}

\includegraphics{03.2-mggpa_files/figure-latex/unnamed-chunk-15-1.pdf}

\section{Método de linhas.}\label{muxe9todo-de-linhas.}

\begin{Shaded}
\begin{Highlighting}[]
\FunctionTok{plotRefToTarget}\NormalTok{(ref,shape[,,}\DecValTok{8}\NormalTok{],}\AttributeTok{links=}\NormalTok{links,}\AttributeTok{method=}\StringTok{"points"}\NormalTok{,}\AttributeTok{gridPars=}\NormalTok{GP1)}
\end{Highlighting}
\end{Shaded}

\includegraphics{03.2-mggpa_files/figure-latex/unnamed-chunk-16-1.pdf}

Aqui, o alvo está em preto e a referência(média) está em cinza.

\section{Método de outlines.}\label{muxe9todo-de-outlines.}

\begin{Shaded}
\begin{Highlighting}[]
\CommentTok{\# Carregando outline}
\NormalTok{outline}\OtherTok{\textless{}{-}}\FunctionTok{warpRefOutline}\NormalTok{(}\StringTok{"dadosmg/outline.mandibula.R.tps"}\NormalTok{, tps[,,}\DecValTok{11}\NormalTok{], ref)}
\end{Highlighting}
\end{Shaded}

\begin{Shaded}
\begin{Highlighting}[]
\CommentTok{\# Gráfico de mudança de forma}
\NormalTok{GP1}\OtherTok{\textless{}{-}}\FunctionTok{gridPar}\NormalTok{(}\AttributeTok{pt.bg=}\StringTok{"gray"}\NormalTok{,}\AttributeTok{tar.out.col =} \StringTok{"red"}\NormalTok{,}\AttributeTok{tar.out.cex =} \FloatTok{0.5}\NormalTok{)}
\FunctionTok{plotRefToTarget}\NormalTok{(ref,shape[,,}\DecValTok{8}\NormalTok{],}\AttributeTok{outline=}\NormalTok{outline}\SpecialCharTok{$}\NormalTok{outline,}\AttributeTok{method=}\StringTok{"TPS"}\NormalTok{)}
\end{Highlighting}
\end{Shaded}

\includegraphics{03.2-mggpa_files/figure-latex/unnamed-chunk-18-1.pdf}

\begin{Shaded}
\begin{Highlighting}[]
\FunctionTok{plotRefToTarget}\NormalTok{(ref,shape[,,}\DecValTok{8}\NormalTok{],}\AttributeTok{outline=}\NormalTok{outline}\SpecialCharTok{$}\NormalTok{outline,}\AttributeTok{method=}\StringTok{"points"}\NormalTok{,}\AttributeTok{gridPars=}\NormalTok{GP1)}
\end{Highlighting}
\end{Shaded}

\includegraphics{03.2-mggpa_files/figure-latex/unnamed-chunk-18-2.pdf}

Com um número suficiente de pontos, dá-se a impressão de um desenho. Veja o exemplo do pacote \emph{geomorph}.

\begin{Shaded}
\begin{Highlighting}[]
\FunctionTok{data}\NormalTok{(}\StringTok{"plethodon"}\NormalTok{) }
\NormalTok{Y.gpa}\OtherTok{\textless{}{-}}\FunctionTok{gpagen}\NormalTok{(plethodon}\SpecialCharTok{$}\NormalTok{land, }\AttributeTok{print.progress =} \ConstantTok{FALSE}\NormalTok{) }
\NormalTok{ref.y}\OtherTok{\textless{}{-}}\FunctionTok{mshape}\NormalTok{(Y.gpa}\SpecialCharTok{$}\NormalTok{coords) }
\NormalTok{GP4}\OtherTok{\textless{}{-}} \FunctionTok{gridPar}\NormalTok{(}\AttributeTok{tar.out.col =} \StringTok{"red"}\NormalTok{, }\AttributeTok{tar.out.cex =} \FloatTok{0.3}\NormalTok{)}
\FunctionTok{plotRefToTarget}\NormalTok{(ref.y,Y.gpa}\SpecialCharTok{$}\NormalTok{coords[,,}\DecValTok{39}\NormalTok{], }\AttributeTok{gridPars=}\NormalTok{GP4, }\AttributeTok{mag=}\DecValTok{3}\NormalTok{, }\AttributeTok{outline=}\NormalTok{plethodon}\SpecialCharTok{$}\NormalTok{outline, }\AttributeTok{method=}\StringTok{"TPS"}\NormalTok{)}
\end{Highlighting}
\end{Shaded}

\includegraphics{03.2-mggpa_files/figure-latex/unnamed-chunk-19-1.pdf}

\section{Magnificação da forma}\label{magnificauxe7uxe3o-da-forma}

Às vezes quando as diferenças de forma são muito sútis, podemos aumentar as diferenças para visualizar melhor. Use com cuidado.

\begin{Shaded}
\begin{Highlighting}[]
\FunctionTok{plotRefToTarget}\NormalTok{(ref,shape[,,}\DecValTok{8}\NormalTok{],}\AttributeTok{outline=}\NormalTok{outline}\SpecialCharTok{$}\NormalTok{outline, }\AttributeTok{method=}\StringTok{"points"}\NormalTok{,}\AttributeTok{gridPars=}\NormalTok{GP1,}\AttributeTok{mag=}\DecValTok{3}\NormalTok{)}
\end{Highlighting}
\end{Shaded}

\includegraphics{03.2-mggpa_files/figure-latex/unnamed-chunk-20-1.pdf}

\chapter{Assimetria e quantificação de erro}\label{assimetria-e-quantificauxe7uxe3o-de-erro}

Duas fontes de adicionais de variação, normalmente bem pequenas, podem ser quantificadas: a assimetria e os erros oridundos da obtenção de dados e da medição de indivíduos. A assimetria é o menor dos efeitos biológicos, e pode ser quantificada através da diferença entre os lados direito e esquerdo de estruturas com simetria bilateral, ou entre estruturas correspondentes em outras formas de simetria, como em asas de insetos ou pétalas de flores. Erros oriundos da obtenção de dados ou de medição dos espécimes (colocação de landmarks) devem também ser muito pequenos para que possam ser ignorados com confiança. Espera-se que esses erros sejam várias vezes menores do que variação entre indivíduos, e mesmo que sejam bastante menores do que a variação vinda da assimetria flutuante.

Para se aprofundar, leia: (\textbf{klingenberg2002?}); (\textbf{klingenberg2015?}).

\section{Assimetria de objeto}\label{assimetria-de-objeto}

Estruturas biológicas com assimetria de objeto, por exemplo, os lados direito e esquerdo do crânio.

\begin{Shaded}
\begin{Highlighting}[]
\CommentTok{\# Carregar dados com landmarks no crânio de tuco{-}tucos}
\FunctionTok{require}\NormalTok{(geomorph)}
\CommentTok{\#\textgreater{} Carregando pacotes exigidos: geomorph}
\CommentTok{\#\textgreater{} Carregando pacotes exigidos: RRPP}
\CommentTok{\#\textgreater{} Carregando pacotes exigidos: rgl}
\CommentTok{\#\textgreater{} Carregando pacotes exigidos: Matrix}
\NormalTok{tps.sim}\OtherTok{\textless{}{-}}\FunctionTok{readland.tps}\NormalTok{(}\StringTok{"dadosmg/Lista tuco dig 2.tps"}\NormalTok{,}\AttributeTok{specID =} \StringTok{"ID"}\NormalTok{, }\AttributeTok{readcurves =} \ConstantTok{FALSE}\NormalTok{)}
\FunctionTok{dim}\NormalTok{(tps.sim)}
\CommentTok{\# Matriz com landmarks simétricos}
\NormalTok{pairs.matrix}\OtherTok{\textless{}{-}}\FunctionTok{matrix}\NormalTok{(}\FunctionTok{c}\NormalTok{(}\DecValTok{2}\NormalTok{,}\DecValTok{3}\NormalTok{,}\DecValTok{5}\NormalTok{,}\DecValTok{6}\NormalTok{,}\DecValTok{7}\NormalTok{,}\DecValTok{8}\NormalTok{,}\DecValTok{10}\NormalTok{,}\DecValTok{11}\NormalTok{,}\DecValTok{12}\NormalTok{,}\DecValTok{13}\NormalTok{,}\DecValTok{14}\NormalTok{,}\DecValTok{15}\NormalTok{,}\DecValTok{16}\NormalTok{,}\DecValTok{17}\NormalTok{,}\DecValTok{18}\NormalTok{,}\DecValTok{19}\NormalTok{,}\DecValTok{20}\NormalTok{,}\DecValTok{21}\NormalTok{,}\DecValTok{23}\NormalTok{,}\DecValTok{24}\NormalTok{,}\DecValTok{25}\NormalTok{,}\DecValTok{26}\NormalTok{,}\DecValTok{27}\NormalTok{,}\DecValTok{28}\NormalTok{),}\AttributeTok{nrow=}\DecValTok{12}\NormalTok{,}\AttributeTok{ncol=}\DecValTok{2}\NormalTok{,}\AttributeTok{byrow=}\NormalTok{T)}
\NormalTok{pairs.matrix}
\CommentTok{\# Construir vetor com rótulo de indivíduos (neste caso, uma marcação por indivíduo)}
\NormalTok{ind}\OtherTok{\textless{}{-}}\FunctionTok{c}\NormalTok{(}\DecValTok{1}\SpecialCharTok{:}\FunctionTok{dim}\NormalTok{(tps.sim)[}\DecValTok{3}\NormalTok{])}
\CommentTok{\# GPA com simetria bilateral}
\NormalTok{b.s}\OtherTok{\textless{}{-}}\FunctionTok{bilat.symmetry}\NormalTok{(tps.sim,}\AttributeTok{ind=}\NormalTok{ind,}\AttributeTok{object.sym=}\ConstantTok{TRUE}\NormalTok{,}\AttributeTok{land.pairs=}\NormalTok{pairs.matrix)}

\CommentTok{\# Componente assimétrico da forma}
\NormalTok{shape.asym}\OtherTok{\textless{}{-}}\NormalTok{b.s}\SpecialCharTok{$}\NormalTok{asymm.shape}
\FunctionTok{plotAllSpecimens}\NormalTok{(shape.asym)}
\end{Highlighting}
\end{Shaded}

\includegraphics{03.3-mgerro_files/figure-latex/unnamed-chunk-1-1.pdf}

\section{Assimetria de correspondência}\label{assimetria-de-corresponduxeancia}

Estruturas biológicas com simetria de correspondência, por exemplo, asas de insetos.

\begin{Shaded}
\begin{Highlighting}[]
\CommentTok{\# Carregar arquivo com landmarks em asas de *Drosophila*}
\NormalTok{fly}\OtherTok{\textless{}{-}}\FunctionTok{read.table}\NormalTok{(}\StringTok{"dadosmg/fly wing replicates Proc.txt"}\NormalTok{,}\AttributeTok{h=}\NormalTok{T,}\AttributeTok{row.names =} \DecValTok{1}\NormalTok{)}
\FunctionTok{head}\NormalTok{(fly)}
\CommentTok{\# Transformar coordenadas em array}
\NormalTok{fly.tps}\OtherTok{\textless{}{-}}\FunctionTok{arrayspecs}\NormalTok{(fly[,}\DecValTok{5}\SpecialCharTok{:}\DecValTok{34}\NormalTok{],}\DecValTok{15}\NormalTok{,}\DecValTok{2}\NormalTok{) }
\NormalTok{fly.tps }
\FunctionTok{plotAllSpecimens}\NormalTok{(fly.tps)}
\end{Highlighting}
\end{Shaded}

\includegraphics{03.3-mgerro_files/figure-latex/unnamed-chunk-2-1.pdf}

\begin{Shaded}
\begin{Highlighting}[]
\CommentTok{\# Vetor com os indivíduos da amostra}
\NormalTok{fly.ind}\OtherTok{\textless{}{-}}\NormalTok{fly}\SpecialCharTok{$}\NormalTok{fly }
\NormalTok{fly.ind}
\CommentTok{\# Vetor informando quais estruturas são asas do lado direito ou do esquerdo}
\NormalTok{fly.side}\OtherTok{\textless{}{-}}\NormalTok{fly}\SpecialCharTok{$}\NormalTok{side }
\NormalTok{fly.side}
\CommentTok{\# GPA com simetria de correspondência}
\NormalTok{fly.b.s}\OtherTok{\textless{}{-}}\FunctionTok{bilat.symmetry}\NormalTok{(fly.tps,}\AttributeTok{ind=}\NormalTok{fly.ind,}\AttributeTok{side=}\NormalTok{fly.side)}
\CommentTok{\# Componente simétrico da forma}
\FunctionTok{plotAllSpecimens}\NormalTok{(fly.b.s}\SpecialCharTok{$}\NormalTok{symm.shape) }
\end{Highlighting}
\end{Shaded}

\includegraphics{03.3-mgerro_files/figure-latex/unnamed-chunk-2-2.pdf}

\begin{Shaded}
\begin{Highlighting}[]
\CommentTok{\# Componente assimétrico da forma}
\FunctionTok{plotAllSpecimens}\NormalTok{(fly.b.s}\SpecialCharTok{$}\NormalTok{asymm.shape)}
\end{Highlighting}
\end{Shaded}

\includegraphics{03.3-mgerro_files/figure-latex/unnamed-chunk-2-3.pdf}

\section{Assimetria direcional}\label{assimetria-direcional}

A forma média dos lados direito e esquerdo.

\begin{Shaded}
\begin{Highlighting}[]
\NormalTok{fly.b.s}\SpecialCharTok{$}\NormalTok{DA.component}
\end{Highlighting}
\end{Shaded}

\section{Assimetria flutuante}\label{assimetria-flutuante}

O desvio específico de cada espécime e índices quantificando o desvio esperado da simetria.

\begin{Shaded}
\begin{Highlighting}[]
\NormalTok{fly.b.s}\SpecialCharTok{$}\NormalTok{FA.component}
\NormalTok{fly.b.s}\SpecialCharTok{$}\NormalTok{signed.AI}
\end{Highlighting}
\end{Shaded}

\section{Quantificação de erro}\label{quantificauxe7uxe3o-de-erro}

A ``Procrustes ANOVA'' (Goodall 1991; Klingengerg \& McIntyre 1998), usada para quantificar diversas fontes de erro em morfometria geométrica, é uma análise multivariada de variância onde as fontes de erro são aninhadas uma dentro da outra. Por exemplo, podemos ter os lados direito/esquerdo de cada indivíduo, mais de uma marcação de landmarks para cada indivíduo, e mais de uma obtenção de imagem para cada indivíduo. A quantificação de variação em cada nível pode ser feita e comparada. No exemplo abaixo, vamos quantificar o erro de medição através de marcações de landmarks repetidas duas vezes em cada indivíduo da amostra. Queremos obter uma medida que represente a confiabilidade da marcação de landmarks por um operador, obtida através da repetição da colocação de landmarks em imagens idênticas, por esse operador; em teoria, os landmarks deveriam estar extamente na mesma posição.

\begin{Shaded}
\begin{Highlighting}[]
\CommentTok{\# Carregar arquivo tps com todos os indivíduos marcados 2 vezes}
\NormalTok{tps.rep}\OtherTok{\textless{}{-}}\FunctionTok{readland.tps}\NormalTok{(}\StringTok{"dadosmg/ventral.dig\_curso\_rep1\_2.tps"}\NormalTok{,}\AttributeTok{specID =} \StringTok{"ID"}\NormalTok{, }\AttributeTok{readcurves =} \ConstantTok{FALSE}\NormalTok{) }
\FunctionTok{dim}\NormalTok{(tps.rep)}
\CommentTok{\# Matriz com os pares de landmarks simétricos}
\NormalTok{pairs.matrix.rep}\OtherTok{\textless{}{-}}\FunctionTok{matrix}\NormalTok{(}\FunctionTok{c}\NormalTok{(}\DecValTok{2}\NormalTok{,}\DecValTok{3}\NormalTok{,}\DecValTok{4}\NormalTok{,}\DecValTok{5}\NormalTok{,}\DecValTok{6}\NormalTok{,}\DecValTok{7}\NormalTok{,}\DecValTok{8}\NormalTok{,}\DecValTok{9}\NormalTok{,}\DecValTok{10}\NormalTok{,}\DecValTok{11}\NormalTok{,}\DecValTok{12}\NormalTok{,}\DecValTok{13}\NormalTok{,}\DecValTok{14}\NormalTok{,}\DecValTok{15}\NormalTok{,}\DecValTok{16}\NormalTok{,}\DecValTok{17}\NormalTok{,}\DecValTok{18}\NormalTok{,}\DecValTok{24}\NormalTok{,}\DecValTok{22}\NormalTok{,}\DecValTok{28}\NormalTok{,}\DecValTok{23}\NormalTok{,}\DecValTok{29}\NormalTok{,}\DecValTok{19}\NormalTok{,}\DecValTok{25}\NormalTok{,}\DecValTok{20}\NormalTok{,}\DecValTok{26}\NormalTok{,}\DecValTok{21}\NormalTok{,}\DecValTok{27}\NormalTok{,}\DecValTok{30}\NormalTok{,}\DecValTok{31}\NormalTok{,}\DecValTok{33}\NormalTok{,}\DecValTok{34}\NormalTok{,}\DecValTok{35}\NormalTok{,}\DecValTok{36}\NormalTok{,}\DecValTok{37}\NormalTok{,}\DecValTok{38}\NormalTok{,}\DecValTok{39}\NormalTok{,}\DecValTok{40}\NormalTok{,}\DecValTok{41}\NormalTok{,}\DecValTok{42}\NormalTok{,}\DecValTok{43}\NormalTok{,}\DecValTok{44}\NormalTok{,}\DecValTok{45}\NormalTok{,}\DecValTok{46}\NormalTok{,}\DecValTok{47}\NormalTok{,}\DecValTok{48}\NormalTok{,}\DecValTok{49}\NormalTok{,}\DecValTok{50}\NormalTok{,}\DecValTok{51}\NormalTok{,}\DecValTok{52}\NormalTok{,}\DecValTok{55}\NormalTok{,}\DecValTok{56}\NormalTok{),}\AttributeTok{nrow=}\DecValTok{26}\NormalTok{,}\AttributeTok{ncol=}\DecValTok{2}\NormalTok{,}\AttributeTok{byrow=}\NormalTok{T) }
\NormalTok{pairs.matrix.rep}
\CommentTok{\# Extração dos rótulos de indivíduos e repetições a partir do ID}
\NormalTok{categories}\OtherTok{\textless{}{-}}\FunctionTok{strsplit}\NormalTok{(}\FunctionTok{dimnames}\NormalTok{(tps.rep)[[}\DecValTok{3}\NormalTok{]], }\StringTok{"\_"}\NormalTok{) }\CommentTok{\# separa os nomes pelo underline}
\NormalTok{classifiers}\OtherTok{\textless{}{-}}\FunctionTok{matrix}\NormalTok{(}\FunctionTok{unlist}\NormalTok{(categories),}\AttributeTok{ncol=}\DecValTok{5}\NormalTok{,}\AttributeTok{byrow=}\NormalTok{T) }\CommentTok{\# transforma em matriz}
\NormalTok{classifiers}\OtherTok{\textless{}{-}}\FunctionTok{cbind}\NormalTok{(}\FunctionTok{dimnames}\NormalTok{(tps.rep)[[}\DecValTok{3}\NormalTok{]],classifiers) }\CommentTok{\# adiciona o ID na primeira coluna}
\FunctionTok{colnames}\NormalTok{(classifiers)}\OtherTok{\textless{}{-}}\FunctionTok{c}\NormalTok{(}\StringTok{"FileID"}\NormalTok{,}\StringTok{"Species"}\NormalTok{,}\StringTok{"Sex"}\NormalTok{,}\StringTok{"View"}\NormalTok{,}\StringTok{"Ind"}\NormalTok{,}\StringTok{"Rep"}\NormalTok{) }\CommentTok{\# renomeia as colunas}
\NormalTok{classifiers}\OtherTok{\textless{}{-}}\FunctionTok{as.data.frame}\NormalTok{(classifiers)}
\CommentTok{\# Indivíduos}
\NormalTok{ind.rep}\OtherTok{\textless{}{-}}\NormalTok{classifiers}\SpecialCharTok{$}\NormalTok{Ind }\CommentTok{\# indivíduos}
\CommentTok{\# Repetições}
\NormalTok{rep.rep}\OtherTok{\textless{}{-}}\NormalTok{classifiers}\SpecialCharTok{$}\NormalTok{Rep }\CommentTok{\# repetições}

\CommentTok{\# Procrustes ANOVA {-} erro de medição}
\NormalTok{b.s.error}\OtherTok{\textless{}{-}}\FunctionTok{bilat.symmetry}\NormalTok{(tps.rep,}\AttributeTok{ind=}\NormalTok{ind.rep,}\AttributeTok{replicate=}\NormalTok{rep.rep,}\AttributeTok{object.sym=}\ConstantTok{TRUE}\NormalTok{,}\AttributeTok{land.pairs=}\NormalTok{pairs.matrix.rep)}
\end{Highlighting}
\end{Shaded}

A tabela de variância informa quanta variação está contida entre as réplicas de um mesmo indivíduo marcados duas vezes (ind:side:replicate).

\begin{Shaded}
\begin{Highlighting}[]
\NormalTok{b.s.error}\SpecialCharTok{$}\NormalTok{shape.anova}
\CommentTok{\#\textgreater{}                     Df      SS        MS     Rsq       F}
\CommentTok{\#\textgreater{} ind                 39 0.34535 0.0088550 0.92935 30.8480}
\CommentTok{\#\textgreater{} side                 1 0.00466 0.0046638 0.01255 16.2470}
\CommentTok{\#\textgreater{} ind:side            39 0.01120 0.0002871 0.03013  2.2096}
\CommentTok{\#\textgreater{} ind:side:replicate  80 0.01039 0.0001299 0.02797        }
\CommentTok{\#\textgreater{} Total              159 0.37160                          }
\CommentTok{\#\textgreater{}                          Z Pr(\textgreater{}F)   }
\CommentTok{\#\textgreater{} ind                 6.8520  0.001 **}
\CommentTok{\#\textgreater{} side                5.2769  0.001 **}
\CommentTok{\#\textgreater{} ind:side           13.2872  0.001 **}
\CommentTok{\#\textgreater{} ind:side:replicate                  }
\CommentTok{\#\textgreater{} Total                               }
\CommentTok{\#\textgreater{} {-}{-}{-}}
\CommentTok{\#\textgreater{} Signif. codes:  }
\CommentTok{\#\textgreater{} 0 \textquotesingle{}***\textquotesingle{} 0.001 \textquotesingle{}**\textquotesingle{} 0.01 \textquotesingle{}*\textquotesingle{} 0.05 \textquotesingle{}.\textquotesingle{} 0.1 \textquotesingle{} \textquotesingle{} 1}
\end{Highlighting}
\end{Shaded}

O erro de medição é pelo menos duas vezes menor do que a assimetria flutuante.

Podemos também calcular um valor de confiança da repetibilidade

\begin{Shaded}
\begin{Highlighting}[]
\NormalTok{R1}\OtherTok{\textless{}{-}}\NormalTok{(((b.s.error}\SpecialCharTok{$}\NormalTok{shape.anova}\SpecialCharTok{$}\NormalTok{MS[}\DecValTok{1}\NormalTok{]}\SpecialCharTok{{-}}\NormalTok{b.s.error}\SpecialCharTok{$}\NormalTok{shape.anova}\SpecialCharTok{$}\NormalTok{MS[}\DecValTok{4}\NormalTok{])}\SpecialCharTok{/}\DecValTok{2}\NormalTok{) }\SpecialCharTok{/}\NormalTok{ ((b.s.error}\SpecialCharTok{$}\NormalTok{shape.anova}\SpecialCharTok{$}\NormalTok{MS[}\DecValTok{4}\NormalTok{]}\SpecialCharTok{+}\NormalTok{((b.s.error}\SpecialCharTok{$}\NormalTok{shape.anova}\SpecialCharTok{$}\NormalTok{MS[}\DecValTok{1}\NormalTok{]}\SpecialCharTok{{-}}\NormalTok{b.s.error}\SpecialCharTok{$}\NormalTok{shape.anova}\SpecialCharTok{$}\NormalTok{MS[}\DecValTok{4}\NormalTok{])}\SpecialCharTok{/}\DecValTok{2}\NormalTok{)))) }
\NormalTok{R1}
\CommentTok{\#\textgreater{} [1] 0.9710822}
\end{Highlighting}
\end{Shaded}

\section{Exercício - Erro e Assimetria}\label{exercuxedcio---erro-e-assimetria}

Usando os dados de asas de Drosophila já carregados (arquivo ``fly wing replicates Proc.txt'') construa um modelo incluindo, além dos lados direto e esquerdo (assimetria) e da repetição das fotos entre indivíduos, também o fator imagem (3° coluna dos dados), que representa a repetição da captura de fotos, i.e.~fotos diferentes do mesmo indivíduo. Uma maneira equivalente de representar a Procrustes ANOVA {[}bilat.symmetry(fly.tps,ind=fly.ind,side=fly.side,replicate=fly.rep{]} com modelos lineares em linguagem geomoph é usando procD.lm: procD.lm(fly.gpa\$coords\textasciitilde fly.ind/fly.side+fly.side/fly.ind). Note que não é necessário usar o fator repetição explicitamente, uma vez que ele está dentro de indivíduos (compare os resíduos da primeira e da segunda função).

\chapter{Semilandmarks}\label{semilandmarks}

\section{Digitalização de semilandmarks}\label{digitalizauxe7uxe3o-de-semilandmarks}

No TPSDig2, abra o arquivo `mandibula.tps' com as coordenadas dos landmarks. Para incluir semilandmarks será usada a ferramenta Draw background curves, que se localiza em `Modes\textgreater Draw curves' ou no ícone parecido com uma caneta amarela. Cliques com o botão esquerdo do mouse atribuem as coordenadas para os semilandmarks, um por vez, e os conectam por uma linha azul. Para finalizar a marcação dos semilandmarks na curva desejada basta clicar com o botão direito. O mesmo botão pode ser usado para clicar sobre a curva e acessar o menu opções, onde é possível, por exemplo, ajustar as distâncias entre os semilandmarks e deixá-las igualmente espaçadas. Ao terminar, salve o arquivo normalmente e abra em um bloco de notas. Note que as coordenadas dos semilandmarks aparecem como pontos em curvas. Neste formato, o arquivo está pronto para ser carregado pelo geomorph do R.
Para abrir o arquivo em outros programas (ex: MorphoJ) é preciso primeiro fazer o deslizamento dos semilandmarks (`sliding') ao longo das curvas. Para isso é preciso: no TPSUtil, (i) transformar os pontos em curva em semilandmarks verdadeiros (operação `Append tps Curve to landmarks'), (ii) informar quais pontos correspondem a semilandmarks em quais curvas (operação `Make sliders file'), e no TPSRelw, (iii) importar o arquivo contendo todas as coordenadas juntas (criado em (i)) e o arquivo slider (criado em (ii)), realizar a GPA com os semilandmarks e exportar o arquivo com as coordenadas alinhadas; esse último arquivo pode ser aberto em qualquer programa. Note que o método para deslizamento dos semilandmarks no TPSRelw pode ser modificado em `Options\textgreater Slide Method'.

\section{Semilandmarks no R}\label{semilandmarks-no-r}

Dados com semilandmarks podem ser carregados diretamente para o R com as coordenadas em formato de curva.

\begin{Shaded}
\begin{Highlighting}[]
\FunctionTok{require}\NormalTok{(geomorph)}
\CommentTok{\#\textgreater{} Carregando pacotes exigidos: geomorph}
\CommentTok{\#\textgreater{} Warning: package \textquotesingle{}geomorph\textquotesingle{} was built under R version 4.3.2}
\CommentTok{\#\textgreater{} Carregando pacotes exigidos: RRPP}
\CommentTok{\#\textgreater{} Warning: package \textquotesingle{}RRPP\textquotesingle{} was built under R version 4.3.2}
\CommentTok{\#\textgreater{} Carregando pacotes exigidos: rgl}
\CommentTok{\#\textgreater{} Carregando pacotes exigidos: Matrix}
\CommentTok{\# Carregar dados usando readcurves = TRUE}
\NormalTok{tps.semi}\OtherTok{\textless{}{-}}\FunctionTok{readland.tps}\NormalTok{(}\StringTok{"dadosmg/mandibula.dig\_curso.tps"}\NormalTok{,}\AttributeTok{specID =} \StringTok{"ID"}\NormalTok{, }\AttributeTok{readcurves =} \ConstantTok{TRUE}\NormalTok{)}
\NormalTok{tps.semi}
\FunctionTok{dim}\NormalTok{(tps.semi)}
\end{Highlighting}
\end{Shaded}

Note que os semilanmarks foram carregados como se fossem landmarks verdadeiros.
Precisamos informar quais deles são semilandmarks e em quais curvas vão deslizar.
Isso pode ser feito em arquivo texto (veja o formato de curves.slide.txt), pelo programa TPSUtil, ou com uma função interativa do pacote \emph{geomorph} chamada define.sliders.

\begin{Shaded}
\begin{Highlighting}[]
\CommentTok{\# Definir semilandmarks}
\NormalTok{curveslide}\OtherTok{\textless{}{-}}\FunctionTok{read.table}\NormalTok{(}\StringTok{"dadosmg/curves.slide.txt"}\NormalTok{,}\AttributeTok{h=}\NormalTok{T) }\CommentTok{\# importa sliders definidos no TPS util }
\NormalTok{curveslide}\OtherTok{\textless{}{-}}\FunctionTok{as.matrix}\NormalTok{(curveslide) }
\NormalTok{curveslide}
\CommentTok{\#\textgreater{}       before slide after}
\CommentTok{\#\textgreater{}  [1,]      2    14    15}
\CommentTok{\#\textgreater{}  [2,]     14    15    16}
\CommentTok{\#\textgreater{}  [3,]     15    16    17}
\CommentTok{\#\textgreater{}  [4,]      3    17    16}
\CommentTok{\#\textgreater{}  [5,]      3    18    19}
\CommentTok{\#\textgreater{}  [6,]     18    19    20}
\CommentTok{\#\textgreater{}  [7,]      4    20    19}
\CommentTok{\#\textgreater{}  [8,]      5    21    22}
\CommentTok{\#\textgreater{}  [9,]     21    22    23}
\CommentTok{\#\textgreater{} [10,]     22    23    24}
\CommentTok{\#\textgreater{} [11,]     23    24    25}
\CommentTok{\#\textgreater{} [12,]     24    25    26}
\CommentTok{\#\textgreater{} [13,]      6    26    25}
\CommentTok{\#\textgreater{} [14,]      6    27    28}
\CommentTok{\#\textgreater{} [15,]     27    28    29}
\CommentTok{\#\textgreater{} [16,]      7    29    28}
\CommentTok{\#\textgreater{} [17,]      7    30    31}
\CommentTok{\#\textgreater{} [18,]     30    31    32}
\CommentTok{\#\textgreater{} [19,]     31    32    33}
\CommentTok{\#\textgreater{} [20,]      8    33    32}
\CommentTok{\#\textgreater{} [21,]      8    34    35}
\CommentTok{\#\textgreater{} [22,]     34    35    36}
\CommentTok{\#\textgreater{} [23,]      9    36    35}
\CommentTok{\#\textgreater{} [24,]      9    37    38}
\CommentTok{\#\textgreater{} [25,]     37    38    39}
\CommentTok{\#\textgreater{} [26,]     38    39    40}
\CommentTok{\#\textgreater{} [27,]     10    40    39}
\CommentTok{\#\textgreater{} [28,]     10    41    42}
\CommentTok{\#\textgreater{} [29,]     41    42    43}
\CommentTok{\#\textgreater{} [30,]     42    43    44}
\CommentTok{\#\textgreater{} [31,]     11    44    43}
\CommentTok{\#\textgreater{} [32,]     11    45    46}
\CommentTok{\#\textgreater{} [33,]     45    46    47}
\CommentTok{\#\textgreater{} [34,]     47    48    49}
\CommentTok{\#\textgreater{} [35,]     46    47    48}
\CommentTok{\#\textgreater{} [36,]     48    49    50}
\CommentTok{\#\textgreater{} [37,]     49    50    51}
\CommentTok{\#\textgreater{} [38,]     50    51    52}
\CommentTok{\#\textgreater{} [39,]     51    52    53}
\CommentTok{\#\textgreater{} [40,]     52    53    54}
\CommentTok{\#\textgreater{} [41,]     12    54    53}
\CommentTok{\#\textgreater{} [42,]     12    55    56}
\CommentTok{\#\textgreater{} [43,]     55    56    57}
\CommentTok{\#\textgreater{} [44,]     56    57    58}
\CommentTok{\#\textgreater{} [45,]     57    58    59}
\CommentTok{\#\textgreater{} [46,]     58    59    60}
\CommentTok{\#\textgreater{} [47,]     13    60    59}
\end{Highlighting}
\end{Shaded}

Os semilandmarks formam a coluna do meio. Cada curva é colocada entre landmarks verdadeiros.

\section{Métodos de deslizamento de semilandmarks}\label{muxe9todos-de-deslizamento-de-semilandmarks}

Agora podemos realizar uma análise de Procrustes deslizando os semilandmarks.
A opção/argumento ``ProcD=FALSE'' vai deslizar os semilandmarks minimizando a energia de deformação, e a opção/argumento ``ProcD=TRUE'' vai deslizar os semilandmarks minizando as distâncias de Procrustes.

\begin{Shaded}
\begin{Highlighting}[]
\CommentTok{\# Deslizamento por Procrustes}
\NormalTok{gpa.semi}\OtherTok{\textless{}{-}}\FunctionTok{gpagen}\NormalTok{(tps.semi,}\AttributeTok{ProcD=}\ConstantTok{TRUE}\NormalTok{,}\AttributeTok{curves=}\NormalTok{curveslide)}
\NormalTok{shape.procD}\OtherTok{\textless{}{-}}\NormalTok{gpa.semi}\SpecialCharTok{$}\NormalTok{coords}
\FunctionTok{plotAllSpecimens}\NormalTok{(shape.procD)}
\end{Highlighting}
\end{Shaded}

\includegraphics{03.4-mgsemi_files/figure-latex/unnamed-chunk-3-1.pdf}

\begin{Shaded}
\begin{Highlighting}[]

\CommentTok{\# Deslizamento por Bending Energy}
\NormalTok{gpa.semi}\OtherTok{\textless{}{-}}\FunctionTok{gpagen}\NormalTok{(tps.semi,}\AttributeTok{ProcD=}\ConstantTok{FALSE}\NormalTok{,}\AttributeTok{curves=}\NormalTok{curveslide)}
\NormalTok{shape.bend}\OtherTok{\textless{}{-}}\NormalTok{gpa.semi}\SpecialCharTok{$}\NormalTok{coords}
\FunctionTok{plotAllSpecimens}\NormalTok{(shape.bend)}
\end{Highlighting}
\end{Shaded}

\includegraphics{03.4-mgsemi_files/figure-latex/unnamed-chunk-3-2.pdf}

Compare os dois métodos.

\begin{Shaded}
\begin{Highlighting}[]
\FunctionTok{par}\NormalTok{(}\AttributeTok{mfrow=}\FunctionTok{c}\NormalTok{(}\DecValTok{1}\NormalTok{,}\DecValTok{2}\NormalTok{))}
\FunctionTok{plotAllSpecimens}\NormalTok{(shape.procD)}
\FunctionTok{plotAllSpecimens}\NormalTok{(shape.bend)}
\end{Highlighting}
\end{Shaded}

\includegraphics{03.4-mgsemi_files/figure-latex/unnamed-chunk-4-1.pdf}

\begin{Shaded}
\begin{Highlighting}[]
\FunctionTok{par}\NormalTok{(}\AttributeTok{mfrow=}\FunctionTok{c}\NormalTok{(}\DecValTok{1}\NormalTok{,}\DecValTok{1}\NormalTok{))}
\end{Highlighting}
\end{Shaded}

\newpage

\chapter{Tamanho do Centróide}\label{tamanho-do-centruxf3ide}

Em morfometria geométrica, o tamanho do centróide é uma medida de tamanho global da estrutura (configuração de landmarks). Ela é obtida como a raiz quadrada da soma das distâncias ao quadrado de cada landmark até o centróide da configuração de landmarks (\textbf{bookstein1991?}). Por definição, essa medida é matematicamente ortogonal à forma. Por ser a única variável de tamanho não correlacionada com a forma em condições de isometria (\textbf{monteiro1999?}; \textbf{zelditch2012?}), ela é a variável de tamanho mais utilizada em morfometria geométrica.

Vamos usar as configurações de landmarks no crânio de roedores para explorar algumas análises de variação de tamanho usando o tamanho do centróide.

\begin{Shaded}
\begin{Highlighting}[]
\FunctionTok{require}\NormalTok{(geomorph)}
\CommentTok{\#\textgreater{} Carregando pacotes exigidos: geomorph}
\CommentTok{\#\textgreater{} Carregando pacotes exigidos: RRPP}
\CommentTok{\#\textgreater{} Carregando pacotes exigidos: rgl}
\CommentTok{\#\textgreater{} Carregando pacotes exigidos: Matrix}
\CommentTok{\# Carregar dados}
\NormalTok{tps}\OtherTok{\textless{}{-}}\FunctionTok{readland.tps}\NormalTok{(}\StringTok{"dadosmg/Lista tuco dig 2.tps"}\NormalTok{,}\AttributeTok{specID =} \StringTok{"ID"}\NormalTok{, }\AttributeTok{readcurves =} \ConstantTok{FALSE}\NormalTok{)}
\FunctionTok{dim}\NormalTok{(tps)}

\CommentTok{\# GPA  }
\NormalTok{gpa}\OtherTok{\textless{}{-}}\FunctionTok{gpagen}\NormalTok{(tps)}

\CommentTok{\# Extrair tamanho do centróide}
\NormalTok{size}\OtherTok{\textless{}{-}}\NormalTok{gpa}\SpecialCharTok{$}\NormalTok{Csize}
\end{Highlighting}
\end{Shaded}

Aqui está o tamanho do centróide.

\begin{Shaded}
\begin{Highlighting}[]
\NormalTok{size}

\FunctionTok{hist}\NormalTok{(size)}
\end{Highlighting}
\end{Shaded}

\includegraphics{03.5-mgcsize_files/figure-latex/unnamed-chunk-2-1.pdf}

\begin{Shaded}
\begin{Highlighting}[]
\FunctionTok{hist}\NormalTok{(}\FunctionTok{log}\NormalTok{(size))}
\end{Highlighting}
\end{Shaded}

\includegraphics{03.5-mgcsize_files/figure-latex/unnamed-chunk-2-2.pdf}

É importante manter os dados organizados e ter uma tabela em formato excel ou similar com as informações de cada espécime na mesma ordem em que eles aparecem na lista .tps. \emph{Organization is key!} Podemos carregar uma tabela assim no R para explorar a variação de tamanho entre classificadores biológicos dos nossos dados.

\begin{Shaded}
\begin{Highlighting}[]
\CommentTok{\# Carregar classificadores a partir de lista externa}
\NormalTok{plan}\OtherTok{\textless{}{-}}\FunctionTok{read.table}\NormalTok{(}\StringTok{"dadosmg/Planilha Ctenomys.txt"}\NormalTok{,}\AttributeTok{h=}\NormalTok{T)}
\NormalTok{plan}
\NormalTok{species}\OtherTok{\textless{}{-}}\NormalTok{plan[,}\DecValTok{1}\NormalTok{]}
\NormalTok{species}\OtherTok{\textless{}{-}}\FunctionTok{as.factor}\NormalTok{(species)}
\NormalTok{sex}\OtherTok{\textless{}{-}}\NormalTok{plan[,}\DecValTok{2}\NormalTok{]}
\NormalTok{sex}\OtherTok{\textless{}{-}}\FunctionTok{as.factor}\NormalTok{(sex)}
\NormalTok{sp.sex}\OtherTok{\textless{}{-}}\NormalTok{plan[,}\DecValTok{3}\NormalTok{]}
\NormalTok{sp.sex}\OtherTok{\textless{}{-}}\FunctionTok{as.factor}\NormalTok{(sp.sex)}
\end{Highlighting}
\end{Shaded}

Variação do tamanho em função de classificadores biológicos.

\begin{Shaded}
\begin{Highlighting}[]
\CommentTok{\# Variação de tamanho entre sexos}
\FunctionTok{boxplot}\NormalTok{(size}\SpecialCharTok{\textasciitilde{}}\NormalTok{sex,}\AttributeTok{ylab=}\StringTok{"CENTROID SIZE"}\NormalTok{)}
\end{Highlighting}
\end{Shaded}

\includegraphics{03.5-mgcsize_files/figure-latex/unnamed-chunk-4-1.pdf}

\begin{Shaded}
\begin{Highlighting}[]
\FunctionTok{boxplot}\NormalTok{(}\FunctionTok{log}\NormalTok{(size)}\SpecialCharTok{\textasciitilde{}}\NormalTok{sex,}\AttributeTok{ylab=}\StringTok{"LOG CENTROID SIZE"}\NormalTok{)}
\end{Highlighting}
\end{Shaded}

\includegraphics{03.5-mgcsize_files/figure-latex/unnamed-chunk-4-2.pdf}

\begin{Shaded}
\begin{Highlighting}[]
\CommentTok{\# Teste t}
\FunctionTok{t.test}\NormalTok{(size}\SpecialCharTok{\textasciitilde{}}\NormalTok{sex) }
\FunctionTok{t.test}\NormalTok{(}\FunctionTok{log}\NormalTok{(size)}\SpecialCharTok{\textasciitilde{}}\NormalTok{sex)}

\CommentTok{\# Variação de tamanho entre espécies}
\FunctionTok{boxplot}\NormalTok{(size}\SpecialCharTok{\textasciitilde{}}\NormalTok{species,}\AttributeTok{ylab=}\StringTok{"CENTROID SIZE"}\NormalTok{)}
\end{Highlighting}
\end{Shaded}

\includegraphics{03.5-mgcsize_files/figure-latex/unnamed-chunk-4-3.pdf}

\begin{Shaded}
\begin{Highlighting}[]
\FunctionTok{boxplot}\NormalTok{(}\FunctionTok{log}\NormalTok{(size)}\SpecialCharTok{\textasciitilde{}}\NormalTok{species,}\AttributeTok{ylab=}\StringTok{"LOG CENTROID SIZE"}\NormalTok{)}
\end{Highlighting}
\end{Shaded}

\includegraphics{03.5-mgcsize_files/figure-latex/unnamed-chunk-4-4.pdf}

\begin{Shaded}
\begin{Highlighting}[]
\CommentTok{\# ANOVA (Analysis of Variance)}
\NormalTok{anova}\OtherTok{\textless{}{-}}\FunctionTok{aov}\NormalTok{(size}\SpecialCharTok{\textasciitilde{}}\NormalTok{species) }
\FunctionTok{summary}\NormalTok{(anova)}
\NormalTok{anova}\OtherTok{\textless{}{-}}\FunctionTok{aov}\NormalTok{(}\FunctionTok{log}\NormalTok{(size)}\SpecialCharTok{\textasciitilde{}}\NormalTok{species) }
\FunctionTok{summary}\NormalTok{(anova)}
\CommentTok{\# Pairwise comparisons}
\FunctionTok{TukeyHSD}\NormalTok{(anova)}

\CommentTok{\# Variação de tamanho entre espécies e sexos}
\FunctionTok{boxplot}\NormalTok{(size}\SpecialCharTok{\textasciitilde{}}\NormalTok{species}\SpecialCharTok{*}\NormalTok{sex,}\AttributeTok{ylab=}\StringTok{"CENTROID SIZE"}\NormalTok{)}
\end{Highlighting}
\end{Shaded}

\includegraphics{03.5-mgcsize_files/figure-latex/unnamed-chunk-4-5.pdf}

\begin{Shaded}
\begin{Highlighting}[]
\FunctionTok{boxplot}\NormalTok{(}\FunctionTok{log}\NormalTok{(size)}\SpecialCharTok{\textasciitilde{}}\NormalTok{species}\SpecialCharTok{*}\NormalTok{sex,}\AttributeTok{ylab=}\StringTok{"LOG CENTROID SIZE"}\NormalTok{)}
\end{Highlighting}
\end{Shaded}

\includegraphics{03.5-mgcsize_files/figure-latex/unnamed-chunk-4-6.pdf}

\begin{Shaded}
\begin{Highlighting}[]
\CommentTok{\# ANOVA com interação}
\NormalTok{anova}\OtherTok{\textless{}{-}}\FunctionTok{aov}\NormalTok{(size}\SpecialCharTok{\textasciitilde{}}\NormalTok{sex}\SpecialCharTok{*}\NormalTok{species)}
\FunctionTok{summary}\NormalTok{(anova)}
\FunctionTok{TukeyHSD}\NormalTok{(anova)}
\NormalTok{anova}\OtherTok{\textless{}{-}}\FunctionTok{aov}\NormalTok{(}\FunctionTok{log}\NormalTok{(size)}\SpecialCharTok{\textasciitilde{}}\NormalTok{sex}\SpecialCharTok{*}\NormalTok{species)}
\FunctionTok{summary}\NormalTok{(anova)}
\FunctionTok{TukeyHSD}\NormalTok{(anova)}
\end{Highlighting}
\end{Shaded}

Um dos usos fundamentais do tamanho do centróide em Morfometria Geométrica é a investigação de alometria, veja o subcapítulo correspondente.

\chapter{Métodos de ordenação}\label{muxe9todos-de-ordenauxe7uxe3o}

Os métodos de ordenação mais usados em dados de morfometria geométrica são a Análise de Componentes Principais (PCA) e a Análise de Variáveis Canônicas (CVA). Ambos os métodos são exploratórios, e providenciam uma descrição dos dados, e não testes de hipóteses. No entanto, a CVA pode ser utilizada como uma função para discriminação de grupos, onde a porcentagem de classificação correta de espécimes dentro de grupos pode ser usada como valor de suporte ao testar a confiabilidade de grupos (ex: espécies).

Para se aprofundar, leia: (\textbf{klingen-monteiro2005?}); (\textbf{zelditch2012?}).

\section{Análise de Componentes Principais}\label{anuxe1lise-de-componentes-principais}

A análise de componentes principais (Principal Component Analysis - PCA) simplifica os complexos padrões de variação e covariação de dados multivariados ao construir componentes principais independentes que resumem a maior parte da variação dos dados em poucos eixos. Assim, também é possível observar diretamente a combinação de dois ou três componentes principais em um gráfico.

Vamos fazer uma PCA com dados de morfometria geométrica. Primeiro, vamos carregar os dados, fazer uma GPA, e carregar os classificadores.

\begin{Shaded}
\begin{Highlighting}[]
\FunctionTok{require}\NormalTok{(geomorph)}
\CommentTok{\#\textgreater{} Carregando pacotes exigidos: geomorph}
\CommentTok{\#\textgreater{} Carregando pacotes exigidos: RRPP}
\CommentTok{\#\textgreater{} Carregando pacotes exigidos: rgl}
\CommentTok{\#\textgreater{} Carregando pacotes exigidos: Matrix}
\CommentTok{\# Carregar dados}
\NormalTok{tps}\OtherTok{\textless{}{-}}\FunctionTok{readland.tps}\NormalTok{(}\StringTok{"dadosmg/Lista tuco dig 2.tps"}\NormalTok{,}\AttributeTok{specID =} \StringTok{"ID"}\NormalTok{, }\AttributeTok{readcurves =} \ConstantTok{FALSE}\NormalTok{)}
\FunctionTok{dim}\NormalTok{(tps)}

\CommentTok{\# Gerar matriz com pares de landmarks simétricos}
\NormalTok{pairs.matrix}\OtherTok{\textless{}{-}}\FunctionTok{matrix}\NormalTok{(}\FunctionTok{c}\NormalTok{(}\DecValTok{2}\NormalTok{,}\DecValTok{3}\NormalTok{,}\DecValTok{5}\NormalTok{,}\DecValTok{6}\NormalTok{,}\DecValTok{7}\NormalTok{,}\DecValTok{8}\NormalTok{,}\DecValTok{10}\NormalTok{,}\DecValTok{11}\NormalTok{,}\DecValTok{12}\NormalTok{,}\DecValTok{13}\NormalTok{,}\DecValTok{14}\NormalTok{,}\DecValTok{15}\NormalTok{,}\DecValTok{16}\NormalTok{,}\DecValTok{17}\NormalTok{,}\DecValTok{18}\NormalTok{,}\DecValTok{19}\NormalTok{,}\DecValTok{20}\NormalTok{,}\DecValTok{21}\NormalTok{,}\DecValTok{23}\NormalTok{,}\DecValTok{24}\NormalTok{,}\DecValTok{25}\NormalTok{,}\DecValTok{26}\NormalTok{,}\DecValTok{27}\NormalTok{,}\DecValTok{28}\NormalTok{),}\AttributeTok{nrow=}\DecValTok{12}\NormalTok{,}\AttributeTok{ncol=}\DecValTok{2}\NormalTok{,}\AttributeTok{byrow=}\NormalTok{T)}
\CommentTok{\# Vetor indivíduos}
\NormalTok{ind}\OtherTok{\textless{}{-}}\FunctionTok{c}\NormalTok{(}\DecValTok{1}\SpecialCharTok{:}\FunctionTok{dim}\NormalTok{(tps)[}\DecValTok{3}\NormalTok{])}

\CommentTok{\# GPA {-} Simetria bilateral }
\NormalTok{b.s}\OtherTok{\textless{}{-}}\FunctionTok{bilat.symmetry}\NormalTok{(tps,}\AttributeTok{ind=}\NormalTok{ind,}\AttributeTok{object.sym=}\ConstantTok{TRUE}\NormalTok{,}\AttributeTok{land.pairs=}\NormalTok{pairs.matrix)}

\CommentTok{\# Matriz de forma}
\NormalTok{shape.sym}\OtherTok{\textless{}{-}}\NormalTok{b.s}\SpecialCharTok{$}\NormalTok{symm.shape}

\CommentTok{\# Carregar classificadores a partir de lista externa}
\NormalTok{plan}\OtherTok{\textless{}{-}}\FunctionTok{read.table}\NormalTok{(}\StringTok{"dadosmg/Planilha Ctenomys.txt"}\NormalTok{,}\AttributeTok{h=}\NormalTok{T)}
\NormalTok{plan}
\NormalTok{species}\OtherTok{\textless{}{-}}\NormalTok{plan[,}\DecValTok{1}\NormalTok{]}
\NormalTok{species}\OtherTok{\textless{}{-}}\FunctionTok{as.factor}\NormalTok{(species)}
\NormalTok{sex}\OtherTok{\textless{}{-}}\NormalTok{plan[,}\DecValTok{2}\NormalTok{]}
\NormalTok{sex}\OtherTok{\textless{}{-}}\FunctionTok{as.factor}\NormalTok{(sex)}
\NormalTok{sp.sex}\OtherTok{\textless{}{-}}\NormalTok{plan[,}\DecValTok{3}\NormalTok{]}
\NormalTok{sp.sex}\OtherTok{\textless{}{-}}\FunctionTok{as.factor}\NormalTok{(sp.sex)}
\end{Highlighting}
\end{Shaded}

Muitas funções realizam a Análise de Componentes Principais no R. Para dados de morfometria geométrica (MG), é melhor usar funções de pacotes desenhados especificamente para MG, já que eles reconhecem a redundância nos graus de liberdade que resultam das transformações de escala, posição e rotação feitas anteriormente na Análise Generalisada de Procrustes, e já retornam o número exato de componentes principais necessários para capturar 100\% da variação da forma. Vamos usar a PCA implementada no pacote \emph{geomorph}.

\begin{Shaded}
\begin{Highlighting}[]
\NormalTok{PCA}\OtherTok{\textless{}{-}}\FunctionTok{gm.prcomp}\NormalTok{(shape.sym)}
\NormalTok{PCA }\CommentTok{\# Porcentagem de explicação dos eixos}
\CommentTok{\#\textgreater{} }
\CommentTok{\#\textgreater{} Ordination type: Principal Component Analysis }
\CommentTok{\#\textgreater{} Centering by OLS mean}
\CommentTok{\#\textgreater{} Orthogonal projection of OLS residuals}
\CommentTok{\#\textgreater{} Number of observations: 30 }
\CommentTok{\#\textgreater{} Number of vectors 28 }
\CommentTok{\#\textgreater{} }
\CommentTok{\#\textgreater{} Importance of Components:}
\CommentTok{\#\textgreater{}                               Comp1        Comp2}
\CommentTok{\#\textgreater{} Eigenvalues            0.0005691898 0.0002349919}
\CommentTok{\#\textgreater{} Proportion of Variance 0.4018322761 0.1658977524}
\CommentTok{\#\textgreater{} Cumulative Proportion  0.4018322761 0.5677300285}
\CommentTok{\#\textgreater{}                               Comp3        Comp4}
\CommentTok{\#\textgreater{} Eigenvalues            0.0001552391 9.049344e{-}05}
\CommentTok{\#\textgreater{} Proportion of Variance 0.1095945175 6.388587e{-}02}
\CommentTok{\#\textgreater{} Cumulative Proportion  0.6773245461 7.412104e{-}01}
\CommentTok{\#\textgreater{}                               Comp5        Comp6}
\CommentTok{\#\textgreater{} Eigenvalues            6.732315e{-}05 5.547785e{-}05}
\CommentTok{\#\textgreater{} Proportion of Variance 4.752828e{-}02 3.916583e{-}02}
\CommentTok{\#\textgreater{} Cumulative Proportion  7.887387e{-}01 8.279045e{-}01}
\CommentTok{\#\textgreater{}                               Comp7        Comp8}
\CommentTok{\#\textgreater{} Eigenvalues            4.232582e{-}05 3.705712e{-}05}
\CommentTok{\#\textgreater{} Proportion of Variance 2.988086e{-}02 2.616131e{-}02}
\CommentTok{\#\textgreater{} Cumulative Proportion  8.577854e{-}01 8.839467e{-}01}
\CommentTok{\#\textgreater{}                               Comp9       Comp10}
\CommentTok{\#\textgreater{} Eigenvalues            3.356797e{-}05 2.770576e{-}05}
\CommentTok{\#\textgreater{} Proportion of Variance 2.369806e{-}02 1.955950e{-}02}
\CommentTok{\#\textgreater{} Cumulative Proportion  9.076448e{-}01 9.272043e{-}01}
\CommentTok{\#\textgreater{}                              Comp11       Comp12}
\CommentTok{\#\textgreater{} Eigenvalues            2.221806e{-}05 1.731913e{-}05}
\CommentTok{\#\textgreater{} Proportion of Variance 1.568534e{-}02 1.222683e{-}02}
\CommentTok{\#\textgreater{} Cumulative Proportion  9.428896e{-}01 9.551164e{-}01}
\CommentTok{\#\textgreater{}                              Comp13       Comp14}
\CommentTok{\#\textgreater{} Eigenvalues            1.405341e{-}05 1.156877e{-}05}
\CommentTok{\#\textgreater{} Proportion of Variance 9.921316e{-}03 8.167232e{-}03}
\CommentTok{\#\textgreater{} Cumulative Proportion  9.650377e{-}01 9.732050e{-}01}
\CommentTok{\#\textgreater{}                              Comp15       Comp16}
\CommentTok{\#\textgreater{} Eigenvalues            1.028007e{-}05 8.308937e{-}06}
\CommentTok{\#\textgreater{} Proportion of Variance 7.257446e{-}03 5.865880e{-}03}
\CommentTok{\#\textgreater{} Cumulative Proportion  9.804624e{-}01 9.863283e{-}01}
\CommentTok{\#\textgreater{}                              Comp17       Comp18}
\CommentTok{\#\textgreater{} Eigenvalues            7.020041e{-}06 3.263883e{-}06}
\CommentTok{\#\textgreater{} Proportion of Variance 4.955955e{-}03 2.304211e{-}03}
\CommentTok{\#\textgreater{} Cumulative Proportion  9.912842e{-}01 9.935885e{-}01}
\CommentTok{\#\textgreater{}                              Comp19       Comp20}
\CommentTok{\#\textgreater{} Eigenvalues            2.679842e{-}06 2.081122e{-}06}
\CommentTok{\#\textgreater{} Proportion of Variance 1.891894e{-}03 1.469214e{-}03}
\CommentTok{\#\textgreater{} Cumulative Proportion  9.954804e{-}01 9.969496e{-}01}
\CommentTok{\#\textgreater{}                              Comp21       Comp22}
\CommentTok{\#\textgreater{} Eigenvalues            1.577657e{-}06 1.186452e{-}06}
\CommentTok{\#\textgreater{} Proportion of Variance 1.113783e{-}03 8.376027e{-}04}
\CommentTok{\#\textgreater{} Cumulative Proportion  9.980634e{-}01 9.989010e{-}01}
\CommentTok{\#\textgreater{}                              Comp23       Comp24}
\CommentTok{\#\textgreater{} Eigenvalues            9.098790e{-}07 4.159837e{-}07}
\CommentTok{\#\textgreater{} Proportion of Variance 6.423494e{-}04 2.936730e{-}04}
\CommentTok{\#\textgreater{} Cumulative Proportion  9.995433e{-}01 9.998370e{-}01}
\CommentTok{\#\textgreater{}                              Comp25       Comp26}
\CommentTok{\#\textgreater{} Eigenvalues            1.598994e{-}07 6.588247e{-}08}
\CommentTok{\#\textgreater{} Proportion of Variance 1.128845e{-}04 4.651120e{-}05}
\CommentTok{\#\textgreater{} Cumulative Proportion  9.999499e{-}01 9.999964e{-}01}
\CommentTok{\#\textgreater{}                              Comp27       Comp28}
\CommentTok{\#\textgreater{} Eigenvalues            5.138116e{-}09 1.613869e{-}19}
\CommentTok{\#\textgreater{} Proportion of Variance 3.627368e{-}06 1.139347e{-}16}
\CommentTok{\#\textgreater{} Cumulative Proportion  1.000000e+00 1.000000e+00}

\FunctionTok{plot}\NormalTok{(PCA,}\AttributeTok{axis1=}\DecValTok{1}\NormalTok{, }\AttributeTok{axis2=}\DecValTok{2}\NormalTok{)}
\end{Highlighting}
\end{Shaded}

\includegraphics{03.6-mgord_files/figure-latex/unnamed-chunk-2-1.pdf}

Antes de observar as mudanças de forma ao longo do gráfico da PCA, vamos criar uma forma média de referência e definir links entre os landmarks.

\begin{Shaded}
\begin{Highlighting}[]
\CommentTok{\# Forma média}
\NormalTok{ref}\OtherTok{\textless{}{-}}\FunctionTok{mshape}\NormalTok{(shape.sym)}

\NormalTok{links.tuco}\OtherTok{\textless{}{-}}\FunctionTok{matrix}\NormalTok{(}\FunctionTok{c}\NormalTok{(}\DecValTok{1}\NormalTok{,}\DecValTok{2}\NormalTok{,}\DecValTok{1}\NormalTok{,}\DecValTok{3}\NormalTok{,}\DecValTok{2}\NormalTok{,}\DecValTok{7}\NormalTok{,}\DecValTok{3}\NormalTok{,}\DecValTok{8}\NormalTok{,}\DecValTok{7}\NormalTok{,}\DecValTok{10}\NormalTok{,}\DecValTok{8}\NormalTok{,}\DecValTok{11}\NormalTok{,}\DecValTok{10}\NormalTok{,}\DecValTok{18}\NormalTok{,}\DecValTok{11}\NormalTok{,}\DecValTok{19}\NormalTok{,}\DecValTok{14}\NormalTok{,}\DecValTok{18}\NormalTok{,}\DecValTok{15}\NormalTok{,}\DecValTok{19}\NormalTok{,}\DecValTok{10}\NormalTok{,}\DecValTok{14}\NormalTok{,}\DecValTok{11}\NormalTok{,}\DecValTok{15}\NormalTok{,}\DecValTok{18}\NormalTok{,}\DecValTok{20}\NormalTok{,}\DecValTok{19}\NormalTok{,}\DecValTok{21}\NormalTok{,}\DecValTok{20}\NormalTok{,}\DecValTok{23}\NormalTok{,}\DecValTok{21}\NormalTok{,}\DecValTok{24}\NormalTok{,}\DecValTok{23}\NormalTok{,}\DecValTok{25}\NormalTok{,}\DecValTok{24}\NormalTok{,}\DecValTok{26}\NormalTok{,}\DecValTok{25}\NormalTok{,}\DecValTok{27}\NormalTok{,}\DecValTok{26}\NormalTok{,}\DecValTok{28}\NormalTok{,}\DecValTok{27}\NormalTok{,}\DecValTok{29}\NormalTok{,}\DecValTok{28}\NormalTok{,}\DecValTok{29}\NormalTok{,}\DecValTok{4}\NormalTok{,}\DecValTok{5}\NormalTok{,}\DecValTok{4}\NormalTok{,}\DecValTok{6}\NormalTok{,}\DecValTok{6}\NormalTok{,}\DecValTok{9}\NormalTok{,}\DecValTok{5}\NormalTok{,}\DecValTok{9}\NormalTok{,}\DecValTok{1}\NormalTok{,}\DecValTok{4}\NormalTok{,}\DecValTok{4}\NormalTok{,}\DecValTok{9}\NormalTok{,}\DecValTok{9}\NormalTok{,}\DecValTok{22}\NormalTok{,}\DecValTok{22}\NormalTok{,}\DecValTok{29}\NormalTok{,}\DecValTok{7}\NormalTok{,}\DecValTok{9}\NormalTok{,}\DecValTok{8}\NormalTok{,}\DecValTok{9}\NormalTok{,}\DecValTok{10}\NormalTok{,}\DecValTok{12}\NormalTok{,}\DecValTok{11}\NormalTok{,}\DecValTok{13}\NormalTok{,}\DecValTok{9}\NormalTok{,}\DecValTok{12}\NormalTok{,}\DecValTok{9}\NormalTok{,}\DecValTok{13}\NormalTok{,}\DecValTok{12}\NormalTok{,}\DecValTok{16}\NormalTok{,}\DecValTok{13}\NormalTok{,}\DecValTok{17}\NormalTok{,}\DecValTok{16}\NormalTok{,}\DecValTok{22}\NormalTok{,}\DecValTok{17}\NormalTok{,}\DecValTok{22}\NormalTok{,}\DecValTok{16}\NormalTok{,}\DecValTok{18}\NormalTok{,}\DecValTok{17}\NormalTok{,}\DecValTok{19}\NormalTok{,}\DecValTok{16}\NormalTok{,}\DecValTok{20}\NormalTok{,}\DecValTok{17}\NormalTok{,}\DecValTok{21}\NormalTok{,}\DecValTok{22}\NormalTok{,}\DecValTok{23}\NormalTok{,}\DecValTok{22}\NormalTok{,}\DecValTok{24}\NormalTok{,}\DecValTok{23}\NormalTok{,}\DecValTok{27}\NormalTok{,}\DecValTok{24}\NormalTok{,}\DecValTok{28}\NormalTok{,}\DecValTok{23}\NormalTok{,}\DecValTok{29}\NormalTok{,}\DecValTok{24}\NormalTok{,}\DecValTok{29}\NormalTok{),}\AttributeTok{nrow=}\DecValTok{50}\NormalTok{,}\AttributeTok{ncol=}\DecValTok{2}\NormalTok{,}\AttributeTok{byrow=}\NormalTok{T)}

\CommentTok{\# Plotar espécimes contra a média}
\NormalTok{GP1}\OtherTok{\textless{}{-}}\FunctionTok{gridPar}\NormalTok{(}\AttributeTok{pt.bg=}\StringTok{"gray"}\NormalTok{,}\AttributeTok{link.col=}\StringTok{"gray"}\NormalTok{,}\AttributeTok{link.lty=}\DecValTok{1}\NormalTok{) }\CommentTok{\# cor dos landmarks e links}
\FunctionTok{plotRefToTarget}\NormalTok{(ref,shape.sym[,,}\DecValTok{8}\NormalTok{],}\AttributeTok{links=}\NormalTok{links.tuco,}\AttributeTok{method=}\StringTok{"TPS"}\NormalTok{) }
\end{Highlighting}
\end{Shaded}

\includegraphics{03.6-mgord_files/figure-latex/unnamed-chunk-3-1.pdf}

\begin{Shaded}
\begin{Highlighting}[]
\FunctionTok{plotRefToTarget}\NormalTok{(ref,shape.sym[,,}\DecValTok{8}\NormalTok{],}\AttributeTok{links=}\NormalTok{links.tuco,}\AttributeTok{method=}\StringTok{"vector"}\NormalTok{) }
\end{Highlighting}
\end{Shaded}

\includegraphics{03.6-mgord_files/figure-latex/unnamed-chunk-3-2.pdf}

\begin{Shaded}
\begin{Highlighting}[]
\FunctionTok{plotRefToTarget}\NormalTok{(ref,shape.sym[,,}\DecValTok{8}\NormalTok{],}\AttributeTok{links=}\NormalTok{links.tuco,}\AttributeTok{method=}\StringTok{"points"}\NormalTok{,}\AttributeTok{gridPars=}\NormalTok{GP1) }\CommentTok{\# alvo = em preto, referência = em cinza}
\end{Highlighting}
\end{Shaded}

\includegraphics{03.6-mgord_files/figure-latex/unnamed-chunk-3-3.pdf}

\begin{Shaded}
\begin{Highlighting}[]

\FunctionTok{plotAllSpecimens}\NormalTok{(shape.sym,}\AttributeTok{mean=}\ConstantTok{TRUE}\NormalTok{,}\AttributeTok{links=}\NormalTok{links.tuco)}
\end{Highlighting}
\end{Shaded}

\includegraphics{03.6-mgord_files/figure-latex/unnamed-chunk-3-4.pdf}

Podemos estilizar a PCA para mostrar, por exemplo, quais indivíduos pertecem a quais grupos. Além disso, é importante observar as formas associadas com porções diferentes do gráfico.

\begin{Shaded}
\begin{Highlighting}[]
\CommentTok{\# PCA Estilizada}
\NormalTok{col.group}\OtherTok{\textless{}{-}}\FunctionTok{rainbow}\NormalTok{(}\FunctionTok{length}\NormalTok{(}\FunctionTok{levels}\NormalTok{(species))) }\CommentTok{\# criar vetor de cores para grupos}
\FunctionTok{names}\NormalTok{(col.group)}\OtherTok{\textless{}{-}}\FunctionTok{levels}\NormalTok{(species)}
\NormalTok{col.group}\OtherTok{\textless{}{-}}\NormalTok{col.group[}\FunctionTok{match}\NormalTok{(species,}\FunctionTok{names}\NormalTok{(col.group))]}
\NormalTok{PCA}\OtherTok{\textless{}{-}}\FunctionTok{gm.prcomp}\NormalTok{(shape.sym) }\CommentTok{\# rodar a PCA}
\NormalTok{xlab}\OtherTok{\textless{}{-}}\FunctionTok{paste}\NormalTok{(}\StringTok{"Principal Component 1 "}\NormalTok{, }\StringTok{"("}\NormalTok{,}\FunctionTok{round}\NormalTok{(PCA}\SpecialCharTok{$}\NormalTok{pc.summary}\SpecialCharTok{$}\NormalTok{importance[}\DecValTok{2}\NormalTok{,}\DecValTok{1}\NormalTok{]}\SpecialCharTok{*}\DecValTok{100}\NormalTok{,}\DecValTok{1}\NormalTok{),}\StringTok{"\%)"}\NormalTok{,}\AttributeTok{sep=}\StringTok{""}\NormalTok{) }\CommentTok{\# Criar label pro eixo 1}
\NormalTok{ylab}\OtherTok{\textless{}{-}}\FunctionTok{paste}\NormalTok{(}\StringTok{"Principal Component 2 "}\NormalTok{, }\StringTok{"("}\NormalTok{,}\FunctionTok{round}\NormalTok{(PCA}\SpecialCharTok{$}\NormalTok{pc.summary}\SpecialCharTok{$}\NormalTok{importance[}\DecValTok{2}\NormalTok{,}\DecValTok{2}\NormalTok{]}\SpecialCharTok{*}\DecValTok{100}\NormalTok{,}\DecValTok{1}\NormalTok{),}\StringTok{"\%)"}\NormalTok{,}\AttributeTok{sep=}\StringTok{""}\NormalTok{) }\CommentTok{\# Criar label pro eixo 2}
\NormalTok{\{mat}\OtherTok{\textless{}{-}}\FunctionTok{matrix}\NormalTok{(}\FunctionTok{c}\NormalTok{(}\DecValTok{4}\NormalTok{,}\DecValTok{5}\NormalTok{,}\DecValTok{0}\NormalTok{,}\DecValTok{1}\NormalTok{,}\DecValTok{1}\NormalTok{,}\DecValTok{2}\NormalTok{,}\DecValTok{1}\NormalTok{,}\DecValTok{1}\NormalTok{,}\DecValTok{3}\NormalTok{),}\DecValTok{3}\NormalTok{) }\CommentTok{\# Dividir a janela do plot}
\FunctionTok{layout}\NormalTok{(mat, }\AttributeTok{widths=}\FunctionTok{c}\NormalTok{(}\DecValTok{1}\NormalTok{,}\DecValTok{1}\NormalTok{,}\DecValTok{1}\NormalTok{), }\AttributeTok{heights=}\FunctionTok{c}\NormalTok{(}\DecValTok{1}\NormalTok{,}\DecValTok{1}\NormalTok{,}\FloatTok{0.6}\NormalTok{))}
\FunctionTok{par}\NormalTok{(}\AttributeTok{mar=}\FunctionTok{c}\NormalTok{(}\DecValTok{4}\NormalTok{, }\DecValTok{4}\NormalTok{, }\DecValTok{1}\NormalTok{, }\DecValTok{1}\NormalTok{))}
\FunctionTok{plot}\NormalTok{(PCA}\SpecialCharTok{$}\NormalTok{x[,}\DecValTok{1}\NormalTok{],PCA}\SpecialCharTok{$}\NormalTok{x[,}\DecValTok{2}\NormalTok{],}\AttributeTok{pch=}\DecValTok{21}\NormalTok{,}\AttributeTok{cex=}\DecValTok{2}\NormalTok{,}\AttributeTok{bg=}\NormalTok{col.group,}\AttributeTok{xlab=}\NormalTok{xlab,}\AttributeTok{ylab=}\NormalTok{ylab,}\AttributeTok{asp=}\NormalTok{T)}
\FunctionTok{legend}\NormalTok{(}\FloatTok{0.025}\NormalTok{,}\SpecialCharTok{{-}}\FloatTok{0.015}\NormalTok{,}\AttributeTok{legend=}\FunctionTok{unique}\NormalTok{(species),}\AttributeTok{pch=}\DecValTok{19}\NormalTok{,}\AttributeTok{col=}\FunctionTok{unique}\NormalTok{(col.group))}
\FunctionTok{plotRefToTarget}\NormalTok{(ref,PCA}\SpecialCharTok{$}\NormalTok{shapes}\SpecialCharTok{$}\NormalTok{shapes.comp1}\SpecialCharTok{$}\NormalTok{min,}\AttributeTok{links=}\NormalTok{links.tuco,}\AttributeTok{method=}\StringTok{"TPS"}\NormalTok{,}\AttributeTok{gridPars=}\NormalTok{GP1)}
\FunctionTok{plotRefToTarget}\NormalTok{(ref,PCA}\SpecialCharTok{$}\NormalTok{shapes}\SpecialCharTok{$}\NormalTok{shapes.comp1}\SpecialCharTok{$}\NormalTok{max,}\AttributeTok{links=}\NormalTok{links.tuco,}\AttributeTok{method=}\StringTok{"TPS"}\NormalTok{,}\AttributeTok{gridPars=}\NormalTok{GP1)}
\FunctionTok{plotRefToTarget}\NormalTok{(ref,PCA}\SpecialCharTok{$}\NormalTok{shapes}\SpecialCharTok{$}\NormalTok{shapes.comp2}\SpecialCharTok{$}\NormalTok{max,}\AttributeTok{links=}\NormalTok{links.tuco,}\AttributeTok{method=}\StringTok{"TPS"}\NormalTok{,}\AttributeTok{gridPars=}\NormalTok{GP1)}
\FunctionTok{plotRefToTarget}\NormalTok{(ref,PCA}\SpecialCharTok{$}\NormalTok{shapes}\SpecialCharTok{$}\NormalTok{shapes.comp2}\SpecialCharTok{$}\NormalTok{min,}\AttributeTok{links=}\NormalTok{links.tuco,}\AttributeTok{method=}\StringTok{"TPS"}\NormalTok{,}\AttributeTok{gridPars=}\NormalTok{GP1)\}}
\end{Highlighting}
\end{Shaded}

\includegraphics{03.6-mgord_files/figure-latex/unnamed-chunk-4-1.pdf}

\begin{Shaded}
\begin{Highlighting}[]
\FunctionTok{par}\NormalTok{(}\AttributeTok{mfrow=}\FunctionTok{c}\NormalTok{(}\DecValTok{1}\NormalTok{,}\DecValTok{1}\NormalTok{)) }\CommentTok{\# janela default do gr?fico  }
\end{Highlighting}
\end{Shaded}

É possível gerar elipses de confiança para os grupos em questão.

\begin{Shaded}
\begin{Highlighting}[]
\FunctionTok{require}\NormalTok{(vegan)}
\CommentTok{\#\textgreater{} Carregando pacotes exigidos: vegan}
\CommentTok{\#\textgreater{} Carregando pacotes exigidos: permute}
\CommentTok{\#\textgreater{} Carregando pacotes exigidos: lattice}
\CommentTok{\#\textgreater{} This is vegan 2.6{-}4}
\NormalTok{\{}\FunctionTok{plot}\NormalTok{(PCA}\SpecialCharTok{$}\NormalTok{x[,}\DecValTok{1}\NormalTok{],PCA}\SpecialCharTok{$}\NormalTok{x[,}\DecValTok{2}\NormalTok{],}\AttributeTok{pch=}\DecValTok{21}\NormalTok{,}\AttributeTok{cex=}\DecValTok{2}\NormalTok{,}\AttributeTok{bg=}\NormalTok{col.group,}\AttributeTok{xlab=}\NormalTok{xlab,}\AttributeTok{ylab=}\NormalTok{ylab,}\AttributeTok{asp=}\NormalTok{T)}
\FunctionTok{ordiellipse}\NormalTok{(PCA}\SpecialCharTok{$}\NormalTok{x,}\AttributeTok{group=}\NormalTok{species, }\AttributeTok{kind=}\StringTok{"sd"}\NormalTok{, }\AttributeTok{conf=}\FloatTok{0.95}\NormalTok{)\}}
\end{Highlighting}
\end{Shaded}

\includegraphics{03.6-mgord_files/figure-latex/unnamed-chunk-5-1.pdf}

\begin{Shaded}
\begin{Highlighting}[]
\NormalTok{\{}\FunctionTok{plot}\NormalTok{(PCA}\SpecialCharTok{$}\NormalTok{x[,}\DecValTok{1}\NormalTok{],PCA}\SpecialCharTok{$}\NormalTok{x[,}\DecValTok{2}\NormalTok{],}\AttributeTok{pch=}\DecValTok{21}\NormalTok{,}\AttributeTok{cex=}\DecValTok{2}\NormalTok{,}\AttributeTok{bg=}\NormalTok{col.group,}\AttributeTok{xlab=}\NormalTok{xlab,}\AttributeTok{ylab=}\NormalTok{ylab,}\AttributeTok{asp=}\NormalTok{T)}
\FunctionTok{ordihull}\NormalTok{(PCA}\SpecialCharTok{$}\NormalTok{x,}\AttributeTok{group=}\NormalTok{species)\}}
\end{Highlighting}
\end{Shaded}

\includegraphics{03.6-mgord_files/figure-latex/unnamed-chunk-5-2.pdf}

As mudanças de forma ao longo dos eixos da PCA podem ser visualizadas comparando-as com a média da amostra. Usualmente, as formas extremas de cada eixo são usadas para visualização. O \emph{geomorph} retorna as formas extremas de maneira simplificada no output da função que faz a PCA.

\begin{Shaded}
\begin{Highlighting}[]
\CommentTok{\# Definir cores dos landmarks e links}
\NormalTok{GP1}\OtherTok{\textless{}{-}}\FunctionTok{gridPar}\NormalTok{(}\AttributeTok{pt.bg=}\StringTok{"gray"}\NormalTok{,}\AttributeTok{link.col=}\StringTok{"gray"}\NormalTok{,}\AttributeTok{link.lty=}\DecValTok{1}\NormalTok{) }\CommentTok{\# cor dos landmarks e links}
\CommentTok{\# Mínimos e máximos contra a média}
\FunctionTok{plotRefToTarget}\NormalTok{(ref,PCA}\SpecialCharTok{$}\NormalTok{shapes}\SpecialCharTok{$}\NormalTok{shapes.comp1}\SpecialCharTok{$}\NormalTok{min,}\AttributeTok{links=}\NormalTok{links.tuco,}\AttributeTok{method=}\StringTok{"TPS"}\NormalTok{,}\AttributeTok{gridPars=}\NormalTok{GP1)}
\end{Highlighting}
\end{Shaded}

\includegraphics{03.6-mgord_files/figure-latex/unnamed-chunk-6-1.pdf}

\begin{Shaded}
\begin{Highlighting}[]
\FunctionTok{plotRefToTarget}\NormalTok{(ref,PCA}\SpecialCharTok{$}\NormalTok{shapes}\SpecialCharTok{$}\NormalTok{shapes.comp1}\SpecialCharTok{$}\NormalTok{max,}\AttributeTok{links=}\NormalTok{links.tuco,}\AttributeTok{method=}\StringTok{"TPS"}\NormalTok{,}\AttributeTok{gridPars=}\NormalTok{GP1)}
\end{Highlighting}
\end{Shaded}

\includegraphics{03.6-mgord_files/figure-latex/unnamed-chunk-6-2.pdf}

\begin{Shaded}
\begin{Highlighting}[]
\CommentTok{\# Mínimo contra máximo}
\FunctionTok{plotRefToTarget}\NormalTok{(PCA}\SpecialCharTok{$}\NormalTok{shapes}\SpecialCharTok{$}\NormalTok{shapes.comp1}\SpecialCharTok{$}\NormalTok{min,PCA}\SpecialCharTok{$}\NormalTok{shapes}\SpecialCharTok{$}\NormalTok{shapes.comp1}\SpecialCharTok{$}\NormalTok{max,}\AttributeTok{links=}\NormalTok{links.tuco,}\AttributeTok{method=}\StringTok{"points"}\NormalTok{,}\AttributeTok{gridPars=}\NormalTok{GP1) }\CommentTok{\# alvo = em preto, referência = em cinza}
\end{Highlighting}
\end{Shaded}

\includegraphics{03.6-mgord_files/figure-latex/unnamed-chunk-6-3.pdf}

\section{Análise de Variação Canônica}\label{anuxe1lise-de-variauxe7uxe3o-canuxf4nica}

Assim como na PCA, a Análise de Variação Canônica (Canonical Variates Analysis - CVA) também cria novas variáveis (eixos de ordenação) que são ortogonais. No entanto, a CVA descreve as diferenças \emph{entre} as médias de \emph{grupos}. Assim, ela pode ser usada para descrever as diferenças de forma entre quaisquer \emph{grupos} biológicos que tenhamos. Contudo, ao contrário da PCA, que simplesmente rotaciona as coordenadas originais, a CVA realiza uma transformação (rescalonamento pela variação dentro de grupos) no sistema de coordenadas original, e por isso não deve ser interpretada livremente para representar variação de forma entre indivíduos. A CVA mostra apenas a direção em que os grupos são melhor discriminados um do outro.

\begin{Shaded}
\begin{Highlighting}[]
\FunctionTok{require}\NormalTok{(Morpho)}
\CommentTok{\#\textgreater{} Carregando pacotes exigidos: Morpho}
\CommentTok{\#\textgreater{} Warning: package \textquotesingle{}Morpho\textquotesingle{} was built under R version 4.3.2}
\CommentTok{\#\textgreater{} }
\CommentTok{\#\textgreater{} Attaching package: \textquotesingle{}Morpho\textquotesingle{}}
\CommentTok{\#\textgreater{} The following object is masked from \textquotesingle{}package:RRPP\textquotesingle{}:}
\CommentTok{\#\textgreater{} }
\CommentTok{\#\textgreater{}     classify}
\NormalTok{cva}\OtherTok{\textless{}{-}}\FunctionTok{CVA}\NormalTok{(PCA}\SpecialCharTok{$}\NormalTok{x[,}\DecValTok{1}\SpecialCharTok{:}\DecValTok{20}\NormalTok{],}\AttributeTok{groups=}\NormalTok{species,}\AttributeTok{cv=}\ConstantTok{TRUE}\NormalTok{) }\CommentTok{\# Jackknife Cross{-}validation}
\NormalTok{cva}\SpecialCharTok{$}\NormalTok{Var}
\CommentTok{\#\textgreater{}      Canonical roots \% Variance Cumulative \%}
\CommentTok{\#\textgreater{} [1,]        88.21224   82.81358     82.81358}
\CommentTok{\#\textgreater{} [2,]        18.30681   17.18642    100.00000}

\NormalTok{\{}\FunctionTok{plot}\NormalTok{(cva}\SpecialCharTok{$}\NormalTok{CVscores,}\AttributeTok{bg=}\NormalTok{col.group,}\AttributeTok{pch=}\DecValTok{21}\NormalTok{,}\AttributeTok{cex=}\DecValTok{2}\NormalTok{,}\AttributeTok{xlab=}\FunctionTok{paste}\NormalTok{(}\StringTok{"CV1  ("}\NormalTok{,}\FunctionTok{paste}\NormalTok{(}\FunctionTok{round}\NormalTok{(cva}\SpecialCharTok{$}\NormalTok{Var[}\DecValTok{1}\NormalTok{,}\DecValTok{2}\NormalTok{],}\DecValTok{1}\NormalTok{),}\StringTok{"\%)"}\NormalTok{),}\AttributeTok{sep=}\StringTok{""}\NormalTok{),}\AttributeTok{ylab=}\FunctionTok{paste}\NormalTok{(}\StringTok{"CV2  ("}\NormalTok{,}\FunctionTok{paste}\NormalTok{(}\FunctionTok{round}\NormalTok{(cva}\SpecialCharTok{$}\NormalTok{Var[}\DecValTok{2}\NormalTok{,}\DecValTok{2}\NormalTok{],}\DecValTok{1}\NormalTok{),}\StringTok{"\%)"}\NormalTok{),}\AttributeTok{sep=}\StringTok{""}\NormalTok{))}
\FunctionTok{legend}\NormalTok{(}\DecValTok{0}\NormalTok{,}\SpecialCharTok{{-}}\DecValTok{2}\NormalTok{,}\AttributeTok{legend=}\FunctionTok{unique}\NormalTok{(species),}\AttributeTok{pch=}\DecValTok{19}\NormalTok{,}\AttributeTok{col=}\FunctionTok{unique}\NormalTok{(col.group))}

\CommentTok{\# gerar elipses de confiança}
\FunctionTok{ordiellipse}\NormalTok{(cva}\SpecialCharTok{$}\NormalTok{CVscores,}\AttributeTok{group=}\NormalTok{species, }\AttributeTok{kind=}\StringTok{"sd"}\NormalTok{, }\AttributeTok{conf=}\FloatTok{0.95}\NormalTok{)\} }
\end{Highlighting}
\end{Shaded}

\includegraphics{03.6-mgord_files/figure-latex/unnamed-chunk-7-1.pdf}

\begin{Shaded}
\begin{Highlighting}[]
\CommentTok{\# ou: ordihull(cva$CVscores,group=species)}
\end{Highlighting}
\end{Shaded}

Uma outra vantagem da CVA é que ela pode ser usada para estimar uma taxa de reclassificação correta de indivíduos dentro dos grupos hipotetizados.

\begin{Shaded}
\begin{Highlighting}[]
\NormalTok{cva}\OtherTok{\textless{}{-}}\FunctionTok{CVA}\NormalTok{(PCA}\SpecialCharTok{$}\NormalTok{x[,}\DecValTok{1}\SpecialCharTok{:}\DecValTok{20}\NormalTok{],}\AttributeTok{groups=}\NormalTok{species,}\AttributeTok{cv=}\ConstantTok{TRUE}\NormalTok{) }\CommentTok{\# Jackknife Cross{-}validation}
\NormalTok{cva}
\CommentTok{\#\textgreater{}  cross{-}validated classification results in frequencies}
\CommentTok{\#\textgreater{}               }
\CommentTok{\#\textgreater{}                C.flamarioni C.minutus C.torquatus}
\CommentTok{\#\textgreater{}   C.flamarioni            9         0           1}
\CommentTok{\#\textgreater{}   C.minutus               0         7           3}
\CommentTok{\#\textgreater{}   C.torquatus             0         0          10}
\CommentTok{\#\textgreater{} }
\CommentTok{\#\textgreater{} }
\CommentTok{\#\textgreater{}  cross{-}validated classification result in \%}
\CommentTok{\#\textgreater{}               }
\CommentTok{\#\textgreater{}                C.flamarioni C.minutus C.torquatus}
\CommentTok{\#\textgreater{}   C.flamarioni           90         0          10}
\CommentTok{\#\textgreater{}   C.minutus               0        70          30}
\CommentTok{\#\textgreater{}   C.torquatus             0         0         100}
\CommentTok{\#\textgreater{} }
\CommentTok{\#\textgreater{} }
\CommentTok{\#\textgreater{}  overall classification accuracy: 86.66667 \%}
\CommentTok{\#\textgreater{} }
\CommentTok{\#\textgreater{}  Kappa statistic: 0.8}
\NormalTok{tab}\OtherTok{\textless{}{-}}\FunctionTok{table}\NormalTok{(species,cva}\SpecialCharTok{$}\NormalTok{class)}
\NormalTok{lda.p}\OtherTok{\textless{}{-}}\FunctionTok{diag}\NormalTok{(tab)}\SpecialCharTok{/}\FunctionTok{summary}\NormalTok{(species)}\SpecialCharTok{*}\DecValTok{100}
\NormalTok{lda.p }\CommentTok{\# proporção de classificação correta para cada grupo}
\CommentTok{\#\textgreater{} C.flamarioni    C.minutus  C.torquatus }
\CommentTok{\#\textgreater{}           90           70          100}
\end{Highlighting}
\end{Shaded}

Visualização da forma ao longo dos eixos da CVA.

\begin{Shaded}
\begin{Highlighting}[]
\CommentTok{\# Mudança de forma ao longo do CV1 (positivo 15, negativo {-}10) e CV2}
\NormalTok{cv1max}\OtherTok{\textless{}{-}} \DecValTok{15}\SpecialCharTok{*}\NormalTok{cva}\SpecialCharTok{$}\NormalTok{CVvis[,}\DecValTok{1}\NormalTok{]}\SpecialCharTok{+}\NormalTok{cva}\SpecialCharTok{$}\NormalTok{Grandm}
\NormalTok{cv1max }\OtherTok{\textless{}{-}} \FunctionTok{restoreShapes}\NormalTok{(cv1max,PCA}\SpecialCharTok{$}\NormalTok{x[,}\DecValTok{1}\SpecialCharTok{:}\DecValTok{20}\NormalTok{],ref)}
\NormalTok{cv1min}\OtherTok{\textless{}{-}} \SpecialCharTok{{-}}\DecValTok{10}\SpecialCharTok{*}\NormalTok{cva}\SpecialCharTok{$}\NormalTok{CVvis[,}\DecValTok{1}\NormalTok{]}\SpecialCharTok{+}\NormalTok{cva}\SpecialCharTok{$}\NormalTok{Grandm}
\NormalTok{cv1min }\OtherTok{\textless{}{-}} \FunctionTok{restoreShapes}\NormalTok{(cv1min,PCA}\SpecialCharTok{$}\NormalTok{x[,}\DecValTok{1}\SpecialCharTok{:}\DecValTok{20}\NormalTok{],ref)}
\NormalTok{cv2max}\OtherTok{\textless{}{-}} \DecValTok{6}\SpecialCharTok{*}\NormalTok{cva}\SpecialCharTok{$}\NormalTok{CVvis[,}\DecValTok{2}\NormalTok{]}\SpecialCharTok{+}\NormalTok{cva}\SpecialCharTok{$}\NormalTok{Grandm}
\NormalTok{cv2max }\OtherTok{\textless{}{-}} \FunctionTok{restoreShapes}\NormalTok{(cv2max,PCA}\SpecialCharTok{$}\NormalTok{x[,}\DecValTok{1}\SpecialCharTok{:}\DecValTok{20}\NormalTok{],ref)}
\NormalTok{cv2min}\OtherTok{\textless{}{-}} \SpecialCharTok{{-}}\DecValTok{6}\SpecialCharTok{*}\NormalTok{cva}\SpecialCharTok{$}\NormalTok{CVvis[,}\DecValTok{2}\NormalTok{]}\SpecialCharTok{+}\NormalTok{cva}\SpecialCharTok{$}\NormalTok{Grandm}
\NormalTok{cv2min }\OtherTok{\textless{}{-}} \FunctionTok{restoreShapes}\NormalTok{(cv2min,PCA}\SpecialCharTok{$}\NormalTok{x[,}\DecValTok{1}\SpecialCharTok{:}\DecValTok{20}\NormalTok{],ref)}
\FunctionTok{plotRefToTarget}\NormalTok{(cv1min,cv1max,}\AttributeTok{links=}\NormalTok{links.tuco,}\AttributeTok{method=}\StringTok{"points"}\NormalTok{) }
\end{Highlighting}
\end{Shaded}

\includegraphics{03.6-mgord_files/figure-latex/unnamed-chunk-9-1.pdf}

\begin{Shaded}
\begin{Highlighting}[]
\FunctionTok{plotRefToTarget}\NormalTok{(cv2min,cv2max,}\AttributeTok{links=}\NormalTok{links.tuco,}\AttributeTok{method=}\StringTok{"points"}\NormalTok{) }
\end{Highlighting}
\end{Shaded}

\includegraphics{03.6-mgord_files/figure-latex/unnamed-chunk-9-2.pdf}

Veja também os modelos linares para discriminação de grupos no subcapítulo correspondente.

\section{Exercício -- Análise de Componentes Principais}\label{exercuxedcio-anuxe1lise-de-componentes-principais}

Usando os mesmos dados deste exercício, faça uma PCA com a função própria do pacote \emph{geomorph}, e uma PCA com a função `prcomp'. Qual a diferença entre as duas PCAs? Como a diferença pode ser explicada?

\chapter{Partial Least Squares}\label{partial-least-squares}

Partial Least Squares (PLS) é um método que explora padrões de covariação entre dois ou mais blocos de variáveis. De modo similar à PCA/CVA, a PLS cria eixos ortogonais que, neste caso, explicam a covariância entre blocos. Em análises de MG, um dos blocos normalmente é a forma. O(s) bloco(s) adicional pode ser composto, por exemplo, por um grupo de variáveis climáticas, geográficas e/ou bióticas ou por outros descritores fenotípicos. Além disso, as variáveis dentro dos blocos não precisam ser independentes uma da outra, e o número de variáveis pode exceder o número de observações. Os blocos de variáveis podem conter variáveis de diversas naturezas (contínuas, categóricas). É importante que a ordem de entrada dos dados corresponda entre os dois blocos.

No exemplo, vamos explorar a covariação entre a forma do crânio e a forma da mandíbula de roedores sigmodontíneos.

Carregar dados.

\begin{Shaded}
\begin{Highlighting}[]
\FunctionTok{require}\NormalTok{(geomorph)}
\CommentTok{\#\textgreater{} Carregando pacotes exigidos: geomorph}
\CommentTok{\#\textgreater{} Carregando pacotes exigidos: RRPP}
\CommentTok{\#\textgreater{} Carregando pacotes exigidos: rgl}
\CommentTok{\#\textgreater{} Carregando pacotes exigidos: Matrix}
\CommentTok{\# Carregar arquivo tps com a vista ventral do crânio}
\NormalTok{tps.v}\OtherTok{\textless{}{-}}\FunctionTok{readland.tps}\NormalTok{(}\StringTok{"dadosmg/ventral.dig\_curso.pls.tps"}\NormalTok{,}\AttributeTok{specID =} \StringTok{"ID"}\NormalTok{, }\AttributeTok{readcurves =} \ConstantTok{FALSE}\NormalTok{)}
\CommentTok{\# Matriz com os pares de landmarks simétricos}
\NormalTok{pairs.matrix}\OtherTok{\textless{}{-}}
\FunctionTok{matrix}\NormalTok{(}\FunctionTok{c}\NormalTok{(}\DecValTok{2}\NormalTok{,}\DecValTok{3}\NormalTok{,}\DecValTok{4}\NormalTok{,}\DecValTok{5}\NormalTok{,}\DecValTok{6}\NormalTok{,}\DecValTok{7}\NormalTok{,}\DecValTok{8}\NormalTok{,}\DecValTok{9}\NormalTok{,}\DecValTok{10}\NormalTok{,}\DecValTok{11}\NormalTok{,}\DecValTok{12}\NormalTok{,}\DecValTok{13}\NormalTok{,}\DecValTok{14}\NormalTok{,}\DecValTok{15}\NormalTok{,}\DecValTok{16}\NormalTok{,}\DecValTok{17}\NormalTok{,}\DecValTok{18}\NormalTok{,}\DecValTok{24}\NormalTok{,}\DecValTok{22}\NormalTok{,}\DecValTok{28}\NormalTok{,}\DecValTok{23}\NormalTok{,}\DecValTok{29}\NormalTok{,}\DecValTok{19}\NormalTok{,}\DecValTok{25}\NormalTok{,}\DecValTok{20}\NormalTok{,}\DecValTok{26}\NormalTok{,}\DecValTok{21}\NormalTok{,}\DecValTok{27}\NormalTok{,}\DecValTok{30}\NormalTok{,}\DecValTok{31}\NormalTok{,}\DecValTok{33}\NormalTok{,}\DecValTok{34}\NormalTok{,}\DecValTok{35}\NormalTok{,}\DecValTok{36}\NormalTok{,}\DecValTok{37}\NormalTok{,}\DecValTok{38}\NormalTok{,}\DecValTok{39}\NormalTok{,}\DecValTok{40}\NormalTok{,}\DecValTok{41}\NormalTok{,}\DecValTok{42}\NormalTok{,}\DecValTok{43}\NormalTok{,}\DecValTok{44}\NormalTok{,}\DecValTok{45}\NormalTok{,}\DecValTok{46}\NormalTok{,}\DecValTok{47}\NormalTok{,}\DecValTok{48}\NormalTok{,}\DecValTok{49}\NormalTok{,}\DecValTok{50}\NormalTok{,}\DecValTok{51}\NormalTok{,}\DecValTok{52}\NormalTok{,}\DecValTok{55}\NormalTok{,}\DecValTok{56}\NormalTok{),}\AttributeTok{nrow=}\DecValTok{26}\NormalTok{,}\AttributeTok{ncol=}\DecValTok{2}\NormalTok{,}\AttributeTok{byrow=}\NormalTok{T)}
\NormalTok{ind.v}\OtherTok{\textless{}{-}}\FunctionTok{c}\NormalTok{(}\DecValTok{1}\SpecialCharTok{:}\FunctionTok{dim}\NormalTok{(tps.v)[}\DecValTok{3}\NormalTok{]) }\CommentTok{\# vetor indivíduos}
\CommentTok{\# GPA {-} Simetria bilateral {-} Extrair componente simetrico da forma}
\NormalTok{b.s}\OtherTok{\textless{}{-}}\FunctionTok{bilat.symmetry}\NormalTok{(tps.v,}\AttributeTok{ind=}\NormalTok{ind.v,}\AttributeTok{object.sym=}\ConstantTok{TRUE}\NormalTok{,}\AttributeTok{land.pairs=}\NormalTok{pairs.matrix)}
\NormalTok{shape.v}\OtherTok{\textless{}{-}}\NormalTok{b.s}\SpecialCharTok{$}\NormalTok{symm.shape }\CommentTok{\# componente simétrico da forma}
\FunctionTok{dimnames}\NormalTok{(shape.v)}\OtherTok{=}\FunctionTok{dimnames}\NormalTok{(tps.v) }\CommentTok{\#manter os nomes dos espécimes}
\NormalTok{shape.v}
\CommentTok{\# Forma média ventral}
\NormalTok{ref.ventral}\OtherTok{\textless{}{-}}\FunctionTok{mshape}\NormalTok{(shape.v)}

\CommentTok{\# Carregar arquivo tps da mandíbula}
\NormalTok{tps}\OtherTok{\textless{}{-}}\FunctionTok{readland.tps}\NormalTok{(}\StringTok{"dadosmg/mandibula.dig\_curso.pls.tps"}\NormalTok{,}\AttributeTok{specID =} \StringTok{"ID"}\NormalTok{, }\AttributeTok{readcurves =} \ConstantTok{FALSE}\NormalTok{)}
\NormalTok{gpa}\OtherTok{\textless{}{-}}\FunctionTok{gpagen}\NormalTok{(tps)}
\NormalTok{shape}\OtherTok{\textless{}{-}}\NormalTok{gpa}\SpecialCharTok{$}\NormalTok{coords}
\NormalTok{size}\OtherTok{\textless{}{-}}\NormalTok{gpa}\SpecialCharTok{$}\NormalTok{Csize}
\NormalTok{ref.mand}\OtherTok{\textless{}{-}}\FunctionTok{mshape}\NormalTok{(shape)}
\end{Highlighting}
\end{Shaded}

Carregar classificadores a partir de lista externa.

\begin{Shaded}
\begin{Highlighting}[]
\NormalTok{plan}\OtherTok{\textless{}{-}}\FunctionTok{read.table}\NormalTok{(}\StringTok{"dadosmg/Planilha\_PLS.txt"}\NormalTok{,}\AttributeTok{h=}\NormalTok{T)}
\NormalTok{plan}
\NormalTok{species.pls}\OtherTok{\textless{}{-}}\NormalTok{plan[,}\DecValTok{2}\NormalTok{]}
\NormalTok{species.pls}
\end{Highlighting}
\end{Shaded}

Partial Least Squares - PLS.

\begin{Shaded}
\begin{Highlighting}[]
\NormalTok{pls}\OtherTok{\textless{}{-}}\FunctionTok{two.b.pls}\NormalTok{(shape,shape.v,}\AttributeTok{iter=}\DecValTok{999}\NormalTok{)}
\CommentTok{\#\textgreater{} }
\CommentTok{\#\textgreater{} Random PLS calculations: 1000 permutations.}
\CommentTok{\#\textgreater{} }
  \SpecialCharTok{|}                                                        
  \ErrorTok{|}                                                  \ErrorTok{|}   \DecValTok{0}\NormalTok{\%}
  \SpecialCharTok{|}                                                        
  \ErrorTok{|}                                                  \ErrorTok{|}   \DecValTok{1}\NormalTok{\%}
  \SpecialCharTok{|}                                                        
  \ErrorTok{|=}                                                 \ErrorTok{|}   \DecValTok{1}\NormalTok{\%}
  \SpecialCharTok{|}                                                        
  \ErrorTok{|=}                                                 \ErrorTok{|}   \DecValTok{2}\NormalTok{\%}
  \SpecialCharTok{|}                                                        
  \ErrorTok{|=}                                                 \ErrorTok{|}   \DecValTok{3}\NormalTok{\%}
  \SpecialCharTok{|}                                                        
  \ErrorTok{|==}                                                \ErrorTok{|}   \DecValTok{3}\NormalTok{\%}
  \SpecialCharTok{|}                                                        
  \ErrorTok{|==}                                                \ErrorTok{|}   \DecValTok{4}\NormalTok{\%}
  \SpecialCharTok{|}                                                        
  \ErrorTok{|==}                                                \ErrorTok{|}   \DecValTok{5}\NormalTok{\%}
  \SpecialCharTok{|}                                                        
  \ErrorTok{|===}                                               \ErrorTok{|}   \DecValTok{5}\NormalTok{\%}
  \SpecialCharTok{|}                                                        
  \ErrorTok{|===}                                               \ErrorTok{|}   \DecValTok{6}\NormalTok{\%}
  \SpecialCharTok{|}                                                        
  \ErrorTok{|===}                                               \ErrorTok{|}   \DecValTok{7}\NormalTok{\%}
  \SpecialCharTok{|}                                                        
  \ErrorTok{|====}                                              \ErrorTok{|}   \DecValTok{7}\NormalTok{\%}
  \SpecialCharTok{|}                                                        
  \ErrorTok{|====}                                              \ErrorTok{|}   \DecValTok{8}\NormalTok{\%}
  \SpecialCharTok{|}                                                        
  \ErrorTok{|====}                                              \ErrorTok{|}   \DecValTok{9}\NormalTok{\%}
  \SpecialCharTok{|}                                                        
  \ErrorTok{|=====}                                             \ErrorTok{|}   \DecValTok{9}\NormalTok{\%}
  \SpecialCharTok{|}                                                        
  \ErrorTok{|=====}                                             \ErrorTok{|}  \DecValTok{10}\NormalTok{\%}
  \SpecialCharTok{|}                                                        
  \ErrorTok{|=====}                                             \ErrorTok{|}  \DecValTok{11}\NormalTok{\%}
  \SpecialCharTok{|}                                                        
  \ErrorTok{|======}                                            \ErrorTok{|}  \DecValTok{11}\NormalTok{\%}
  \SpecialCharTok{|}                                                        
  \ErrorTok{|======}                                            \ErrorTok{|}  \DecValTok{12}\NormalTok{\%}
  \SpecialCharTok{|}                                                        
  \ErrorTok{|======}                                            \ErrorTok{|}  \DecValTok{13}\NormalTok{\%}
  \SpecialCharTok{|}                                                        
  \ErrorTok{|=======}                                           \ErrorTok{|}  \DecValTok{13}\NormalTok{\%}
  \SpecialCharTok{|}                                                        
  \ErrorTok{|=======}                                           \ErrorTok{|}  \DecValTok{14}\NormalTok{\%}
  \SpecialCharTok{|}                                                        
  \ErrorTok{|=======}                                           \ErrorTok{|}  \DecValTok{15}\NormalTok{\%}
  \SpecialCharTok{|}                                                        
  \ErrorTok{|========}                                          \ErrorTok{|}  \DecValTok{15}\NormalTok{\%}
  \SpecialCharTok{|}                                                        
  \ErrorTok{|========}                                          \ErrorTok{|}  \DecValTok{16}\NormalTok{\%}
  \SpecialCharTok{|}                                                        
  \ErrorTok{|========}                                          \ErrorTok{|}  \DecValTok{17}\NormalTok{\%}
  \SpecialCharTok{|}                                                        
  \ErrorTok{|=========}                                         \ErrorTok{|}  \DecValTok{17}\NormalTok{\%}
  \SpecialCharTok{|}                                                        
  \ErrorTok{|=========}                                         \ErrorTok{|}  \DecValTok{18}\NormalTok{\%}
  \SpecialCharTok{|}                                                        
  \ErrorTok{|=========}                                         \ErrorTok{|}  \DecValTok{19}\NormalTok{\%}
  \SpecialCharTok{|}                                                        
  \ErrorTok{|==========}                                        \ErrorTok{|}  \DecValTok{19}\NormalTok{\%}
  \SpecialCharTok{|}                                                        
  \ErrorTok{|==========}                                        \ErrorTok{|}  \DecValTok{20}\NormalTok{\%}
  \SpecialCharTok{|}                                                        
  \ErrorTok{|==========}                                        \ErrorTok{|}  \DecValTok{21}\NormalTok{\%}
  \SpecialCharTok{|}                                                        
  \ErrorTok{|===========}                                       \ErrorTok{|}  \DecValTok{21}\NormalTok{\%}
  \SpecialCharTok{|}                                                        
  \ErrorTok{|===========}                                       \ErrorTok{|}  \DecValTok{22}\NormalTok{\%}
  \SpecialCharTok{|}                                                        
  \ErrorTok{|===========}                                       \ErrorTok{|}  \DecValTok{23}\NormalTok{\%}
  \SpecialCharTok{|}                                                        
  \ErrorTok{|============}                                      \ErrorTok{|}  \DecValTok{23}\NormalTok{\%}
  \SpecialCharTok{|}                                                        
  \ErrorTok{|============}                                      \ErrorTok{|}  \DecValTok{24}\NormalTok{\%}
  \SpecialCharTok{|}                                                        
  \ErrorTok{|============}                                      \ErrorTok{|}  \DecValTok{25}\NormalTok{\%}
  \SpecialCharTok{|}                                                        
  \ErrorTok{|=============}                                     \ErrorTok{|}  \DecValTok{25}\NormalTok{\%}
  \SpecialCharTok{|}                                                        
  \ErrorTok{|=============}                                     \ErrorTok{|}  \DecValTok{26}\NormalTok{\%}
  \SpecialCharTok{|}                                                        
  \ErrorTok{|=============}                                     \ErrorTok{|}  \DecValTok{27}\NormalTok{\%}
  \SpecialCharTok{|}                                                        
  \ErrorTok{|==============}                                    \ErrorTok{|}  \DecValTok{27}\NormalTok{\%}
  \SpecialCharTok{|}                                                        
  \ErrorTok{|==============}                                    \ErrorTok{|}  \DecValTok{28}\NormalTok{\%}
  \SpecialCharTok{|}                                                        
  \ErrorTok{|==============}                                    \ErrorTok{|}  \DecValTok{29}\NormalTok{\%}
  \SpecialCharTok{|}                                                        
  \ErrorTok{|===============}                                   \ErrorTok{|}  \DecValTok{29}\NormalTok{\%}
  \SpecialCharTok{|}                                                        
  \ErrorTok{|===============}                                   \ErrorTok{|}  \DecValTok{30}\NormalTok{\%}
  \SpecialCharTok{|}                                                        
  \ErrorTok{|===============}                                   \ErrorTok{|}  \DecValTok{31}\NormalTok{\%}
  \SpecialCharTok{|}                                                        
  \ErrorTok{|================}                                  \ErrorTok{|}  \DecValTok{31}\NormalTok{\%}
  \SpecialCharTok{|}                                                        
  \ErrorTok{|================}                                  \ErrorTok{|}  \DecValTok{32}\NormalTok{\%}
  \SpecialCharTok{|}                                                        
  \ErrorTok{|================}                                  \ErrorTok{|}  \DecValTok{33}\NormalTok{\%}
  \SpecialCharTok{|}                                                        
  \ErrorTok{|=================}                                 \ErrorTok{|}  \DecValTok{33}\NormalTok{\%}
  \SpecialCharTok{|}                                                        
  \ErrorTok{|=================}                                 \ErrorTok{|}  \DecValTok{34}\NormalTok{\%}
  \SpecialCharTok{|}                                                        
  \ErrorTok{|=================}                                 \ErrorTok{|}  \DecValTok{35}\NormalTok{\%}
  \SpecialCharTok{|}                                                        
  \ErrorTok{|==================}                                \ErrorTok{|}  \DecValTok{35}\NormalTok{\%}
  \SpecialCharTok{|}                                                        
  \ErrorTok{|==================}                                \ErrorTok{|}  \DecValTok{36}\NormalTok{\%}
  \SpecialCharTok{|}                                                        
  \ErrorTok{|==================}                                \ErrorTok{|}  \DecValTok{37}\NormalTok{\%}
  \SpecialCharTok{|}                                                        
  \ErrorTok{|===================}                               \ErrorTok{|}  \DecValTok{37}\NormalTok{\%}
  \SpecialCharTok{|}                                                        
  \ErrorTok{|===================}                               \ErrorTok{|}  \DecValTok{38}\NormalTok{\%}
  \SpecialCharTok{|}                                                        
  \ErrorTok{|===================}                               \ErrorTok{|}  \DecValTok{39}\NormalTok{\%}
  \SpecialCharTok{|}                                                        
  \ErrorTok{|====================}                              \ErrorTok{|}  \DecValTok{39}\NormalTok{\%}
  \SpecialCharTok{|}                                                        
  \ErrorTok{|====================}                              \ErrorTok{|}  \DecValTok{40}\NormalTok{\%}
  \SpecialCharTok{|}                                                        
  \ErrorTok{|====================}                              \ErrorTok{|}  \DecValTok{41}\NormalTok{\%}
  \SpecialCharTok{|}                                                        
  \ErrorTok{|=====================}                             \ErrorTok{|}  \DecValTok{41}\NormalTok{\%}
  \SpecialCharTok{|}                                                        
  \ErrorTok{|=====================}                             \ErrorTok{|}  \DecValTok{42}\NormalTok{\%}
  \SpecialCharTok{|}                                                        
  \ErrorTok{|=====================}                             \ErrorTok{|}  \DecValTok{43}\NormalTok{\%}
  \SpecialCharTok{|}                                                        
  \ErrorTok{|======================}                            \ErrorTok{|}  \DecValTok{43}\NormalTok{\%}
  \SpecialCharTok{|}                                                        
  \ErrorTok{|======================}                            \ErrorTok{|}  \DecValTok{44}\NormalTok{\%}
  \SpecialCharTok{|}                                                        
  \ErrorTok{|======================}                            \ErrorTok{|}  \DecValTok{45}\NormalTok{\%}
  \SpecialCharTok{|}                                                        
  \ErrorTok{|=======================}                           \ErrorTok{|}  \DecValTok{45}\NormalTok{\%}
  \SpecialCharTok{|}                                                        
  \ErrorTok{|=======================}                           \ErrorTok{|}  \DecValTok{46}\NormalTok{\%}
  \SpecialCharTok{|}                                                        
  \ErrorTok{|=======================}                           \ErrorTok{|}  \DecValTok{47}\NormalTok{\%}
  \SpecialCharTok{|}                                                        
  \ErrorTok{|========================}                          \ErrorTok{|}  \DecValTok{47}\NormalTok{\%}
  \SpecialCharTok{|}                                                        
  \ErrorTok{|========================}                          \ErrorTok{|}  \DecValTok{48}\NormalTok{\%}
  \SpecialCharTok{|}                                                        
  \ErrorTok{|========================}                          \ErrorTok{|}  \DecValTok{49}\NormalTok{\%}
  \SpecialCharTok{|}                                                        
  \ErrorTok{|=========================}                         \ErrorTok{|}  \DecValTok{49}\NormalTok{\%}
  \SpecialCharTok{|}                                                        
  \ErrorTok{|=========================}                         \ErrorTok{|}  \DecValTok{50}\NormalTok{\%}
  \SpecialCharTok{|}                                                        
  \ErrorTok{|=========================}                         \ErrorTok{|}  \DecValTok{51}\NormalTok{\%}
  \SpecialCharTok{|}                                                        
  \ErrorTok{|==========================}                        \ErrorTok{|}  \DecValTok{51}\NormalTok{\%}
  \SpecialCharTok{|}                                                        
  \ErrorTok{|==========================}                        \ErrorTok{|}  \DecValTok{52}\NormalTok{\%}
  \SpecialCharTok{|}                                                        
  \ErrorTok{|==========================}                        \ErrorTok{|}  \DecValTok{53}\NormalTok{\%}
  \SpecialCharTok{|}                                                        
  \ErrorTok{|===========================}                       \ErrorTok{|}  \DecValTok{53}\NormalTok{\%}
  \SpecialCharTok{|}                                                        
  \ErrorTok{|===========================}                       \ErrorTok{|}  \DecValTok{54}\NormalTok{\%}
  \SpecialCharTok{|}                                                        
  \ErrorTok{|===========================}                       \ErrorTok{|}  \DecValTok{55}\NormalTok{\%}
  \SpecialCharTok{|}                                                        
  \ErrorTok{|============================}                      \ErrorTok{|}  \DecValTok{55}\NormalTok{\%}
  \SpecialCharTok{|}                                                        
  \ErrorTok{|============================}                      \ErrorTok{|}  \DecValTok{56}\NormalTok{\%}
  \SpecialCharTok{|}                                                        
  \ErrorTok{|============================}                      \ErrorTok{|}  \DecValTok{57}\NormalTok{\%}
  \SpecialCharTok{|}                                                        
  \ErrorTok{|=============================}                     \ErrorTok{|}  \DecValTok{57}\NormalTok{\%}
  \SpecialCharTok{|}                                                        
  \ErrorTok{|=============================}                     \ErrorTok{|}  \DecValTok{58}\NormalTok{\%}
  \SpecialCharTok{|}                                                        
  \ErrorTok{|=============================}                     \ErrorTok{|}  \DecValTok{59}\NormalTok{\%}
  \SpecialCharTok{|}                                                        
  \ErrorTok{|==============================}                    \ErrorTok{|}  \DecValTok{59}\NormalTok{\%}
  \SpecialCharTok{|}                                                        
  \ErrorTok{|==============================}                    \ErrorTok{|}  \DecValTok{60}\NormalTok{\%}
  \SpecialCharTok{|}                                                        
  \ErrorTok{|==============================}                    \ErrorTok{|}  \DecValTok{61}\NormalTok{\%}
  \SpecialCharTok{|}                                                        
  \ErrorTok{|===============================}                   \ErrorTok{|}  \DecValTok{61}\NormalTok{\%}
  \SpecialCharTok{|}                                                        
  \ErrorTok{|===============================}                   \ErrorTok{|}  \DecValTok{62}\NormalTok{\%}
  \SpecialCharTok{|}                                                        
  \ErrorTok{|===============================}                   \ErrorTok{|}  \DecValTok{63}\NormalTok{\%}
  \SpecialCharTok{|}                                                        
  \ErrorTok{|================================}                  \ErrorTok{|}  \DecValTok{63}\NormalTok{\%}
  \SpecialCharTok{|}                                                        
  \ErrorTok{|================================}                  \ErrorTok{|}  \DecValTok{64}\NormalTok{\%}
  \SpecialCharTok{|}                                                        
  \ErrorTok{|================================}                  \ErrorTok{|}  \DecValTok{65}\NormalTok{\%}
  \SpecialCharTok{|}                                                        
  \ErrorTok{|=================================}                 \ErrorTok{|}  \DecValTok{65}\NormalTok{\%}
  \SpecialCharTok{|}                                                        
  \ErrorTok{|=================================}                 \ErrorTok{|}  \DecValTok{66}\NormalTok{\%}
  \SpecialCharTok{|}                                                        
  \ErrorTok{|=================================}                 \ErrorTok{|}  \DecValTok{67}\NormalTok{\%}
  \SpecialCharTok{|}                                                        
  \ErrorTok{|==================================}                \ErrorTok{|}  \DecValTok{67}\NormalTok{\%}
  \SpecialCharTok{|}                                                        
  \ErrorTok{|==================================}                \ErrorTok{|}  \DecValTok{68}\NormalTok{\%}
  \SpecialCharTok{|}                                                        
  \ErrorTok{|==================================}                \ErrorTok{|}  \DecValTok{69}\NormalTok{\%}
  \SpecialCharTok{|}                                                        
  \ErrorTok{|===================================}               \ErrorTok{|}  \DecValTok{69}\NormalTok{\%}
  \SpecialCharTok{|}                                                        
  \ErrorTok{|===================================}               \ErrorTok{|}  \DecValTok{70}\NormalTok{\%}
  \SpecialCharTok{|}                                                        
  \ErrorTok{|===================================}               \ErrorTok{|}  \DecValTok{71}\NormalTok{\%}
  \SpecialCharTok{|}                                                        
  \ErrorTok{|====================================}              \ErrorTok{|}  \DecValTok{71}\NormalTok{\%}
  \SpecialCharTok{|}                                                        
  \ErrorTok{|====================================}              \ErrorTok{|}  \DecValTok{72}\NormalTok{\%}
  \SpecialCharTok{|}                                                        
  \ErrorTok{|====================================}              \ErrorTok{|}  \DecValTok{73}\NormalTok{\%}
  \SpecialCharTok{|}                                                        
  \ErrorTok{|=====================================}             \ErrorTok{|}  \DecValTok{73}\NormalTok{\%}
  \SpecialCharTok{|}                                                        
  \ErrorTok{|=====================================}             \ErrorTok{|}  \DecValTok{74}\NormalTok{\%}
  \SpecialCharTok{|}                                                        
  \ErrorTok{|=====================================}             \ErrorTok{|}  \DecValTok{75}\NormalTok{\%}
  \SpecialCharTok{|}                                                        
  \ErrorTok{|======================================}            \ErrorTok{|}  \DecValTok{75}\NormalTok{\%}
  \SpecialCharTok{|}                                                        
  \ErrorTok{|======================================}            \ErrorTok{|}  \DecValTok{76}\NormalTok{\%}
  \SpecialCharTok{|}                                                        
  \ErrorTok{|======================================}            \ErrorTok{|}  \DecValTok{77}\NormalTok{\%}
  \SpecialCharTok{|}                                                        
  \ErrorTok{|=======================================}           \ErrorTok{|}  \DecValTok{77}\NormalTok{\%}
  \SpecialCharTok{|}                                                        
  \ErrorTok{|=======================================}           \ErrorTok{|}  \DecValTok{78}\NormalTok{\%}
  \SpecialCharTok{|}                                                        
  \ErrorTok{|=======================================}           \ErrorTok{|}  \DecValTok{79}\NormalTok{\%}
  \SpecialCharTok{|}                                                        
  \ErrorTok{|========================================}          \ErrorTok{|}  \DecValTok{79}\NormalTok{\%}
  \SpecialCharTok{|}                                                        
  \ErrorTok{|========================================}          \ErrorTok{|}  \DecValTok{80}\NormalTok{\%}
  \SpecialCharTok{|}                                                        
  \ErrorTok{|========================================}          \ErrorTok{|}  \DecValTok{81}\NormalTok{\%}
  \SpecialCharTok{|}                                                        
  \ErrorTok{|=========================================}         \ErrorTok{|}  \DecValTok{81}\NormalTok{\%}
  \SpecialCharTok{|}                                                        
  \ErrorTok{|=========================================}         \ErrorTok{|}  \DecValTok{82}\NormalTok{\%}
  \SpecialCharTok{|}                                                        
  \ErrorTok{|=========================================}         \ErrorTok{|}  \DecValTok{83}\NormalTok{\%}
  \SpecialCharTok{|}                                                        
  \ErrorTok{|==========================================}        \ErrorTok{|}  \DecValTok{83}\NormalTok{\%}
  \SpecialCharTok{|}                                                        
  \ErrorTok{|==========================================}        \ErrorTok{|}  \DecValTok{84}\NormalTok{\%}
  \SpecialCharTok{|}                                                        
  \ErrorTok{|==========================================}        \ErrorTok{|}  \DecValTok{85}\NormalTok{\%}
  \SpecialCharTok{|}                                                        
  \ErrorTok{|===========================================}       \ErrorTok{|}  \DecValTok{85}\NormalTok{\%}
  \SpecialCharTok{|}                                                        
  \ErrorTok{|===========================================}       \ErrorTok{|}  \DecValTok{86}\NormalTok{\%}
  \SpecialCharTok{|}                                                        
  \ErrorTok{|===========================================}       \ErrorTok{|}  \DecValTok{87}\NormalTok{\%}
  \SpecialCharTok{|}                                                        
  \ErrorTok{|============================================}      \ErrorTok{|}  \DecValTok{87}\NormalTok{\%}
  \SpecialCharTok{|}                                                        
  \ErrorTok{|============================================}      \ErrorTok{|}  \DecValTok{88}\NormalTok{\%}
  \SpecialCharTok{|}                                                        
  \ErrorTok{|============================================}      \ErrorTok{|}  \DecValTok{89}\NormalTok{\%}
  \SpecialCharTok{|}                                                        
  \ErrorTok{|=============================================}     \ErrorTok{|}  \DecValTok{89}\NormalTok{\%}
  \SpecialCharTok{|}                                                        
  \ErrorTok{|=============================================}     \ErrorTok{|}  \DecValTok{90}\NormalTok{\%}
  \SpecialCharTok{|}                                                        
  \ErrorTok{|=============================================}     \ErrorTok{|}  \DecValTok{91}\NormalTok{\%}
  \SpecialCharTok{|}                                                        
  \ErrorTok{|==============================================}    \ErrorTok{|}  \DecValTok{91}\NormalTok{\%}
  \SpecialCharTok{|}                                                        
  \ErrorTok{|==============================================}    \ErrorTok{|}  \DecValTok{92}\NormalTok{\%}
  \SpecialCharTok{|}                                                        
  \ErrorTok{|==============================================}    \ErrorTok{|}  \DecValTok{93}\NormalTok{\%}
  \SpecialCharTok{|}                                                        
  \ErrorTok{|===============================================}   \ErrorTok{|}  \DecValTok{93}\NormalTok{\%}
  \SpecialCharTok{|}                                                        
  \ErrorTok{|===============================================}   \ErrorTok{|}  \DecValTok{94}\NormalTok{\%}
  \SpecialCharTok{|}                                                        
  \ErrorTok{|===============================================}   \ErrorTok{|}  \DecValTok{95}\NormalTok{\%}
  \SpecialCharTok{|}                                                        
  \ErrorTok{|================================================}  \ErrorTok{|}  \DecValTok{95}\NormalTok{\%}
  \SpecialCharTok{|}                                                        
  \ErrorTok{|================================================}  \ErrorTok{|}  \DecValTok{96}\NormalTok{\%}
  \SpecialCharTok{|}                                                        
  \ErrorTok{|================================================}  \ErrorTok{|}  \DecValTok{97}\NormalTok{\%}
  \SpecialCharTok{|}                                                        
  \ErrorTok{|=================================================} \ErrorTok{|}  \DecValTok{97}\NormalTok{\%}
  \SpecialCharTok{|}                                                        
  \ErrorTok{|=================================================} \ErrorTok{|}  \DecValTok{98}\NormalTok{\%}
  \SpecialCharTok{|}                                                        
  \ErrorTok{|=================================================} \ErrorTok{|}  \DecValTok{99}\NormalTok{\%}
  \SpecialCharTok{|}                                                        
  \ErrorTok{|==================================================|}  \DecValTok{99}\NormalTok{\%}
  \SpecialCharTok{|}                                                        
  \ErrorTok{|==================================================|} \DecValTok{100}\NormalTok{\%}
\NormalTok{pls}
\CommentTok{\#\textgreater{} }
\CommentTok{\#\textgreater{} Call:}
\CommentTok{\#\textgreater{} two.b.pls(A1 = shape, A2 = shape.v, iter = 999) }
\CommentTok{\#\textgreater{} }
\CommentTok{\#\textgreater{} }
\CommentTok{\#\textgreater{} }
\CommentTok{\#\textgreater{} r{-}PLS: 0.897}
\CommentTok{\#\textgreater{} }
\CommentTok{\#\textgreater{} Effect Size (Z): 4.257}
\CommentTok{\#\textgreater{} }
\CommentTok{\#\textgreater{} P{-}value: 0.001}
\CommentTok{\#\textgreater{} }
\CommentTok{\#\textgreater{} Based on 1000 random permutations}

\FunctionTok{plot}\NormalTok{(pls,}\AttributeTok{label=}\NormalTok{species.pls)}
\end{Highlighting}
\end{Shaded}

\includegraphics{03.7-mgpls_files/figure-latex/unnamed-chunk-3-1.pdf}

É possível visualizar as mudanças de forma ao longo dos eixos da PLS de diferentes maneiras.
Uma delas é usando uma função interativa do pacote \emph{geomorph}

\begin{Shaded}
\begin{Highlighting}[]
\NormalTok{pls.plot}\OtherTok{\textless{}{-}}\FunctionTok{plot}\NormalTok{(pls)}
\end{Highlighting}
\end{Shaded}

\includegraphics{03.7-mgpls_files/figure-latex/unnamed-chunk-4-1.pdf}

\begin{Shaded}
\begin{Highlighting}[]
\CommentTok{\#picknplot.shape(pls.plot)}
\end{Highlighting}
\end{Shaded}

A maneira mais tradicional é encontrar as formas a partir dos valores nos eixos da PLS.

\begin{Shaded}
\begin{Highlighting}[]
\NormalTok{pls}\OtherTok{\textless{}{-}}\FunctionTok{two.b.pls}\NormalTok{(shape,shape.v,}\AttributeTok{iter=}\DecValTok{999}\NormalTok{)}
\FunctionTok{plot}\NormalTok{(pls,}\AttributeTok{label=}\NormalTok{species.pls)}

\CommentTok{\# Valores mínimos e máximos do pls plot}
\NormalTok{preds.v}\OtherTok{\textless{}{-}}\FunctionTok{shape.predictor}\NormalTok{(shape.v,}\FunctionTok{two.d.array}\NormalTok{(shape),}\AttributeTok{method=}\StringTok{"PLS"}\NormalTok{,}
                         \AttributeTok{Vmin=}\FunctionTok{min}\NormalTok{(pls}\SpecialCharTok{$}\NormalTok{XScores[,}\DecValTok{1}\NormalTok{]),}
                         \AttributeTok{Vmax=}\FunctionTok{max}\NormalTok{(pls}\SpecialCharTok{$}\NormalTok{XScores[,}\DecValTok{1}\NormalTok{])) }
\NormalTok{preds.v}

\CommentTok{\# Valores mínimos e máximos do pls2}
\NormalTok{pls2}\OtherTok{\textless{}{-}}\FunctionTok{two.b.pls}\NormalTok{(shape.v,shape,}\AttributeTok{iter=}\DecValTok{999}\NormalTok{)}
\FunctionTok{plot}\NormalTok{(pls2,}\AttributeTok{label=}\NormalTok{species.pls)}

\NormalTok{preds.m}\OtherTok{\textless{}{-}}\FunctionTok{shape.predictor}\NormalTok{(shape,}\FunctionTok{two.d.array}\NormalTok{(shape.v),}\AttributeTok{method=}\StringTok{"PLS"}\NormalTok{,}
                         \AttributeTok{Mmin=}\FunctionTok{min}\NormalTok{(pls2}\SpecialCharTok{$}\NormalTok{XScores[,}\DecValTok{1}\NormalTok{]),}
                         \AttributeTok{Mmax=}\FunctionTok{max}\NormalTok{(pls2}\SpecialCharTok{$}\NormalTok{XScores[,}\DecValTok{1}\NormalTok{])) }
\NormalTok{preds.m}
\end{Highlighting}
\end{Shaded}

Visualização das formas mínimas e máximas dos eixos da PLS.

\begin{Shaded}
\begin{Highlighting}[]
\NormalTok{\{}\FunctionTok{plot}\NormalTok{(pls,}\AttributeTok{label=}\NormalTok{species.pls)}
\FunctionTok{par}\NormalTok{(}\AttributeTok{mfrow=}\FunctionTok{c}\NormalTok{(}\DecValTok{2}\NormalTok{,}\DecValTok{2}\NormalTok{),}\AttributeTok{mar=}\FunctionTok{c}\NormalTok{(}\FloatTok{0.01}\NormalTok{,}\FloatTok{0.01}\NormalTok{,}\FloatTok{0.01}\NormalTok{,}\FloatTok{0.01}\NormalTok{))}
\FunctionTok{plotRefToTarget}\NormalTok{(ref.mand,preds.m}\SpecialCharTok{$}\NormalTok{Mmax)}
\FunctionTok{text}\NormalTok{(}\DecValTok{0}\NormalTok{,}\FloatTok{0.4}\NormalTok{,}\AttributeTok{labels=}\StringTok{"Valores negativos PLS1 Bloco1"}\NormalTok{,}\AttributeTok{cex=}\FloatTok{1.3}\NormalTok{)}
\FunctionTok{plotRefToTarget}\NormalTok{(ref.ventral,preds.v}\SpecialCharTok{$}\NormalTok{Vmin)}
\FunctionTok{text}\NormalTok{(}\DecValTok{0}\NormalTok{,}\FloatTok{0.25}\NormalTok{,}\AttributeTok{labels=}\StringTok{"Valores negativos PLS1 Bloco2"}\NormalTok{,}\AttributeTok{cex=}\FloatTok{1.3}\NormalTok{)}
\FunctionTok{plotRefToTarget}\NormalTok{(ref.mand,preds.m}\SpecialCharTok{$}\NormalTok{Mmin)}
\FunctionTok{text}\NormalTok{(}\DecValTok{0}\NormalTok{,}\FloatTok{0.4}\NormalTok{,}\AttributeTok{labels=}\StringTok{"Valores positivos PLS1 Bloco1"}\NormalTok{,}\AttributeTok{cex=}\FloatTok{1.3}\NormalTok{)}
\FunctionTok{plotRefToTarget}\NormalTok{(ref.ventral,preds.v}\SpecialCharTok{$}\NormalTok{Vmax)}
\FunctionTok{text}\NormalTok{(}\DecValTok{0}\NormalTok{,}\FloatTok{0.25}\NormalTok{,}\AttributeTok{labels=}\StringTok{"Valores positivos PLS1 Bloco2"}\NormalTok{,}\AttributeTok{cex=}\FloatTok{1.3}\NormalTok{)\}}
\end{Highlighting}
\end{Shaded}

\includegraphics{03.7-mgpls_files/figure-latex/unnamed-chunk-6-1.pdf} \includegraphics{03.7-mgpls_files/figure-latex/unnamed-chunk-6-2.pdf}

\begin{Shaded}
\begin{Highlighting}[]
\FunctionTok{par}\NormalTok{(}\AttributeTok{mfrow=}\FunctionTok{c}\NormalTok{(}\DecValTok{1}\NormalTok{,}\DecValTok{1}\NormalTok{))}
\end{Highlighting}
\end{Shaded}

\chapter{Modelos lineares para discriminação de grupos}\label{modelos-lineares-para-discriminauxe7uxe3o-de-grupos}

Uma classe ampla de modelos preditivos está abarcada pelo nome de Modelos Lineares Generalizados, incluindo regressão múltipla e análises de variância/covariância. Uma vasta gama de modelos pode ser ajustada com a função `procD.lm' e outras do pacote \emph{geomorph}, que usa distância de Procrustes entre espécimes para calcular tamanhos de efeito e randomização para calcular valores de \emph{P} (de modo equivalente à MANOVA não paramétrica -- (\textbf{goodall1991?}; \textbf{pillar1996?}; \textbf{anderson2001?}); assim, o número de variáveis de forma, normalmente muito alto em dados de MG, não é uma restrição ao ajuste de modelos (\textbf{collyer2015?}). Podemos usar esses modelos baseados em distância para responder diversas questões quando a variável resposta é composta por variáveis de forma. Aqui, vamos aplicá-los no contexto de discriminação de grupos, quando a variável preditora é categórica.

\section{MANOVA}\label{manova}

Vamos testar se existem diferenças de forma entre espécies de tuco-tucos.

\begin{Shaded}
\begin{Highlighting}[]
\FunctionTok{require}\NormalTok{(geomorph)}
\CommentTok{\#\textgreater{} Carregando pacotes exigidos: geomorph}
\CommentTok{\#\textgreater{} Carregando pacotes exigidos: RRPP}
\CommentTok{\#\textgreater{} Carregando pacotes exigidos: rgl}
\CommentTok{\#\textgreater{} Carregando pacotes exigidos: Matrix}
\CommentTok{\# Carregar dados}
\NormalTok{tps}\OtherTok{\textless{}{-}}\FunctionTok{readland.tps}\NormalTok{(}\StringTok{"dadosmg/Lista tuco dig 2.tps"}\NormalTok{,}\AttributeTok{specID =} \StringTok{"ID"}\NormalTok{, }\AttributeTok{readcurves =} \ConstantTok{FALSE}\NormalTok{)}
\FunctionTok{dim}\NormalTok{(tps)}

\CommentTok{\# Gerar matriz com pares de landmarks simétricos}
\NormalTok{pairs.matrix}\OtherTok{\textless{}{-}}\FunctionTok{matrix}\NormalTok{(}\FunctionTok{c}\NormalTok{(}\DecValTok{2}\NormalTok{,}\DecValTok{3}\NormalTok{,}\DecValTok{5}\NormalTok{,}\DecValTok{6}\NormalTok{,}\DecValTok{7}\NormalTok{,}\DecValTok{8}\NormalTok{,}\DecValTok{10}\NormalTok{,}\DecValTok{11}\NormalTok{,}\DecValTok{12}\NormalTok{,}\DecValTok{13}\NormalTok{,}\DecValTok{14}\NormalTok{,}\DecValTok{15}\NormalTok{,}\DecValTok{16}\NormalTok{,}\DecValTok{17}\NormalTok{,}\DecValTok{18}\NormalTok{,}\DecValTok{19}\NormalTok{,}\DecValTok{20}\NormalTok{,}\DecValTok{21}\NormalTok{,}\DecValTok{23}\NormalTok{,}\DecValTok{24}\NormalTok{,}\DecValTok{25}\NormalTok{,}\DecValTok{26}\NormalTok{,}\DecValTok{27}\NormalTok{,}\DecValTok{28}\NormalTok{),}\AttributeTok{nrow=}\DecValTok{12}\NormalTok{,}\AttributeTok{ncol=}\DecValTok{2}\NormalTok{,}\AttributeTok{byrow=}\NormalTok{T)}
\CommentTok{\# Vetor indivíduos}
\NormalTok{ind}\OtherTok{\textless{}{-}}\FunctionTok{c}\NormalTok{(}\DecValTok{1}\SpecialCharTok{:}\FunctionTok{dim}\NormalTok{(tps)[}\DecValTok{3}\NormalTok{])}

\CommentTok{\# GPA {-} Simetria bilateral }
\NormalTok{b.s}\OtherTok{\textless{}{-}}\FunctionTok{bilat.symmetry}\NormalTok{(tps,}\AttributeTok{ind=}\NormalTok{ind,}\AttributeTok{object.sym=}\ConstantTok{TRUE}\NormalTok{,}\AttributeTok{land.pairs=}\NormalTok{pairs.matrix)}

\CommentTok{\# Matriz de forma}
\NormalTok{shape.sym}\OtherTok{\textless{}{-}}\NormalTok{b.s}\SpecialCharTok{$}\NormalTok{symm.shape}

\CommentTok{\# Carregar classificadores a partir de lista externa}
\NormalTok{plan}\OtherTok{\textless{}{-}}\FunctionTok{read.table}\NormalTok{(}\StringTok{"dadosmg/Planilha Ctenomys.txt"}\NormalTok{,}\AttributeTok{h=}\NormalTok{T)}
\NormalTok{plan}
\NormalTok{species}\OtherTok{\textless{}{-}}\NormalTok{plan[,}\DecValTok{1}\NormalTok{]}
\NormalTok{species}\OtherTok{\textless{}{-}}\FunctionTok{as.factor}\NormalTok{(species)}
\NormalTok{sex}\OtherTok{\textless{}{-}}\NormalTok{plan[,}\DecValTok{2}\NormalTok{]}
\NormalTok{sex}\OtherTok{\textless{}{-}}\FunctionTok{as.factor}\NormalTok{(sex)}
\NormalTok{sp.sex}\OtherTok{\textless{}{-}}\NormalTok{plan[,}\DecValTok{3}\NormalTok{]}
\NormalTok{sp.sex}\OtherTok{\textless{}{-}}\FunctionTok{as.factor}\NormalTok{(sp.sex)}
\end{Highlighting}
\end{Shaded}

MANOVA.

\begin{Shaded}
\begin{Highlighting}[]
\CommentTok{\# MANOVA (Multivariate Analysis of Variance)}
\NormalTok{manova}\OtherTok{\textless{}{-}}\FunctionTok{procD.lm}\NormalTok{(shape.sym}\SpecialCharTok{\textasciitilde{}}\NormalTok{species, }\AttributeTok{iter=}\DecValTok{999}\NormalTok{, }\AttributeTok{RRPP=} \ConstantTok{TRUE}\NormalTok{)}
\FunctionTok{summary}\NormalTok{(manova)}
\CommentTok{\#\textgreater{} }
\CommentTok{\#\textgreater{} Analysis of Variance, using Residual Randomization}
\CommentTok{\#\textgreater{} Permutation procedure: Randomization of null model residuals }
\CommentTok{\#\textgreater{} Number of permutations: 1000 }
\CommentTok{\#\textgreater{} Estimation method: Ordinary Least Squares }
\CommentTok{\#\textgreater{} Sums of Squares and Cross{-}products: Type I }
\CommentTok{\#\textgreater{} Effect sizes (Z) based on F distributions}
\CommentTok{\#\textgreater{} }
\CommentTok{\#\textgreater{}           Df       SS        MS     Rsq      F      Z}
\CommentTok{\#\textgreater{} species    2 0.019283 0.0096414 0.46942 11.944 4.7202}
\CommentTok{\#\textgreater{} Residuals 27 0.021795 0.0008072 0.53058              }
\CommentTok{\#\textgreater{} Total     29 0.041078                                }
\CommentTok{\#\textgreater{}           Pr(\textgreater{}F)   }
\CommentTok{\#\textgreater{} species    0.001 **}
\CommentTok{\#\textgreater{} Residuals          }
\CommentTok{\#\textgreater{} Total              }
\CommentTok{\#\textgreater{} {-}{-}{-}}
\CommentTok{\#\textgreater{} Signif. codes:  }
\CommentTok{\#\textgreater{} 0 \textquotesingle{}***\textquotesingle{} 0.001 \textquotesingle{}**\textquotesingle{} 0.01 \textquotesingle{}*\textquotesingle{} 0.05 \textquotesingle{}.\textquotesingle{} 0.1 \textquotesingle{} \textquotesingle{} 1}
\CommentTok{\#\textgreater{} }
\CommentTok{\#\textgreater{} Call: procD.lm(f1 = shape.sym \textasciitilde{} species, iter = 999, RRPP = TRUE)}

\NormalTok{manova.sex}\OtherTok{\textless{}{-}}\FunctionTok{procD.lm}\NormalTok{(shape.sym}\SpecialCharTok{\textasciitilde{}}\NormalTok{sex, }\AttributeTok{iter=}\DecValTok{999}\NormalTok{, }\AttributeTok{RRPP=} \ConstantTok{TRUE}\NormalTok{)}
\FunctionTok{summary}\NormalTok{(manova.sex)}
\CommentTok{\#\textgreater{} }
\CommentTok{\#\textgreater{} Analysis of Variance, using Residual Randomization}
\CommentTok{\#\textgreater{} Permutation procedure: Randomization of null model residuals }
\CommentTok{\#\textgreater{} Number of permutations: 1000 }
\CommentTok{\#\textgreater{} Estimation method: Ordinary Least Squares }
\CommentTok{\#\textgreater{} Sums of Squares and Cross{-}products: Type I }
\CommentTok{\#\textgreater{} Effect sizes (Z) based on F distributions}
\CommentTok{\#\textgreater{} }
\CommentTok{\#\textgreater{}           Df       SS        MS     Rsq      F      Z}
\CommentTok{\#\textgreater{} sex        1 0.002892 0.0028915 0.07039 2.1202 1.4855}
\CommentTok{\#\textgreater{} Residuals 28 0.038187 0.0013638 0.92961              }
\CommentTok{\#\textgreater{} Total     29 0.041078                                }
\CommentTok{\#\textgreater{}           Pr(\textgreater{}F)  }
\CommentTok{\#\textgreater{} sex        0.077 .}
\CommentTok{\#\textgreater{} Residuals         }
\CommentTok{\#\textgreater{} Total             }
\CommentTok{\#\textgreater{} {-}{-}{-}}
\CommentTok{\#\textgreater{} Signif. codes:  }
\CommentTok{\#\textgreater{} 0 \textquotesingle{}***\textquotesingle{} 0.001 \textquotesingle{}**\textquotesingle{} 0.01 \textquotesingle{}*\textquotesingle{} 0.05 \textquotesingle{}.\textquotesingle{} 0.1 \textquotesingle{} \textquotesingle{} 1}
\CommentTok{\#\textgreater{} }
\CommentTok{\#\textgreater{} Call: procD.lm(f1 = shape.sym \textasciitilde{} sex, iter = 999, RRPP = TRUE)}

\NormalTok{manova.completa}\OtherTok{\textless{}{-}}\FunctionTok{procD.lm}\NormalTok{(shape.sym}\SpecialCharTok{\textasciitilde{}}\NormalTok{species}\SpecialCharTok{*}\NormalTok{sex, }\AttributeTok{iter=}\DecValTok{999}\NormalTok{, }\AttributeTok{RRPP=} \ConstantTok{TRUE}\NormalTok{)}
\FunctionTok{summary}\NormalTok{(manova.completa)}
\CommentTok{\#\textgreater{} }
\CommentTok{\#\textgreater{} Analysis of Variance, using Residual Randomization}
\CommentTok{\#\textgreater{} Permutation procedure: Randomization of null model residuals }
\CommentTok{\#\textgreater{} Number of permutations: 1000 }
\CommentTok{\#\textgreater{} Estimation method: Ordinary Least Squares }
\CommentTok{\#\textgreater{} Sums of Squares and Cross{-}products: Type I }
\CommentTok{\#\textgreater{} Effect sizes (Z) based on F distributions}
\CommentTok{\#\textgreater{} }
\CommentTok{\#\textgreater{}             Df       SS        MS     Rsq       F      Z}
\CommentTok{\#\textgreater{} species      2 0.019283 0.0096414 0.46942 13.9065 4.9897}
\CommentTok{\#\textgreater{} sex          1 0.002892 0.0028915 0.07039  4.1707 3.5268}
\CommentTok{\#\textgreater{} species:sex  2 0.002264 0.0011322 0.05512  1.6331 1.7627}
\CommentTok{\#\textgreater{} Residuals   24 0.016639 0.0006933 0.40506               }
\CommentTok{\#\textgreater{} Total       29 0.041078                                 }
\CommentTok{\#\textgreater{}             Pr(\textgreater{}F)   }
\CommentTok{\#\textgreater{} species      0.001 **}
\CommentTok{\#\textgreater{} sex          0.001 **}
\CommentTok{\#\textgreater{} species:sex  0.036 * }
\CommentTok{\#\textgreater{} Residuals            }
\CommentTok{\#\textgreater{} Total                }
\CommentTok{\#\textgreater{} {-}{-}{-}}
\CommentTok{\#\textgreater{} Signif. codes:  }
\CommentTok{\#\textgreater{} 0 \textquotesingle{}***\textquotesingle{} 0.001 \textquotesingle{}**\textquotesingle{} 0.01 \textquotesingle{}*\textquotesingle{} 0.05 \textquotesingle{}.\textquotesingle{} 0.1 \textquotesingle{} \textquotesingle{} 1}
\CommentTok{\#\textgreater{} }
\CommentTok{\#\textgreater{} Call: procD.lm(f1 = shape.sym \textasciitilde{} species * sex, iter = 999, RRPP = TRUE)}

\CommentTok{\# MANOVA Wilks\textquotesingle{}s lambda}
\NormalTok{PCA}\OtherTok{\textless{}{-}}\FunctionTok{gm.prcomp}\NormalTok{(shape.sym)}
\NormalTok{manova.w}\OtherTok{\textless{}{-}}\FunctionTok{manova}\NormalTok{(PCA}\SpecialCharTok{$}\NormalTok{x[,}\DecValTok{1}\SpecialCharTok{:}\DecValTok{26}\NormalTok{]}\SpecialCharTok{\textasciitilde{}}\NormalTok{species)}
\FunctionTok{summary}\NormalTok{(manova.w,}\AttributeTok{test=}\StringTok{"Wilks"}\NormalTok{)}
\CommentTok{\#\textgreater{}           Df      Wilks approx F num Df den Df  Pr(\textgreater{}F)  }
\CommentTok{\#\textgreater{} species    2 7.9746e{-}05    8.537     52      4 0.02414 *}
\CommentTok{\#\textgreater{} Residuals 27                                            }
\CommentTok{\#\textgreater{} {-}{-}{-}}
\CommentTok{\#\textgreater{} Signif. codes:  }
\CommentTok{\#\textgreater{} 0 \textquotesingle{}***\textquotesingle{} 0.001 \textquotesingle{}**\textquotesingle{} 0.01 \textquotesingle{}*\textquotesingle{} 0.05 \textquotesingle{}.\textquotesingle{} 0.1 \textquotesingle{} \textquotesingle{} 1}

\NormalTok{manova.w.sex}\OtherTok{\textless{}{-}}\FunctionTok{manova}\NormalTok{(PCA}\SpecialCharTok{$}\NormalTok{x[,}\DecValTok{1}\SpecialCharTok{:}\DecValTok{26}\NormalTok{]}\SpecialCharTok{\textasciitilde{}}\NormalTok{sex)}
\FunctionTok{summary}\NormalTok{(manova.w.sex,}\AttributeTok{test=}\StringTok{"Wilks"}\NormalTok{)}
\CommentTok{\#\textgreater{}           Df    Wilks approx F num Df den Df  Pr(\textgreater{}F)  }
\CommentTok{\#\textgreater{} sex        1 0.021083   5.3575     26      3 0.09549 .}
\CommentTok{\#\textgreater{} Residuals 28                                          }
\CommentTok{\#\textgreater{} {-}{-}{-}}
\CommentTok{\#\textgreater{} Signif. codes:  }
\CommentTok{\#\textgreater{} 0 \textquotesingle{}***\textquotesingle{} 0.001 \textquotesingle{}**\textquotesingle{} 0.01 \textquotesingle{}*\textquotesingle{} 0.05 \textquotesingle{}.\textquotesingle{} 0.1 \textquotesingle{} \textquotesingle{} 1}
\end{Highlighting}
\end{Shaded}

Além da resposta geral, é interessante obter comparações par a par.

\begin{Shaded}
\begin{Highlighting}[]
\CommentTok{\# Pairwise comparisons}
\FunctionTok{require}\NormalTok{(RRPP)}
\NormalTok{manova.pairwise}\OtherTok{\textless{}{-}}\FunctionTok{pairwise}\NormalTok{(manova,}\AttributeTok{groups=}\NormalTok{species)}
\FunctionTok{summary}\NormalTok{(manova.pairwise)}
\CommentTok{\#\textgreater{} }
\CommentTok{\#\textgreater{} Pairwise comparisons}
\CommentTok{\#\textgreater{} }
\CommentTok{\#\textgreater{} Groups: C.flamarioni C.minutus C.torquatus }
\CommentTok{\#\textgreater{} }
\CommentTok{\#\textgreater{} RRPP: 1000 permutations}
\CommentTok{\#\textgreater{} }
\CommentTok{\#\textgreater{} LS means:}
\CommentTok{\#\textgreater{} Vectors hidden (use show.vectors = TRUE to view)}
\CommentTok{\#\textgreater{} }
\CommentTok{\#\textgreater{} Pairwise distances between means, plus statistics}
\CommentTok{\#\textgreater{}                                   d  UCL (95\%)        Z}
\CommentTok{\#\textgreater{} C.flamarioni:C.minutus   0.03831876 0.02502431 3.065210}
\CommentTok{\#\textgreater{} C.flamarioni:C.torquatus 0.05438892 0.02542406 4.093932}
\CommentTok{\#\textgreater{} C.minutus:C.torquatus    0.03685612 0.02512571 2.914480}
\CommentTok{\#\textgreater{}                          Pr \textgreater{} d}
\CommentTok{\#\textgreater{} C.flamarioni:C.minutus    0.003}
\CommentTok{\#\textgreater{} C.flamarioni:C.torquatus  0.001}
\CommentTok{\#\textgreater{} C.minutus:C.torquatus     0.001}
\end{Highlighting}
\end{Shaded}

Manova par a par ``manual''.

\begin{Shaded}
\begin{Highlighting}[]
\CommentTok{\# Manova par a par}
\NormalTok{gdf}\FloatTok{.1}\OtherTok{\textless{}{-}} \FunctionTok{geomorph.data.frame}\NormalTok{(}\AttributeTok{coords=}\NormalTok{shape.sym[,,}\FunctionTok{which}\NormalTok{(species}\SpecialCharTok{==}\StringTok{"C.flamarioni"}\SpecialCharTok{|}\NormalTok{species}\SpecialCharTok{==}\StringTok{"C.minutus"}\NormalTok{)],}\AttributeTok{species=}\NormalTok{species[species}\SpecialCharTok{==}\StringTok{"C.flamarioni"} \SpecialCharTok{|}\NormalTok{ species}\SpecialCharTok{==}\StringTok{"C.minutus"}\NormalTok{])}
\NormalTok{manova.par1}\OtherTok{\textless{}{-}}\FunctionTok{procD.lm}\NormalTok{(coords}\SpecialCharTok{\textasciitilde{}}\NormalTok{species, }\AttributeTok{iter=}\DecValTok{999}\NormalTok{, }\AttributeTok{data =}\NormalTok{ gdf}\FloatTok{.1}\NormalTok{)}
\NormalTok{manova.par1}
\FunctionTok{summary}\NormalTok{(manova.par1)}
\NormalTok{gdf}\FloatTok{.2}\OtherTok{\textless{}{-}}
  \FunctionTok{geomorph.data.frame}\NormalTok{(}\AttributeTok{coords=}\NormalTok{shape.sym[,,}\FunctionTok{which}\NormalTok{(species}\SpecialCharTok{==}\StringTok{"C.flamarioni"}\SpecialCharTok{|}\NormalTok{species}\SpecialCharTok{==}\StringTok{"C.torquatus"}\NormalTok{)],}\AttributeTok{species=}\NormalTok{species[species}\SpecialCharTok{==}\StringTok{"C.flamarioni"} \SpecialCharTok{|}\NormalTok{ species}\SpecialCharTok{==}\StringTok{"C.torquatus"}\NormalTok{])}
\NormalTok{manova.par2}\OtherTok{\textless{}{-}}\FunctionTok{procD.lm}\NormalTok{(coords}\SpecialCharTok{\textasciitilde{}}\NormalTok{species, }\AttributeTok{iter=}\DecValTok{999}\NormalTok{, }\AttributeTok{data =}\NormalTok{ gdf}\FloatTok{.2}\NormalTok{)}
\NormalTok{manova.par2}
\FunctionTok{summary}\NormalTok{(manova.par2)}
\NormalTok{gdf}\FloatTok{.3}\OtherTok{\textless{}{-}}
  \FunctionTok{geomorph.data.frame}\NormalTok{(}\AttributeTok{coords=}\NormalTok{shape.sym[,,}\FunctionTok{which}\NormalTok{(species}\SpecialCharTok{==}\StringTok{"C.minutus"}\SpecialCharTok{|}\NormalTok{species}\SpecialCharTok{==}\StringTok{"C.torquatus"}\NormalTok{)],}\AttributeTok{species=}\NormalTok{species[species}\SpecialCharTok{==}\StringTok{"C.minutus"} \SpecialCharTok{|}\NormalTok{ species}\SpecialCharTok{==}\StringTok{"C.torquatus"}\NormalTok{])}
\NormalTok{manova.par3}\OtherTok{\textless{}{-}}\FunctionTok{procD.lm}\NormalTok{(coords}\SpecialCharTok{\textasciitilde{}}\NormalTok{species, }\AttributeTok{iter=}\DecValTok{999}\NormalTok{, }\AttributeTok{data =}\NormalTok{ gdf}\FloatTok{.3}\NormalTok{)}
\NormalTok{manova.par3}
\FunctionTok{summary}\NormalTok{(manova.par3)}
\end{Highlighting}
\end{Shaded}

\section{Fenograma}\label{fenograma}

Dependendo da pergunta de estudo, fenogramas podem ser úteis para visualizar os padrões de diferenciação de forma em um formato de árvore de distâncias morfológicas. Aqui, vamos criar um fenograma com distâncias de Procrustes e agrupamento de Neighbot-Joining.

\begin{Shaded}
\begin{Highlighting}[]
\NormalTok{obj}\OtherTok{\textless{}{-}}\FunctionTok{summary}\NormalTok{(manova.pairwise)}
\NormalTok{obj}\SpecialCharTok{$}\NormalTok{pairwise.tables}\SpecialCharTok{$}\NormalTok{D}
\NormalTok{means.dist}\OtherTok{\textless{}{-}}\NormalTok{obj}\SpecialCharTok{$}\NormalTok{pairwise.tables}\SpecialCharTok{$}\NormalTok{D}
\FunctionTok{require}\NormalTok{(ape)}
\CommentTok{\#\textgreater{} Carregando pacotes exigidos: ape}
\NormalTok{fen}\OtherTok{\textless{}{-}}\FunctionTok{nj}\NormalTok{(means.dist)}
\FunctionTok{plot}\NormalTok{(fen)}
\end{Highlighting}
\end{Shaded}

\includegraphics{03.8-mgaov_files/figure-latex/unnamed-chunk-5-1.pdf}

\begin{Shaded}
\begin{Highlighting}[]
\FunctionTok{plot}\NormalTok{(fen,}\AttributeTok{type=}\StringTok{"unrooted"}\NormalTok{)}
\end{Highlighting}
\end{Shaded}

\includegraphics{03.8-mgaov_files/figure-latex/unnamed-chunk-5-2.pdf}

\chapter{Alometria}\label{alometria}

De acordo com a escola Gould-Mosimann, alometria é covariação entre tamanho e forma. Esse conceito faz sentido em Morfometria Geométrica (MG), já que há uma distinção formal entre tamanho e forma. Assim, podemos testar se existe uma associação entre forma e tamanho, ou seja, alometria, ou se não existe tal associação, ou seja, isometria. Em MG, isso é tradicionalmente feito através de uma regressão multivariada da forma sobre o tamanho do centróide.

A alometria é classicamente dividida em três níveis: \emph{estática} (dentro de uma amostra homogênea, ex: adultos de uma população), \emph{ontogenética} (ao longo do desenvolvimento, ex: indivíduos de diferentes idades em uma população) e \emph{evolutiva} (entre várias espécies). A forma e a magnitude das relações alométricas dentro e entre grupos oferecem oportunidades interessantes para entender processos biológicos. Em alguns casos, pode-se também estar interessado em correção de tamanho/alometria nos dados, casos estes onde é fundamental investigar se há uma tendência alométrica comum aos dados (ou seja, quando as relações alométricas não diferem entre grupos).

Para textos abrangentes sobre alometria, veja: (\textbf{zelditch2012?}); (\textbf{klingenberg2016?}).

\section{Alometria simples}\label{alometria-simples}

Vamos investigar se a forma e o tamanho do crânio de tuco-tucos estão associados.

\begin{Shaded}
\begin{Highlighting}[]
\FunctionTok{require}\NormalTok{(geomorph)}
\CommentTok{\#\textgreater{} Carregando pacotes exigidos: geomorph}
\CommentTok{\#\textgreater{} Carregando pacotes exigidos: RRPP}
\CommentTok{\#\textgreater{} Carregando pacotes exigidos: rgl}
\CommentTok{\#\textgreater{} Carregando pacotes exigidos: Matrix}
\CommentTok{\# Carregar dados}
\NormalTok{tps}\OtherTok{\textless{}{-}}\FunctionTok{readland.tps}\NormalTok{(}\StringTok{"dadosmg/Lista tuco dig 2.tps"}\NormalTok{,}\AttributeTok{specID =} \StringTok{"ID"}\NormalTok{, }\AttributeTok{readcurves =} \ConstantTok{FALSE}\NormalTok{)}
\FunctionTok{dim}\NormalTok{(tps)}

\CommentTok{\# Gerar matriz com pares de landmarks simétricos}
\NormalTok{pairs.matrix}\OtherTok{\textless{}{-}}\FunctionTok{matrix}\NormalTok{(}\FunctionTok{c}\NormalTok{(}\DecValTok{2}\NormalTok{,}\DecValTok{3}\NormalTok{,}\DecValTok{5}\NormalTok{,}\DecValTok{6}\NormalTok{,}\DecValTok{7}\NormalTok{,}\DecValTok{8}\NormalTok{,}\DecValTok{10}\NormalTok{,}\DecValTok{11}\NormalTok{,}\DecValTok{12}\NormalTok{,}\DecValTok{13}\NormalTok{,}\DecValTok{14}\NormalTok{,}\DecValTok{15}\NormalTok{,}\DecValTok{16}\NormalTok{,}\DecValTok{17}\NormalTok{,}\DecValTok{18}\NormalTok{,}\DecValTok{19}\NormalTok{,}\DecValTok{20}\NormalTok{,}\DecValTok{21}\NormalTok{,}\DecValTok{23}\NormalTok{,}\DecValTok{24}\NormalTok{,}\DecValTok{25}\NormalTok{,}\DecValTok{26}\NormalTok{,}\DecValTok{27}\NormalTok{,}\DecValTok{28}\NormalTok{),}\AttributeTok{nrow=}\DecValTok{12}\NormalTok{,}\AttributeTok{ncol=}\DecValTok{2}\NormalTok{,}\AttributeTok{byrow=}\NormalTok{T)}
\CommentTok{\# Vetor indivíduos}
\NormalTok{ind}\OtherTok{\textless{}{-}}\FunctionTok{c}\NormalTok{(}\DecValTok{1}\SpecialCharTok{:}\FunctionTok{dim}\NormalTok{(tps)[}\DecValTok{3}\NormalTok{])}

\CommentTok{\# GPA {-} Extrair tamanho do centróide}
\NormalTok{gpa}\OtherTok{\textless{}{-}}\FunctionTok{gpagen}\NormalTok{(tps)}
\NormalTok{size}\OtherTok{\textless{}{-}}\NormalTok{gpa}\SpecialCharTok{$}\NormalTok{Csize}

\CommentTok{\# GPA {-} Simetria bilateral }
\NormalTok{b.s}\OtherTok{\textless{}{-}}\FunctionTok{bilat.symmetry}\NormalTok{(tps,}\AttributeTok{ind=}\NormalTok{ind,}\AttributeTok{object.sym=}\ConstantTok{TRUE}\NormalTok{,}\AttributeTok{land.pairs=}\NormalTok{pairs.matrix)}

\CommentTok{\# Matriz de forma}
\NormalTok{shape.sym}\OtherTok{\textless{}{-}}\NormalTok{b.s}\SpecialCharTok{$}\NormalTok{symm.shape}

\CommentTok{\# Carregar classificadores a partir de lista externa}
\NormalTok{plan}\OtherTok{\textless{}{-}}\FunctionTok{read.table}\NormalTok{(}\StringTok{"dadosmg/Planilha Ctenomys.txt"}\NormalTok{,}\AttributeTok{h=}\NormalTok{T)}
\NormalTok{plan}
\NormalTok{species}\OtherTok{\textless{}{-}}\NormalTok{plan[,}\DecValTok{1}\NormalTok{]}
\NormalTok{species}\OtherTok{\textless{}{-}}\FunctionTok{as.factor}\NormalTok{(species)}
\NormalTok{sex}\OtherTok{\textless{}{-}}\NormalTok{plan[,}\DecValTok{2}\NormalTok{]}
\NormalTok{sex}\OtherTok{\textless{}{-}}\FunctionTok{as.factor}\NormalTok{(sex)}
\NormalTok{sp.sex}\OtherTok{\textless{}{-}}\NormalTok{plan[,}\DecValTok{3}\NormalTok{]}
\NormalTok{sp.sex}\OtherTok{\textless{}{-}}\FunctionTok{as.factor}\NormalTok{(sp.sex)}
\end{Highlighting}
\end{Shaded}

Usaremos um modelo linear multivariado para testar a associação entre tamanho e forma. Também podemos visualizar a relação entre tamanho e forma em um espaço bivariado. Pra isso, precisamos resumir a variação de forma em um único vetor, o que pode ser feito de algumas maneiras. A maneira mais interessante é usar \emph{regresion scores} (Drake \& Klingenberg 2008) que retornam uma variável de forma que tem covariação máxima com o tamanho do centróide.

\begin{Shaded}
\begin{Highlighting}[]
\CommentTok{\# Alometria simples}
\NormalTok{alometria}\OtherTok{\textless{}{-}}\FunctionTok{procD.lm}\NormalTok{(shape.sym}\SpecialCharTok{\textasciitilde{}}\FunctionTok{log}\NormalTok{(size),}\AttributeTok{iter=}\DecValTok{999}\NormalTok{)}
\FunctionTok{summary}\NormalTok{(alometria)}
\CommentTok{\#\textgreater{} }
\CommentTok{\#\textgreater{} Analysis of Variance, using Residual Randomization}
\CommentTok{\#\textgreater{} Permutation procedure: Randomization of null model residuals }
\CommentTok{\#\textgreater{} Number of permutations: 1000 }
\CommentTok{\#\textgreater{} Estimation method: Ordinary Least Squares }
\CommentTok{\#\textgreater{} Sums of Squares and Cross{-}products: Type I }
\CommentTok{\#\textgreater{} Effect sizes (Z) based on F distributions}
\CommentTok{\#\textgreater{} }
\CommentTok{\#\textgreater{}           Df       SS        MS     Rsq      F      Z}
\CommentTok{\#\textgreater{} log(size)  1 0.002667 0.0026672 0.06493 1.9443 1.4121}
\CommentTok{\#\textgreater{} Residuals 28 0.038411 0.0013718 0.93507              }
\CommentTok{\#\textgreater{} Total     29 0.041078                                }
\CommentTok{\#\textgreater{}           Pr(\textgreater{}F)  }
\CommentTok{\#\textgreater{} log(size)  0.082 .}
\CommentTok{\#\textgreater{} Residuals         }
\CommentTok{\#\textgreater{} Total             }
\CommentTok{\#\textgreater{} {-}{-}{-}}
\CommentTok{\#\textgreater{} Signif. codes:  }
\CommentTok{\#\textgreater{} 0 \textquotesingle{}***\textquotesingle{} 0.001 \textquotesingle{}**\textquotesingle{} 0.01 \textquotesingle{}*\textquotesingle{} 0.05 \textquotesingle{}.\textquotesingle{} 0.1 \textquotesingle{} \textquotesingle{} 1}
\CommentTok{\#\textgreater{} }
\CommentTok{\#\textgreater{} Call: procD.lm(f1 = shape.sym \textasciitilde{} log(size), iter = 999)}
\FunctionTok{plot}\NormalTok{(alometria,}\AttributeTok{type=}\StringTok{"regression"}\NormalTok{,}\AttributeTok{predictor=}\FunctionTok{log}\NormalTok{(size),}
     \AttributeTok{reg.type =} \StringTok{"RegScore"}\NormalTok{, }\AttributeTok{pch=}\DecValTok{19}\NormalTok{, }\AttributeTok{xlab=}\StringTok{"Log Size"}\NormalTok{)}
\end{Highlighting}
\end{Shaded}

\includegraphics{03.9-mgalo_files/figure-latex/unnamed-chunk-2-1.pdf}

Finalmente, mas não menos importante, vamos visualizar as diferenças de forma entre valores diferentes de tamanho.

\begin{Shaded}
\begin{Highlighting}[]
\NormalTok{plot.alometria}\OtherTok{\textless{}{-}}\FunctionTok{plot}\NormalTok{(alometria,}\AttributeTok{type=}\StringTok{"regression"}\NormalTok{,}\AttributeTok{predictor=}\FunctionTok{log}\NormalTok{(size),}
     \AttributeTok{reg.type =} \StringTok{"RegScore"}\NormalTok{, }\AttributeTok{pch=}\DecValTok{19}\NormalTok{, }\AttributeTok{xlab=}\StringTok{"Log Size"}\NormalTok{)}
\end{Highlighting}
\end{Shaded}

\includegraphics{03.9-mgalo_files/figure-latex/unnamed-chunk-3-1.pdf}

\begin{Shaded}
\begin{Highlighting}[]
\NormalTok{preds}\OtherTok{\textless{}{-}}\FunctionTok{shape.predictor}\NormalTok{(alometria}\SpecialCharTok{$}\NormalTok{GM}\SpecialCharTok{$}\NormalTok{fitted, }\AttributeTok{x =}\NormalTok{ plot.alometria}\SpecialCharTok{$}\NormalTok{RegScore, }
                        \AttributeTok{predmin =} \FunctionTok{min}\NormalTok{(plot.alometria}\SpecialCharTok{$}\NormalTok{RegScore), }
                        \AttributeTok{predmax =} \FunctionTok{max}\NormalTok{(plot.alometria}\SpecialCharTok{$}\NormalTok{RegScore))}
\NormalTok{ref }\OtherTok{\textless{}{-}}\FunctionTok{mshape}\NormalTok{(shape.sym)}
\CommentTok{\# Menores tamanhos}
\FunctionTok{plotRefToTarget}\NormalTok{(ref, preds}\SpecialCharTok{$}\NormalTok{predmin, }\AttributeTok{mag=}\DecValTok{2}\NormalTok{)}
\end{Highlighting}
\end{Shaded}

\includegraphics{03.9-mgalo_files/figure-latex/unnamed-chunk-3-2.pdf}

\begin{Shaded}
\begin{Highlighting}[]
\CommentTok{\# Maiores tamanhos}
\FunctionTok{plotRefToTarget}\NormalTok{(ref, preds}\SpecialCharTok{$}\NormalTok{predmax, }\AttributeTok{mag=}\DecValTok{2}\NormalTok{)}
\end{Highlighting}
\end{Shaded}

\includegraphics{03.9-mgalo_files/figure-latex/unnamed-chunk-3-3.pdf}

\section{Correção de efeitos alométricos}\label{correuxe7uxe3o-de-efeitos-alomuxe9tricos}

A correção de tamanho/alometria em estudos de MG pode ser feita usando os resíduos da regressão multivariada da forma sobre o tamanho do centróide em análises posteriores. Se a amostra for composta por grupos biológicos, é importante primeiro investigar se há uma tendência alométrica comum entre os grupos.

\begin{Shaded}
\begin{Highlighting}[]
\NormalTok{shape.resid}\OtherTok{\textless{}{-}}
  \FunctionTok{arrayspecs}\NormalTok{(alometria}\SpecialCharTok{$}\NormalTok{residuals,}\AttributeTok{p=}\FunctionTok{dim}\NormalTok{(shape.sym)[}\DecValTok{1}\NormalTok{],}\AttributeTok{k=}\FunctionTok{dim}\NormalTok{(shape.sym)[}\DecValTok{2}\NormalTok{]) }\CommentTok{\#size adjusted residuals}
\NormalTok{adj.shape}\OtherTok{\textless{}{-}}\NormalTok{shape.resid}\SpecialCharTok{+}\FunctionTok{array}\NormalTok{(gpa}\SpecialCharTok{$}\NormalTok{consensus, }\FunctionTok{dim}\NormalTok{(shape.resid)) }\CommentTok{\# allometry{-}free shapes{-}{-}somar a média retorna os números para o espaço de forma original e permite a visualização de mudanças de forma}

\CommentTok{\# PCA of allometry{-}free shape space}
\NormalTok{PCA}\OtherTok{\textless{}{-}}\FunctionTok{gm.prcomp}\NormalTok{(adj.shape)}
\FunctionTok{plot}\NormalTok{(PCA,}\AttributeTok{axis1=}\DecValTok{1}\NormalTok{, }\AttributeTok{axis2=}\DecValTok{2}\NormalTok{)}
\end{Highlighting}
\end{Shaded}

\includegraphics{03.9-mgalo_files/figure-latex/unnamed-chunk-4-1.pdf}

\section{Comparação de trajetórias alométricas entre grupos}\label{comparauxe7uxe3o-de-trajetuxf3rias-alomuxe9tricas-entre-grupos}

A análise de alometria considerando grupos (e.g., espécies, habitats) pode revelar se as relações alométricas (e.g., inclinação, ângulo, e comprimento das retas) são diferentes entre os grupos. Em outras palavras, é possível testar a homogeneidade ou heterogeneidade dos \emph{slopes}/inclinações. Se as relações alométricas diferirem, a pergunta a ser respondida é quais fatores causam as diferenças alométricas.

\begin{Shaded}
\begin{Highlighting}[]
\CommentTok{\# Teste de Homogeneidade de Slopes}
\NormalTok{alometria.HOS}\OtherTok{\textless{}{-}}\FunctionTok{procD.lm}\NormalTok{(shape.sym}\SpecialCharTok{\textasciitilde{}}\FunctionTok{log}\NormalTok{(size)}\SpecialCharTok{*}\NormalTok{species,}\AttributeTok{iter=}\DecValTok{999}\NormalTok{)}
\FunctionTok{summary}\NormalTok{(alometria.HOS)}
\CommentTok{\#\textgreater{} }
\CommentTok{\#\textgreater{} Analysis of Variance, using Residual Randomization}
\CommentTok{\#\textgreater{} Permutation procedure: Randomization of null model residuals }
\CommentTok{\#\textgreater{} Number of permutations: 1000 }
\CommentTok{\#\textgreater{} Estimation method: Ordinary Least Squares }
\CommentTok{\#\textgreater{} Sums of Squares and Cross{-}products: Type I }
\CommentTok{\#\textgreater{} Effect sizes (Z) based on F distributions}
\CommentTok{\#\textgreater{} }
\CommentTok{\#\textgreater{}                   Df       SS        MS     Rsq       F}
\CommentTok{\#\textgreater{} log(size)          1 0.002667 0.0026672 0.06493  3.7930}
\CommentTok{\#\textgreater{} species            2 0.019233 0.0096167 0.46821 13.6760}
\CommentTok{\#\textgreater{} log(size):species  2 0.002301 0.0011507 0.05602  1.6364}
\CommentTok{\#\textgreater{} Residuals         24 0.016876 0.0007032 0.41083        }
\CommentTok{\#\textgreater{} Total             29 0.041078                          }
\CommentTok{\#\textgreater{}                        Z Pr(\textgreater{}F)   }
\CommentTok{\#\textgreater{} log(size)         2.4779  0.006 **}
\CommentTok{\#\textgreater{} species           4.6694  0.001 **}
\CommentTok{\#\textgreater{} log(size):species 1.7715  0.040 * }
\CommentTok{\#\textgreater{} Residuals                         }
\CommentTok{\#\textgreater{} Total                             }
\CommentTok{\#\textgreater{} {-}{-}{-}}
\CommentTok{\#\textgreater{} Signif. codes:  }
\CommentTok{\#\textgreater{} 0 \textquotesingle{}***\textquotesingle{} 0.001 \textquotesingle{}**\textquotesingle{} 0.01 \textquotesingle{}*\textquotesingle{} 0.05 \textquotesingle{}.\textquotesingle{} 0.1 \textquotesingle{} \textquotesingle{} 1}
\CommentTok{\#\textgreater{} }
\CommentTok{\#\textgreater{} Call: procD.lm(f1 = shape.sym \textasciitilde{} log(size) * species, iter = 999)}
\CommentTok{\# Comparando com modelo nulo}
\NormalTok{alometria.nulo}\OtherTok{\textless{}{-}}\FunctionTok{procD.lm}\NormalTok{(shape.sym}\SpecialCharTok{\textasciitilde{}}\NormalTok{size}\SpecialCharTok{+}\NormalTok{species,}\AttributeTok{iter=}\DecValTok{999}\NormalTok{)}
\FunctionTok{anova}\NormalTok{(alometria.nulo,alometria.HOS) }\CommentTok{\# R2 do modelo full é a diferença entre modelos}
\CommentTok{\#\textgreater{} }
\CommentTok{\#\textgreater{} Sums of Squares calculations for 1 model: 1000 permutations.}
\CommentTok{\#\textgreater{} }
  \SpecialCharTok{|}                                                        
  \ErrorTok{|}                                                  \ErrorTok{|}   \DecValTok{0}\NormalTok{\%}
  \SpecialCharTok{|}                                                        
  \ErrorTok{|}                                                  \ErrorTok{|}   \DecValTok{1}\NormalTok{\%}
  \SpecialCharTok{|}                                                        
  \ErrorTok{|=}                                                 \ErrorTok{|}   \DecValTok{1}\NormalTok{\%}
  \SpecialCharTok{|}                                                        
  \ErrorTok{|=}                                                 \ErrorTok{|}   \DecValTok{2}\NormalTok{\%}
  \SpecialCharTok{|}                                                        
  \ErrorTok{|=}                                                 \ErrorTok{|}   \DecValTok{3}\NormalTok{\%}
  \SpecialCharTok{|}                                                        
  \ErrorTok{|==}                                                \ErrorTok{|}   \DecValTok{3}\NormalTok{\%}
  \SpecialCharTok{|}                                                        
  \ErrorTok{|==}                                                \ErrorTok{|}   \DecValTok{4}\NormalTok{\%}
  \SpecialCharTok{|}                                                        
  \ErrorTok{|==}                                                \ErrorTok{|}   \DecValTok{5}\NormalTok{\%}
  \SpecialCharTok{|}                                                        
  \ErrorTok{|===}                                               \ErrorTok{|}   \DecValTok{5}\NormalTok{\%}
  \SpecialCharTok{|}                                                        
  \ErrorTok{|===}                                               \ErrorTok{|}   \DecValTok{6}\NormalTok{\%}
  \SpecialCharTok{|}                                                        
  \ErrorTok{|===}                                               \ErrorTok{|}   \DecValTok{7}\NormalTok{\%}
  \SpecialCharTok{|}                                                        
  \ErrorTok{|====}                                              \ErrorTok{|}   \DecValTok{7}\NormalTok{\%}
  \SpecialCharTok{|}                                                        
  \ErrorTok{|====}                                              \ErrorTok{|}   \DecValTok{8}\NormalTok{\%}
  \SpecialCharTok{|}                                                        
  \ErrorTok{|====}                                              \ErrorTok{|}   \DecValTok{9}\NormalTok{\%}
  \SpecialCharTok{|}                                                        
  \ErrorTok{|=====}                                             \ErrorTok{|}   \DecValTok{9}\NormalTok{\%}
  \SpecialCharTok{|}                                                        
  \ErrorTok{|=====}                                             \ErrorTok{|}  \DecValTok{10}\NormalTok{\%}
  \SpecialCharTok{|}                                                        
  \ErrorTok{|=====}                                             \ErrorTok{|}  \DecValTok{11}\NormalTok{\%}
  \SpecialCharTok{|}                                                        
  \ErrorTok{|======}                                            \ErrorTok{|}  \DecValTok{11}\NormalTok{\%}
  \SpecialCharTok{|}                                                        
  \ErrorTok{|======}                                            \ErrorTok{|}  \DecValTok{12}\NormalTok{\%}
  \SpecialCharTok{|}                                                        
  \ErrorTok{|======}                                            \ErrorTok{|}  \DecValTok{13}\NormalTok{\%}
  \SpecialCharTok{|}                                                        
  \ErrorTok{|=======}                                           \ErrorTok{|}  \DecValTok{13}\NormalTok{\%}
  \SpecialCharTok{|}                                                        
  \ErrorTok{|=======}                                           \ErrorTok{|}  \DecValTok{14}\NormalTok{\%}
  \SpecialCharTok{|}                                                        
  \ErrorTok{|=======}                                           \ErrorTok{|}  \DecValTok{15}\NormalTok{\%}
  \SpecialCharTok{|}                                                        
  \ErrorTok{|========}                                          \ErrorTok{|}  \DecValTok{15}\NormalTok{\%}
  \SpecialCharTok{|}                                                        
  \ErrorTok{|========}                                          \ErrorTok{|}  \DecValTok{16}\NormalTok{\%}
  \SpecialCharTok{|}                                                        
  \ErrorTok{|========}                                          \ErrorTok{|}  \DecValTok{17}\NormalTok{\%}
  \SpecialCharTok{|}                                                        
  \ErrorTok{|=========}                                         \ErrorTok{|}  \DecValTok{17}\NormalTok{\%}
  \SpecialCharTok{|}                                                        
  \ErrorTok{|=========}                                         \ErrorTok{|}  \DecValTok{18}\NormalTok{\%}
  \SpecialCharTok{|}                                                        
  \ErrorTok{|=========}                                         \ErrorTok{|}  \DecValTok{19}\NormalTok{\%}
  \SpecialCharTok{|}                                                        
  \ErrorTok{|==========}                                        \ErrorTok{|}  \DecValTok{19}\NormalTok{\%}
  \SpecialCharTok{|}                                                        
  \ErrorTok{|==========}                                        \ErrorTok{|}  \DecValTok{20}\NormalTok{\%}
  \SpecialCharTok{|}                                                        
  \ErrorTok{|==========}                                        \ErrorTok{|}  \DecValTok{21}\NormalTok{\%}
  \SpecialCharTok{|}                                                        
  \ErrorTok{|===========}                                       \ErrorTok{|}  \DecValTok{21}\NormalTok{\%}
  \SpecialCharTok{|}                                                        
  \ErrorTok{|===========}                                       \ErrorTok{|}  \DecValTok{22}\NormalTok{\%}
  \SpecialCharTok{|}                                                        
  \ErrorTok{|===========}                                       \ErrorTok{|}  \DecValTok{23}\NormalTok{\%}
  \SpecialCharTok{|}                                                        
  \ErrorTok{|============}                                      \ErrorTok{|}  \DecValTok{23}\NormalTok{\%}
  \SpecialCharTok{|}                                                        
  \ErrorTok{|============}                                      \ErrorTok{|}  \DecValTok{24}\NormalTok{\%}
  \SpecialCharTok{|}                                                        
  \ErrorTok{|============}                                      \ErrorTok{|}  \DecValTok{25}\NormalTok{\%}
  \SpecialCharTok{|}                                                        
  \ErrorTok{|=============}                                     \ErrorTok{|}  \DecValTok{25}\NormalTok{\%}
  \SpecialCharTok{|}                                                        
  \ErrorTok{|=============}                                     \ErrorTok{|}  \DecValTok{26}\NormalTok{\%}
  \SpecialCharTok{|}                                                        
  \ErrorTok{|=============}                                     \ErrorTok{|}  \DecValTok{27}\NormalTok{\%}
  \SpecialCharTok{|}                                                        
  \ErrorTok{|==============}                                    \ErrorTok{|}  \DecValTok{27}\NormalTok{\%}
  \SpecialCharTok{|}                                                        
  \ErrorTok{|==============}                                    \ErrorTok{|}  \DecValTok{28}\NormalTok{\%}
  \SpecialCharTok{|}                                                        
  \ErrorTok{|==============}                                    \ErrorTok{|}  \DecValTok{29}\NormalTok{\%}
  \SpecialCharTok{|}                                                        
  \ErrorTok{|===============}                                   \ErrorTok{|}  \DecValTok{29}\NormalTok{\%}
  \SpecialCharTok{|}                                                        
  \ErrorTok{|===============}                                   \ErrorTok{|}  \DecValTok{30}\NormalTok{\%}
  \SpecialCharTok{|}                                                        
  \ErrorTok{|===============}                                   \ErrorTok{|}  \DecValTok{31}\NormalTok{\%}
  \SpecialCharTok{|}                                                        
  \ErrorTok{|================}                                  \ErrorTok{|}  \DecValTok{31}\NormalTok{\%}
  \SpecialCharTok{|}                                                        
  \ErrorTok{|================}                                  \ErrorTok{|}  \DecValTok{32}\NormalTok{\%}
  \SpecialCharTok{|}                                                        
  \ErrorTok{|================}                                  \ErrorTok{|}  \DecValTok{33}\NormalTok{\%}
  \SpecialCharTok{|}                                                        
  \ErrorTok{|=================}                                 \ErrorTok{|}  \DecValTok{33}\NormalTok{\%}
  \SpecialCharTok{|}                                                        
  \ErrorTok{|=================}                                 \ErrorTok{|}  \DecValTok{34}\NormalTok{\%}
  \SpecialCharTok{|}                                                        
  \ErrorTok{|=================}                                 \ErrorTok{|}  \DecValTok{35}\NormalTok{\%}
  \SpecialCharTok{|}                                                        
  \ErrorTok{|==================}                                \ErrorTok{|}  \DecValTok{35}\NormalTok{\%}
  \SpecialCharTok{|}                                                        
  \ErrorTok{|==================}                                \ErrorTok{|}  \DecValTok{36}\NormalTok{\%}
  \SpecialCharTok{|}                                                        
  \ErrorTok{|==================}                                \ErrorTok{|}  \DecValTok{37}\NormalTok{\%}
  \SpecialCharTok{|}                                                        
  \ErrorTok{|===================}                               \ErrorTok{|}  \DecValTok{37}\NormalTok{\%}
  \SpecialCharTok{|}                                                        
  \ErrorTok{|===================}                               \ErrorTok{|}  \DecValTok{38}\NormalTok{\%}
  \SpecialCharTok{|}                                                        
  \ErrorTok{|===================}                               \ErrorTok{|}  \DecValTok{39}\NormalTok{\%}
  \SpecialCharTok{|}                                                        
  \ErrorTok{|====================}                              \ErrorTok{|}  \DecValTok{39}\NormalTok{\%}
  \SpecialCharTok{|}                                                        
  \ErrorTok{|====================}                              \ErrorTok{|}  \DecValTok{40}\NormalTok{\%}
  \SpecialCharTok{|}                                                        
  \ErrorTok{|====================}                              \ErrorTok{|}  \DecValTok{41}\NormalTok{\%}
  \SpecialCharTok{|}                                                        
  \ErrorTok{|=====================}                             \ErrorTok{|}  \DecValTok{41}\NormalTok{\%}
  \SpecialCharTok{|}                                                        
  \ErrorTok{|=====================}                             \ErrorTok{|}  \DecValTok{42}\NormalTok{\%}
  \SpecialCharTok{|}                                                        
  \ErrorTok{|=====================}                             \ErrorTok{|}  \DecValTok{43}\NormalTok{\%}
  \SpecialCharTok{|}                                                        
  \ErrorTok{|======================}                            \ErrorTok{|}  \DecValTok{43}\NormalTok{\%}
  \SpecialCharTok{|}                                                        
  \ErrorTok{|======================}                            \ErrorTok{|}  \DecValTok{44}\NormalTok{\%}
  \SpecialCharTok{|}                                                        
  \ErrorTok{|======================}                            \ErrorTok{|}  \DecValTok{45}\NormalTok{\%}
  \SpecialCharTok{|}                                                        
  \ErrorTok{|=======================}                           \ErrorTok{|}  \DecValTok{45}\NormalTok{\%}
  \SpecialCharTok{|}                                                        
  \ErrorTok{|=======================}                           \ErrorTok{|}  \DecValTok{46}\NormalTok{\%}
  \SpecialCharTok{|}                                                        
  \ErrorTok{|=======================}                           \ErrorTok{|}  \DecValTok{47}\NormalTok{\%}
  \SpecialCharTok{|}                                                        
  \ErrorTok{|========================}                          \ErrorTok{|}  \DecValTok{47}\NormalTok{\%}
  \SpecialCharTok{|}                                                        
  \ErrorTok{|========================}                          \ErrorTok{|}  \DecValTok{48}\NormalTok{\%}
  \SpecialCharTok{|}                                                        
  \ErrorTok{|========================}                          \ErrorTok{|}  \DecValTok{49}\NormalTok{\%}
  \SpecialCharTok{|}                                                        
  \ErrorTok{|=========================}                         \ErrorTok{|}  \DecValTok{49}\NormalTok{\%}
  \SpecialCharTok{|}                                                        
  \ErrorTok{|=========================}                         \ErrorTok{|}  \DecValTok{50}\NormalTok{\%}
  \SpecialCharTok{|}                                                        
  \ErrorTok{|=========================}                         \ErrorTok{|}  \DecValTok{51}\NormalTok{\%}
  \SpecialCharTok{|}                                                        
  \ErrorTok{|==========================}                        \ErrorTok{|}  \DecValTok{51}\NormalTok{\%}
  \SpecialCharTok{|}                                                        
  \ErrorTok{|==========================}                        \ErrorTok{|}  \DecValTok{52}\NormalTok{\%}
  \SpecialCharTok{|}                                                        
  \ErrorTok{|==========================}                        \ErrorTok{|}  \DecValTok{53}\NormalTok{\%}
  \SpecialCharTok{|}                                                        
  \ErrorTok{|===========================}                       \ErrorTok{|}  \DecValTok{53}\NormalTok{\%}
  \SpecialCharTok{|}                                                        
  \ErrorTok{|===========================}                       \ErrorTok{|}  \DecValTok{54}\NormalTok{\%}
  \SpecialCharTok{|}                                                        
  \ErrorTok{|===========================}                       \ErrorTok{|}  \DecValTok{55}\NormalTok{\%}
  \SpecialCharTok{|}                                                        
  \ErrorTok{|============================}                      \ErrorTok{|}  \DecValTok{55}\NormalTok{\%}
  \SpecialCharTok{|}                                                        
  \ErrorTok{|============================}                      \ErrorTok{|}  \DecValTok{56}\NormalTok{\%}
  \SpecialCharTok{|}                                                        
  \ErrorTok{|============================}                      \ErrorTok{|}  \DecValTok{57}\NormalTok{\%}
  \SpecialCharTok{|}                                                        
  \ErrorTok{|=============================}                     \ErrorTok{|}  \DecValTok{57}\NormalTok{\%}
  \SpecialCharTok{|}                                                        
  \ErrorTok{|=============================}                     \ErrorTok{|}  \DecValTok{58}\NormalTok{\%}
  \SpecialCharTok{|}                                                        
  \ErrorTok{|=============================}                     \ErrorTok{|}  \DecValTok{59}\NormalTok{\%}
  \SpecialCharTok{|}                                                        
  \ErrorTok{|==============================}                    \ErrorTok{|}  \DecValTok{59}\NormalTok{\%}
  \SpecialCharTok{|}                                                        
  \ErrorTok{|==============================}                    \ErrorTok{|}  \DecValTok{60}\NormalTok{\%}
  \SpecialCharTok{|}                                                        
  \ErrorTok{|==============================}                    \ErrorTok{|}  \DecValTok{61}\NormalTok{\%}
  \SpecialCharTok{|}                                                        
  \ErrorTok{|===============================}                   \ErrorTok{|}  \DecValTok{61}\NormalTok{\%}
  \SpecialCharTok{|}                                                        
  \ErrorTok{|===============================}                   \ErrorTok{|}  \DecValTok{62}\NormalTok{\%}
  \SpecialCharTok{|}                                                        
  \ErrorTok{|===============================}                   \ErrorTok{|}  \DecValTok{63}\NormalTok{\%}
  \SpecialCharTok{|}                                                        
  \ErrorTok{|================================}                  \ErrorTok{|}  \DecValTok{63}\NormalTok{\%}
  \SpecialCharTok{|}                                                        
  \ErrorTok{|================================}                  \ErrorTok{|}  \DecValTok{64}\NormalTok{\%}
  \SpecialCharTok{|}                                                        
  \ErrorTok{|================================}                  \ErrorTok{|}  \DecValTok{65}\NormalTok{\%}
  \SpecialCharTok{|}                                                        
  \ErrorTok{|=================================}                 \ErrorTok{|}  \DecValTok{65}\NormalTok{\%}
  \SpecialCharTok{|}                                                        
  \ErrorTok{|=================================}                 \ErrorTok{|}  \DecValTok{66}\NormalTok{\%}
  \SpecialCharTok{|}                                                        
  \ErrorTok{|=================================}                 \ErrorTok{|}  \DecValTok{67}\NormalTok{\%}
  \SpecialCharTok{|}                                                        
  \ErrorTok{|==================================}                \ErrorTok{|}  \DecValTok{67}\NormalTok{\%}
  \SpecialCharTok{|}                                                        
  \ErrorTok{|==================================}                \ErrorTok{|}  \DecValTok{68}\NormalTok{\%}
  \SpecialCharTok{|}                                                        
  \ErrorTok{|==================================}                \ErrorTok{|}  \DecValTok{69}\NormalTok{\%}
  \SpecialCharTok{|}                                                        
  \ErrorTok{|===================================}               \ErrorTok{|}  \DecValTok{69}\NormalTok{\%}
  \SpecialCharTok{|}                                                        
  \ErrorTok{|===================================}               \ErrorTok{|}  \DecValTok{70}\NormalTok{\%}
  \SpecialCharTok{|}                                                        
  \ErrorTok{|===================================}               \ErrorTok{|}  \DecValTok{71}\NormalTok{\%}
  \SpecialCharTok{|}                                                        
  \ErrorTok{|====================================}              \ErrorTok{|}  \DecValTok{71}\NormalTok{\%}
  \SpecialCharTok{|}                                                        
  \ErrorTok{|====================================}              \ErrorTok{|}  \DecValTok{72}\NormalTok{\%}
  \SpecialCharTok{|}                                                        
  \ErrorTok{|====================================}              \ErrorTok{|}  \DecValTok{73}\NormalTok{\%}
  \SpecialCharTok{|}                                                        
  \ErrorTok{|=====================================}             \ErrorTok{|}  \DecValTok{73}\NormalTok{\%}
  \SpecialCharTok{|}                                                        
  \ErrorTok{|=====================================}             \ErrorTok{|}  \DecValTok{74}\NormalTok{\%}
  \SpecialCharTok{|}                                                        
  \ErrorTok{|=====================================}             \ErrorTok{|}  \DecValTok{75}\NormalTok{\%}
  \SpecialCharTok{|}                                                        
  \ErrorTok{|======================================}            \ErrorTok{|}  \DecValTok{75}\NormalTok{\%}
  \SpecialCharTok{|}                                                        
  \ErrorTok{|======================================}            \ErrorTok{|}  \DecValTok{76}\NormalTok{\%}
  \SpecialCharTok{|}                                                        
  \ErrorTok{|======================================}            \ErrorTok{|}  \DecValTok{77}\NormalTok{\%}
  \SpecialCharTok{|}                                                        
  \ErrorTok{|=======================================}           \ErrorTok{|}  \DecValTok{77}\NormalTok{\%}
  \SpecialCharTok{|}                                                        
  \ErrorTok{|=======================================}           \ErrorTok{|}  \DecValTok{78}\NormalTok{\%}
  \SpecialCharTok{|}                                                        
  \ErrorTok{|=======================================}           \ErrorTok{|}  \DecValTok{79}\NormalTok{\%}
  \SpecialCharTok{|}                                                        
  \ErrorTok{|========================================}          \ErrorTok{|}  \DecValTok{79}\NormalTok{\%}
  \SpecialCharTok{|}                                                        
  \ErrorTok{|========================================}          \ErrorTok{|}  \DecValTok{80}\NormalTok{\%}
  \SpecialCharTok{|}                                                        
  \ErrorTok{|========================================}          \ErrorTok{|}  \DecValTok{81}\NormalTok{\%}
  \SpecialCharTok{|}                                                        
  \ErrorTok{|=========================================}         \ErrorTok{|}  \DecValTok{81}\NormalTok{\%}
  \SpecialCharTok{|}                                                        
  \ErrorTok{|=========================================}         \ErrorTok{|}  \DecValTok{82}\NormalTok{\%}
  \SpecialCharTok{|}                                                        
  \ErrorTok{|=========================================}         \ErrorTok{|}  \DecValTok{83}\NormalTok{\%}
  \SpecialCharTok{|}                                                        
  \ErrorTok{|==========================================}        \ErrorTok{|}  \DecValTok{83}\NormalTok{\%}
  \SpecialCharTok{|}                                                        
  \ErrorTok{|==========================================}        \ErrorTok{|}  \DecValTok{84}\NormalTok{\%}
  \SpecialCharTok{|}                                                        
  \ErrorTok{|==========================================}        \ErrorTok{|}  \DecValTok{85}\NormalTok{\%}
  \SpecialCharTok{|}                                                        
  \ErrorTok{|===========================================}       \ErrorTok{|}  \DecValTok{85}\NormalTok{\%}
  \SpecialCharTok{|}                                                        
  \ErrorTok{|===========================================}       \ErrorTok{|}  \DecValTok{86}\NormalTok{\%}
  \SpecialCharTok{|}                                                        
  \ErrorTok{|===========================================}       \ErrorTok{|}  \DecValTok{87}\NormalTok{\%}
  \SpecialCharTok{|}                                                        
  \ErrorTok{|============================================}      \ErrorTok{|}  \DecValTok{87}\NormalTok{\%}
  \SpecialCharTok{|}                                                        
  \ErrorTok{|============================================}      \ErrorTok{|}  \DecValTok{88}\NormalTok{\%}
  \SpecialCharTok{|}                                                        
  \ErrorTok{|============================================}      \ErrorTok{|}  \DecValTok{89}\NormalTok{\%}
  \SpecialCharTok{|}                                                        
  \ErrorTok{|=============================================}     \ErrorTok{|}  \DecValTok{89}\NormalTok{\%}
  \SpecialCharTok{|}                                                        
  \ErrorTok{|=============================================}     \ErrorTok{|}  \DecValTok{90}\NormalTok{\%}
  \SpecialCharTok{|}                                                        
  \ErrorTok{|=============================================}     \ErrorTok{|}  \DecValTok{91}\NormalTok{\%}
  \SpecialCharTok{|}                                                        
  \ErrorTok{|==============================================}    \ErrorTok{|}  \DecValTok{91}\NormalTok{\%}
  \SpecialCharTok{|}                                                        
  \ErrorTok{|==============================================}    \ErrorTok{|}  \DecValTok{92}\NormalTok{\%}
  \SpecialCharTok{|}                                                        
  \ErrorTok{|==============================================}    \ErrorTok{|}  \DecValTok{93}\NormalTok{\%}
  \SpecialCharTok{|}                                                        
  \ErrorTok{|===============================================}   \ErrorTok{|}  \DecValTok{93}\NormalTok{\%}
  \SpecialCharTok{|}                                                        
  \ErrorTok{|===============================================}   \ErrorTok{|}  \DecValTok{94}\NormalTok{\%}
  \SpecialCharTok{|}                                                        
  \ErrorTok{|===============================================}   \ErrorTok{|}  \DecValTok{95}\NormalTok{\%}
  \SpecialCharTok{|}                                                        
  \ErrorTok{|================================================}  \ErrorTok{|}  \DecValTok{95}\NormalTok{\%}
  \SpecialCharTok{|}                                                        
  \ErrorTok{|================================================}  \ErrorTok{|}  \DecValTok{96}\NormalTok{\%}
  \SpecialCharTok{|}                                                        
  \ErrorTok{|================================================}  \ErrorTok{|}  \DecValTok{97}\NormalTok{\%}
  \SpecialCharTok{|}                                                        
  \ErrorTok{|=================================================} \ErrorTok{|}  \DecValTok{97}\NormalTok{\%}
  \SpecialCharTok{|}                                                        
  \ErrorTok{|=================================================} \ErrorTok{|}  \DecValTok{98}\NormalTok{\%}
  \SpecialCharTok{|}                                                        
  \ErrorTok{|=================================================} \ErrorTok{|}  \DecValTok{99}\NormalTok{\%}
  \SpecialCharTok{|}                                                        
  \ErrorTok{|==================================================|}  \DecValTok{99}\NormalTok{\%}
  \SpecialCharTok{|}                                                        
  \ErrorTok{|==================================================|} \DecValTok{100}\NormalTok{\%}
\CommentTok{\#\textgreater{} }
\CommentTok{\#\textgreater{} Analysis of Variance, using Residual Randomization}
\CommentTok{\#\textgreater{} Permutation procedure: Randomization of null model residuals }
\CommentTok{\#\textgreater{} Number of permutations: 1000 }
\CommentTok{\#\textgreater{} Estimation method: Ordinary Least Squares }
\CommentTok{\#\textgreater{} Effect sizes (Z) based on F distributions}
\CommentTok{\#\textgreater{} }
\CommentTok{\#\textgreater{}                                   ResDf Df      RSS}
\CommentTok{\#\textgreater{} shape.sym \textasciitilde{} size + species (Null)    26  1 0.019228}
\CommentTok{\#\textgreater{} shape.sym \textasciitilde{} log(size) * species      24  2 0.016876}
\CommentTok{\#\textgreater{} Total                                29    0.041078}
\CommentTok{\#\textgreater{}                                          SS        MS}
\CommentTok{\#\textgreater{} shape.sym \textasciitilde{} size + species (Null)                    }
\CommentTok{\#\textgreater{} shape.sym \textasciitilde{} log(size) * species   0.0023522 0.0011761}
\CommentTok{\#\textgreater{} Total                                                }
\CommentTok{\#\textgreater{}                                        Rsq      F      Z}
\CommentTok{\#\textgreater{} shape.sym \textasciitilde{} size + species (Null) 0.000000              }
\CommentTok{\#\textgreater{} shape.sym \textasciitilde{} log(size) * species   0.057263 1.6726 1.8423}
\CommentTok{\#\textgreater{} Total                                                   }
\CommentTok{\#\textgreater{}                                       P Pr(\textgreater{}F)}
\CommentTok{\#\textgreater{} shape.sym \textasciitilde{} size + species (Null)             }
\CommentTok{\#\textgreater{} shape.sym \textasciitilde{} log(size) * species   0.034       }
\CommentTok{\#\textgreater{} Total}
\end{Highlighting}
\end{Shaded}

Visualizando trajetórias alométricas.

\begin{Shaded}
\begin{Highlighting}[]
\CommentTok{\# Definindo cores para grupos}
\NormalTok{col.group}\OtherTok{\textless{}{-}}\FunctionTok{rainbow}\NormalTok{(}\FunctionTok{length}\NormalTok{(}\FunctionTok{levels}\NormalTok{(species))) }\CommentTok{\# criar vetor de cores para grupos}
\FunctionTok{names}\NormalTok{(col.group)}\OtherTok{\textless{}{-}}\FunctionTok{levels}\NormalTok{(species)}
\NormalTok{col.group}\OtherTok{\textless{}{-}}\NormalTok{col.group[}\FunctionTok{match}\NormalTok{(species,}\FunctionTok{names}\NormalTok{(col.group))]}

\CommentTok{\# Plot RegScore}
\NormalTok{\{}\FunctionTok{plotAllometry}\NormalTok{(alometria.HOS,}\AttributeTok{size=}\NormalTok{size,}\AttributeTok{logsz=}\ConstantTok{TRUE}\NormalTok{,}
              \AttributeTok{method=}\StringTok{"RegScore"}\NormalTok{,}\AttributeTok{pch=}\DecValTok{19}\NormalTok{,}\AttributeTok{col=}\NormalTok{col.group)}
\FunctionTok{legend}\NormalTok{(}\FloatTok{7.125}\NormalTok{,}\SpecialCharTok{{-}}\FloatTok{0.007}\NormalTok{,}\AttributeTok{legend=}\FunctionTok{unique}\NormalTok{(species),}\AttributeTok{pch=}\DecValTok{19}\NormalTok{,}\AttributeTok{col=}\FunctionTok{unique}\NormalTok{(col.group))\}}
\end{Highlighting}
\end{Shaded}

\includegraphics{03.9-mgalo_files/figure-latex/unnamed-chunk-6-1.pdf}

\begin{Shaded}
\begin{Highlighting}[]

\CommentTok{\# Plot fitted values}
\NormalTok{\{}\FunctionTok{plotAllometry}\NormalTok{(alometria.HOS,}\AttributeTok{size=}\NormalTok{size,}\AttributeTok{logsz=}\ConstantTok{TRUE}\NormalTok{,}
              \AttributeTok{method=}\StringTok{"PredLine"}\NormalTok{,}\AttributeTok{pch=}\DecValTok{19}\NormalTok{,}\AttributeTok{col=}\NormalTok{col.group)}
\FunctionTok{legend}\NormalTok{(}\FloatTok{7.11}\NormalTok{,}\FloatTok{0.00}\NormalTok{,}\AttributeTok{legend=}\FunctionTok{unique}\NormalTok{(species),}\AttributeTok{pch=}\DecValTok{19}\NormalTok{,}\AttributeTok{col=}\FunctionTok{unique}\NormalTok{(col.group))\}}
\end{Highlighting}
\end{Shaded}

\includegraphics{03.9-mgalo_files/figure-latex/unnamed-chunk-6-2.pdf}

Diferenças de magnitude e orientação das retas.

\begin{Shaded}
\begin{Highlighting}[]
\CommentTok{\# Diferenças de magnitude (comprimento das retas)}
\NormalTok{allo.mancova}\OtherTok{\textless{}{-}}\FunctionTok{pairwise}\NormalTok{(alometria.HOS,}\AttributeTok{groups=}\NormalTok{species)}
\FunctionTok{summary}\NormalTok{(allo.mancova)}
\FunctionTok{names}\NormalTok{(allo.mancova)}
\FunctionTok{summary}\NormalTok{(allo.mancova, }\AttributeTok{confidence =} \FloatTok{0.95}\NormalTok{, }\AttributeTok{test.type =} \StringTok{"DL"}\NormalTok{)}

\CommentTok{\# Diferenças de orientação/direção dos vetores (ângulos, vector correlation)}
\FunctionTok{summary}\NormalTok{(allo.mancova,}\AttributeTok{confidence =} \FloatTok{0.95}\NormalTok{,}\AttributeTok{test.type =} \StringTok{"VC"}\NormalTok{, }\AttributeTok{angle.type =} \StringTok{"deg"}\NormalTok{)}
\end{Highlighting}
\end{Shaded}

\chapter{Análise de Trajetória}\label{anuxe1lise-de-trajetuxf3ria}

Análises de trajetória fenotípica (\emph{Phenotypic Trajectory Analysis} - PTA) descrevem mudanças de forma ao longo de um gradiente, por exemplo, ecológico. Diferente de análises multivariadas padrão (MANOVAs) que permitem avaliar diferenças de forma média, com PTA é possível também investigar diferenças nas trajetórias: localização, comprimento, orientação e forma das trajetórias.
Alguns tipos de perguntas biológicas são adequadamente respondidas com análises de trajetória.
No exemplo, vamos testar a hipótese de deslocamento de caracteres causada por competição assumida. Duas espécies de tuco-tucos vivem em regiões de alopatria (só uma das espécies é encontrada) e regiões de simpatria (as duas são encontradas). Vamos investigar se a morfologia de ambas na região de simpatria evoluiu para se tornar mais diferente uma da outra do que nas regiões de alopatria.

\begin{Shaded}
\begin{Highlighting}[]
\CommentTok{\# Carregar arquivo .tps}
\FunctionTok{require}\NormalTok{(geomorph)}
\CommentTok{\#\textgreater{} Carregando pacotes exigidos: geomorph}
\CommentTok{\#\textgreater{} Carregando pacotes exigidos: RRPP}
\CommentTok{\#\textgreater{} Carregando pacotes exigidos: rgl}
\CommentTok{\#\textgreater{} Carregando pacotes exigidos: Matrix}
\NormalTok{tps}\OtherTok{\textless{}{-}}\FunctionTok{readland.tps}\NormalTok{(}\StringTok{"dadosmg/tucodeslocamento.tps"}\NormalTok{,}\AttributeTok{specID =} \StringTok{"ID"}\NormalTok{, }\AttributeTok{readcurves =} \ConstantTok{FALSE}\NormalTok{)}
\FunctionTok{dim}\NormalTok{(tps)}
\CommentTok{\# GPA}
\CommentTok{\# gerar matriz com os pares de landmarks simétricos}
\NormalTok{pairs.matrix}\OtherTok{\textless{}{-}}
  \FunctionTok{matrix}\NormalTok{(}\FunctionTok{c}\NormalTok{(}\DecValTok{2}\NormalTok{,}\DecValTok{3}\NormalTok{,}\DecValTok{4}\NormalTok{,}\DecValTok{5}\NormalTok{,}\DecValTok{6}\NormalTok{,}\DecValTok{7}\NormalTok{,}\DecValTok{8}\NormalTok{,}\DecValTok{9}\NormalTok{,}\DecValTok{10}\NormalTok{,}\DecValTok{11}\NormalTok{,}\DecValTok{12}\NormalTok{,}\DecValTok{13}\NormalTok{,}\DecValTok{15}\NormalTok{,}\DecValTok{16}\NormalTok{,}\DecValTok{17}\NormalTok{,}\DecValTok{18}\NormalTok{,}\DecValTok{19}\NormalTok{,}\DecValTok{20}\NormalTok{,}\DecValTok{21}\NormalTok{,}\DecValTok{22}\NormalTok{,}\DecValTok{23}\NormalTok{,}\DecValTok{24}\NormalTok{,}\DecValTok{25}\NormalTok{,}\DecValTok{26}\NormalTok{,}\DecValTok{29}\NormalTok{,}\DecValTok{30}\NormalTok{),}\AttributeTok{nrow=}\DecValTok{13}\NormalTok{,}\AttributeTok{ncol=}\DecValTok{2}\NormalTok{,}\AttributeTok{byrow=}\NormalTok{T)}
\NormalTok{pairs.matrix}
\NormalTok{ind}\OtherTok{\textless{}{-}}\FunctionTok{c}\NormalTok{(}\DecValTok{1}\SpecialCharTok{:}\FunctionTok{dim}\NormalTok{(tps)[}\DecValTok{3}\NormalTok{]) }\CommentTok{\# vetor indivíduos}
\CommentTok{\# GPA {-} Simetria bilateral}
\NormalTok{gpa.object}\OtherTok{\textless{}{-}}\FunctionTok{gpagen}\NormalTok{(tps)}
\NormalTok{size}\OtherTok{\textless{}{-}}\NormalTok{gpa.object}\SpecialCharTok{$}\NormalTok{Csize}
\NormalTok{b.s}\OtherTok{\textless{}{-}}\FunctionTok{bilat.symmetry}\NormalTok{(tps,}\AttributeTok{ind=}\NormalTok{ind,}\AttributeTok{object.sym=}\ConstantTok{TRUE}\NormalTok{,}\AttributeTok{land.pairs=}\NormalTok{pairs.matrix)}
\NormalTok{shape.sym}\OtherTok{\textless{}{-}}\NormalTok{b.s}\SpecialCharTok{$}\NormalTok{symm.shape }\CommentTok{\# componente simétrico da forma}
\FunctionTok{plotAllSpecimens}\NormalTok{(shape.sym)}
\end{Highlighting}
\end{Shaded}

\includegraphics{03.10-mgtraj_files/figure-latex/unnamed-chunk-1-1.pdf}

\begin{Shaded}
\begin{Highlighting}[]

\CommentTok{\# Carregar classificadores a partir de lista externa}
\NormalTok{plan}\OtherTok{\textless{}{-}}\FunctionTok{read.table}\NormalTok{(}\StringTok{"dadosmg/tucodeslocamento.txt"}\NormalTok{,}\AttributeTok{h=}\NormalTok{T)}
\FunctionTok{attach}\NormalTok{(plan)}
\NormalTok{Species}
\NormalTok{Local}

\CommentTok{\# Deslocamento de Caracteres}
\CommentTok{\# Análises de Trajetória}
\NormalTok{modelo.traje}\OtherTok{\textless{}{-}}\FunctionTok{procD.lm}\NormalTok{(shape.sym}\SpecialCharTok{\textasciitilde{}}\NormalTok{Species}\SpecialCharTok{*}\NormalTok{Local)}
\NormalTok{trajetoria}\OtherTok{\textless{}{-}}\FunctionTok{trajectory.analysis}\NormalTok{(modelo.traje,}\AttributeTok{groups =}\NormalTok{ Species,}\AttributeTok{traj.pts =}\NormalTok{ Local)}

\CommentTok{\# Diferenças de magnitude (comprimento das trajetórias)}
\FunctionTok{summary}\NormalTok{(trajetoria,}\AttributeTok{attribute=}\StringTok{"MD"}\NormalTok{)}

\CommentTok{\# Diferenças de ângulo (correlação entre trajetórias)}
\FunctionTok{summary}\NormalTok{(trajetoria, }\AttributeTok{attribute=}\StringTok{"TC"}\NormalTok{, }\AttributeTok{angle.type =} \StringTok{"deg"}\NormalTok{)}

\CommentTok{\# Visualização}
\CommentTok{\# Plot {-} Símbolos=Espécies; Cores=Local (preto=1, vermelho=2)}
\NormalTok{\{TP}\OtherTok{\textless{}{-}}\FunctionTok{plot}\NormalTok{(trajetoria,}\AttributeTok{pch=}\FunctionTok{as.numeric}\NormalTok{(}\FunctionTok{as.factor}\NormalTok{(Species))}\SpecialCharTok{+}\DecValTok{20}\NormalTok{,}\AttributeTok{bg=}\FunctionTok{as.numeric}\NormalTok{(}\FunctionTok{as.factor}\NormalTok{(Local)))}
\FunctionTok{add.trajectories}\NormalTok{(TP, }\AttributeTok{traj.pch =} \FunctionTok{c}\NormalTok{(}\DecValTok{21}\NormalTok{, }\DecValTok{22}\NormalTok{),}
                 \AttributeTok{start.bg =} \DecValTok{1}\NormalTok{, }\AttributeTok{end.bg =} \DecValTok{2}\NormalTok{)}
\FunctionTok{legend}\NormalTok{(}\StringTok{"topright"}\NormalTok{, }\FunctionTok{levels}\NormalTok{(}\FunctionTok{as.factor}\NormalTok{(Species)), }\AttributeTok{pch =} \FunctionTok{c}\NormalTok{(}\DecValTok{21}\NormalTok{, }\DecValTok{22}\NormalTok{))}
\FunctionTok{legend}\NormalTok{(}\StringTok{"bottomleft"}\NormalTok{, }\FunctionTok{levels}\NormalTok{(}\FunctionTok{as.factor}\NormalTok{(Local)), }\AttributeTok{pch =} \DecValTok{19}\NormalTok{,}
       \AttributeTok{col =} \FunctionTok{c}\NormalTok{(}\StringTok{"black"}\NormalTok{,}\StringTok{"red"}\NormalTok{))\}}
\end{Highlighting}
\end{Shaded}

\includegraphics{03.10-mgtraj_files/figure-latex/unnamed-chunk-1-2.pdf}

\chapter{Integração morfológica e modularidade}\label{integrauxe7uxe3o-morfoluxf3gica-e-modularidade}

Integração morfológica e modularidade são conceitos próximos, mas que dão ênfase para aspectos diferentes da morfologia: associação e dissociação entre partes. A arquitetura do mapa genótipo-fenótipo dá origem aos padrões e magnitudes de integração morfológica. Desenvolvimento e função compartilhadas são também aspectos fundamentais para entendermos as origens da integração morfológica e da modularidade.

Existem vários métodos para acessar a estrutura de integração morfológica em dados de morfometria. Alguns testam hipóteses \emph{a priori} de modularidade, revisados em (\textbf{zelditch2012?}). No exemplo, vamos usar um método que compara a variância entre módulos com a variância dentro de módulos.

Para textos abrangentes sobre integração morfológica e modularidade no contexto de morfometria geométrica, veja: (\textbf{zelditch2012?}); (\textbf{klingenberg2014?}).

Vamos investigar integração morfológica e modularidade sobre uma hipótese de módulos crânianos usando landmarks em crânios de roedores sigmodontíneos.

\begin{Shaded}
\begin{Highlighting}[]
\FunctionTok{require}\NormalTok{(geomorph)}
\CommentTok{\#\textgreater{} Carregando pacotes exigidos: geomorph}
\CommentTok{\#\textgreater{} Carregando pacotes exigidos: RRPP}
\CommentTok{\#\textgreater{} Carregando pacotes exigidos: rgl}
\CommentTok{\#\textgreater{} Carregando pacotes exigidos: Matrix}
\CommentTok{\# Carregar arquivo tps com a vista ventral do crânio }
\NormalTok{tps.v}\OtherTok{\textless{}{-}}\FunctionTok{readland.tps}\NormalTok{(}\StringTok{"dadosmg/ventral.dig\_curso.pls.tps"}\NormalTok{,}\AttributeTok{specID =} \StringTok{"ID"}\NormalTok{, }\AttributeTok{readcurves =} \ConstantTok{FALSE}\NormalTok{)}

\CommentTok{\# Gerar matriz com os pares de landmarks simétricos}
\NormalTok{pairs.matrix}\OtherTok{\textless{}{-}}\FunctionTok{matrix}\NormalTok{(}\FunctionTok{c}\NormalTok{(}\DecValTok{2}\NormalTok{,}\DecValTok{3}\NormalTok{,}\DecValTok{4}\NormalTok{,}\DecValTok{5}\NormalTok{,}\DecValTok{6}\NormalTok{,}\DecValTok{7}\NormalTok{,}\DecValTok{8}\NormalTok{,}\DecValTok{9}\NormalTok{,}\DecValTok{10}\NormalTok{,}\DecValTok{11}\NormalTok{,}\DecValTok{12}\NormalTok{,}\DecValTok{13}\NormalTok{,}\DecValTok{14}\NormalTok{,}\DecValTok{15}\NormalTok{,}\DecValTok{16}\NormalTok{,}\DecValTok{17}\NormalTok{,}\DecValTok{18}\NormalTok{,}\DecValTok{24}\NormalTok{,}\DecValTok{22}\NormalTok{,}\DecValTok{28}\NormalTok{,}\DecValTok{23}\NormalTok{,}\DecValTok{29}\NormalTok{,}\DecValTok{19}\NormalTok{,}\DecValTok{25}\NormalTok{,}\DecValTok{20}\NormalTok{,}\DecValTok{26}\NormalTok{,}\DecValTok{21}\NormalTok{,}\DecValTok{27}\NormalTok{,}\DecValTok{30}\NormalTok{,}\DecValTok{31}\NormalTok{,}\DecValTok{33}\NormalTok{,}\DecValTok{34}\NormalTok{,}\DecValTok{35}\NormalTok{,}\DecValTok{36}\NormalTok{,}\DecValTok{37}\NormalTok{,}\DecValTok{38}\NormalTok{,}\DecValTok{39}\NormalTok{,}\DecValTok{40}\NormalTok{,}\DecValTok{41}\NormalTok{,}\DecValTok{42}\NormalTok{,}\DecValTok{43}\NormalTok{,}\DecValTok{44}\NormalTok{,}\DecValTok{45}\NormalTok{,}\DecValTok{46}\NormalTok{,}\DecValTok{47}\NormalTok{,}\DecValTok{48}\NormalTok{,}\DecValTok{49}\NormalTok{,}\DecValTok{50}\NormalTok{,}\DecValTok{51}\NormalTok{,}\DecValTok{52}\NormalTok{,}\DecValTok{55}\NormalTok{,}\DecValTok{56}\NormalTok{),}\AttributeTok{nrow=}\DecValTok{26}\NormalTok{,}\AttributeTok{ncol=}\DecValTok{2}\NormalTok{,}\AttributeTok{byrow=}\NormalTok{T)}
\NormalTok{ind.v}\OtherTok{\textless{}{-}}\FunctionTok{c}\NormalTok{(}\DecValTok{1}\SpecialCharTok{:}\FunctionTok{dim}\NormalTok{(tps.v)[}\DecValTok{3}\NormalTok{]) }\CommentTok{\# vetor indivíduos}
\CommentTok{\# GPA {-} Simetria bilateral {-} Extrair componente simetrico da forma}
\NormalTok{b.s}\OtherTok{\textless{}{-}}\FunctionTok{bilat.symmetry}\NormalTok{(tps.v,}\AttributeTok{ind=}\NormalTok{ind.v,}\AttributeTok{object.sym=}\ConstantTok{TRUE}\NormalTok{,}\AttributeTok{land.pairs=}\NormalTok{pairs.matrix)}
\NormalTok{shape.v}\OtherTok{\textless{}{-}}\NormalTok{b.s}\SpecialCharTok{$}\NormalTok{symm.shape }\CommentTok{\# componente sim?trico da forma}
\NormalTok{ref.ventral}\OtherTok{\textless{}{-}}\FunctionTok{mshape}\NormalTok{(shape.v)}
\end{Highlighting}
\end{Shaded}

Primeiro, precisamos definir a hipótese a ser testada. A função `define.modules' permite desenharmos a hipótese de maneira interativa. Selecione os landmarks de um módulo, pressione esc, selecione os landmarks do segundo módulo, etc.

\begin{Shaded}
\begin{Highlighting}[]
\NormalTok{partition}\OtherTok{\textless{}{-}}\FunctionTok{define.modules}\NormalTok{(ref.ventral,}\DecValTok{2}\NormalTok{)}
\end{Highlighting}
\end{Shaded}

Testando a hipótese nula de ausência de modularidade.

\begin{Shaded}
\begin{Highlighting}[]
\NormalTok{partition}\OtherTok{\textless{}{-}}\FunctionTok{c}\NormalTok{(}\FunctionTok{rep}\NormalTok{(}\DecValTok{1}\NormalTok{,}\DecValTok{36}\NormalTok{),}\FunctionTok{rep}\NormalTok{(}\DecValTok{2}\NormalTok{,}\DecValTok{20}\NormalTok{))}
\NormalTok{mod}\OtherTok{\textless{}{-}}\FunctionTok{modularity.test}\NormalTok{(shape.v,partition,}\AttributeTok{iter=}\DecValTok{999}\NormalTok{) }
\NormalTok{mod}
\FunctionTok{plot}\NormalTok{(mod)}
\end{Highlighting}
\end{Shaded}

Testando a hipótese nula de ausência de integração morfológica.

\begin{Shaded}
\begin{Highlighting}[]
\NormalTok{int}\OtherTok{\textless{}{-}}\FunctionTok{integration.test}\NormalTok{(shape.v,}\AttributeTok{partition.gp=}\NormalTok{partition,}\AttributeTok{iter=}\DecValTok{999}\NormalTok{) }
\NormalTok{int}
\FunctionTok{plot}\NormalTok{(int)}
\end{Highlighting}
\end{Shaded}

Função interativa para visualização da forma dos módulos.

\begin{Shaded}
\begin{Highlighting}[]
\FunctionTok{plot}\NormalTok{(int)}
\NormalTok{int.plot}\OtherTok{\textless{}{-}}\FunctionTok{plot}\NormalTok{(int,}\AttributeTok{shape=}\ConstantTok{TRUE}\NormalTok{)}
\CommentTok{\#picknplot.shape(int.plot)}
\end{Highlighting}
\end{Shaded}

\chapter{Disparidade e variação morfológica}\label{disparidade-e-variauxe7uxe3o-morfoluxf3gica}

Disparidade e variação são conceitos próximos, ambos relacionados à ideia de ``variedade'' ou ``diversidade'' morfológica; as mesmas métricas podem ser utilizadas para ambos. No entanto, uma distinção conceitual é aplicada: (i) Disparidade tem relação com a variedade de um grupo de espécies, e é o resultado do processo evolutivo de longo prazo; (ii) Variação tem relação com a variedade de indivíduos dentro de uma única população.

Vamos investigar a disparidade total e dentro/entre espécies de roedores sigmodontíneos.

\begin{Shaded}
\begin{Highlighting}[]
\FunctionTok{require}\NormalTok{(geomorph)}
\CommentTok{\#\textgreater{} Carregando pacotes exigidos: geomorph}
\CommentTok{\#\textgreater{} Carregando pacotes exigidos: RRPP}
\CommentTok{\#\textgreater{} Carregando pacotes exigidos: rgl}
\CommentTok{\#\textgreater{} Carregando pacotes exigidos: Matrix}
\CommentTok{\# Carregar arquivo tps com a vista ventral do crânio }
\NormalTok{tps.v}\OtherTok{\textless{}{-}}\FunctionTok{readland.tps}\NormalTok{(}\StringTok{"dadosmg/ventral.dig\_curso.pls.tps"}\NormalTok{,}\AttributeTok{specID =} \StringTok{"ID"}\NormalTok{, }\AttributeTok{readcurves =} \ConstantTok{FALSE}\NormalTok{)}

\CommentTok{\# Gerar matriz com os pares de landmarks simétricos}
\NormalTok{pairs.matrix}\OtherTok{\textless{}{-}}\FunctionTok{matrix}\NormalTok{(}\FunctionTok{c}\NormalTok{(}\DecValTok{2}\NormalTok{,}\DecValTok{3}\NormalTok{,}\DecValTok{4}\NormalTok{,}\DecValTok{5}\NormalTok{,}\DecValTok{6}\NormalTok{,}\DecValTok{7}\NormalTok{,}\DecValTok{8}\NormalTok{,}\DecValTok{9}\NormalTok{,}\DecValTok{10}\NormalTok{,}\DecValTok{11}\NormalTok{,}\DecValTok{12}\NormalTok{,}\DecValTok{13}\NormalTok{,}\DecValTok{14}\NormalTok{,}\DecValTok{15}\NormalTok{,}\DecValTok{16}\NormalTok{,}\DecValTok{17}\NormalTok{,}\DecValTok{18}\NormalTok{,}\DecValTok{24}\NormalTok{,}\DecValTok{22}\NormalTok{,}\DecValTok{28}\NormalTok{,}\DecValTok{23}\NormalTok{,}\DecValTok{29}\NormalTok{,}\DecValTok{19}\NormalTok{,}\DecValTok{25}\NormalTok{,}\DecValTok{20}\NormalTok{,}\DecValTok{26}\NormalTok{,}\DecValTok{21}\NormalTok{,}\DecValTok{27}\NormalTok{,}\DecValTok{30}\NormalTok{,}\DecValTok{31}\NormalTok{,}\DecValTok{33}\NormalTok{,}\DecValTok{34}\NormalTok{,}\DecValTok{35}\NormalTok{,}\DecValTok{36}\NormalTok{,}\DecValTok{37}\NormalTok{,}\DecValTok{38}\NormalTok{,}\DecValTok{39}\NormalTok{,}\DecValTok{40}\NormalTok{,}\DecValTok{41}\NormalTok{,}\DecValTok{42}\NormalTok{,}\DecValTok{43}\NormalTok{,}\DecValTok{44}\NormalTok{,}\DecValTok{45}\NormalTok{,}\DecValTok{46}\NormalTok{,}\DecValTok{47}\NormalTok{,}\DecValTok{48}\NormalTok{,}\DecValTok{49}\NormalTok{,}\DecValTok{50}\NormalTok{,}\DecValTok{51}\NormalTok{,}\DecValTok{52}\NormalTok{,}\DecValTok{55}\NormalTok{,}\DecValTok{56}\NormalTok{),}\AttributeTok{nrow=}\DecValTok{26}\NormalTok{,}\AttributeTok{ncol=}\DecValTok{2}\NormalTok{,}\AttributeTok{byrow=}\NormalTok{T)}
\NormalTok{ind.v}\OtherTok{\textless{}{-}}\FunctionTok{c}\NormalTok{(}\DecValTok{1}\SpecialCharTok{:}\FunctionTok{dim}\NormalTok{(tps.v)[}\DecValTok{3}\NormalTok{]) }\CommentTok{\# vetor indivíduos}
\CommentTok{\# GPA {-} Simetria bilateral {-} Extrair componente simetrico da forma}
\NormalTok{b.s}\OtherTok{\textless{}{-}}\FunctionTok{bilat.symmetry}\NormalTok{(tps.v,}\AttributeTok{ind=}\NormalTok{ind.v,}\AttributeTok{object.sym=}\ConstantTok{TRUE}\NormalTok{,}\AttributeTok{land.pairs=}\NormalTok{pairs.matrix)}
\NormalTok{shape.v}\OtherTok{\textless{}{-}}\NormalTok{b.s}\SpecialCharTok{$}\NormalTok{symm.shape }\CommentTok{\# componente sim?trico da forma}
\NormalTok{ref.ventral}\OtherTok{\textless{}{-}}\FunctionTok{mshape}\NormalTok{(shape.v)}

\CommentTok{\# Carregar classificadores a partir de lista externa}
\NormalTok{plan}\OtherTok{\textless{}{-}}\FunctionTok{read.table}\NormalTok{(}\StringTok{"dadosmg/Planilha\_PLS.txt"}\NormalTok{,}\AttributeTok{h=}\NormalTok{T)}
\NormalTok{plan}
\NormalTok{species}\OtherTok{\textless{}{-}}\FunctionTok{as.factor}\NormalTok{(plan[,}\DecValTok{2}\NormalTok{])}
\NormalTok{species}
\end{Highlighting}
\end{Shaded}

Disparidade morfológica.

\begin{Shaded}
\begin{Highlighting}[]
\CommentTok{\# Disparidade Morfológica}
\NormalTok{div.total}\OtherTok{\textless{}{-}}\FunctionTok{morphol.disparity}\NormalTok{(shape.v}\SpecialCharTok{\textasciitilde{}}\DecValTok{1}\NormalTok{,}\AttributeTok{groups=}\ConstantTok{NULL}\NormalTok{,}\AttributeTok{iter=}\DecValTok{999}\NormalTok{)}
\CommentTok{\#\textgreater{} }
\CommentTok{\#\textgreater{} Please be aware that printing progress slows down the analysis (perhaps slightly).}
\CommentTok{\#\textgreater{} }
\CommentTok{\#\textgreater{} Preliminary Model Fit...}
\CommentTok{\#\textgreater{} }
\CommentTok{\#\textgreater{} Sums of Squares calculations: 3 permutations.}
\CommentTok{\#\textgreater{} }
  \SpecialCharTok{|}                                                        
  \ErrorTok{|}                                                  \ErrorTok{|}   \DecValTok{0}\NormalTok{\%}
  \SpecialCharTok{|}                                                        
  \ErrorTok{|=================}                                 \ErrorTok{|}  \DecValTok{33}\NormalTok{\%}
  \SpecialCharTok{|}                                                        
  \ErrorTok{|=================================}                 \ErrorTok{|}  \DecValTok{67}\NormalTok{\%}
  \SpecialCharTok{|}                                                        
  \ErrorTok{|==================================================|} \DecValTok{100}\NormalTok{\%}
\CommentTok{\#\textgreater{} }
\CommentTok{\#\textgreater{} No terms for ANOVA; only RSS calculated in each permutation}
\CommentTok{\#\textgreater{} No factor in formula or model terms from which to define groups.}
\CommentTok{\#\textgreater{} Procrustes variance:}
\CommentTok{\#\textgreater{} 0.002195294}
\NormalTok{div.total}
\CommentTok{\#\textgreater{} [1] 0.002195294}
\NormalTok{div.sp}\OtherTok{\textless{}{-}}\FunctionTok{morphol.disparity}\NormalTok{(shape.v}\SpecialCharTok{\textasciitilde{}}\NormalTok{species,}\AttributeTok{groups=}\SpecialCharTok{\textasciitilde{}}\NormalTok{species,}\AttributeTok{iter=}\DecValTok{999}\NormalTok{)}
\CommentTok{\#\textgreater{} }
\CommentTok{\#\textgreater{} Please be aware that printing progress slows down the analysis (perhaps slightly).}
\CommentTok{\#\textgreater{} }
\CommentTok{\#\textgreater{} Preliminary Model Fit...}
\CommentTok{\#\textgreater{} }
\CommentTok{\#\textgreater{} Sums of Squares calculations: 3 permutations.}
\CommentTok{\#\textgreater{} }
  \SpecialCharTok{|}                                                        
  \ErrorTok{|}                                                  \ErrorTok{|}   \DecValTok{0}\NormalTok{\%}
  \SpecialCharTok{|}                                                        
  \ErrorTok{|=================}                                 \ErrorTok{|}  \DecValTok{33}\NormalTok{\%}
  \SpecialCharTok{|}                                                        
  \ErrorTok{|=================================}                 \ErrorTok{|}  \DecValTok{67}\NormalTok{\%}
  \SpecialCharTok{|}                                                        
  \ErrorTok{|==================================================|} \DecValTok{100}\NormalTok{\%}
\CommentTok{\#\textgreater{} }
\CommentTok{\#\textgreater{} }
\CommentTok{\#\textgreater{} Performimg pairwise comparisons of disparity}
\CommentTok{\#\textgreater{} }
  \SpecialCharTok{|}                                                        
  \ErrorTok{|}                                                  \ErrorTok{|}   \DecValTok{0}\NormalTok{\%}
  \SpecialCharTok{|}                                                        
  \ErrorTok{|}                                                  \ErrorTok{|}   \DecValTok{1}\NormalTok{\%}
  \SpecialCharTok{|}                                                        
  \ErrorTok{|=}                                                 \ErrorTok{|}   \DecValTok{1}\NormalTok{\%}
  \SpecialCharTok{|}                                                        
  \ErrorTok{|=}                                                 \ErrorTok{|}   \DecValTok{2}\NormalTok{\%}
  \SpecialCharTok{|}                                                        
  \ErrorTok{|=}                                                 \ErrorTok{|}   \DecValTok{3}\NormalTok{\%}
  \SpecialCharTok{|}                                                        
  \ErrorTok{|==}                                                \ErrorTok{|}   \DecValTok{3}\NormalTok{\%}
  \SpecialCharTok{|}                                                        
  \ErrorTok{|==}                                                \ErrorTok{|}   \DecValTok{4}\NormalTok{\%}
  \SpecialCharTok{|}                                                        
  \ErrorTok{|==}                                                \ErrorTok{|}   \DecValTok{5}\NormalTok{\%}
  \SpecialCharTok{|}                                                        
  \ErrorTok{|===}                                               \ErrorTok{|}   \DecValTok{5}\NormalTok{\%}
  \SpecialCharTok{|}                                                        
  \ErrorTok{|===}                                               \ErrorTok{|}   \DecValTok{6}\NormalTok{\%}
  \SpecialCharTok{|}                                                        
  \ErrorTok{|===}                                               \ErrorTok{|}   \DecValTok{7}\NormalTok{\%}
  \SpecialCharTok{|}                                                        
  \ErrorTok{|====}                                              \ErrorTok{|}   \DecValTok{7}\NormalTok{\%}
  \SpecialCharTok{|}                                                        
  \ErrorTok{|====}                                              \ErrorTok{|}   \DecValTok{8}\NormalTok{\%}
  \SpecialCharTok{|}                                                        
  \ErrorTok{|====}                                              \ErrorTok{|}   \DecValTok{9}\NormalTok{\%}
  \SpecialCharTok{|}                                                        
  \ErrorTok{|=====}                                             \ErrorTok{|}   \DecValTok{9}\NormalTok{\%}
  \SpecialCharTok{|}                                                        
  \ErrorTok{|=====}                                             \ErrorTok{|}  \DecValTok{10}\NormalTok{\%}
  \SpecialCharTok{|}                                                        
  \ErrorTok{|=====}                                             \ErrorTok{|}  \DecValTok{11}\NormalTok{\%}
  \SpecialCharTok{|}                                                        
  \ErrorTok{|======}                                            \ErrorTok{|}  \DecValTok{11}\NormalTok{\%}
  \SpecialCharTok{|}                                                        
  \ErrorTok{|======}                                            \ErrorTok{|}  \DecValTok{12}\NormalTok{\%}
  \SpecialCharTok{|}                                                        
  \ErrorTok{|======}                                            \ErrorTok{|}  \DecValTok{13}\NormalTok{\%}
  \SpecialCharTok{|}                                                        
  \ErrorTok{|=======}                                           \ErrorTok{|}  \DecValTok{13}\NormalTok{\%}
  \SpecialCharTok{|}                                                        
  \ErrorTok{|=======}                                           \ErrorTok{|}  \DecValTok{14}\NormalTok{\%}
  \SpecialCharTok{|}                                                        
  \ErrorTok{|=======}                                           \ErrorTok{|}  \DecValTok{15}\NormalTok{\%}
  \SpecialCharTok{|}                                                        
  \ErrorTok{|========}                                          \ErrorTok{|}  \DecValTok{15}\NormalTok{\%}
  \SpecialCharTok{|}                                                        
  \ErrorTok{|========}                                          \ErrorTok{|}  \DecValTok{16}\NormalTok{\%}
  \SpecialCharTok{|}                                                        
  \ErrorTok{|========}                                          \ErrorTok{|}  \DecValTok{17}\NormalTok{\%}
  \SpecialCharTok{|}                                                        
  \ErrorTok{|=========}                                         \ErrorTok{|}  \DecValTok{17}\NormalTok{\%}
  \SpecialCharTok{|}                                                        
  \ErrorTok{|=========}                                         \ErrorTok{|}  \DecValTok{18}\NormalTok{\%}
  \SpecialCharTok{|}                                                        
  \ErrorTok{|=========}                                         \ErrorTok{|}  \DecValTok{19}\NormalTok{\%}
  \SpecialCharTok{|}                                                        
  \ErrorTok{|==========}                                        \ErrorTok{|}  \DecValTok{19}\NormalTok{\%}
  \SpecialCharTok{|}                                                        
  \ErrorTok{|==========}                                        \ErrorTok{|}  \DecValTok{20}\NormalTok{\%}
  \SpecialCharTok{|}                                                        
  \ErrorTok{|==========}                                        \ErrorTok{|}  \DecValTok{21}\NormalTok{\%}
  \SpecialCharTok{|}                                                        
  \ErrorTok{|===========}                                       \ErrorTok{|}  \DecValTok{21}\NormalTok{\%}
  \SpecialCharTok{|}                                                        
  \ErrorTok{|===========}                                       \ErrorTok{|}  \DecValTok{22}\NormalTok{\%}
  \SpecialCharTok{|}                                                        
  \ErrorTok{|===========}                                       \ErrorTok{|}  \DecValTok{23}\NormalTok{\%}
  \SpecialCharTok{|}                                                        
  \ErrorTok{|============}                                      \ErrorTok{|}  \DecValTok{23}\NormalTok{\%}
  \SpecialCharTok{|}                                                        
  \ErrorTok{|============}                                      \ErrorTok{|}  \DecValTok{24}\NormalTok{\%}
  \SpecialCharTok{|}                                                        
  \ErrorTok{|============}                                      \ErrorTok{|}  \DecValTok{25}\NormalTok{\%}
  \SpecialCharTok{|}                                                        
  \ErrorTok{|=============}                                     \ErrorTok{|}  \DecValTok{25}\NormalTok{\%}
  \SpecialCharTok{|}                                                        
  \ErrorTok{|=============}                                     \ErrorTok{|}  \DecValTok{26}\NormalTok{\%}
  \SpecialCharTok{|}                                                        
  \ErrorTok{|=============}                                     \ErrorTok{|}  \DecValTok{27}\NormalTok{\%}
  \SpecialCharTok{|}                                                        
  \ErrorTok{|==============}                                    \ErrorTok{|}  \DecValTok{27}\NormalTok{\%}
  \SpecialCharTok{|}                                                        
  \ErrorTok{|==============}                                    \ErrorTok{|}  \DecValTok{28}\NormalTok{\%}
  \SpecialCharTok{|}                                                        
  \ErrorTok{|==============}                                    \ErrorTok{|}  \DecValTok{29}\NormalTok{\%}
  \SpecialCharTok{|}                                                        
  \ErrorTok{|===============}                                   \ErrorTok{|}  \DecValTok{29}\NormalTok{\%}
  \SpecialCharTok{|}                                                        
  \ErrorTok{|===============}                                   \ErrorTok{|}  \DecValTok{30}\NormalTok{\%}
  \SpecialCharTok{|}                                                        
  \ErrorTok{|===============}                                   \ErrorTok{|}  \DecValTok{31}\NormalTok{\%}
  \SpecialCharTok{|}                                                        
  \ErrorTok{|================}                                  \ErrorTok{|}  \DecValTok{31}\NormalTok{\%}
  \SpecialCharTok{|}                                                        
  \ErrorTok{|================}                                  \ErrorTok{|}  \DecValTok{32}\NormalTok{\%}
  \SpecialCharTok{|}                                                        
  \ErrorTok{|================}                                  \ErrorTok{|}  \DecValTok{33}\NormalTok{\%}
  \SpecialCharTok{|}                                                        
  \ErrorTok{|=================}                                 \ErrorTok{|}  \DecValTok{33}\NormalTok{\%}
  \SpecialCharTok{|}                                                        
  \ErrorTok{|=================}                                 \ErrorTok{|}  \DecValTok{34}\NormalTok{\%}
  \SpecialCharTok{|}                                                        
  \ErrorTok{|=================}                                 \ErrorTok{|}  \DecValTok{35}\NormalTok{\%}
  \SpecialCharTok{|}                                                        
  \ErrorTok{|==================}                                \ErrorTok{|}  \DecValTok{35}\NormalTok{\%}
  \SpecialCharTok{|}                                                        
  \ErrorTok{|==================}                                \ErrorTok{|}  \DecValTok{36}\NormalTok{\%}
  \SpecialCharTok{|}                                                        
  \ErrorTok{|==================}                                \ErrorTok{|}  \DecValTok{37}\NormalTok{\%}
  \SpecialCharTok{|}                                                        
  \ErrorTok{|===================}                               \ErrorTok{|}  \DecValTok{37}\NormalTok{\%}
  \SpecialCharTok{|}                                                        
  \ErrorTok{|===================}                               \ErrorTok{|}  \DecValTok{38}\NormalTok{\%}
  \SpecialCharTok{|}                                                        
  \ErrorTok{|===================}                               \ErrorTok{|}  \DecValTok{39}\NormalTok{\%}
  \SpecialCharTok{|}                                                        
  \ErrorTok{|====================}                              \ErrorTok{|}  \DecValTok{39}\NormalTok{\%}
  \SpecialCharTok{|}                                                        
  \ErrorTok{|====================}                              \ErrorTok{|}  \DecValTok{40}\NormalTok{\%}
  \SpecialCharTok{|}                                                        
  \ErrorTok{|====================}                              \ErrorTok{|}  \DecValTok{41}\NormalTok{\%}
  \SpecialCharTok{|}                                                        
  \ErrorTok{|=====================}                             \ErrorTok{|}  \DecValTok{41}\NormalTok{\%}
  \SpecialCharTok{|}                                                        
  \ErrorTok{|=====================}                             \ErrorTok{|}  \DecValTok{42}\NormalTok{\%}
  \SpecialCharTok{|}                                                        
  \ErrorTok{|=====================}                             \ErrorTok{|}  \DecValTok{43}\NormalTok{\%}
  \SpecialCharTok{|}                                                        
  \ErrorTok{|======================}                            \ErrorTok{|}  \DecValTok{43}\NormalTok{\%}
  \SpecialCharTok{|}                                                        
  \ErrorTok{|======================}                            \ErrorTok{|}  \DecValTok{44}\NormalTok{\%}
  \SpecialCharTok{|}                                                        
  \ErrorTok{|======================}                            \ErrorTok{|}  \DecValTok{45}\NormalTok{\%}
  \SpecialCharTok{|}                                                        
  \ErrorTok{|=======================}                           \ErrorTok{|}  \DecValTok{45}\NormalTok{\%}
  \SpecialCharTok{|}                                                        
  \ErrorTok{|=======================}                           \ErrorTok{|}  \DecValTok{46}\NormalTok{\%}
  \SpecialCharTok{|}                                                        
  \ErrorTok{|=======================}                           \ErrorTok{|}  \DecValTok{47}\NormalTok{\%}
  \SpecialCharTok{|}                                                        
  \ErrorTok{|========================}                          \ErrorTok{|}  \DecValTok{47}\NormalTok{\%}
  \SpecialCharTok{|}                                                        
  \ErrorTok{|========================}                          \ErrorTok{|}  \DecValTok{48}\NormalTok{\%}
  \SpecialCharTok{|}                                                        
  \ErrorTok{|========================}                          \ErrorTok{|}  \DecValTok{49}\NormalTok{\%}
  \SpecialCharTok{|}                                                        
  \ErrorTok{|=========================}                         \ErrorTok{|}  \DecValTok{49}\NormalTok{\%}
  \SpecialCharTok{|}                                                        
  \ErrorTok{|=========================}                         \ErrorTok{|}  \DecValTok{50}\NormalTok{\%}
  \SpecialCharTok{|}                                                        
  \ErrorTok{|=========================}                         \ErrorTok{|}  \DecValTok{51}\NormalTok{\%}
  \SpecialCharTok{|}                                                        
  \ErrorTok{|==========================}                        \ErrorTok{|}  \DecValTok{51}\NormalTok{\%}
  \SpecialCharTok{|}                                                        
  \ErrorTok{|==========================}                        \ErrorTok{|}  \DecValTok{52}\NormalTok{\%}
  \SpecialCharTok{|}                                                        
  \ErrorTok{|==========================}                        \ErrorTok{|}  \DecValTok{53}\NormalTok{\%}
  \SpecialCharTok{|}                                                        
  \ErrorTok{|===========================}                       \ErrorTok{|}  \DecValTok{53}\NormalTok{\%}
  \SpecialCharTok{|}                                                        
  \ErrorTok{|===========================}                       \ErrorTok{|}  \DecValTok{54}\NormalTok{\%}
  \SpecialCharTok{|}                                                        
  \ErrorTok{|===========================}                       \ErrorTok{|}  \DecValTok{55}\NormalTok{\%}
  \SpecialCharTok{|}                                                        
  \ErrorTok{|============================}                      \ErrorTok{|}  \DecValTok{55}\NormalTok{\%}
  \SpecialCharTok{|}                                                        
  \ErrorTok{|============================}                      \ErrorTok{|}  \DecValTok{56}\NormalTok{\%}
  \SpecialCharTok{|}                                                        
  \ErrorTok{|============================}                      \ErrorTok{|}  \DecValTok{57}\NormalTok{\%}
  \SpecialCharTok{|}                                                        
  \ErrorTok{|=============================}                     \ErrorTok{|}  \DecValTok{57}\NormalTok{\%}
  \SpecialCharTok{|}                                                        
  \ErrorTok{|=============================}                     \ErrorTok{|}  \DecValTok{58}\NormalTok{\%}
  \SpecialCharTok{|}                                                        
  \ErrorTok{|=============================}                     \ErrorTok{|}  \DecValTok{59}\NormalTok{\%}
  \SpecialCharTok{|}                                                        
  \ErrorTok{|==============================}                    \ErrorTok{|}  \DecValTok{59}\NormalTok{\%}
  \SpecialCharTok{|}                                                        
  \ErrorTok{|==============================}                    \ErrorTok{|}  \DecValTok{60}\NormalTok{\%}
  \SpecialCharTok{|}                                                        
  \ErrorTok{|==============================}                    \ErrorTok{|}  \DecValTok{61}\NormalTok{\%}
  \SpecialCharTok{|}                                                        
  \ErrorTok{|===============================}                   \ErrorTok{|}  \DecValTok{61}\NormalTok{\%}
  \SpecialCharTok{|}                                                        
  \ErrorTok{|===============================}                   \ErrorTok{|}  \DecValTok{62}\NormalTok{\%}
  \SpecialCharTok{|}                                                        
  \ErrorTok{|===============================}                   \ErrorTok{|}  \DecValTok{63}\NormalTok{\%}
  \SpecialCharTok{|}                                                        
  \ErrorTok{|================================}                  \ErrorTok{|}  \DecValTok{63}\NormalTok{\%}
  \SpecialCharTok{|}                                                        
  \ErrorTok{|================================}                  \ErrorTok{|}  \DecValTok{64}\NormalTok{\%}
  \SpecialCharTok{|}                                                        
  \ErrorTok{|================================}                  \ErrorTok{|}  \DecValTok{65}\NormalTok{\%}
  \SpecialCharTok{|}                                                        
  \ErrorTok{|=================================}                 \ErrorTok{|}  \DecValTok{65}\NormalTok{\%}
  \SpecialCharTok{|}                                                        
  \ErrorTok{|=================================}                 \ErrorTok{|}  \DecValTok{66}\NormalTok{\%}
  \SpecialCharTok{|}                                                        
  \ErrorTok{|=================================}                 \ErrorTok{|}  \DecValTok{67}\NormalTok{\%}
  \SpecialCharTok{|}                                                        
  \ErrorTok{|==================================}                \ErrorTok{|}  \DecValTok{67}\NormalTok{\%}
  \SpecialCharTok{|}                                                        
  \ErrorTok{|==================================}                \ErrorTok{|}  \DecValTok{68}\NormalTok{\%}
  \SpecialCharTok{|}                                                        
  \ErrorTok{|==================================}                \ErrorTok{|}  \DecValTok{69}\NormalTok{\%}
  \SpecialCharTok{|}                                                        
  \ErrorTok{|===================================}               \ErrorTok{|}  \DecValTok{69}\NormalTok{\%}
  \SpecialCharTok{|}                                                        
  \ErrorTok{|===================================}               \ErrorTok{|}  \DecValTok{70}\NormalTok{\%}
  \SpecialCharTok{|}                                                        
  \ErrorTok{|===================================}               \ErrorTok{|}  \DecValTok{71}\NormalTok{\%}
  \SpecialCharTok{|}                                                        
  \ErrorTok{|====================================}              \ErrorTok{|}  \DecValTok{71}\NormalTok{\%}
  \SpecialCharTok{|}                                                        
  \ErrorTok{|====================================}              \ErrorTok{|}  \DecValTok{72}\NormalTok{\%}
  \SpecialCharTok{|}                                                        
  \ErrorTok{|====================================}              \ErrorTok{|}  \DecValTok{73}\NormalTok{\%}
  \SpecialCharTok{|}                                                        
  \ErrorTok{|=====================================}             \ErrorTok{|}  \DecValTok{73}\NormalTok{\%}
  \SpecialCharTok{|}                                                        
  \ErrorTok{|=====================================}             \ErrorTok{|}  \DecValTok{74}\NormalTok{\%}
  \SpecialCharTok{|}                                                        
  \ErrorTok{|=====================================}             \ErrorTok{|}  \DecValTok{75}\NormalTok{\%}
  \SpecialCharTok{|}                                                        
  \ErrorTok{|======================================}            \ErrorTok{|}  \DecValTok{75}\NormalTok{\%}
  \SpecialCharTok{|}                                                        
  \ErrorTok{|======================================}            \ErrorTok{|}  \DecValTok{76}\NormalTok{\%}
  \SpecialCharTok{|}                                                        
  \ErrorTok{|======================================}            \ErrorTok{|}  \DecValTok{77}\NormalTok{\%}
  \SpecialCharTok{|}                                                        
  \ErrorTok{|=======================================}           \ErrorTok{|}  \DecValTok{77}\NormalTok{\%}
  \SpecialCharTok{|}                                                        
  \ErrorTok{|=======================================}           \ErrorTok{|}  \DecValTok{78}\NormalTok{\%}
  \SpecialCharTok{|}                                                        
  \ErrorTok{|=======================================}           \ErrorTok{|}  \DecValTok{79}\NormalTok{\%}
  \SpecialCharTok{|}                                                        
  \ErrorTok{|========================================}          \ErrorTok{|}  \DecValTok{79}\NormalTok{\%}
  \SpecialCharTok{|}                                                        
  \ErrorTok{|========================================}          \ErrorTok{|}  \DecValTok{80}\NormalTok{\%}
  \SpecialCharTok{|}                                                        
  \ErrorTok{|========================================}          \ErrorTok{|}  \DecValTok{81}\NormalTok{\%}
  \SpecialCharTok{|}                                                        
  \ErrorTok{|=========================================}         \ErrorTok{|}  \DecValTok{81}\NormalTok{\%}
  \SpecialCharTok{|}                                                        
  \ErrorTok{|=========================================}         \ErrorTok{|}  \DecValTok{82}\NormalTok{\%}
  \SpecialCharTok{|}                                                        
  \ErrorTok{|=========================================}         \ErrorTok{|}  \DecValTok{83}\NormalTok{\%}
  \SpecialCharTok{|}                                                        
  \ErrorTok{|==========================================}        \ErrorTok{|}  \DecValTok{83}\NormalTok{\%}
  \SpecialCharTok{|}                                                        
  \ErrorTok{|==========================================}        \ErrorTok{|}  \DecValTok{84}\NormalTok{\%}
  \SpecialCharTok{|}                                                        
  \ErrorTok{|==========================================}        \ErrorTok{|}  \DecValTok{85}\NormalTok{\%}
  \SpecialCharTok{|}                                                        
  \ErrorTok{|===========================================}       \ErrorTok{|}  \DecValTok{85}\NormalTok{\%}
  \SpecialCharTok{|}                                                        
  \ErrorTok{|===========================================}       \ErrorTok{|}  \DecValTok{86}\NormalTok{\%}
  \SpecialCharTok{|}                                                        
  \ErrorTok{|===========================================}       \ErrorTok{|}  \DecValTok{87}\NormalTok{\%}
  \SpecialCharTok{|}                                                        
  \ErrorTok{|============================================}      \ErrorTok{|}  \DecValTok{87}\NormalTok{\%}
  \SpecialCharTok{|}                                                        
  \ErrorTok{|============================================}      \ErrorTok{|}  \DecValTok{88}\NormalTok{\%}
  \SpecialCharTok{|}                                                        
  \ErrorTok{|============================================}      \ErrorTok{|}  \DecValTok{89}\NormalTok{\%}
  \SpecialCharTok{|}                                                        
  \ErrorTok{|=============================================}     \ErrorTok{|}  \DecValTok{89}\NormalTok{\%}
  \SpecialCharTok{|}                                                        
  \ErrorTok{|=============================================}     \ErrorTok{|}  \DecValTok{90}\NormalTok{\%}
  \SpecialCharTok{|}                                                        
  \ErrorTok{|=============================================}     \ErrorTok{|}  \DecValTok{91}\NormalTok{\%}
  \SpecialCharTok{|}                                                        
  \ErrorTok{|==============================================}    \ErrorTok{|}  \DecValTok{91}\NormalTok{\%}
  \SpecialCharTok{|}                                                        
  \ErrorTok{|==============================================}    \ErrorTok{|}  \DecValTok{92}\NormalTok{\%}
  \SpecialCharTok{|}                                                        
  \ErrorTok{|==============================================}    \ErrorTok{|}  \DecValTok{93}\NormalTok{\%}
  \SpecialCharTok{|}                                                        
  \ErrorTok{|===============================================}   \ErrorTok{|}  \DecValTok{93}\NormalTok{\%}
  \SpecialCharTok{|}                                                        
  \ErrorTok{|===============================================}   \ErrorTok{|}  \DecValTok{94}\NormalTok{\%}
  \SpecialCharTok{|}                                                        
  \ErrorTok{|===============================================}   \ErrorTok{|}  \DecValTok{95}\NormalTok{\%}
  \SpecialCharTok{|}                                                        
  \ErrorTok{|================================================}  \ErrorTok{|}  \DecValTok{95}\NormalTok{\%}
  \SpecialCharTok{|}                                                        
  \ErrorTok{|================================================}  \ErrorTok{|}  \DecValTok{96}\NormalTok{\%}
  \SpecialCharTok{|}                                                        
  \ErrorTok{|================================================}  \ErrorTok{|}  \DecValTok{97}\NormalTok{\%}
  \SpecialCharTok{|}                                                        
  \ErrorTok{|=================================================} \ErrorTok{|}  \DecValTok{97}\NormalTok{\%}
  \SpecialCharTok{|}                                                        
  \ErrorTok{|=================================================} \ErrorTok{|}  \DecValTok{98}\NormalTok{\%}
  \SpecialCharTok{|}                                                        
  \ErrorTok{|=================================================} \ErrorTok{|}  \DecValTok{99}\NormalTok{\%}
  \SpecialCharTok{|}                                                        
  \ErrorTok{|==================================================|}  \DecValTok{99}\NormalTok{\%}
  \SpecialCharTok{|}                                                        
  \ErrorTok{|==================================================|} \DecValTok{100}\NormalTok{\%}
\NormalTok{div.sp}
\CommentTok{\#\textgreater{} }
\CommentTok{\#\textgreater{} Call:}
\CommentTok{\#\textgreater{} morphol.disparity(f1 = shape.v \textasciitilde{} species, groups = \textasciitilde{}species,  }
\CommentTok{\#\textgreater{}     iter = 999) }
\CommentTok{\#\textgreater{} }
\CommentTok{\#\textgreater{} }
\CommentTok{\#\textgreater{} }
\CommentTok{\#\textgreater{} Randomized Residual Permutation Procedure Used}
\CommentTok{\#\textgreater{} 1000 Permutations}
\CommentTok{\#\textgreater{} }
\CommentTok{\#\textgreater{} Procrustes variances for defined groups}
\CommentTok{\#\textgreater{}         Aded         Aega         Aexa         Aupi }
\CommentTok{\#\textgreater{} 0.0004630470 0.0002893775 0.0006065943 0.0003949857 }
\CommentTok{\#\textgreater{}         Elmo         Irta         Jshu         Lomi }
\CommentTok{\#\textgreater{} 0.0001739802 0.0003795813 0.0004866557 0.0003683766 }
\CommentTok{\#\textgreater{} }
\CommentTok{\#\textgreater{} }
\CommentTok{\#\textgreater{} Pairwise absolute differences between variances}
\CommentTok{\#\textgreater{}              Aded         Aega         Aexa         Aupi}
\CommentTok{\#\textgreater{} Aded 0.000000e+00 1.736695e{-}04 0.0001435473 6.806128e{-}05}
\CommentTok{\#\textgreater{} Aega 1.736695e{-}04 0.000000e+00 0.0003172168 1.056082e{-}04}
\CommentTok{\#\textgreater{} Aexa 1.435473e{-}04 3.172168e{-}04 0.0000000000 2.116086e{-}04}
\CommentTok{\#\textgreater{} Aupi 6.806128e{-}05 1.056082e{-}04 0.0002116086 0.000000e+00}
\CommentTok{\#\textgreater{} Elmo 2.890668e{-}04 1.153973e{-}04 0.0004326141 2.210055e{-}04}
\CommentTok{\#\textgreater{} Irta 8.346564e{-}05 9.020385e{-}05 0.0002270129 1.540436e{-}05}
\CommentTok{\#\textgreater{} Jshu 2.360875e{-}05 1.972782e{-}04 0.0001199386 9.167002e{-}05}
\CommentTok{\#\textgreater{} Lomi 9.467042e{-}05 7.899906e{-}05 0.0002382177 2.660915e{-}05}
\CommentTok{\#\textgreater{}              Elmo         Irta         Jshu         Lomi}
\CommentTok{\#\textgreater{} Aded 0.0002890668 8.346564e{-}05 2.360875e{-}05 9.467042e{-}05}
\CommentTok{\#\textgreater{} Aega 0.0001153973 9.020385e{-}05 1.972782e{-}04 7.899906e{-}05}
\CommentTok{\#\textgreater{} Aexa 0.0004326141 2.270129e{-}04 1.199386e{-}04 2.382177e{-}04}
\CommentTok{\#\textgreater{} Aupi 0.0002210055 1.540436e{-}05 9.167002e{-}05 2.660915e{-}05}
\CommentTok{\#\textgreater{} Elmo 0.0000000000 2.056012e{-}04 3.126755e{-}04 1.943964e{-}04}
\CommentTok{\#\textgreater{} Irta 0.0002056012 0.000000e+00 1.070744e{-}04 1.120479e{-}05}
\CommentTok{\#\textgreater{} Jshu 0.0003126755 1.070744e{-}04 0.000000e+00 1.182792e{-}04}
\CommentTok{\#\textgreater{} Lomi 0.0001943964 1.120479e{-}05 1.182792e{-}04 0.000000e+00}
\CommentTok{\#\textgreater{} }
\CommentTok{\#\textgreater{} }
\CommentTok{\#\textgreater{} P{-}Values}
\CommentTok{\#\textgreater{}       Aded  Aega  Aexa  Aupi  Elmo  Irta  Jshu  Lomi}
\CommentTok{\#\textgreater{} Aded 1.000 0.175 0.247 0.600 0.021 0.520 0.844 0.468}
\CommentTok{\#\textgreater{} Aega 0.175 1.000 0.009 0.407 0.343 0.477 0.118 0.527}
\CommentTok{\#\textgreater{} Aexa 0.247 0.009 1.000 0.089 0.001 0.062 0.353 0.058}
\CommentTok{\#\textgreater{} Aupi 0.600 0.407 0.089 1.000 0.072 0.917 0.469 0.842}
\CommentTok{\#\textgreater{} Elmo 0.021 0.343 0.001 0.072 1.000 0.097 0.011 0.146}
\CommentTok{\#\textgreater{} Irta 0.520 0.477 0.062 0.917 0.097 1.000 0.418 0.934}
\CommentTok{\#\textgreater{} Jshu 0.844 0.118 0.353 0.469 0.011 0.418 1.000 0.361}
\CommentTok{\#\textgreater{} Lomi 0.468 0.527 0.058 0.842 0.146 0.934 0.361 1.000}
\CommentTok{\# Comparar com dispersão de pontos nos primeiros eixos da PCA:}
\NormalTok{col.group}\OtherTok{\textless{}{-}}\FunctionTok{rainbow}\NormalTok{(}\FunctionTok{length}\NormalTok{(}\FunctionTok{levels}\NormalTok{(species))) }\CommentTok{\# criar vetor de cores para grupos}
\FunctionTok{names}\NormalTok{(col.group)}\OtherTok{\textless{}{-}}\FunctionTok{levels}\NormalTok{(species)}
\NormalTok{col.group}\OtherTok{\textless{}{-}}\NormalTok{col.group[}\FunctionTok{match}\NormalTok{(species,}\FunctionTok{names}\NormalTok{(col.group))]}
\NormalTok{PCA}\OtherTok{\textless{}{-}}\FunctionTok{gm.prcomp}\NormalTok{(shape.v) }\CommentTok{\# rodar a PCA}
\FunctionTok{plot}\NormalTok{(PCA}\SpecialCharTok{$}\NormalTok{x[,}\DecValTok{1}\NormalTok{],PCA}\SpecialCharTok{$}\NormalTok{x[,}\DecValTok{2}\NormalTok{],}\AttributeTok{pch=}\DecValTok{21}\NormalTok{,}\AttributeTok{cex=}\DecValTok{2}\NormalTok{,}\AttributeTok{bg=}\NormalTok{col.group,}\AttributeTok{asp=}\NormalTok{T)}
\end{Highlighting}
\end{Shaded}

\includegraphics{03.12-mgdisp_files/figure-latex/unnamed-chunk-2-1.pdf}

\begin{Shaded}
\begin{Highlighting}[]
\CommentTok{\# Disparidade considerando alometria}
\NormalTok{gpa.v}\OtherTok{\textless{}{-}}\FunctionTok{gpagen}\NormalTok{(tps.v)}
\CommentTok{\#\textgreater{} }
\CommentTok{\#\textgreater{} Performing GPA}
\CommentTok{\#\textgreater{} }
  \SpecialCharTok{|}                                                        
  \ErrorTok{|}                                                  \ErrorTok{|}   \DecValTok{0}\NormalTok{\%}
  \SpecialCharTok{|}                                                        
  \ErrorTok{|============}                                      \ErrorTok{|}  \DecValTok{25}\NormalTok{\%}
  \SpecialCharTok{|}                                                        
  \ErrorTok{|=========================}                         \ErrorTok{|}  \DecValTok{50}\NormalTok{\%}
  \SpecialCharTok{|}                                                        
  \ErrorTok{|==================================================|} \DecValTok{100}\NormalTok{\%}
\CommentTok{\#\textgreater{} }
\CommentTok{\#\textgreater{} Making projections... Finished!}
\NormalTok{size}\OtherTok{\textless{}{-}}\NormalTok{gpa.v}\SpecialCharTok{$}\NormalTok{Csize}
\NormalTok{div.sp}\OtherTok{\textless{}{-}}\FunctionTok{morphol.disparity}\NormalTok{(shape.v}\SpecialCharTok{\textasciitilde{}}\FunctionTok{log}\NormalTok{(size)}\SpecialCharTok{*}\NormalTok{species,}\AttributeTok{groups=}\SpecialCharTok{\textasciitilde{}}\NormalTok{species,}\AttributeTok{iter=}\DecValTok{999}\NormalTok{)}
\CommentTok{\#\textgreater{} }
\CommentTok{\#\textgreater{} Please be aware that printing progress slows down the analysis (perhaps slightly).}
\CommentTok{\#\textgreater{} }
\CommentTok{\#\textgreater{} Preliminary Model Fit...}
\CommentTok{\#\textgreater{} }
\CommentTok{\#\textgreater{} Sums of Squares calculations: 3 permutations.}
\CommentTok{\#\textgreater{} }
  \SpecialCharTok{|}                                                        
  \ErrorTok{|}                                                  \ErrorTok{|}   \DecValTok{0}\NormalTok{\%}
  \SpecialCharTok{|}                                                        
  \ErrorTok{|=================}                                 \ErrorTok{|}  \DecValTok{33}\NormalTok{\%}
  \SpecialCharTok{|}                                                        
  \ErrorTok{|=================================}                 \ErrorTok{|}  \DecValTok{67}\NormalTok{\%}
  \SpecialCharTok{|}                                                        
  \ErrorTok{|==================================================|} \DecValTok{100}\NormalTok{\%}
\CommentTok{\#\textgreater{} }
\CommentTok{\#\textgreater{} }
\CommentTok{\#\textgreater{} Performimg pairwise comparisons of disparity}
\CommentTok{\#\textgreater{} }
  \SpecialCharTok{|}                                                        
  \ErrorTok{|}                                                  \ErrorTok{|}   \DecValTok{0}\NormalTok{\%}
  \SpecialCharTok{|}                                                        
  \ErrorTok{|}                                                  \ErrorTok{|}   \DecValTok{1}\NormalTok{\%}
  \SpecialCharTok{|}                                                        
  \ErrorTok{|=}                                                 \ErrorTok{|}   \DecValTok{1}\NormalTok{\%}
  \SpecialCharTok{|}                                                        
  \ErrorTok{|=}                                                 \ErrorTok{|}   \DecValTok{2}\NormalTok{\%}
  \SpecialCharTok{|}                                                        
  \ErrorTok{|=}                                                 \ErrorTok{|}   \DecValTok{3}\NormalTok{\%}
  \SpecialCharTok{|}                                                        
  \ErrorTok{|==}                                                \ErrorTok{|}   \DecValTok{3}\NormalTok{\%}
  \SpecialCharTok{|}                                                        
  \ErrorTok{|==}                                                \ErrorTok{|}   \DecValTok{4}\NormalTok{\%}
  \SpecialCharTok{|}                                                        
  \ErrorTok{|==}                                                \ErrorTok{|}   \DecValTok{5}\NormalTok{\%}
  \SpecialCharTok{|}                                                        
  \ErrorTok{|===}                                               \ErrorTok{|}   \DecValTok{5}\NormalTok{\%}
  \SpecialCharTok{|}                                                        
  \ErrorTok{|===}                                               \ErrorTok{|}   \DecValTok{6}\NormalTok{\%}
  \SpecialCharTok{|}                                                        
  \ErrorTok{|===}                                               \ErrorTok{|}   \DecValTok{7}\NormalTok{\%}
  \SpecialCharTok{|}                                                        
  \ErrorTok{|====}                                              \ErrorTok{|}   \DecValTok{7}\NormalTok{\%}
  \SpecialCharTok{|}                                                        
  \ErrorTok{|====}                                              \ErrorTok{|}   \DecValTok{8}\NormalTok{\%}
  \SpecialCharTok{|}                                                        
  \ErrorTok{|====}                                              \ErrorTok{|}   \DecValTok{9}\NormalTok{\%}
  \SpecialCharTok{|}                                                        
  \ErrorTok{|=====}                                             \ErrorTok{|}   \DecValTok{9}\NormalTok{\%}
  \SpecialCharTok{|}                                                        
  \ErrorTok{|=====}                                             \ErrorTok{|}  \DecValTok{10}\NormalTok{\%}
  \SpecialCharTok{|}                                                        
  \ErrorTok{|=====}                                             \ErrorTok{|}  \DecValTok{11}\NormalTok{\%}
  \SpecialCharTok{|}                                                        
  \ErrorTok{|======}                                            \ErrorTok{|}  \DecValTok{11}\NormalTok{\%}
  \SpecialCharTok{|}                                                        
  \ErrorTok{|======}                                            \ErrorTok{|}  \DecValTok{12}\NormalTok{\%}
  \SpecialCharTok{|}                                                        
  \ErrorTok{|======}                                            \ErrorTok{|}  \DecValTok{13}\NormalTok{\%}
  \SpecialCharTok{|}                                                        
  \ErrorTok{|=======}                                           \ErrorTok{|}  \DecValTok{13}\NormalTok{\%}
  \SpecialCharTok{|}                                                        
  \ErrorTok{|=======}                                           \ErrorTok{|}  \DecValTok{14}\NormalTok{\%}
  \SpecialCharTok{|}                                                        
  \ErrorTok{|=======}                                           \ErrorTok{|}  \DecValTok{15}\NormalTok{\%}
  \SpecialCharTok{|}                                                        
  \ErrorTok{|========}                                          \ErrorTok{|}  \DecValTok{15}\NormalTok{\%}
  \SpecialCharTok{|}                                                        
  \ErrorTok{|========}                                          \ErrorTok{|}  \DecValTok{16}\NormalTok{\%}
  \SpecialCharTok{|}                                                        
  \ErrorTok{|========}                                          \ErrorTok{|}  \DecValTok{17}\NormalTok{\%}
  \SpecialCharTok{|}                                                        
  \ErrorTok{|=========}                                         \ErrorTok{|}  \DecValTok{17}\NormalTok{\%}
  \SpecialCharTok{|}                                                        
  \ErrorTok{|=========}                                         \ErrorTok{|}  \DecValTok{18}\NormalTok{\%}
  \SpecialCharTok{|}                                                        
  \ErrorTok{|=========}                                         \ErrorTok{|}  \DecValTok{19}\NormalTok{\%}
  \SpecialCharTok{|}                                                        
  \ErrorTok{|==========}                                        \ErrorTok{|}  \DecValTok{19}\NormalTok{\%}
  \SpecialCharTok{|}                                                        
  \ErrorTok{|==========}                                        \ErrorTok{|}  \DecValTok{20}\NormalTok{\%}
  \SpecialCharTok{|}                                                        
  \ErrorTok{|==========}                                        \ErrorTok{|}  \DecValTok{21}\NormalTok{\%}
  \SpecialCharTok{|}                                                        
  \ErrorTok{|===========}                                       \ErrorTok{|}  \DecValTok{21}\NormalTok{\%}
  \SpecialCharTok{|}                                                        
  \ErrorTok{|===========}                                       \ErrorTok{|}  \DecValTok{22}\NormalTok{\%}
  \SpecialCharTok{|}                                                        
  \ErrorTok{|===========}                                       \ErrorTok{|}  \DecValTok{23}\NormalTok{\%}
  \SpecialCharTok{|}                                                        
  \ErrorTok{|============}                                      \ErrorTok{|}  \DecValTok{23}\NormalTok{\%}
  \SpecialCharTok{|}                                                        
  \ErrorTok{|============}                                      \ErrorTok{|}  \DecValTok{24}\NormalTok{\%}
  \SpecialCharTok{|}                                                        
  \ErrorTok{|============}                                      \ErrorTok{|}  \DecValTok{25}\NormalTok{\%}
  \SpecialCharTok{|}                                                        
  \ErrorTok{|=============}                                     \ErrorTok{|}  \DecValTok{25}\NormalTok{\%}
  \SpecialCharTok{|}                                                        
  \ErrorTok{|=============}                                     \ErrorTok{|}  \DecValTok{26}\NormalTok{\%}
  \SpecialCharTok{|}                                                        
  \ErrorTok{|=============}                                     \ErrorTok{|}  \DecValTok{27}\NormalTok{\%}
  \SpecialCharTok{|}                                                        
  \ErrorTok{|==============}                                    \ErrorTok{|}  \DecValTok{27}\NormalTok{\%}
  \SpecialCharTok{|}                                                        
  \ErrorTok{|==============}                                    \ErrorTok{|}  \DecValTok{28}\NormalTok{\%}
  \SpecialCharTok{|}                                                        
  \ErrorTok{|==============}                                    \ErrorTok{|}  \DecValTok{29}\NormalTok{\%}
  \SpecialCharTok{|}                                                        
  \ErrorTok{|===============}                                   \ErrorTok{|}  \DecValTok{29}\NormalTok{\%}
  \SpecialCharTok{|}                                                        
  \ErrorTok{|===============}                                   \ErrorTok{|}  \DecValTok{30}\NormalTok{\%}
  \SpecialCharTok{|}                                                        
  \ErrorTok{|===============}                                   \ErrorTok{|}  \DecValTok{31}\NormalTok{\%}
  \SpecialCharTok{|}                                                        
  \ErrorTok{|================}                                  \ErrorTok{|}  \DecValTok{31}\NormalTok{\%}
  \SpecialCharTok{|}                                                        
  \ErrorTok{|================}                                  \ErrorTok{|}  \DecValTok{32}\NormalTok{\%}
  \SpecialCharTok{|}                                                        
  \ErrorTok{|================}                                  \ErrorTok{|}  \DecValTok{33}\NormalTok{\%}
  \SpecialCharTok{|}                                                        
  \ErrorTok{|=================}                                 \ErrorTok{|}  \DecValTok{33}\NormalTok{\%}
  \SpecialCharTok{|}                                                        
  \ErrorTok{|=================}                                 \ErrorTok{|}  \DecValTok{34}\NormalTok{\%}
  \SpecialCharTok{|}                                                        
  \ErrorTok{|=================}                                 \ErrorTok{|}  \DecValTok{35}\NormalTok{\%}
  \SpecialCharTok{|}                                                        
  \ErrorTok{|==================}                                \ErrorTok{|}  \DecValTok{35}\NormalTok{\%}
  \SpecialCharTok{|}                                                        
  \ErrorTok{|==================}                                \ErrorTok{|}  \DecValTok{36}\NormalTok{\%}
  \SpecialCharTok{|}                                                        
  \ErrorTok{|==================}                                \ErrorTok{|}  \DecValTok{37}\NormalTok{\%}
  \SpecialCharTok{|}                                                        
  \ErrorTok{|===================}                               \ErrorTok{|}  \DecValTok{37}\NormalTok{\%}
  \SpecialCharTok{|}                                                        
  \ErrorTok{|===================}                               \ErrorTok{|}  \DecValTok{38}\NormalTok{\%}
  \SpecialCharTok{|}                                                        
  \ErrorTok{|===================}                               \ErrorTok{|}  \DecValTok{39}\NormalTok{\%}
  \SpecialCharTok{|}                                                        
  \ErrorTok{|====================}                              \ErrorTok{|}  \DecValTok{39}\NormalTok{\%}
  \SpecialCharTok{|}                                                        
  \ErrorTok{|====================}                              \ErrorTok{|}  \DecValTok{40}\NormalTok{\%}
  \SpecialCharTok{|}                                                        
  \ErrorTok{|====================}                              \ErrorTok{|}  \DecValTok{41}\NormalTok{\%}
  \SpecialCharTok{|}                                                        
  \ErrorTok{|=====================}                             \ErrorTok{|}  \DecValTok{41}\NormalTok{\%}
  \SpecialCharTok{|}                                                        
  \ErrorTok{|=====================}                             \ErrorTok{|}  \DecValTok{42}\NormalTok{\%}
  \SpecialCharTok{|}                                                        
  \ErrorTok{|=====================}                             \ErrorTok{|}  \DecValTok{43}\NormalTok{\%}
  \SpecialCharTok{|}                                                        
  \ErrorTok{|======================}                            \ErrorTok{|}  \DecValTok{43}\NormalTok{\%}
  \SpecialCharTok{|}                                                        
  \ErrorTok{|======================}                            \ErrorTok{|}  \DecValTok{44}\NormalTok{\%}
  \SpecialCharTok{|}                                                        
  \ErrorTok{|======================}                            \ErrorTok{|}  \DecValTok{45}\NormalTok{\%}
  \SpecialCharTok{|}                                                        
  \ErrorTok{|=======================}                           \ErrorTok{|}  \DecValTok{45}\NormalTok{\%}
  \SpecialCharTok{|}                                                        
  \ErrorTok{|=======================}                           \ErrorTok{|}  \DecValTok{46}\NormalTok{\%}
  \SpecialCharTok{|}                                                        
  \ErrorTok{|=======================}                           \ErrorTok{|}  \DecValTok{47}\NormalTok{\%}
  \SpecialCharTok{|}                                                        
  \ErrorTok{|========================}                          \ErrorTok{|}  \DecValTok{47}\NormalTok{\%}
  \SpecialCharTok{|}                                                        
  \ErrorTok{|========================}                          \ErrorTok{|}  \DecValTok{48}\NormalTok{\%}
  \SpecialCharTok{|}                                                        
  \ErrorTok{|========================}                          \ErrorTok{|}  \DecValTok{49}\NormalTok{\%}
  \SpecialCharTok{|}                                                        
  \ErrorTok{|=========================}                         \ErrorTok{|}  \DecValTok{49}\NormalTok{\%}
  \SpecialCharTok{|}                                                        
  \ErrorTok{|=========================}                         \ErrorTok{|}  \DecValTok{50}\NormalTok{\%}
  \SpecialCharTok{|}                                                        
  \ErrorTok{|=========================}                         \ErrorTok{|}  \DecValTok{51}\NormalTok{\%}
  \SpecialCharTok{|}                                                        
  \ErrorTok{|==========================}                        \ErrorTok{|}  \DecValTok{51}\NormalTok{\%}
  \SpecialCharTok{|}                                                        
  \ErrorTok{|==========================}                        \ErrorTok{|}  \DecValTok{52}\NormalTok{\%}
  \SpecialCharTok{|}                                                        
  \ErrorTok{|==========================}                        \ErrorTok{|}  \DecValTok{53}\NormalTok{\%}
  \SpecialCharTok{|}                                                        
  \ErrorTok{|===========================}                       \ErrorTok{|}  \DecValTok{53}\NormalTok{\%}
  \SpecialCharTok{|}                                                        
  \ErrorTok{|===========================}                       \ErrorTok{|}  \DecValTok{54}\NormalTok{\%}
  \SpecialCharTok{|}                                                        
  \ErrorTok{|===========================}                       \ErrorTok{|}  \DecValTok{55}\NormalTok{\%}
  \SpecialCharTok{|}                                                        
  \ErrorTok{|============================}                      \ErrorTok{|}  \DecValTok{55}\NormalTok{\%}
  \SpecialCharTok{|}                                                        
  \ErrorTok{|============================}                      \ErrorTok{|}  \DecValTok{56}\NormalTok{\%}
  \SpecialCharTok{|}                                                        
  \ErrorTok{|============================}                      \ErrorTok{|}  \DecValTok{57}\NormalTok{\%}
  \SpecialCharTok{|}                                                        
  \ErrorTok{|=============================}                     \ErrorTok{|}  \DecValTok{57}\NormalTok{\%}
  \SpecialCharTok{|}                                                        
  \ErrorTok{|=============================}                     \ErrorTok{|}  \DecValTok{58}\NormalTok{\%}
  \SpecialCharTok{|}                                                        
  \ErrorTok{|=============================}                     \ErrorTok{|}  \DecValTok{59}\NormalTok{\%}
  \SpecialCharTok{|}                                                        
  \ErrorTok{|==============================}                    \ErrorTok{|}  \DecValTok{59}\NormalTok{\%}
  \SpecialCharTok{|}                                                        
  \ErrorTok{|==============================}                    \ErrorTok{|}  \DecValTok{60}\NormalTok{\%}
  \SpecialCharTok{|}                                                        
  \ErrorTok{|==============================}                    \ErrorTok{|}  \DecValTok{61}\NormalTok{\%}
  \SpecialCharTok{|}                                                        
  \ErrorTok{|===============================}                   \ErrorTok{|}  \DecValTok{61}\NormalTok{\%}
  \SpecialCharTok{|}                                                        
  \ErrorTok{|===============================}                   \ErrorTok{|}  \DecValTok{62}\NormalTok{\%}
  \SpecialCharTok{|}                                                        
  \ErrorTok{|===============================}                   \ErrorTok{|}  \DecValTok{63}\NormalTok{\%}
  \SpecialCharTok{|}                                                        
  \ErrorTok{|================================}                  \ErrorTok{|}  \DecValTok{63}\NormalTok{\%}
  \SpecialCharTok{|}                                                        
  \ErrorTok{|================================}                  \ErrorTok{|}  \DecValTok{64}\NormalTok{\%}
  \SpecialCharTok{|}                                                        
  \ErrorTok{|================================}                  \ErrorTok{|}  \DecValTok{65}\NormalTok{\%}
  \SpecialCharTok{|}                                                        
  \ErrorTok{|=================================}                 \ErrorTok{|}  \DecValTok{65}\NormalTok{\%}
  \SpecialCharTok{|}                                                        
  \ErrorTok{|=================================}                 \ErrorTok{|}  \DecValTok{66}\NormalTok{\%}
  \SpecialCharTok{|}                                                        
  \ErrorTok{|=================================}                 \ErrorTok{|}  \DecValTok{67}\NormalTok{\%}
  \SpecialCharTok{|}                                                        
  \ErrorTok{|==================================}                \ErrorTok{|}  \DecValTok{67}\NormalTok{\%}
  \SpecialCharTok{|}                                                        
  \ErrorTok{|==================================}                \ErrorTok{|}  \DecValTok{68}\NormalTok{\%}
  \SpecialCharTok{|}                                                        
  \ErrorTok{|==================================}                \ErrorTok{|}  \DecValTok{69}\NormalTok{\%}
  \SpecialCharTok{|}                                                        
  \ErrorTok{|===================================}               \ErrorTok{|}  \DecValTok{69}\NormalTok{\%}
  \SpecialCharTok{|}                                                        
  \ErrorTok{|===================================}               \ErrorTok{|}  \DecValTok{70}\NormalTok{\%}
  \SpecialCharTok{|}                                                        
  \ErrorTok{|===================================}               \ErrorTok{|}  \DecValTok{71}\NormalTok{\%}
  \SpecialCharTok{|}                                                        
  \ErrorTok{|====================================}              \ErrorTok{|}  \DecValTok{71}\NormalTok{\%}
  \SpecialCharTok{|}                                                        
  \ErrorTok{|====================================}              \ErrorTok{|}  \DecValTok{72}\NormalTok{\%}
  \SpecialCharTok{|}                                                        
  \ErrorTok{|====================================}              \ErrorTok{|}  \DecValTok{73}\NormalTok{\%}
  \SpecialCharTok{|}                                                        
  \ErrorTok{|=====================================}             \ErrorTok{|}  \DecValTok{73}\NormalTok{\%}
  \SpecialCharTok{|}                                                        
  \ErrorTok{|=====================================}             \ErrorTok{|}  \DecValTok{74}\NormalTok{\%}
  \SpecialCharTok{|}                                                        
  \ErrorTok{|=====================================}             \ErrorTok{|}  \DecValTok{75}\NormalTok{\%}
  \SpecialCharTok{|}                                                        
  \ErrorTok{|======================================}            \ErrorTok{|}  \DecValTok{75}\NormalTok{\%}
  \SpecialCharTok{|}                                                        
  \ErrorTok{|======================================}            \ErrorTok{|}  \DecValTok{76}\NormalTok{\%}
  \SpecialCharTok{|}                                                        
  \ErrorTok{|======================================}            \ErrorTok{|}  \DecValTok{77}\NormalTok{\%}
  \SpecialCharTok{|}                                                        
  \ErrorTok{|=======================================}           \ErrorTok{|}  \DecValTok{77}\NormalTok{\%}
  \SpecialCharTok{|}                                                        
  \ErrorTok{|=======================================}           \ErrorTok{|}  \DecValTok{78}\NormalTok{\%}
  \SpecialCharTok{|}                                                        
  \ErrorTok{|=======================================}           \ErrorTok{|}  \DecValTok{79}\NormalTok{\%}
  \SpecialCharTok{|}                                                        
  \ErrorTok{|========================================}          \ErrorTok{|}  \DecValTok{79}\NormalTok{\%}
  \SpecialCharTok{|}                                                        
  \ErrorTok{|========================================}          \ErrorTok{|}  \DecValTok{80}\NormalTok{\%}
  \SpecialCharTok{|}                                                        
  \ErrorTok{|========================================}          \ErrorTok{|}  \DecValTok{81}\NormalTok{\%}
  \SpecialCharTok{|}                                                        
  \ErrorTok{|=========================================}         \ErrorTok{|}  \DecValTok{81}\NormalTok{\%}
  \SpecialCharTok{|}                                                        
  \ErrorTok{|=========================================}         \ErrorTok{|}  \DecValTok{82}\NormalTok{\%}
  \SpecialCharTok{|}                                                        
  \ErrorTok{|=========================================}         \ErrorTok{|}  \DecValTok{83}\NormalTok{\%}
  \SpecialCharTok{|}                                                        
  \ErrorTok{|==========================================}        \ErrorTok{|}  \DecValTok{83}\NormalTok{\%}
  \SpecialCharTok{|}                                                        
  \ErrorTok{|==========================================}        \ErrorTok{|}  \DecValTok{84}\NormalTok{\%}
  \SpecialCharTok{|}                                                        
  \ErrorTok{|==========================================}        \ErrorTok{|}  \DecValTok{85}\NormalTok{\%}
  \SpecialCharTok{|}                                                        
  \ErrorTok{|===========================================}       \ErrorTok{|}  \DecValTok{85}\NormalTok{\%}
  \SpecialCharTok{|}                                                        
  \ErrorTok{|===========================================}       \ErrorTok{|}  \DecValTok{86}\NormalTok{\%}
  \SpecialCharTok{|}                                                        
  \ErrorTok{|===========================================}       \ErrorTok{|}  \DecValTok{87}\NormalTok{\%}
  \SpecialCharTok{|}                                                        
  \ErrorTok{|============================================}      \ErrorTok{|}  \DecValTok{87}\NormalTok{\%}
  \SpecialCharTok{|}                                                        
  \ErrorTok{|============================================}      \ErrorTok{|}  \DecValTok{88}\NormalTok{\%}
  \SpecialCharTok{|}                                                        
  \ErrorTok{|============================================}      \ErrorTok{|}  \DecValTok{89}\NormalTok{\%}
  \SpecialCharTok{|}                                                        
  \ErrorTok{|=============================================}     \ErrorTok{|}  \DecValTok{89}\NormalTok{\%}
  \SpecialCharTok{|}                                                        
  \ErrorTok{|=============================================}     \ErrorTok{|}  \DecValTok{90}\NormalTok{\%}
  \SpecialCharTok{|}                                                        
  \ErrorTok{|=============================================}     \ErrorTok{|}  \DecValTok{91}\NormalTok{\%}
  \SpecialCharTok{|}                                                        
  \ErrorTok{|==============================================}    \ErrorTok{|}  \DecValTok{91}\NormalTok{\%}
  \SpecialCharTok{|}                                                        
  \ErrorTok{|==============================================}    \ErrorTok{|}  \DecValTok{92}\NormalTok{\%}
  \SpecialCharTok{|}                                                        
  \ErrorTok{|==============================================}    \ErrorTok{|}  \DecValTok{93}\NormalTok{\%}
  \SpecialCharTok{|}                                                        
  \ErrorTok{|===============================================}   \ErrorTok{|}  \DecValTok{93}\NormalTok{\%}
  \SpecialCharTok{|}                                                        
  \ErrorTok{|===============================================}   \ErrorTok{|}  \DecValTok{94}\NormalTok{\%}
  \SpecialCharTok{|}                                                        
  \ErrorTok{|===============================================}   \ErrorTok{|}  \DecValTok{95}\NormalTok{\%}
  \SpecialCharTok{|}                                                        
  \ErrorTok{|================================================}  \ErrorTok{|}  \DecValTok{95}\NormalTok{\%}
  \SpecialCharTok{|}                                                        
  \ErrorTok{|================================================}  \ErrorTok{|}  \DecValTok{96}\NormalTok{\%}
  \SpecialCharTok{|}                                                        
  \ErrorTok{|================================================}  \ErrorTok{|}  \DecValTok{97}\NormalTok{\%}
  \SpecialCharTok{|}                                                        
  \ErrorTok{|=================================================} \ErrorTok{|}  \DecValTok{97}\NormalTok{\%}
  \SpecialCharTok{|}                                                        
  \ErrorTok{|=================================================} \ErrorTok{|}  \DecValTok{98}\NormalTok{\%}
  \SpecialCharTok{|}                                                        
  \ErrorTok{|=================================================} \ErrorTok{|}  \DecValTok{99}\NormalTok{\%}
  \SpecialCharTok{|}                                                        
  \ErrorTok{|==================================================|}  \DecValTok{99}\NormalTok{\%}
  \SpecialCharTok{|}                                                        
  \ErrorTok{|==================================================|} \DecValTok{100}\NormalTok{\%}
\NormalTok{div.sp}
\CommentTok{\#\textgreater{} }
\CommentTok{\#\textgreater{} Call:}
\CommentTok{\#\textgreater{} morphol.disparity(f1 = shape.v \textasciitilde{} log(size) * species, groups = \textasciitilde{}species,  }
\CommentTok{\#\textgreater{}     iter = 999) }
\CommentTok{\#\textgreater{} }
\CommentTok{\#\textgreater{} }
\CommentTok{\#\textgreater{} }
\CommentTok{\#\textgreater{} Randomized Residual Permutation Procedure Used}
\CommentTok{\#\textgreater{} 1000 Permutations}
\CommentTok{\#\textgreater{} }
\CommentTok{\#\textgreater{} Procrustes variances for defined groups}
\CommentTok{\#\textgreater{}         Aded         Aega         Aexa         Aupi }
\CommentTok{\#\textgreater{} 0.0003260927 0.0001331069 0.0004333711 0.0001910877 }
\CommentTok{\#\textgreater{}         Elmo         Irta         Jshu         Lomi }
\CommentTok{\#\textgreater{} 0.0001310704 0.0002895081 0.0003230759 0.0002042844 }
\CommentTok{\#\textgreater{} }
\CommentTok{\#\textgreater{} }
\CommentTok{\#\textgreater{} Pairwise absolute differences between variances}
\CommentTok{\#\textgreater{}              Aded         Aega         Aexa         Aupi}
\CommentTok{\#\textgreater{} Aded 0.000000e+00 1.929858e{-}04 0.0001072784 1.350051e{-}04}
\CommentTok{\#\textgreater{} Aega 1.929858e{-}04 0.000000e+00 0.0003002643 5.798077e{-}05}
\CommentTok{\#\textgreater{} Aexa 1.072784e{-}04 3.002643e{-}04 0.0000000000 2.422835e{-}04}
\CommentTok{\#\textgreater{} Aupi 1.350051e{-}04 5.798077e{-}05 0.0002422835 0.000000e+00}
\CommentTok{\#\textgreater{} Elmo 1.950223e{-}04 2.036498e{-}06 0.0003023008 6.001727e{-}05}
\CommentTok{\#\textgreater{} Irta 3.658457e{-}05 1.564013e{-}04 0.0001438630 9.842049e{-}05}
\CommentTok{\#\textgreater{} Jshu 3.016813e{-}06 1.899690e{-}04 0.0001102952 1.319882e{-}04}
\CommentTok{\#\textgreater{} Lomi 1.218084e{-}04 7.117747e{-}05 0.0002290868 1.319670e{-}05}
\CommentTok{\#\textgreater{}              Elmo         Irta         Jshu         Lomi}
\CommentTok{\#\textgreater{} Aded 1.950223e{-}04 3.658457e{-}05 3.016813e{-}06 1.218084e{-}04}
\CommentTok{\#\textgreater{} Aega 2.036498e{-}06 1.564013e{-}04 1.899690e{-}04 7.117747e{-}05}
\CommentTok{\#\textgreater{} Aexa 3.023008e{-}04 1.438630e{-}04 1.102952e{-}04 2.290868e{-}04}
\CommentTok{\#\textgreater{} Aupi 6.001727e{-}05 9.842049e{-}05 1.319882e{-}04 1.319670e{-}05}
\CommentTok{\#\textgreater{} Elmo 0.000000e+00 1.584378e{-}04 1.920055e{-}04 7.321397e{-}05}
\CommentTok{\#\textgreater{} Irta 1.584378e{-}04 0.000000e+00 3.356776e{-}05 8.522379e{-}05}
\CommentTok{\#\textgreater{} Jshu 1.920055e{-}04 3.356776e{-}05 0.000000e+00 1.187915e{-}04}
\CommentTok{\#\textgreater{} Lomi 7.321397e{-}05 8.522379e{-}05 1.187915e{-}04 0.000000e+00}
\CommentTok{\#\textgreater{} }
\CommentTok{\#\textgreater{} }
\CommentTok{\#\textgreater{} P{-}Values}
\CommentTok{\#\textgreater{}       Aded  Aega  Aexa  Aupi  Elmo  Irta  Jshu  Lomi}
\CommentTok{\#\textgreater{} Aded 1.000 0.096 0.350 0.241 0.098 0.760 0.980 0.306}
\CommentTok{\#\textgreater{} Aega 0.096 1.000 0.012 0.614 0.980 0.162 0.104 0.521}
\CommentTok{\#\textgreater{} Aexa 0.350 0.012 1.000 0.043 0.009 0.252 0.331 0.043}
\CommentTok{\#\textgreater{} Aupi 0.241 0.614 0.043 1.000 0.581 0.398 0.236 0.915}
\CommentTok{\#\textgreater{} Elmo 0.098 0.980 0.009 0.581 1.000 0.158 0.092 0.527}
\CommentTok{\#\textgreater{} Irta 0.760 0.162 0.252 0.398 0.158 1.000 0.760 0.474}
\CommentTok{\#\textgreater{} Jshu 0.980 0.104 0.331 0.236 0.092 0.760 1.000 0.302}
\CommentTok{\#\textgreater{} Lomi 0.306 0.521 0.043 0.915 0.527 0.474 0.302 1.000}
\end{Highlighting}
\end{Shaded}

\chapter{Morfometria geométrica e filogenias}\label{morfometria-geomuxe9trica-e-filogenias}

Quando as morfologias são de espécies diferentes, os dados não são independentes, já que as espécies são conectadas por uma árvore filogenética.
A não-independência filogenética precisa ser levada em consideração em análises de correlação.
Pra além disso, podemos entender muito sobre padrões e processos evolutivos da morfologia quando temos dados morfológicos de diferentes espécies e hipóteses filogenéticas.
A parte deste manual sobre Macroevolução explora muitas perguntas/análises que podem ser feitas em um contexto comparativo.
Aqui, os exemplos vão apresentar algumas análises comparativas no contexto de morfometria geométrica.

Para textos abrangentes sobre macroevolução e métodos filogenéticos comparativos no contexto de morfometria geométrica, veja: (\textbf{zelditch2012?}); (\textbf{monteiro2013?}); (\textbf{adams\&collyer2018?}).

Para literatura geral, e muitas outras opções analíticas, veja o material e as referências na parte sobre Macroevolução deste manual.

Vamos carregar dados morfológicos obtidos por morfometria geométrica e uma hipótese filogenética para explorar algumas análises filogenéticas comparativas.

\begin{Shaded}
\begin{Highlighting}[]
\FunctionTok{require}\NormalTok{(geomorph)}
\CommentTok{\#\textgreater{} Carregando pacotes exigidos: geomorph}
\CommentTok{\#\textgreater{} Carregando pacotes exigidos: RRPP}
\CommentTok{\#\textgreater{} Carregando pacotes exigidos: rgl}
\CommentTok{\#\textgreater{} Carregando pacotes exigidos: Matrix}
\CommentTok{\# Carregar dados de forma média por espécie em formato de matriz}
\NormalTok{tps}\OtherTok{\textless{}{-}}\FunctionTok{readland.tps}\NormalTok{(}\StringTok{"dadosmg/mandibula.dig\_curso.pls.tps"}\NormalTok{,}\AttributeTok{specID =} \StringTok{"ID"}\NormalTok{, }\AttributeTok{readcurves =} \ConstantTok{FALSE}\NormalTok{)}
\NormalTok{gpa}\OtherTok{\textless{}{-}}\FunctionTok{gpagen}\NormalTok{(tps)}
\NormalTok{shape}\OtherTok{\textless{}{-}}\NormalTok{gpa}\SpecialCharTok{$}\NormalTok{coords}
\NormalTok{size}\OtherTok{\textless{}{-}}\NormalTok{gpa}\SpecialCharTok{$}\NormalTok{Csize}
\NormalTok{ref.mand}\OtherTok{\textless{}{-}}\FunctionTok{mshape}\NormalTok{(shape)}

\CommentTok{\# Carregar classificadores a partir de lista externa}
\NormalTok{plan}\OtherTok{\textless{}{-}}\FunctionTok{read.table}\NormalTok{(}\StringTok{"dadosmg/Planilha\_PLS.txt"}\NormalTok{,}\AttributeTok{h=}\NormalTok{T)}
\NormalTok{plan}
\NormalTok{species}\OtherTok{\textless{}{-}}\FunctionTok{as.factor}\NormalTok{(plan[,}\DecValTok{2}\NormalTok{])}
\NormalTok{species}

\CommentTok{\# Calcular forma média por espécie}
\NormalTok{shape}\FloatTok{.2}\NormalTok{d}\OtherTok{\textless{}{-}}\FunctionTok{two.d.array}\NormalTok{(shape)}
\NormalTok{shape}\FloatTok{.2}\NormalTok{d.means}\OtherTok{\textless{}{-}}\FunctionTok{rowsum}\NormalTok{(shape}\FloatTok{.2}\NormalTok{d,species)}\SpecialCharTok{/}\FunctionTok{as.vector}\NormalTok{(}\FunctionTok{table}\NormalTok{(species))}
\NormalTok{shape.means}\OtherTok{\textless{}{-}}\FunctionTok{arrayspecs}\NormalTok{(shape}\FloatTok{.2}\NormalTok{d.means,}\FunctionTok{dim}\NormalTok{(shape)[}\DecValTok{1}\NormalTok{],}\FunctionTok{dim}\NormalTok{(shape)[}\DecValTok{2}\NormalTok{])}

\NormalTok{size.means}\OtherTok{\textless{}{-}}\FunctionTok{rowsum}\NormalTok{(size,species)}\SpecialCharTok{/}\FunctionTok{as.vector}\NormalTok{(}\FunctionTok{table}\NormalTok{(species))}
\end{Highlighting}
\end{Shaded}

Hipótese filogenética.

\begin{Shaded}
\begin{Highlighting}[]
\FunctionTok{require}\NormalTok{(ape)}
\CommentTok{\#\textgreater{} Carregando pacotes exigidos: ape}
\NormalTok{tree}\OtherTok{\textless{}{-}}\FunctionTok{read.nexus}\NormalTok{(}\StringTok{"dadosmg/nexus.tree.sig.nex"}\NormalTok{)}

\FunctionTok{plot}\NormalTok{(tree)}
\end{Highlighting}
\end{Shaded}

\includegraphics{03.13-mgpcm_files/figure-latex/unnamed-chunk-2-1.pdf}

\begin{Shaded}
\begin{Highlighting}[]
\NormalTok{tree}\OtherTok{\textless{}{-}}\FunctionTok{compute.brlen}\NormalTok{(tree,}\DecValTok{1}\NormalTok{) }\CommentTok{\# definir comprimento dos ramos = 1}
\end{Highlighting}
\end{Shaded}

\section{Filomorfoespaço}\label{filomorfoespauxe7o}

Podemos representar de maneira gráfica a filogenia no espaço de forma criado com uma análise de componentes principais.
Essa é a PCA padrão, inalterada, com a projeção da filogenia e estimativa de formas ancestrais.

\begin{Shaded}
\begin{Highlighting}[]
\NormalTok{PCA}\OtherTok{\textless{}{-}}\FunctionTok{gm.prcomp}\NormalTok{(shape.means,}\AttributeTok{phy=}\NormalTok{tree)}
\FunctionTok{summary}\NormalTok{(PCA)}
\FunctionTok{plot}\NormalTok{(PCA,}\AttributeTok{phylo=}\ConstantTok{TRUE}\NormalTok{)}
\end{Highlighting}
\end{Shaded}

\includegraphics{03.13-mgpcm_files/figure-latex/unnamed-chunk-3-1.pdf}

Filomorfoespaço 3D.

\begin{Shaded}
\begin{Highlighting}[]
\FunctionTok{require}\NormalTok{(phytools)}
\CommentTok{\#\textgreater{} Carregando pacotes exigidos: phytools}
\CommentTok{\#\textgreater{} Carregando pacotes exigidos: maps}
\CommentTok{\#phylomorphospace3d(tree,PCA$x[,1:3])}
\end{Highlighting}
\end{Shaded}

A visualização de formas pode ser feita ao longo dos eixos, como na PCA tradicional, bem como podem ser visualizadas as formas ancestrais estimadas, formas médias de espécies atuais, etc.

\begin{Shaded}
\begin{Highlighting}[]
\NormalTok{ref}\OtherTok{\textless{}{-}}\FunctionTok{mshape}\NormalTok{(shape.means)}
\NormalTok{ancestral.shapes}\OtherTok{\textless{}{-}}\FunctionTok{arrayspecs}\NormalTok{(PCA}\SpecialCharTok{$}\NormalTok{ancestors,}\DecValTok{13}\NormalTok{,}\DecValTok{2}\NormalTok{)}
\CommentTok{\# Forma ancestral contra a média}
\FunctionTok{plotRefToTarget}\NormalTok{(ref,ancestral.shapes[,,}\DecValTok{1}\NormalTok{],}\AttributeTok{method=}\StringTok{"TPS"}\NormalTok{)}
\end{Highlighting}
\end{Shaded}

\includegraphics{03.13-mgpcm_files/figure-latex/unnamed-chunk-5-1.pdf}

\begin{Shaded}
\begin{Highlighting}[]
\CommentTok{\# Forma ancestral contra espécie atual}
\FunctionTok{plotRefToTarget}\NormalTok{(ancestral.shapes[,,}\DecValTok{1}\NormalTok{],shape.means[,,}\StringTok{"Aexa"}\NormalTok{],}\AttributeTok{method=}\StringTok{"TPS"}\NormalTok{)}
\end{Highlighting}
\end{Shaded}

\includegraphics{03.13-mgpcm_files/figure-latex/unnamed-chunk-5-2.pdf}

Visualizar tamanho ao longo da filogenia.

\begin{Shaded}
\begin{Highlighting}[]
\NormalTok{size.means1}\OtherTok{\textless{}{-}}\FunctionTok{as.vector}\NormalTok{(size.means)}
\FunctionTok{names}\NormalTok{(size.means1)}\OtherTok{=}\FunctionTok{rownames}\NormalTok{(size.means)}
\NormalTok{tree1}\OtherTok{\textless{}{-}}\FunctionTok{compute.brlen}\NormalTok{(tree,}\AttributeTok{method=}\StringTok{"Grafen"}\NormalTok{) }\CommentTok{\# make ultrametric}
\CommentTok{\# cores}
\FunctionTok{contMap}\NormalTok{(tree1,size.means1)}
\end{Highlighting}
\end{Shaded}

\includegraphics{03.13-mgpcm_files/figure-latex/unnamed-chunk-6-1.pdf}

\begin{Shaded}
\begin{Highlighting}[]
\CommentTok{\# dots}
\FunctionTok{dotTree}\NormalTok{(tree1,size.means1,}\AttributeTok{length=}\DecValTok{10}\NormalTok{)}
\end{Highlighting}
\end{Shaded}

\includegraphics{03.13-mgpcm_files/figure-latex/unnamed-chunk-6-2.pdf}

\begin{Shaded}
\begin{Highlighting}[]
\CommentTok{\# barras}
\FunctionTok{plotTree.barplot}\NormalTok{(tree1,size.means1)}
\end{Highlighting}
\end{Shaded}

\includegraphics{03.13-mgpcm_files/figure-latex/unnamed-chunk-6-3.pdf}

\section{PCA filogenética}\label{pca-filogenuxe9tica}

A PCA filogenética \emph{sensu} (\textbf{revell2009?}).

\begin{Shaded}
\begin{Highlighting}[]
\NormalTok{phyloPCA}\OtherTok{\textless{}{-}}\FunctionTok{gm.prcomp}\NormalTok{(shape.means,}\AttributeTok{phy=}\NormalTok{tree,}\AttributeTok{GLS=}\ConstantTok{TRUE}\NormalTok{)}
\FunctionTok{summary}\NormalTok{(phyloPCA)}
\FunctionTok{plot}\NormalTok{(phyloPCA,}\AttributeTok{phylo=}\ConstantTok{TRUE}\NormalTok{)}
\end{Highlighting}
\end{Shaded}

\includegraphics{03.13-mgpcm_files/figure-latex/unnamed-chunk-7-1.pdf}

\section{Sinal filogenético multivariado}\label{sinal-filogenuxe9tico-multivariado}

Sinal filogenético é um padrão que mostra se espécies mais próximas filogeneticamente têm morfologias mais similares do que espécies mais distantes filogeneticamente.
Existem vários métodos para o cálculo de sinal filogenético, o mais comum é a estatística K (Blomberg, Garland JR., and Ives 2003; \textbf{adams2014?}).
K=1 indica que espécies próximas são tão parecidas quanto prevê o modelo Browniano K\textless1 indica sinal filogenético menor do que esperado pelo modelo Browniano K\textgreater1 indica sinal filogenético maior do que esperado pelo modelo Browniano

\begin{Shaded}
\begin{Highlighting}[]
\CommentTok{\# Sinal Filogenético para forma do crânio}
\NormalTok{sinal.k}\OtherTok{\textless{}{-}}\FunctionTok{physignal}\NormalTok{(shape.means,tree,}\AttributeTok{iter=}\DecValTok{999}\NormalTok{)}
\NormalTok{sinal.k}
\CommentTok{\#\textgreater{} }
\CommentTok{\#\textgreater{} Call:}
\CommentTok{\#\textgreater{} physignal(A = shape.means, phy = tree, iter = 999) }
\CommentTok{\#\textgreater{} }
\CommentTok{\#\textgreater{} }
\CommentTok{\#\textgreater{} }
\CommentTok{\#\textgreater{} Observed Phylogenetic Signal (K): 0.6146}
\CommentTok{\#\textgreater{} }
\CommentTok{\#\textgreater{} P{-}value: 0.158}
\CommentTok{\#\textgreater{} }
\CommentTok{\#\textgreater{} Based on 1000 random permutations}
\CommentTok{\#\textgreater{} }
\CommentTok{\#\textgreater{}  Use physignal.z to estimate effect size.}
\CommentTok{\# Sinal Filogenético para tamanho do crânio}
\NormalTok{sinal.k.size}\OtherTok{\textless{}{-}}\FunctionTok{physignal}\NormalTok{(size.means,tree,}\AttributeTok{iter=}\DecValTok{999}\NormalTok{)}
\NormalTok{sinal.k.size}
\CommentTok{\#\textgreater{} }
\CommentTok{\#\textgreater{} Call:}
\CommentTok{\#\textgreater{} physignal(A = size.means, phy = tree, iter = 999) }
\CommentTok{\#\textgreater{} }
\CommentTok{\#\textgreater{} }
\CommentTok{\#\textgreater{} }
\CommentTok{\#\textgreater{} Observed Phylogenetic Signal (K): 0.5612}
\CommentTok{\#\textgreater{} }
\CommentTok{\#\textgreater{} P{-}value: 0.416}
\CommentTok{\#\textgreater{} }
\CommentTok{\#\textgreater{} Based on 1000 random permutations}
\CommentTok{\#\textgreater{} }
\CommentTok{\#\textgreater{}  Use physignal.z to estimate effect size.}
\end{Highlighting}
\end{Shaded}

\section{Análise filogenética generalizada por quadrados mínimos (PGLS multivariada)}\label{anuxe1lise-filogenuxe9tica-generalizada-por-quadrados-muxednimos-pgls-multivariada}

\begin{Shaded}
\begin{Highlighting}[]
\CommentTok{\# Variáveis hipotéticas para usar como exemplo}
\NormalTok{nomes}\OtherTok{\textless{}{-}}\FunctionTok{c}\NormalTok{(}\StringTok{"Aded"}\NormalTok{,}\StringTok{"Aega"}\NormalTok{,}\StringTok{"Aexa"}\NormalTok{,}\StringTok{"Aupi"}\NormalTok{,}\StringTok{"Elmo"}\NormalTok{,}\StringTok{"Irta"}\NormalTok{,}\StringTok{"Jshu"}\NormalTok{,}\StringTok{"Lomi"}\NormalTok{)}
\NormalTok{temp}\OtherTok{\textless{}{-}}\FunctionTok{c}\NormalTok{(}\FloatTok{12.15}\NormalTok{,}\DecValTok{22}\NormalTok{,}\FloatTok{16.02}\NormalTok{,}\FloatTok{10.42}\NormalTok{,}\FloatTok{8.21}\NormalTok{,}\FloatTok{8.39}\NormalTok{,}\DecValTok{25}\NormalTok{,}\FloatTok{7.86}\NormalTok{)}
\NormalTok{alt}\OtherTok{\textless{}{-}}\FunctionTok{c}\NormalTok{(}\DecValTok{2951}\NormalTok{,}\DecValTok{300}\NormalTok{,}\DecValTok{1931}\NormalTok{,}\DecValTok{3222}\NormalTok{,}\FloatTok{753.96}\NormalTok{,}\FloatTok{719.94}\NormalTok{,}\DecValTok{700}\NormalTok{,}\FloatTok{866.65}\NormalTok{)}
\FunctionTok{names}\NormalTok{(temp)}\OtherTok{\textless{}{-}}\NormalTok{nomes}
\FunctionTok{names}\NormalTok{(alt)}\OtherTok{\textless{}{-}}\NormalTok{nomes}

\NormalTok{gdf}\OtherTok{\textless{}{-}}\FunctionTok{geomorph.data.frame}\NormalTok{(}\AttributeTok{shape.means=}\NormalTok{shape.means,}\AttributeTok{temp=}\NormalTok{temp,}\AttributeTok{tree=}\NormalTok{tree)}
\NormalTok{fit.pgls}\OtherTok{\textless{}{-}}\FunctionTok{procD.pgls}\NormalTok{(shape.means}\SpecialCharTok{\textasciitilde{}}\NormalTok{temp,tree,}\AttributeTok{data=}\NormalTok{gdf)}
\CommentTok{\#\textgreater{} }
\CommentTok{\#\textgreater{} Please be aware that printing progress slows down the analysis (perhaps slightly).}
\CommentTok{\#\textgreater{} }
\CommentTok{\#\textgreater{} Preliminary Model Fit...}
\CommentTok{\#\textgreater{} }
\CommentTok{\#\textgreater{} Sums of Squares calculations: 1000 permutations.}
\CommentTok{\#\textgreater{} }
  \SpecialCharTok{|}                                                        
  \ErrorTok{|}                                                  \ErrorTok{|}   \DecValTok{0}\NormalTok{\%}
  \SpecialCharTok{|}                                                        
  \ErrorTok{|}                                                  \ErrorTok{|}   \DecValTok{1}\NormalTok{\%}
  \SpecialCharTok{|}                                                        
  \ErrorTok{|=}                                                 \ErrorTok{|}   \DecValTok{1}\NormalTok{\%}
  \SpecialCharTok{|}                                                        
  \ErrorTok{|=}                                                 \ErrorTok{|}   \DecValTok{2}\NormalTok{\%}
  \SpecialCharTok{|}                                                        
  \ErrorTok{|=}                                                 \ErrorTok{|}   \DecValTok{3}\NormalTok{\%}
  \SpecialCharTok{|}                                                        
  \ErrorTok{|==}                                                \ErrorTok{|}   \DecValTok{3}\NormalTok{\%}
  \SpecialCharTok{|}                                                        
  \ErrorTok{|==}                                                \ErrorTok{|}   \DecValTok{4}\NormalTok{\%}
  \SpecialCharTok{|}                                                        
  \ErrorTok{|==}                                                \ErrorTok{|}   \DecValTok{5}\NormalTok{\%}
  \SpecialCharTok{|}                                                        
  \ErrorTok{|===}                                               \ErrorTok{|}   \DecValTok{5}\NormalTok{\%}
  \SpecialCharTok{|}                                                        
  \ErrorTok{|===}                                               \ErrorTok{|}   \DecValTok{6}\NormalTok{\%}
  \SpecialCharTok{|}                                                        
  \ErrorTok{|===}                                               \ErrorTok{|}   \DecValTok{7}\NormalTok{\%}
  \SpecialCharTok{|}                                                        
  \ErrorTok{|====}                                              \ErrorTok{|}   \DecValTok{7}\NormalTok{\%}
  \SpecialCharTok{|}                                                        
  \ErrorTok{|====}                                              \ErrorTok{|}   \DecValTok{8}\NormalTok{\%}
  \SpecialCharTok{|}                                                        
  \ErrorTok{|====}                                              \ErrorTok{|}   \DecValTok{9}\NormalTok{\%}
  \SpecialCharTok{|}                                                        
  \ErrorTok{|=====}                                             \ErrorTok{|}   \DecValTok{9}\NormalTok{\%}
  \SpecialCharTok{|}                                                        
  \ErrorTok{|=====}                                             \ErrorTok{|}  \DecValTok{10}\NormalTok{\%}
  \SpecialCharTok{|}                                                        
  \ErrorTok{|=====}                                             \ErrorTok{|}  \DecValTok{11}\NormalTok{\%}
  \SpecialCharTok{|}                                                        
  \ErrorTok{|======}                                            \ErrorTok{|}  \DecValTok{11}\NormalTok{\%}
  \SpecialCharTok{|}                                                        
  \ErrorTok{|======}                                            \ErrorTok{|}  \DecValTok{12}\NormalTok{\%}
  \SpecialCharTok{|}                                                        
  \ErrorTok{|======}                                            \ErrorTok{|}  \DecValTok{13}\NormalTok{\%}
  \SpecialCharTok{|}                                                        
  \ErrorTok{|=======}                                           \ErrorTok{|}  \DecValTok{13}\NormalTok{\%}
  \SpecialCharTok{|}                                                        
  \ErrorTok{|=======}                                           \ErrorTok{|}  \DecValTok{14}\NormalTok{\%}
  \SpecialCharTok{|}                                                        
  \ErrorTok{|=======}                                           \ErrorTok{|}  \DecValTok{15}\NormalTok{\%}
  \SpecialCharTok{|}                                                        
  \ErrorTok{|========}                                          \ErrorTok{|}  \DecValTok{15}\NormalTok{\%}
  \SpecialCharTok{|}                                                        
  \ErrorTok{|========}                                          \ErrorTok{|}  \DecValTok{16}\NormalTok{\%}
  \SpecialCharTok{|}                                                        
  \ErrorTok{|========}                                          \ErrorTok{|}  \DecValTok{17}\NormalTok{\%}
  \SpecialCharTok{|}                                                        
  \ErrorTok{|=========}                                         \ErrorTok{|}  \DecValTok{17}\NormalTok{\%}
  \SpecialCharTok{|}                                                        
  \ErrorTok{|=========}                                         \ErrorTok{|}  \DecValTok{18}\NormalTok{\%}
  \SpecialCharTok{|}                                                        
  \ErrorTok{|=========}                                         \ErrorTok{|}  \DecValTok{19}\NormalTok{\%}
  \SpecialCharTok{|}                                                        
  \ErrorTok{|==========}                                        \ErrorTok{|}  \DecValTok{19}\NormalTok{\%}
  \SpecialCharTok{|}                                                        
  \ErrorTok{|==========}                                        \ErrorTok{|}  \DecValTok{20}\NormalTok{\%}
  \SpecialCharTok{|}                                                        
  \ErrorTok{|==========}                                        \ErrorTok{|}  \DecValTok{21}\NormalTok{\%}
  \SpecialCharTok{|}                                                        
  \ErrorTok{|===========}                                       \ErrorTok{|}  \DecValTok{21}\NormalTok{\%}
  \SpecialCharTok{|}                                                        
  \ErrorTok{|===========}                                       \ErrorTok{|}  \DecValTok{22}\NormalTok{\%}
  \SpecialCharTok{|}                                                        
  \ErrorTok{|===========}                                       \ErrorTok{|}  \DecValTok{23}\NormalTok{\%}
  \SpecialCharTok{|}                                                        
  \ErrorTok{|============}                                      \ErrorTok{|}  \DecValTok{23}\NormalTok{\%}
  \SpecialCharTok{|}                                                        
  \ErrorTok{|============}                                      \ErrorTok{|}  \DecValTok{24}\NormalTok{\%}
  \SpecialCharTok{|}                                                        
  \ErrorTok{|============}                                      \ErrorTok{|}  \DecValTok{25}\NormalTok{\%}
  \SpecialCharTok{|}                                                        
  \ErrorTok{|=============}                                     \ErrorTok{|}  \DecValTok{25}\NormalTok{\%}
  \SpecialCharTok{|}                                                        
  \ErrorTok{|=============}                                     \ErrorTok{|}  \DecValTok{26}\NormalTok{\%}
  \SpecialCharTok{|}                                                        
  \ErrorTok{|=============}                                     \ErrorTok{|}  \DecValTok{27}\NormalTok{\%}
  \SpecialCharTok{|}                                                        
  \ErrorTok{|==============}                                    \ErrorTok{|}  \DecValTok{27}\NormalTok{\%}
  \SpecialCharTok{|}                                                        
  \ErrorTok{|==============}                                    \ErrorTok{|}  \DecValTok{28}\NormalTok{\%}
  \SpecialCharTok{|}                                                        
  \ErrorTok{|==============}                                    \ErrorTok{|}  \DecValTok{29}\NormalTok{\%}
  \SpecialCharTok{|}                                                        
  \ErrorTok{|===============}                                   \ErrorTok{|}  \DecValTok{29}\NormalTok{\%}
  \SpecialCharTok{|}                                                        
  \ErrorTok{|===============}                                   \ErrorTok{|}  \DecValTok{30}\NormalTok{\%}
  \SpecialCharTok{|}                                                        
  \ErrorTok{|===============}                                   \ErrorTok{|}  \DecValTok{31}\NormalTok{\%}
  \SpecialCharTok{|}                                                        
  \ErrorTok{|================}                                  \ErrorTok{|}  \DecValTok{31}\NormalTok{\%}
  \SpecialCharTok{|}                                                        
  \ErrorTok{|================}                                  \ErrorTok{|}  \DecValTok{32}\NormalTok{\%}
  \SpecialCharTok{|}                                                        
  \ErrorTok{|================}                                  \ErrorTok{|}  \DecValTok{33}\NormalTok{\%}
  \SpecialCharTok{|}                                                        
  \ErrorTok{|=================}                                 \ErrorTok{|}  \DecValTok{33}\NormalTok{\%}
  \SpecialCharTok{|}                                                        
  \ErrorTok{|=================}                                 \ErrorTok{|}  \DecValTok{34}\NormalTok{\%}
  \SpecialCharTok{|}                                                        
  \ErrorTok{|=================}                                 \ErrorTok{|}  \DecValTok{35}\NormalTok{\%}
  \SpecialCharTok{|}                                                        
  \ErrorTok{|==================}                                \ErrorTok{|}  \DecValTok{35}\NormalTok{\%}
  \SpecialCharTok{|}                                                        
  \ErrorTok{|==================}                                \ErrorTok{|}  \DecValTok{36}\NormalTok{\%}
  \SpecialCharTok{|}                                                        
  \ErrorTok{|==================}                                \ErrorTok{|}  \DecValTok{37}\NormalTok{\%}
  \SpecialCharTok{|}                                                        
  \ErrorTok{|===================}                               \ErrorTok{|}  \DecValTok{37}\NormalTok{\%}
  \SpecialCharTok{|}                                                        
  \ErrorTok{|===================}                               \ErrorTok{|}  \DecValTok{38}\NormalTok{\%}
  \SpecialCharTok{|}                                                        
  \ErrorTok{|===================}                               \ErrorTok{|}  \DecValTok{39}\NormalTok{\%}
  \SpecialCharTok{|}                                                        
  \ErrorTok{|====================}                              \ErrorTok{|}  \DecValTok{39}\NormalTok{\%}
  \SpecialCharTok{|}                                                        
  \ErrorTok{|====================}                              \ErrorTok{|}  \DecValTok{40}\NormalTok{\%}
  \SpecialCharTok{|}                                                        
  \ErrorTok{|====================}                              \ErrorTok{|}  \DecValTok{41}\NormalTok{\%}
  \SpecialCharTok{|}                                                        
  \ErrorTok{|=====================}                             \ErrorTok{|}  \DecValTok{41}\NormalTok{\%}
  \SpecialCharTok{|}                                                        
  \ErrorTok{|=====================}                             \ErrorTok{|}  \DecValTok{42}\NormalTok{\%}
  \SpecialCharTok{|}                                                        
  \ErrorTok{|=====================}                             \ErrorTok{|}  \DecValTok{43}\NormalTok{\%}
  \SpecialCharTok{|}                                                        
  \ErrorTok{|======================}                            \ErrorTok{|}  \DecValTok{43}\NormalTok{\%}
  \SpecialCharTok{|}                                                        
  \ErrorTok{|======================}                            \ErrorTok{|}  \DecValTok{44}\NormalTok{\%}
  \SpecialCharTok{|}                                                        
  \ErrorTok{|======================}                            \ErrorTok{|}  \DecValTok{45}\NormalTok{\%}
  \SpecialCharTok{|}                                                        
  \ErrorTok{|=======================}                           \ErrorTok{|}  \DecValTok{45}\NormalTok{\%}
  \SpecialCharTok{|}                                                        
  \ErrorTok{|=======================}                           \ErrorTok{|}  \DecValTok{46}\NormalTok{\%}
  \SpecialCharTok{|}                                                        
  \ErrorTok{|=======================}                           \ErrorTok{|}  \DecValTok{47}\NormalTok{\%}
  \SpecialCharTok{|}                                                        
  \ErrorTok{|========================}                          \ErrorTok{|}  \DecValTok{47}\NormalTok{\%}
  \SpecialCharTok{|}                                                        
  \ErrorTok{|========================}                          \ErrorTok{|}  \DecValTok{48}\NormalTok{\%}
  \SpecialCharTok{|}                                                        
  \ErrorTok{|========================}                          \ErrorTok{|}  \DecValTok{49}\NormalTok{\%}
  \SpecialCharTok{|}                                                        
  \ErrorTok{|=========================}                         \ErrorTok{|}  \DecValTok{49}\NormalTok{\%}
  \SpecialCharTok{|}                                                        
  \ErrorTok{|=========================}                         \ErrorTok{|}  \DecValTok{50}\NormalTok{\%}
  \SpecialCharTok{|}                                                        
  \ErrorTok{|=========================}                         \ErrorTok{|}  \DecValTok{51}\NormalTok{\%}
  \SpecialCharTok{|}                                                        
  \ErrorTok{|==========================}                        \ErrorTok{|}  \DecValTok{51}\NormalTok{\%}
  \SpecialCharTok{|}                                                        
  \ErrorTok{|==========================}                        \ErrorTok{|}  \DecValTok{52}\NormalTok{\%}
  \SpecialCharTok{|}                                                        
  \ErrorTok{|==========================}                        \ErrorTok{|}  \DecValTok{53}\NormalTok{\%}
  \SpecialCharTok{|}                                                        
  \ErrorTok{|===========================}                       \ErrorTok{|}  \DecValTok{53}\NormalTok{\%}
  \SpecialCharTok{|}                                                        
  \ErrorTok{|===========================}                       \ErrorTok{|}  \DecValTok{54}\NormalTok{\%}
  \SpecialCharTok{|}                                                        
  \ErrorTok{|===========================}                       \ErrorTok{|}  \DecValTok{55}\NormalTok{\%}
  \SpecialCharTok{|}                                                        
  \ErrorTok{|============================}                      \ErrorTok{|}  \DecValTok{55}\NormalTok{\%}
  \SpecialCharTok{|}                                                        
  \ErrorTok{|============================}                      \ErrorTok{|}  \DecValTok{56}\NormalTok{\%}
  \SpecialCharTok{|}                                                        
  \ErrorTok{|============================}                      \ErrorTok{|}  \DecValTok{57}\NormalTok{\%}
  \SpecialCharTok{|}                                                        
  \ErrorTok{|=============================}                     \ErrorTok{|}  \DecValTok{57}\NormalTok{\%}
  \SpecialCharTok{|}                                                        
  \ErrorTok{|=============================}                     \ErrorTok{|}  \DecValTok{58}\NormalTok{\%}
  \SpecialCharTok{|}                                                        
  \ErrorTok{|=============================}                     \ErrorTok{|}  \DecValTok{59}\NormalTok{\%}
  \SpecialCharTok{|}                                                        
  \ErrorTok{|==============================}                    \ErrorTok{|}  \DecValTok{59}\NormalTok{\%}
  \SpecialCharTok{|}                                                        
  \ErrorTok{|==============================}                    \ErrorTok{|}  \DecValTok{60}\NormalTok{\%}
  \SpecialCharTok{|}                                                        
  \ErrorTok{|==============================}                    \ErrorTok{|}  \DecValTok{61}\NormalTok{\%}
  \SpecialCharTok{|}                                                        
  \ErrorTok{|===============================}                   \ErrorTok{|}  \DecValTok{61}\NormalTok{\%}
  \SpecialCharTok{|}                                                        
  \ErrorTok{|===============================}                   \ErrorTok{|}  \DecValTok{62}\NormalTok{\%}
  \SpecialCharTok{|}                                                        
  \ErrorTok{|===============================}                   \ErrorTok{|}  \DecValTok{63}\NormalTok{\%}
  \SpecialCharTok{|}                                                        
  \ErrorTok{|================================}                  \ErrorTok{|}  \DecValTok{63}\NormalTok{\%}
  \SpecialCharTok{|}                                                        
  \ErrorTok{|================================}                  \ErrorTok{|}  \DecValTok{64}\NormalTok{\%}
  \SpecialCharTok{|}                                                        
  \ErrorTok{|================================}                  \ErrorTok{|}  \DecValTok{65}\NormalTok{\%}
  \SpecialCharTok{|}                                                        
  \ErrorTok{|=================================}                 \ErrorTok{|}  \DecValTok{65}\NormalTok{\%}
  \SpecialCharTok{|}                                                        
  \ErrorTok{|=================================}                 \ErrorTok{|}  \DecValTok{66}\NormalTok{\%}
  \SpecialCharTok{|}                                                        
  \ErrorTok{|=================================}                 \ErrorTok{|}  \DecValTok{67}\NormalTok{\%}
  \SpecialCharTok{|}                                                        
  \ErrorTok{|==================================}                \ErrorTok{|}  \DecValTok{67}\NormalTok{\%}
  \SpecialCharTok{|}                                                        
  \ErrorTok{|==================================}                \ErrorTok{|}  \DecValTok{68}\NormalTok{\%}
  \SpecialCharTok{|}                                                        
  \ErrorTok{|==================================}                \ErrorTok{|}  \DecValTok{69}\NormalTok{\%}
  \SpecialCharTok{|}                                                        
  \ErrorTok{|===================================}               \ErrorTok{|}  \DecValTok{69}\NormalTok{\%}
  \SpecialCharTok{|}                                                        
  \ErrorTok{|===================================}               \ErrorTok{|}  \DecValTok{70}\NormalTok{\%}
  \SpecialCharTok{|}                                                        
  \ErrorTok{|===================================}               \ErrorTok{|}  \DecValTok{71}\NormalTok{\%}
  \SpecialCharTok{|}                                                        
  \ErrorTok{|====================================}              \ErrorTok{|}  \DecValTok{71}\NormalTok{\%}
  \SpecialCharTok{|}                                                        
  \ErrorTok{|====================================}              \ErrorTok{|}  \DecValTok{72}\NormalTok{\%}
  \SpecialCharTok{|}                                                        
  \ErrorTok{|====================================}              \ErrorTok{|}  \DecValTok{73}\NormalTok{\%}
  \SpecialCharTok{|}                                                        
  \ErrorTok{|=====================================}             \ErrorTok{|}  \DecValTok{73}\NormalTok{\%}
  \SpecialCharTok{|}                                                        
  \ErrorTok{|=====================================}             \ErrorTok{|}  \DecValTok{74}\NormalTok{\%}
  \SpecialCharTok{|}                                                        
  \ErrorTok{|=====================================}             \ErrorTok{|}  \DecValTok{75}\NormalTok{\%}
  \SpecialCharTok{|}                                                        
  \ErrorTok{|======================================}            \ErrorTok{|}  \DecValTok{75}\NormalTok{\%}
  \SpecialCharTok{|}                                                        
  \ErrorTok{|======================================}            \ErrorTok{|}  \DecValTok{76}\NormalTok{\%}
  \SpecialCharTok{|}                                                        
  \ErrorTok{|======================================}            \ErrorTok{|}  \DecValTok{77}\NormalTok{\%}
  \SpecialCharTok{|}                                                        
  \ErrorTok{|=======================================}           \ErrorTok{|}  \DecValTok{77}\NormalTok{\%}
  \SpecialCharTok{|}                                                        
  \ErrorTok{|=======================================}           \ErrorTok{|}  \DecValTok{78}\NormalTok{\%}
  \SpecialCharTok{|}                                                        
  \ErrorTok{|=======================================}           \ErrorTok{|}  \DecValTok{79}\NormalTok{\%}
  \SpecialCharTok{|}                                                        
  \ErrorTok{|========================================}          \ErrorTok{|}  \DecValTok{79}\NormalTok{\%}
  \SpecialCharTok{|}                                                        
  \ErrorTok{|========================================}          \ErrorTok{|}  \DecValTok{80}\NormalTok{\%}
  \SpecialCharTok{|}                                                        
  \ErrorTok{|========================================}          \ErrorTok{|}  \DecValTok{81}\NormalTok{\%}
  \SpecialCharTok{|}                                                        
  \ErrorTok{|=========================================}         \ErrorTok{|}  \DecValTok{81}\NormalTok{\%}
  \SpecialCharTok{|}                                                        
  \ErrorTok{|=========================================}         \ErrorTok{|}  \DecValTok{82}\NormalTok{\%}
  \SpecialCharTok{|}                                                        
  \ErrorTok{|=========================================}         \ErrorTok{|}  \DecValTok{83}\NormalTok{\%}
  \SpecialCharTok{|}                                                        
  \ErrorTok{|==========================================}        \ErrorTok{|}  \DecValTok{83}\NormalTok{\%}
  \SpecialCharTok{|}                                                        
  \ErrorTok{|==========================================}        \ErrorTok{|}  \DecValTok{84}\NormalTok{\%}
  \SpecialCharTok{|}                                                        
  \ErrorTok{|==========================================}        \ErrorTok{|}  \DecValTok{85}\NormalTok{\%}
  \SpecialCharTok{|}                                                        
  \ErrorTok{|===========================================}       \ErrorTok{|}  \DecValTok{85}\NormalTok{\%}
  \SpecialCharTok{|}                                                        
  \ErrorTok{|===========================================}       \ErrorTok{|}  \DecValTok{86}\NormalTok{\%}
  \SpecialCharTok{|}                                                        
  \ErrorTok{|===========================================}       \ErrorTok{|}  \DecValTok{87}\NormalTok{\%}
  \SpecialCharTok{|}                                                        
  \ErrorTok{|============================================}      \ErrorTok{|}  \DecValTok{87}\NormalTok{\%}
  \SpecialCharTok{|}                                                        
  \ErrorTok{|============================================}      \ErrorTok{|}  \DecValTok{88}\NormalTok{\%}
  \SpecialCharTok{|}                                                        
  \ErrorTok{|============================================}      \ErrorTok{|}  \DecValTok{89}\NormalTok{\%}
  \SpecialCharTok{|}                                                        
  \ErrorTok{|=============================================}     \ErrorTok{|}  \DecValTok{89}\NormalTok{\%}
  \SpecialCharTok{|}                                                        
  \ErrorTok{|=============================================}     \ErrorTok{|}  \DecValTok{90}\NormalTok{\%}
  \SpecialCharTok{|}                                                        
  \ErrorTok{|=============================================}     \ErrorTok{|}  \DecValTok{91}\NormalTok{\%}
  \SpecialCharTok{|}                                                        
  \ErrorTok{|==============================================}    \ErrorTok{|}  \DecValTok{91}\NormalTok{\%}
  \SpecialCharTok{|}                                                        
  \ErrorTok{|==============================================}    \ErrorTok{|}  \DecValTok{92}\NormalTok{\%}
  \SpecialCharTok{|}                                                        
  \ErrorTok{|==============================================}    \ErrorTok{|}  \DecValTok{93}\NormalTok{\%}
  \SpecialCharTok{|}                                                        
  \ErrorTok{|===============================================}   \ErrorTok{|}  \DecValTok{93}\NormalTok{\%}
  \SpecialCharTok{|}                                                        
  \ErrorTok{|===============================================}   \ErrorTok{|}  \DecValTok{94}\NormalTok{\%}
  \SpecialCharTok{|}                                                        
  \ErrorTok{|===============================================}   \ErrorTok{|}  \DecValTok{95}\NormalTok{\%}
  \SpecialCharTok{|}                                                        
  \ErrorTok{|================================================}  \ErrorTok{|}  \DecValTok{95}\NormalTok{\%}
  \SpecialCharTok{|}                                                        
  \ErrorTok{|================================================}  \ErrorTok{|}  \DecValTok{96}\NormalTok{\%}
  \SpecialCharTok{|}                                                        
  \ErrorTok{|================================================}  \ErrorTok{|}  \DecValTok{97}\NormalTok{\%}
  \SpecialCharTok{|}                                                        
  \ErrorTok{|=================================================} \ErrorTok{|}  \DecValTok{97}\NormalTok{\%}
  \SpecialCharTok{|}                                                        
  \ErrorTok{|=================================================} \ErrorTok{|}  \DecValTok{98}\NormalTok{\%}
  \SpecialCharTok{|}                                                        
  \ErrorTok{|=================================================} \ErrorTok{|}  \DecValTok{99}\NormalTok{\%}
  \SpecialCharTok{|}                                                        
  \ErrorTok{|==================================================|}  \DecValTok{99}\NormalTok{\%}
  \SpecialCharTok{|}                                                        
  \ErrorTok{|==================================================|} \DecValTok{100}\NormalTok{\%}
\FunctionTok{summary}\NormalTok{(fit.pgls)}
\CommentTok{\#\textgreater{} }
\CommentTok{\#\textgreater{} Analysis of Variance, using Residual Randomization}
\CommentTok{\#\textgreater{} Permutation procedure: Randomization of null model residuals }
\CommentTok{\#\textgreater{} Number of permutations: 1000 }
\CommentTok{\#\textgreater{} Estimation method: Generalized Least{-}Squares (via OLS projection) }
\CommentTok{\#\textgreater{} Sums of Squares and Cross{-}products: Type I }
\CommentTok{\#\textgreater{} Effect sizes (Z) based on F distributions}
\CommentTok{\#\textgreater{} }
\CommentTok{\#\textgreater{}           Df       SS       MS     Rsq      F       Z}
\CommentTok{\#\textgreater{} temp       1 0.072892 0.072892 0.26685 2.1839 0.93768}
\CommentTok{\#\textgreater{} Residuals  6 0.200262 0.033377 0.73315               }
\CommentTok{\#\textgreater{} Total      7 0.273153                                }
\CommentTok{\#\textgreater{}           Pr(\textgreater{}F)}
\CommentTok{\#\textgreater{} temp       0.202}
\CommentTok{\#\textgreater{} Residuals       }
\CommentTok{\#\textgreater{} Total           }
\CommentTok{\#\textgreater{} }
\CommentTok{\#\textgreater{} Call: procD.lm(f1 = shape.means \textasciitilde{} temp, iter = iter, seed = seed, RRPP = TRUE,  }
\CommentTok{\#\textgreater{}     SS.type = SS.type, effect.type = effect.type, int.first = int.first,  }
\CommentTok{\#\textgreater{}     Cov = Cov, data = data, print.progress = print.progress)}
\end{Highlighting}
\end{Shaded}

\section{Contrastes filogenéticos independentes}\label{contrastes-filogenuxe9ticos-independentes}

\begin{Shaded}
\begin{Highlighting}[]
\NormalTok{pic.size}\OtherTok{\textless{}{-}}\FunctionTok{pic}\NormalTok{(size.means1,tree)}
\NormalTok{pic.temp}\OtherTok{\textless{}{-}}\FunctionTok{pic}\NormalTok{(temp,tree)}
\NormalTok{fit.pic}\OtherTok{\textless{}{-}}\FunctionTok{lm}\NormalTok{(pic.size}\SpecialCharTok{\textasciitilde{}}\NormalTok{pic.temp}\SpecialCharTok{+}\DecValTok{0}\NormalTok{) }\CommentTok{\# regressão pela origem, a direção da subtração foi arbitrária no cálculo dos PICs}
\FunctionTok{summary.aov}\NormalTok{(fit.pic) }\CommentTok{\# mesmo resultado da PGLS}
\FunctionTok{plot}\NormalTok{(pic.size}\SpecialCharTok{\textasciitilde{}}\NormalTok{pic.temp)}
\end{Highlighting}
\end{Shaded}

\includegraphics{03.13-mgpcm_files/figure-latex/unnamed-chunk-10-1.pdf}

\section{Exercício - Morfometria e filogenias}\label{exercuxedcio---morfometria-e-filogenias}

Carregue o arquivo .tps `Roedores\_Digitalizados.tps' contendo coordenadas digitalizadas na mandíbula de 143 indivíduos de 19 espécies de roedores.
Carregue também a informação taxonômica de cada indivíduo (`Roedores\_classificador.txt'), e uma hipótese filogenética para as 19 espécies (`Roedores\_tree\_nexus.nex').
Existe um arquivo adicional que informa a força da mordida média de cada uma das 19 espécies (`Roedores\_BiteForce.txt').
Investigue a força do sinal filogenético para o tamanho e a forma do crânio, e descubra se existe relação entre a força da mordida e a forma (e tamanho) do crânio.

\part{PARTE IV - Macroevolução}\label{part-parte-iv---macroevoluuxe7uxe3o}

\chapter{Macroevolução}\label{macroevoluuxe7uxe3o}

Em construção.

\chapter{Lendo e manipulando árvores filogenéticas no R}\label{lendo-e-manipulando-uxe1rvores-filogenuxe9ticas-no-r}

É aconselhável criar um único diretório contendo a(s) árvore(s) e os dados fenotípicos, e informar ao R de qual diretório os arquivos devem ser lidos. Árvores filogenéticas podem ser importadas facilmente para dentro do ambiente R através de funções como read.tree e read.nexus do pacote \emph{ape}.

Neste capítulo, vamos explorar algumas maneiras de representar graficamente as árvores filogenéticas e manipulá-las. Para entender mais sobre a estrutura da árvore em formato phylo volte ao \hyperref[cap2]{Capítulo 2} do manual.

Vamos ler uma árvore filogenética do gênero de roedores neotropicais \emph{Akodon} e representá-la graficamente.

\begin{Shaded}
\begin{Highlighting}[]
\FunctionTok{require}\NormalTok{(ape)}
\CommentTok{\#\textgreater{} Carregando pacotes exigidos: ape}
\NormalTok{akodon.tree}\OtherTok{\textless{}{-}}\FunctionTok{read.tree}\NormalTok{(}\AttributeTok{file=}\StringTok{"dadospcm/akodon.tree"}\NormalTok{)}
\NormalTok{akodon.tree}
\CommentTok{\#\textgreater{} }
\CommentTok{\#\textgreater{} Phylogenetic tree with 31 tips and 30 internal nodes.}
\CommentTok{\#\textgreater{} }
\CommentTok{\#\textgreater{} Tip labels:}
\CommentTok{\#\textgreater{}   Castoria\_angustidens, Akodon\_mimus, Akodon\_lindberghi, Akodon\_subfuscus, Akodon\_lutescens, Akodon\_sylvanus, ...}
\CommentTok{\#\textgreater{} Node labels:}
\CommentTok{\#\textgreater{}   474, 475, 476, 477, 479, 480, ...}
\CommentTok{\#\textgreater{} }
\CommentTok{\#\textgreater{} Rooted; includes branch lengths.}
\CommentTok{\# Existem muitas formas de plotar a árvore}
\FunctionTok{plot}\NormalTok{(akodon.tree)}
\end{Highlighting}
\end{Shaded}

\includegraphics{04.1-pcmtrees_files/figure-latex/unnamed-chunk-1-1.pdf}

\begin{Shaded}
\begin{Highlighting}[]
\CommentTok{\# cladograma orientado pra esquerda}
\FunctionTok{plot}\NormalTok{(akodon.tree, }\AttributeTok{type =} \StringTok{"c"}\NormalTok{, }\AttributeTok{use.edge.length =} \ConstantTok{FALSE}\NormalTok{,}\AttributeTok{direction=}\StringTok{"l"}\NormalTok{)}
\end{Highlighting}
\end{Shaded}

\includegraphics{04.1-pcmtrees_files/figure-latex/unnamed-chunk-1-2.pdf}

\begin{Shaded}
\begin{Highlighting}[]
\CommentTok{\# nomes alinhados à esquerda}
\FunctionTok{plot}\NormalTok{(akodon.tree, }\AttributeTok{type =} \StringTok{"c"}\NormalTok{, }\AttributeTok{use.edge.length =} \ConstantTok{FALSE}\NormalTok{,}\AttributeTok{direction=}\StringTok{"l"}\NormalTok{,}\AttributeTok{adj=}\DecValTok{0}\NormalTok{)}
\end{Highlighting}
\end{Shaded}

\includegraphics{04.1-pcmtrees_files/figure-latex/unnamed-chunk-1-3.pdf}

\begin{Shaded}
\begin{Highlighting}[]
\CommentTok{\# visualizar partes da árvore com zoom}
\FunctionTok{zoom}\NormalTok{(akodon.tree, }\DecValTok{13}\SpecialCharTok{:}\DecValTok{17}\NormalTok{, }\AttributeTok{subtree =} \ConstantTok{FALSE}\NormalTok{, }\AttributeTok{cex =} \FloatTok{0.8}\NormalTok{)}
\end{Highlighting}
\end{Shaded}

\includegraphics{04.1-pcmtrees_files/figure-latex/unnamed-chunk-1-4.pdf}

\begin{Shaded}
\begin{Highlighting}[]
\FunctionTok{par}\NormalTok{(}\AttributeTok{mfrow=}\FunctionTok{c}\NormalTok{(}\DecValTok{1}\NormalTok{,}\DecValTok{1}\NormalTok{))}
\CommentTok{\# mudando cores}
\FunctionTok{plot}\NormalTok{(akodon.tree, }\AttributeTok{type =} \StringTok{"fan"}\NormalTok{, }\AttributeTok{edge.color =} \StringTok{"violetred"}\NormalTok{, }\AttributeTok{tip.color =} \StringTok{"turquoise1"}\NormalTok{,}
 \AttributeTok{cex =} \FloatTok{0.5}\NormalTok{)}
\end{Highlighting}
\end{Shaded}

\includegraphics{04.1-pcmtrees_files/figure-latex/unnamed-chunk-1-5.pdf}

\begin{Shaded}
\begin{Highlighting}[]
\CommentTok{\# com o pacote phytools}
\FunctionTok{require}\NormalTok{(phytools)}
\CommentTok{\#\textgreater{} Carregando pacotes exigidos: phytools}
\CommentTok{\#\textgreater{} Warning: package \textquotesingle{}phytools\textquotesingle{} was built under R version 4.3.2}
\CommentTok{\#\textgreater{} Carregando pacotes exigidos: maps}
\FunctionTok{plotTree}\NormalTok{(akodon.tree,}\AttributeTok{ftype=}\StringTok{"i"}\NormalTok{,}\AttributeTok{fsize=}\FloatTok{0.8}\NormalTok{,}\AttributeTok{lwd=}\DecValTok{1}\NormalTok{)}
\end{Highlighting}
\end{Shaded}

\includegraphics{04.1-pcmtrees_files/figure-latex/unnamed-chunk-1-6.pdf}

\begin{Shaded}
\begin{Highlighting}[]
\CommentTok{\# filograma arredondado}
\FunctionTok{roundPhylogram}\NormalTok{(akodon.tree)}
\end{Highlighting}
\end{Shaded}

\includegraphics{04.1-pcmtrees_files/figure-latex/unnamed-chunk-1-7.pdf}

\begin{Shaded}
\begin{Highlighting}[]
\CommentTok{\# cladograma desenraizado (unrooted)}
\FunctionTok{plot}\NormalTok{(}\FunctionTok{unroot}\NormalTok{(akodon.tree),}\AttributeTok{type=}\StringTok{"unrooted"}\NormalTok{,}\AttributeTok{cex=}\FloatTok{0.8}\NormalTok{,}\AttributeTok{use.edge.length=}\ConstantTok{FALSE}\NormalTok{,}\AttributeTok{lab4ut=}\StringTok{"axial"}\NormalTok{)}
\end{Highlighting}
\end{Shaded}

\includegraphics{04.1-pcmtrees_files/figure-latex/unnamed-chunk-1-8.pdf}

\begin{Shaded}
\begin{Highlighting}[]
\CommentTok{\# árvore circular (fan tree)}
\FunctionTok{plotTree}\NormalTok{(akodon.tree,}\AttributeTok{type=}\StringTok{"fan"}\NormalTok{,}\AttributeTok{fsize=}\FloatTok{0.8}\NormalTok{,}\AttributeTok{lwd=}\DecValTok{1}\NormalTok{,}\AttributeTok{ftype=}\StringTok{"i"}\NormalTok{)}
\end{Highlighting}
\end{Shaded}

\includegraphics{04.1-pcmtrees_files/figure-latex/unnamed-chunk-1-9.pdf}

\begin{Shaded}
\begin{Highlighting}[]
\CommentTok{\# Legenda dos comprimento de ramos}
\NormalTok{\{}\FunctionTok{plotTree}\NormalTok{(akodon.tree,}\AttributeTok{ftype=}\StringTok{"i"}\NormalTok{,}\AttributeTok{fsize=}\FloatTok{0.8}\NormalTok{,}\AttributeTok{lwd=}\DecValTok{1}\NormalTok{,}\AttributeTok{mar=}\FunctionTok{c}\NormalTok{(}\FloatTok{3.0}\NormalTok{,}\FloatTok{0.1}\NormalTok{,}\FloatTok{0.1}\NormalTok{,}\FloatTok{0.1}\NormalTok{))}
\FunctionTok{axisPhylo}\NormalTok{()\}}
\end{Highlighting}
\end{Shaded}

\includegraphics{04.1-pcmtrees_files/figure-latex/unnamed-chunk-1-10.pdf}

\section{Extraindo subclados ou espécies da árvore}\label{extraindo-subclados-ou-espuxe9cies-da-uxe1rvore}

Podemos remover ou manter apenas determinadas espécies ou subclados da árvore. Pra isso, as espécies devem estar identificadas pelos seus nomes (tip labels) ou os subclados pelo nome ou número do nó que identifica o ancestral comum mais recente do subclado. Abaixo estão exemplificadas as duas formas.

\begin{Shaded}
\begin{Highlighting}[]
\NormalTok{akodon.tree}\SpecialCharTok{$}\NormalTok{tip.label}
\CommentTok{\# Exemplo: identificar as espécies do gênero com ocorrência no sul do Brasil}
\NormalTok{sul.especies}\OtherTok{\textless{}{-}}
\FunctionTok{c}\NormalTok{(}\StringTok{"Akodon\_paranaensis"}\NormalTok{,}\StringTok{"Akodon\_azarae"}\NormalTok{,}\StringTok{"Akodon\_reigi"}\NormalTok{,}\StringTok{"Akodon\_montensis"}\NormalTok{,}\StringTok{"Cast}
\StringTok{oria\_angustidens"}\NormalTok{)}
\CommentTok{\# remover essas espécies da árvore}
\NormalTok{akodon.drop}\OtherTok{\textless{}{-}}\FunctionTok{drop.tip}\NormalTok{(akodon.tree,sul.especies)}
\NormalTok{akodon.drop}
\FunctionTok{plotTree}\NormalTok{(akodon.tree,}\AttributeTok{ftype=}\StringTok{"i"}\NormalTok{,}\AttributeTok{fsize=}\FloatTok{0.8}\NormalTok{,}\AttributeTok{lwd=}\DecValTok{1}\NormalTok{)}
\end{Highlighting}
\end{Shaded}

\includegraphics{04.1-pcmtrees_files/figure-latex/unnamed-chunk-2-1.pdf}

\begin{Shaded}
\begin{Highlighting}[]
\FunctionTok{plotTree}\NormalTok{(akodon.drop,}\AttributeTok{ftype=}\StringTok{"i"}\NormalTok{,}\AttributeTok{fsize=}\FloatTok{0.8}\NormalTok{,}\AttributeTok{lwd=}\DecValTok{1}\NormalTok{)}
\end{Highlighting}
\end{Shaded}

\includegraphics{04.1-pcmtrees_files/figure-latex/unnamed-chunk-2-2.pdf}

\begin{Shaded}
\begin{Highlighting}[]
\CommentTok{\# ou manter apenas essas espécies na árvore}
\NormalTok{akodon.drop1}\OtherTok{\textless{}{-}}\FunctionTok{drop.tip}\NormalTok{(akodon.tree,}\FunctionTok{setdiff}\NormalTok{(akodon.tree}\SpecialCharTok{$}\NormalTok{tip.label,sul.especies))}
\NormalTok{akodon.drop1}
\FunctionTok{plotTree}\NormalTok{(akodon.drop1,}\AttributeTok{ftype=}\StringTok{"i"}\NormalTok{,}\AttributeTok{fsize=}\FloatTok{0.8}\NormalTok{,}\AttributeTok{lwd=}\DecValTok{1}\NormalTok{)}
\end{Highlighting}
\end{Shaded}

\includegraphics{04.1-pcmtrees_files/figure-latex/unnamed-chunk-2-3.pdf}

\begin{Shaded}
\begin{Highlighting}[]
\CommentTok{\# extrair um subclado da árvore}
\NormalTok{\{}\FunctionTok{plotTree}\NormalTok{(akodon.tree,}\AttributeTok{ftype=}\StringTok{"i"}\NormalTok{,}\AttributeTok{fsize=}\FloatTok{0.8}\NormalTok{,}\AttributeTok{lwd=}\DecValTok{1}\NormalTok{)}
\FunctionTok{nodelabels}\NormalTok{()\}}
\end{Highlighting}
\end{Shaded}

\includegraphics{04.1-pcmtrees_files/figure-latex/unnamed-chunk-2-4.pdf}

\begin{Shaded}
\begin{Highlighting}[]
\NormalTok{clado1}\OtherTok{\textless{}{-}}\FunctionTok{extract.clade}\NormalTok{(akodon.tree,}\AttributeTok{node=}\DecValTok{35}\NormalTok{)}
\FunctionTok{plotTree}\NormalTok{(clado1,}\AttributeTok{ftype=}\StringTok{"i"}\NormalTok{,}\AttributeTok{fsize=}\FloatTok{0.8}\NormalTok{,}\AttributeTok{lwd=}\DecValTok{1}\NormalTok{)}
\end{Highlighting}
\end{Shaded}

\includegraphics{04.1-pcmtrees_files/figure-latex/unnamed-chunk-2-5.pdf}

\begin{Shaded}
\begin{Highlighting}[]
\CommentTok{\# remover o clado}
\NormalTok{akodon.noclado1}\OtherTok{\textless{}{-}}\FunctionTok{drop.tip}\NormalTok{(akodon.tree,clado1}\SpecialCharTok{$}\NormalTok{tip.label)}
\NormalTok{akodon.noclado1}
\FunctionTok{plotTree}\NormalTok{(akodon.noclado1,}\AttributeTok{ftype=}\StringTok{"i"}\NormalTok{,}\AttributeTok{fsize=}\FloatTok{0.8}\NormalTok{,}\AttributeTok{lwd=}\DecValTok{1}\NormalTok{)}
\end{Highlighting}
\end{Shaded}

\includegraphics{04.1-pcmtrees_files/figure-latex/unnamed-chunk-2-6.pdf}

\begin{Shaded}
\begin{Highlighting}[]
\CommentTok{\# função interativa do phytools para remover clados}
\NormalTok{akodon.col}\OtherTok{\textless{}{-}}\FunctionTok{collapseTree}\NormalTok{(akodon.tree)}
\end{Highlighting}
\end{Shaded}

\includegraphics{04.1-pcmtrees_files/figure-latex/unnamed-chunk-2-7.pdf}

\begin{Shaded}
\begin{Highlighting}[]
\FunctionTok{plotTree}\NormalTok{(akodon.col,}\AttributeTok{ftype=}\StringTok{"i"}\NormalTok{,}\AttributeTok{fsize=}\FloatTok{0.8}\NormalTok{,}\AttributeTok{lwd=}\DecValTok{1}\NormalTok{)}
\end{Highlighting}
\end{Shaded}

\includegraphics{04.1-pcmtrees_files/figure-latex/unnamed-chunk-2-8.pdf}

\section{Incerteza filogenética}\label{incerteza-filogenuxe9tica}

Uma árvore filogenética é uma hipótese sobre as relações de parentesco entre aquelas espécies. Assim, é recomendável avaliar a incerteza resultante do uso de hipóteses alternativas. Uma das maneiras de fazer isso é repetindo as análises de interesse para um conjunto de árvores (hipóteses).

\subsection{Politomias}\label{politomias}

Uma das fontes de incerteza filogenética são as politomias. Algumas das funções implementando métodos comparativos no R, que veremos, vão aceitar ávores com politomias, outras funções não vão. Podemos considerar incerteza filogenética ao resolver aleatoriamente as politomias um grande número de vezes, e considerar todas as hipóteses alternativas nos nossos testes estatísticos. Uma das maneiras mais simples de resolver politomias no R é usando a função multi2di do pacote \emph{ape}.

\begin{Shaded}
\begin{Highlighting}[]
\NormalTok{ex}\OtherTok{\textless{}{-}}\FunctionTok{read.tree}\NormalTok{(}\AttributeTok{text=}\StringTok{"((A,B,C),(D,E));"}\NormalTok{)}
\FunctionTok{plot}\NormalTok{(ex)}
\end{Highlighting}
\end{Shaded}

\includegraphics{04.1-pcmtrees_files/figure-latex/unnamed-chunk-3-1.pdf}

\begin{Shaded}
\begin{Highlighting}[]
\FunctionTok{is.binary}\NormalTok{(ex)}
\CommentTok{\# resolver ao acaso uma única vez}
\NormalTok{ex1}\OtherTok{\textless{}{-}}\FunctionTok{multi2di}\NormalTok{(ex)}
\FunctionTok{plot}\NormalTok{(ex1)}
\end{Highlighting}
\end{Shaded}

\includegraphics{04.1-pcmtrees_files/figure-latex/unnamed-chunk-3-2.pdf}

\begin{Shaded}
\begin{Highlighting}[]
\FunctionTok{is.binary}\NormalTok{(ex1)}
\end{Highlighting}
\end{Shaded}

Lembre-se que ao resolver apenas 1 vez ao acaso não existe avaliação de incerteza filogenética, é preciso repetir o processo um determinado número de vezes, que vai depender de quantas instâncias de politomia estiverem presentes na árvore, e quantas espécies estiverem nas politomias. Existem outros métodos que permitem gerar múltiplas árvores com as politomias resolvidas ao acaso. Por exemplo: \href{https://bioinfo.inf.ufg.br/sunplin/}{sunplin}- Rangel et al. (2015).

\subsection{Multiplas árvores}\label{multiplas-uxe1rvores}

Múltiplas filogenias são úteis quando queremos incorporar incerteza filogenética, por exemplo, replicando as mesmas análises em múltiplas árvores.
Árvores múltiplas pertencem ao objeto de classe multiPhylo, que é uma lista de filogenias da classe phylo. Muitas funções podem prontamente ser aplicadas tanto para phylo como para multiPhylo. Algumas vão exigir a replicação em cada filogenia separada.
Vamos juntar as árvores que criamos até aqui para exemplificar um objeto multiPhylo.

\begin{Shaded}
\begin{Highlighting}[]
\NormalTok{akodon.trees}\OtherTok{\textless{}{-}}\FunctionTok{c}\NormalTok{(akodon.tree,akodon.drop,akodon.drop1,clado1,akodon.noclado1)}
\NormalTok{akodon.trees}
\FunctionTok{class}\NormalTok{(akodon.trees)}
\FunctionTok{print}\NormalTok{(akodon.trees,}\AttributeTok{details=}\ConstantTok{TRUE}\NormalTok{)}
\FunctionTok{str}\NormalTok{(akodon.trees[}\DecValTok{1}\NormalTok{])}
\end{Highlighting}
\end{Shaded}

Um exemplo para uso de funções em múltiplas árvores, uma por vez, exemplificado com o recorte de uma parte da árvore.

\begin{Shaded}
\begin{Highlighting}[]
\CommentTok{\# Extraindo clados de múltiplas árvores}
\NormalTok{trees}\OtherTok{\textless{}{-}}\FunctionTok{pbtree}\NormalTok{(}\AttributeTok{n=}\DecValTok{50}\NormalTok{,}\AttributeTok{nsim=}\DecValTok{100}\NormalTok{)}
\NormalTok{trees}
\FunctionTok{plot}\NormalTok{(trees[[}\DecValTok{1}\NormalTok{]])}
\end{Highlighting}
\end{Shaded}

\includegraphics{04.1-pcmtrees_files/figure-latex/unnamed-chunk-5-1.pdf}

\begin{Shaded}
\begin{Highlighting}[]
\NormalTok{tree\_cortado}\OtherTok{\textless{}{-}} \FunctionTok{vector}\NormalTok{(}\AttributeTok{mode =} \StringTok{"list"}\NormalTok{, }\FunctionTok{length}\NormalTok{(trees))}
\ControlFlowTok{for}\NormalTok{(i }\ControlFlowTok{in} \DecValTok{1}\SpecialCharTok{:}\FunctionTok{length}\NormalTok{(trees))\{}
\NormalTok{ tree\_cortado[[i]]}\OtherTok{\textless{}{-}}\FunctionTok{keep.tip}\NormalTok{(trees[[i]], }\AttributeTok{tip=}\NormalTok{trees[[}\DecValTok{1}\NormalTok{]]}\SpecialCharTok{$}\NormalTok{tip.label[}\DecValTok{1}\SpecialCharTok{:}\DecValTok{12}\NormalTok{])}
\NormalTok{\}}
\NormalTok{tree\_cortado}
\FunctionTok{plot}\NormalTok{(tree\_cortado[[}\DecValTok{1}\NormalTok{]])}
\end{Highlighting}
\end{Shaded}

\includegraphics{04.1-pcmtrees_files/figure-latex/unnamed-chunk-5-2.pdf}

\chapter{Sinal filogenético}\label{sinal-filogenuxe9tico}

Sinal filogenético é meramente um padrão que mostra se espécies mais próximas filogenéticamente têm fenótipos mais similares do que espécies mais distantes filogenéticamente. A magnitude do sinal filogenético é, portanto, útil para percebermos quão forte é essa tendência de espécies similares fenotipicamente estarem próximas filogenéticamente.

Existem vários métodos para calcular sinal filogenético. Vamos explorar aqui a estatística K Blomberg, Garland JR., and Ives (2003), o lambda de Pagel Pagel (1999), e a curva PSR Diniz-Filho et al. (2012a).

Vamos investigar a magnitude do sinal filogenético no comprimento do crânio de um gênero de roedores sul-americanos.

\begin{Shaded}
\begin{Highlighting}[]
\CommentTok{\# Carregar dados fenotípicos}
\NormalTok{dados}\OtherTok{\textless{}{-}}\FunctionTok{read.table}\NormalTok{(}\StringTok{"dadospcm/akodon.txt"}\NormalTok{,}\AttributeTok{h=}\NormalTok{T)}
\NormalTok{dados}
\CommentTok{\# A função attach permite obter cada coluna dos dados}
\FunctionTok{attach}\NormalTok{(dados)}
\CommentTok{\# Vamos (re)atribuir os nomes das espécies}
\NormalTok{skull\_length}\OtherTok{\textless{}{-}}\FunctionTok{setNames}\NormalTok{(skull\_length,species)}
\NormalTok{skull\_length}

\CommentTok{\# Carregar árvore}
\FunctionTok{require}\NormalTok{(ape)}
\CommentTok{\#\textgreater{} Carregando pacotes exigidos: ape}
\NormalTok{tree}\OtherTok{\textless{}{-}}\FunctionTok{read.tree}\NormalTok{(}\StringTok{"dadospcm/akodon.tree"}\NormalTok{)}
\NormalTok{\{}\FunctionTok{plot}\NormalTok{(tree,}\AttributeTok{cex=}\FloatTok{0.4}\NormalTok{)}
\FunctionTok{axisPhylo}\NormalTok{()\}}
\end{Highlighting}
\end{Shaded}

\includegraphics{04.2-pcmsinal_files/figure-latex/unnamed-chunk-1-1.pdf}

É importante conferir a correspondência entre as espécies na filogenia e nos dados. Todas as espécies precisam estar tanto na árvore como nos dados, e também tem que existir correspondência exata nos nomes. Existem algumas funções, como a name.check e a treedata, do pacote \emph{geiger}, que fazem essa conferência.

\begin{Shaded}
\begin{Highlighting}[]
\FunctionTok{require}\NormalTok{(geiger)}
\CommentTok{\#\textgreater{} Carregando pacotes exigidos: geiger}
\CommentTok{\#\textgreater{} Carregando pacotes exigidos: phytools}
\CommentTok{\#\textgreater{} Carregando pacotes exigidos: maps}
\FunctionTok{name.check}\NormalTok{(tree,dados,}\AttributeTok{data.names=}\NormalTok{species)}
\CommentTok{\#\textgreater{} [1] "OK"}
\end{Highlighting}
\end{Shaded}

É importante visualizar a distribuição no atributo na filogenia.

\begin{Shaded}
\begin{Highlighting}[]
\FunctionTok{require}\NormalTok{(phytools)}
\FunctionTok{plotTree.barplot}\NormalTok{(tree,skull\_length,}
 \AttributeTok{args.plotTree=}\FunctionTok{list}\NormalTok{(}\AttributeTok{fsize=}\FloatTok{0.5}\NormalTok{,}\AttributeTok{ftype=}\StringTok{"i"}\NormalTok{))}
\end{Highlighting}
\end{Shaded}

\includegraphics{04.2-pcmsinal_files/figure-latex/unnamed-chunk-3-1.pdf}

\section{Estatística K}\label{estatuxedstica-k}

K=1 indica que espécies próximas são tão parecidas quanto prevê o modelo Browniano de evolução K\textless1 indica sinal filogenético menor do que esperado pelo modelo Browniano de evolução K\textgreater1 indica sinal filogenético maior do que esperado pelo modelo Browniano de evolução

\begin{Shaded}
\begin{Highlighting}[]
\FunctionTok{require}\NormalTok{(phytools)}
\FunctionTok{phylosig}\NormalTok{(tree,skull\_length,}\AttributeTok{method =} \StringTok{"K"}\NormalTok{,}\AttributeTok{test =} \ConstantTok{TRUE}\NormalTok{)}
\CommentTok{\#\textgreater{} }
\CommentTok{\#\textgreater{} Phylogenetic signal K : 0.336983 }
\CommentTok{\#\textgreater{} P{-}value (based on 1000 randomizations) : 0.646}
\end{Highlighting}
\end{Shaded}

\section{Lambda de Pagel}\label{lambda-de-pagel}

λ=0 indica ausência de sinal filogenético (sinal esperado em uma filogenia estrela) λ=1 indica sinal filogenético correspondente ao modelo Browniano

\begin{Shaded}
\begin{Highlighting}[]
\FunctionTok{phylosig}\NormalTok{(tree,skull\_length,}\AttributeTok{method =} \StringTok{"lambda"}\NormalTok{,}\AttributeTok{test =} \ConstantTok{TRUE}\NormalTok{)}
\CommentTok{\#\textgreater{} }
\CommentTok{\#\textgreater{} Phylogenetic signal lambda : 7.95528e{-}05 }
\CommentTok{\#\textgreater{} logL(lambda) : {-}42.2205 }
\CommentTok{\#\textgreater{} LR(lambda=0) : {-}0.000914493 }
\CommentTok{\#\textgreater{} P{-}value (based on LR test) : 1}
\end{Highlighting}
\end{Shaded}

\section{Curva PSR}\label{curva-psr}

A regressão de autovetores filogenéticos (\emph{phylogenetic eigenvector regression} - PVR) (Diniz-Filho, de Sant'Ana, and Bini 1998) pode ser usada para estimar a quantidade de sinal filogenético nos dados, além de permitir a partição da variância entre a filogenia e demais causas, como veremos no capítulo sobre regressão filogenética. A PVR funciona com uma regressão do fenótipo contra um conjunto de autovetores filogenéticos extraídos por uma análise de coordenadas principais de uma matriz de distâncias filogenéticas. O R² da regressão entre o fenótipo e os autovetores filogenéticos funciona como uma medida de sinal filogenético. No entanto, se todos os autovetores forem usados, o R² inevitavelmente alcança 1, e se alguns autovetores forem omitidos, parte da filogenia é ignorada (Rohlf 2001). No entanto, Diniz-Filho et al. (2012a) descobriu que se PVRs sucessivas forem feitas, adicionando um autovetor filogenético por vez como preditor, e depois os R²s de cada regressão forem plotados contra os autovalores cumulativos (\%), uma relação linear (linha de 45º) é obtida sob um modelo Browniano de evolução. Essa abordagem, chamada de PSR (phylogenetic signal-representation) permite estimar o sinal filogenético e oferece a habilidade de visualizar o padrão de acúmulo de variação do fenótipo (acelaração e desaceleração de taxas) ao longo do tempo, em comparação com o modelo Browniano de evolução (Diniz-Filho et al. 2012b) e outros modelos, como OU (Maestri et al. 2017).

O exemplo abaixo calcula e representa graficamente a curva PSR.

\begin{Shaded}
\begin{Highlighting}[]
\CommentTok{\#install.packages("PVR")}
\FunctionTok{require}\NormalTok{(PVR)}
\CommentTok{\#\textgreater{} Carregando pacotes exigidos: PVR}
\CommentTok{\#\textgreater{} PVR 0.3 loaded}
\CommentTok{\# Decomposição de matriz de distância filogenética}
\NormalTok{pvr\_obj}\OtherTok{\textless{}{-}}\FunctionTok{PVRdecomp}\NormalTok{(tree)}
\CommentTok{\# Regressão de autovetores filogenéticos}
\NormalTok{pvr\_reg}\OtherTok{\textless{}{-}}\FunctionTok{PVR}\NormalTok{(pvr\_obj,tree,skull\_length)}
\CommentTok{\# Sinal Filogenético com Cuva PSR}
\NormalTok{psr\_obj}\OtherTok{\textless{}{-}}\FunctionTok{PSR}\NormalTok{(pvr\_obj,}\AttributeTok{trait=}\NormalTok{skull\_length,}
\AttributeTok{null.model=}\ConstantTok{TRUE}\NormalTok{,}\AttributeTok{Brownian.model=}\ConstantTok{TRUE}\NormalTok{,}\AttributeTok{times=}\DecValTok{10}\NormalTok{)}
\NormalTok{psr\_obj}
\CommentTok{\#PSRplot(psr\_obj,info="both")}
\end{Highlighting}
\end{Shaded}

É comum encontrar menções de que regressões ou outras análises foram corrigidas pela filogenia porque um sinal filogenético foi encontrado nas variáveis brutas. No entanto, o uso de regressão filogenética deve estar associado aos resíduos do modelo de interesse apresentarem sinal filogenético (Revell 2010).

\section{Exercício - Árvores e sinal filogenético}\label{exercuxedcio---uxe1rvores-e-sinal-filogenuxe9tico}

Carregue uma árvore filogenética para o seu grupo de estudo e observe a estrutura dela. Represente a árvore graficamente de diferentes formas. Se você tiver variáveis fenotípicas para o seu grupo de estudo, tente carregar no R e testar o sinal filogenético.

\chapter{Modelos para evolução de fenótipos contínuos}\label{modelos-para-evoluuxe7uxe3o-de-fenuxf3tipos-contuxednuos}

Existem vários modelos de evolução de atributos contínuos univariados. Os modelos macroevolutivos mais comuns são o Browniano (BM), o modelo de Ornstein-Uhlenbeck (OU) e o modelo Early-Burst (EB). Sobre ajuste de modelos evolutivos, leia Butler and King (2004).

No exemplo, a função fitContinuous do pacote \emph{geiger} vai ser usada para ajustar os modelos. O método é baseado em verossimilhança (\emph{likelihood}), e os resultados vão retornar o valor de máxima verossimilhança (\emph{maximum likelihood}) como estimativa para os parâmetros. Diferentes modelos podem ser comparados, desde que sobre os mesmos dados. Vamos carregar vários atributos fenotípicos para 77 espécies de Primatas.

\begin{Shaded}
\begin{Highlighting}[]
\FunctionTok{require}\NormalTok{(phytools)}
\CommentTok{\#\textgreater{} Carregando pacotes exigidos: phytools}
\CommentTok{\#\textgreater{} Carregando pacotes exigidos: ape}
\CommentTok{\#\textgreater{} Carregando pacotes exigidos: maps}
\FunctionTok{require}\NormalTok{(geiger)}
\CommentTok{\#\textgreater{} Carregando pacotes exigidos: geiger}
\CommentTok{\# Carregar dados}
\NormalTok{dados}\OtherTok{\textless{}{-}}\FunctionTok{read.table}\NormalTok{(}\StringTok{"dadospcm/primate{-}data.txt"}\NormalTok{,}\AttributeTok{h=}\NormalTok{T,}\AttributeTok{row.names=}\DecValTok{1}\NormalTok{)}
\NormalTok{dados}
\CommentTok{\# Carregar filogenia}
\NormalTok{tree}\OtherTok{\textless{}{-}}\FunctionTok{read.tree}\NormalTok{(}\StringTok{"dadospcm/primate{-}tree.txt"}\NormalTok{)}
\FunctionTok{plot}\NormalTok{(tree)}
\end{Highlighting}
\end{Shaded}

\includegraphics{04.3-pcmcont_files/figure-latex/unnamed-chunk-1-1.pdf}

\begin{Shaded}
\begin{Highlighting}[]
\FunctionTok{plotTree}\NormalTok{(tree,}\AttributeTok{fsize=}\FloatTok{0.4}\NormalTok{,}\AttributeTok{ftype=}\StringTok{"i"}\NormalTok{,}\AttributeTok{type=}\StringTok{"fan"}\NormalTok{,}\AttributeTok{lwd=}\DecValTok{1}\NormalTok{)}
\end{Highlighting}
\end{Shaded}

\includegraphics{04.3-pcmcont_files/figure-latex/unnamed-chunk-1-2.pdf}

Conferir correspondência entre espécies nos dados e na filogenia.

\begin{Shaded}
\begin{Highlighting}[]
\NormalTok{match.species}\OtherTok{\textless{}{-}}\FunctionTok{treedata}\NormalTok{(tree,dados)}
\FunctionTok{names}\NormalTok{(match.species)}
\NormalTok{tree}\OtherTok{\textless{}{-}}\NormalTok{match.species}\SpecialCharTok{$}\NormalTok{phy}
\FunctionTok{str}\NormalTok{(tree)}
\FunctionTok{plotTree}\NormalTok{(tree,}\AttributeTok{fsize=}\FloatTok{0.5}\NormalTok{,}\AttributeTok{ftype=}\StringTok{"i"}\NormalTok{,}\AttributeTok{type=}\StringTok{"fan"}\NormalTok{,}\AttributeTok{lwd=}\DecValTok{1}\NormalTok{)}
\end{Highlighting}
\end{Shaded}

\includegraphics{04.3-pcmcont_files/figure-latex/unnamed-chunk-2-1.pdf}

\begin{Shaded}
\begin{Highlighting}[]
\NormalTok{bodymass}\OtherTok{\textless{}{-}}\NormalTok{dados[,}\DecValTok{1}\NormalTok{]}
\NormalTok{bodymass}\OtherTok{\textless{}{-}}\FunctionTok{setNames}\NormalTok{(bodymass,}\FunctionTok{rownames}\NormalTok{(dados))}
\NormalTok{bodymass}
\end{Highlighting}
\end{Shaded}

\section{Modelos evolutivos contínuos}\label{modelos-evolutivos-contuxednuos}

\subsection{Modelo Browniano}\label{modelo-browniano}

O modelo Browniano estima a taxa de evolução σ2 (sigsq no output), e o valor do atributo na raiz X(0) (z0 no output). Outras estimativas do output são o likelihood (lnL) do modelo, o AIC/AICc e o número de parâmetros (k na literatura) do modelo.

\begin{Shaded}
\begin{Highlighting}[]
\NormalTok{BM}\OtherTok{\textless{}{-}}\FunctionTok{fitContinuous}\NormalTok{(tree,}\FunctionTok{log}\NormalTok{(bodymass),}\AttributeTok{model=}\StringTok{"BM"}\NormalTok{)}
\NormalTok{BM}
\end{Highlighting}
\end{Shaded}

\subsection{Modelo de Ornstein-Uhlenbeck}\label{modelo-de-ornstein-uhlenbeck}

O modelo OU estima um parâmetro adicional, α ou alpha, que é a força de atração para a média do ótimo, μ ou θ.

\begin{Shaded}
\begin{Highlighting}[]
\NormalTok{OU}\OtherTok{\textless{}{-}}\FunctionTok{fitContinuous}\NormalTok{(tree,}\FunctionTok{log}\NormalTok{(bodymass),}\AttributeTok{model=}\StringTok{"OU"}\NormalTok{)}
\NormalTok{OU}
\end{Highlighting}
\end{Shaded}

\subsection{Modelo de Early-Burst/Late-Burst}\label{modelo-de-early-burstlate-burst}

No modelo EB o parâmetro \emph{a} é a taxa de mudança através do tempo. O valor de \emph{a} é normalmente menor ou igual a 0. Quando a é negativo, as taxas de evolução decrescem com o tempo. Se a é próximo de 0, a evolução é aproximadamente Browniana.

\begin{Shaded}
\begin{Highlighting}[]
\NormalTok{EB}\OtherTok{\textless{}{-}}\FunctionTok{fitContinuous}\NormalTok{(tree,}\FunctionTok{log}\NormalTok{(bodymass),}\AttributeTok{model=}\StringTok{"EB"}\NormalTok{)}
\NormalTok{EB}
\end{Highlighting}
\end{Shaded}

\subsection{Comparação de modelos}\label{comparauxe7uxe3o-de-modelos}

Nós podemos comparar o ajuste dos modelos usando um critério de informação, como o AIC. Existe debate sobre quão grande deve ser a diferença no AIC para definir um único ``melhor'' modelo, variando de 2-10 unidades. É comum usar 4 unidades, então EB não tem um ajuste muito melhor do que BM ou OU nos dados do exemplo.

\begin{Shaded}
\begin{Highlighting}[]
\NormalTok{BM}\SpecialCharTok{$}\NormalTok{opt}\SpecialCharTok{$}\NormalTok{aicc}
\NormalTok{OU}\SpecialCharTok{$}\NormalTok{opt}\SpecialCharTok{$}\NormalTok{aicc}
\NormalTok{EB}\SpecialCharTok{$}\NormalTok{opt}\SpecialCharTok{$}\NormalTok{aicc}
\NormalTok{aic.scores}\OtherTok{\textless{}{-}}\FunctionTok{setNames}\NormalTok{(}\FunctionTok{c}\NormalTok{(BM}\SpecialCharTok{$}\NormalTok{opt}\SpecialCharTok{$}\NormalTok{aicc,OU}\SpecialCharTok{$}\NormalTok{opt}\SpecialCharTok{$}\NormalTok{aicc,EB}\SpecialCharTok{$}\NormalTok{opt}\SpecialCharTok{$}\NormalTok{aicc),}
 \FunctionTok{c}\NormalTok{(}\StringTok{"BM"}\NormalTok{,}\StringTok{"OU"}\NormalTok{,}\StringTok{"EB"}\NormalTok{))}
\NormalTok{aic.scores}
\FunctionTok{aicw}\NormalTok{(aic.scores) }\CommentTok{\# Delta AICc e AICc weights}
\end{Highlighting}
\end{Shaded}

\section{Reconstrução ancestral de dados contínuos}\label{reconstruuxe7uxe3o-ancestral-de-dados-contuxednuos}

Vamos estimar o tamanho corporal ancestral para primatas, sob um modelo Browniano de evolução.
As funções ace (ape) e fastAnc (phytools), entre outras, implementam reconstrução ancestral via \emph{maximum likelihood} (encontrando os estados com verossimilhança máxima de terem surgido sob o modelo escolhido).

\begin{Shaded}
\begin{Highlighting}[]
\NormalTok{rec}\OtherTok{\textless{}{-}}\FunctionTok{fastAnc}\NormalTok{(tree,}\FunctionTok{log}\NormalTok{(bodymass),}\AttributeTok{vars=}\ConstantTok{TRUE}\NormalTok{,}\AttributeTok{CI=}\ConstantTok{TRUE}\NormalTok{)}
\NormalTok{rec}
\FunctionTok{print}\NormalTok{(rec,}\AttributeTok{printlen=}\DecValTok{10}\NormalTok{)}
\NormalTok{\{}\FunctionTok{plotTree}\NormalTok{(tree,}\AttributeTok{fsize=}\FloatTok{0.5}\NormalTok{,}\AttributeTok{ftype=}\StringTok{"i"}\NormalTok{,}\AttributeTok{type=}\StringTok{"fan"}\NormalTok{,}\AttributeTok{lwd=}\DecValTok{1}\NormalTok{)}
\FunctionTok{nodelabels}\NormalTok{()\}}
\end{Highlighting}
\end{Shaded}

\includegraphics{04.3-pcmcont_files/figure-latex/unnamed-chunk-7-1.pdf}

\begin{Shaded}
\begin{Highlighting}[]
\CommentTok{\# Visualizando a reconstrução}
\NormalTok{rec.map}\OtherTok{\textless{}{-}}\FunctionTok{contMap}\NormalTok{(tree,}\FunctionTok{log}\NormalTok{(bodymass))}
\end{Highlighting}
\end{Shaded}

\includegraphics{04.3-pcmcont_files/figure-latex/unnamed-chunk-7-2.pdf}

\begin{Shaded}
\begin{Highlighting}[]
\NormalTok{rec.map}
\CommentTok{\# customizando}
\NormalTok{rec.map}\OtherTok{\textless{}{-}}\FunctionTok{setMap}\NormalTok{(rec.map,}\AttributeTok{invert=}\ConstantTok{TRUE}\NormalTok{) }\CommentTok{\# inverte cores}
\FunctionTok{plot}\NormalTok{(rec.map,}\AttributeTok{type=}\StringTok{"fan"}\NormalTok{,}\AttributeTok{fsize=}\FloatTok{0.5}\NormalTok{)}
\end{Highlighting}
\end{Shaded}

\includegraphics{04.3-pcmcont_files/figure-latex/unnamed-chunk-7-3.pdf}

\begin{Shaded}
\begin{Highlighting}[]
\CommentTok{\# Modo fenograma \textquotesingle{}traitgram\textquotesingle{}}
\FunctionTok{par}\NormalTok{(}\AttributeTok{mfrow =} \FunctionTok{c}\NormalTok{(}\DecValTok{1}\NormalTok{, }\DecValTok{1}\NormalTok{),}\AttributeTok{mar=}\FunctionTok{c}\NormalTok{(}\DecValTok{4}\NormalTok{, }\DecValTok{4}\NormalTok{, }\DecValTok{2}\NormalTok{, }\DecValTok{2}\NormalTok{))}
\NormalTok{phe}\OtherTok{\textless{}{-}}\FunctionTok{phenogram}\NormalTok{(tree,}\FunctionTok{log}\NormalTok{(bodymass),}\AttributeTok{ftype=}\StringTok{"i"}\NormalTok{,}\AttributeTok{fsize=}\FloatTok{0.5}\NormalTok{)}
\end{Highlighting}
\end{Shaded}

\includegraphics{04.3-pcmcont_files/figure-latex/unnamed-chunk-7-4.pdf}

\begin{Shaded}
\begin{Highlighting}[]
\CommentTok{\# Incerteza (CI 95\%) nas estimativas dos nós}
\FunctionTok{fancyTree}\NormalTok{(tree,}\AttributeTok{type=}\StringTok{"phenogram95"}\NormalTok{,}\AttributeTok{x=}\FunctionTok{log}\NormalTok{(bodymass),}\AttributeTok{ftype=}\StringTok{"i"}\NormalTok{,}\AttributeTok{fsize=}\FloatTok{0.5}\NormalTok{) }\CommentTok{\# pode demorar}
\end{Highlighting}
\end{Shaded}

\includegraphics{04.3-pcmcont_files/figure-latex/unnamed-chunk-7-5.pdf}

\begin{Shaded}
\begin{Highlighting}[]
\CommentTok{\# Adicionando barra de erros com cores}
\NormalTok{rec.map}\OtherTok{\textless{}{-}}\FunctionTok{contMap}\NormalTok{(tree,}\FunctionTok{log}\NormalTok{(bodymass),}\AttributeTok{fsize=}\FloatTok{0.5}\NormalTok{,}\AttributeTok{ftype=}\StringTok{"i"}\NormalTok{)}
\FunctionTok{errorbar.contMap}\NormalTok{(rec.map)}
\end{Highlighting}
\end{Shaded}

\includegraphics{04.3-pcmcont_files/figure-latex/unnamed-chunk-7-6.pdf}

Uma crítica comum a reconstrução ancestral é que a incerteza ao redor do nós pode ser grande. No entanto, se o modelo evolutivo está correto, as estimativas são confiáveis (95\% de chance de o valor estimado incluir o valor `real'). Em teoria, é possível corrigir a estimativa se o valor de algum nó é conhecido (e.g., fóssil), como no exemplo abaixo.

\begin{Shaded}
\begin{Highlighting}[]
\CommentTok{\# Reconstrução com valores de nós conhecidos}
\NormalTok{rec}\SpecialCharTok{$}\NormalTok{ace }\CommentTok{\# estados ancestrais estimados}
\NormalTok{anc.states}\OtherTok{\textless{}{-}}\NormalTok{rec}\SpecialCharTok{$}\NormalTok{ace[}\FunctionTok{as.character}\NormalTok{(}\FunctionTok{c}\NormalTok{(}\DecValTok{78}\NormalTok{,}\DecValTok{80}\NormalTok{,}\DecValTok{82}\NormalTok{))]}
\NormalTok{anc.states }\CommentTok{\# valores estimados para nós 78 (raiz), 80 e 82}
\NormalTok{anc.a}\OtherTok{\textless{}{-}}\FunctionTok{c}\NormalTok{(}\FloatTok{4.2}\NormalTok{,}\FloatTok{5.1}\NormalTok{,}\FloatTok{5.3}\NormalTok{) }\CommentTok{\# valores conhecidos}
\FunctionTok{names}\NormalTok{(anc.a)}\OtherTok{=}\FunctionTok{names}\NormalTok{(anc.states)}
\NormalTok{anc.a}
\NormalTok{rec}\FloatTok{.1}\OtherTok{\textless{}{-}}\FunctionTok{fastAnc}\NormalTok{(tree,}\AttributeTok{x=}\FunctionTok{log}\NormalTok{(bodymass),}\AttributeTok{anc.states=}\NormalTok{anc.a,}\AttributeTok{vars=}\ConstantTok{TRUE}\NormalTok{,}\AttributeTok{CI=}\ConstantTok{TRUE}\NormalTok{)}
\FunctionTok{print}\NormalTok{(rec}\FloatTok{.1}\NormalTok{,}\AttributeTok{printlen=}\DecValTok{10}\NormalTok{)}
\CommentTok{\# Comparando com a reconstrução original}
\FunctionTok{print}\NormalTok{(rec,}\AttributeTok{printlen=}\DecValTok{10}\NormalTok{)}
\CommentTok{\# Visualizando reconstrução com nós conhecidos}
\CommentTok{\# lista contendo valores terminais + estimados com fastAnc}
\NormalTok{X}\OtherTok{\textless{}{-}}\FunctionTok{c}\NormalTok{(}\FunctionTok{log}\NormalTok{(bodymass),rec}\FloatTok{.1}\SpecialCharTok{$}\NormalTok{ace)}
\FunctionTok{par}\NormalTok{(}\AttributeTok{mfrow =} \FunctionTok{c}\NormalTok{(}\DecValTok{1}\NormalTok{, }\DecValTok{1}\NormalTok{),}\AttributeTok{mar=}\FunctionTok{c}\NormalTok{(}\DecValTok{4}\NormalTok{, }\DecValTok{4}\NormalTok{, }\DecValTok{2}\NormalTok{, }\DecValTok{2}\NormalTok{))}
\FunctionTok{phenogram}\NormalTok{(tree,X,}\AttributeTok{ftype=}\StringTok{"i"}\NormalTok{,}\AttributeTok{fsize=}\FloatTok{0.5}\NormalTok{)}
\end{Highlighting}
\end{Shaded}

\includegraphics{04.3-pcmcont_files/figure-latex/unnamed-chunk-8-1.pdf}

\section{Exercício - Modelos contínuos e variação intraespecífica}\label{exercuxedcio---modelos-contuxednuos-e-variauxe7uxe3o-intraespecuxedfica}

Carregue os dados:
primate-data-erro.txt
primate-tree.txt
Estime o sinal filogenético com a estatística K, e ajuste e compare os modelos BM, OU e EB para os dados considerando a variação intraespecífica.
O arquivo de dados contém log de massa corporal para dez indivíduos de cada espécie de primata. Faça uma média de tamanho corporal por espécie e calcule o erro padrão da
média.
Descubra (leia a documentação das funções ?phylosig e ?fitContinuous) como incorporar o erro padrão no cálculo de sinal filogenético e ajuste de modelos.
Você pode usar a função apply para calcular a média e o desvio padrão da média (SD). O erro padrão da média (uma medida de quão precisa é a estimativa da média) pode ser obtido através da fórmula SD/√N (N=tamanho da amostra) = se\textless-sd/sqrt(10).

\chapter{Modelos para evolução de fenótipos discretos}\label{modelos-para-evoluuxe7uxe3o-de-fenuxf3tipos-discretos}

Sobre o ajuste de modelos evolutivos para dados discretos leia (Yang 2006; Revell and Harmon 2022).

Vamos carregar os mesmos dados de primatas que usamos no capítulo anterior. Os ajustes de modelos seguem a mesma lógica. Aqui, vamos aplicar para fenótipos categóricos/discretos.

\begin{Shaded}
\begin{Highlighting}[]
\FunctionTok{require}\NormalTok{(ape)}
\CommentTok{\#\textgreater{} Carregando pacotes exigidos: ape}
\FunctionTok{require}\NormalTok{(phytools)}
\CommentTok{\#\textgreater{} Carregando pacotes exigidos: phytools}
\CommentTok{\#\textgreater{} Carregando pacotes exigidos: maps}
\FunctionTok{require}\NormalTok{(geiger)}
\CommentTok{\#\textgreater{} Carregando pacotes exigidos: geiger}
\CommentTok{\# Filogenia}
\NormalTok{tree}\OtherTok{\textless{}{-}}\FunctionTok{read.tree}\NormalTok{(}\StringTok{"dadospcm/primate{-}tree.txt"}\NormalTok{)}
\FunctionTok{plotTree}\NormalTok{(tree,}\AttributeTok{fsize=}\FloatTok{0.4}\NormalTok{,}\AttributeTok{ftype=}\StringTok{"i"}\NormalTok{,}\AttributeTok{type=}\StringTok{"fan"}\NormalTok{,}\AttributeTok{lwd=}\DecValTok{1}\NormalTok{)}
\end{Highlighting}
\end{Shaded}

\includegraphics{04.4-pcmdisc_files/figure-latex/unnamed-chunk-1-1.pdf}

\begin{Shaded}
\begin{Highlighting}[]
\CommentTok{\# Atributos}
\NormalTok{dados}\OtherTok{\textless{}{-}}\FunctionTok{read.table}\NormalTok{(}\StringTok{"dadospcm/primate{-}data.txt"}\NormalTok{,}\AttributeTok{h=}\NormalTok{T,}\AttributeTok{row.names=}\DecValTok{1}\NormalTok{)}
\FunctionTok{head}\NormalTok{(dados)}
\CommentTok{\# Recortar filogenia}
\NormalTok{match.species}\OtherTok{\textless{}{-}}\FunctionTok{treedata}\NormalTok{(tree,dados)}
\NormalTok{tree}\OtherTok{\textless{}{-}}\NormalTok{match.species}\SpecialCharTok{$}\NormalTok{phy}
\FunctionTok{plotTree}\NormalTok{(tree,}\AttributeTok{fsize=}\FloatTok{0.5}\NormalTok{,}\AttributeTok{ftype=}\StringTok{"i"}\NormalTok{,}\AttributeTok{type=}\StringTok{"fan"}\NormalTok{,}\AttributeTok{lwd=}\DecValTok{1}\NormalTok{)}
\end{Highlighting}
\end{Shaded}

\includegraphics{04.4-pcmdisc_files/figure-latex/unnamed-chunk-1-2.pdf}

O conjunto de dados contém atributos binários e categóricos informando o tamanho do grupo social em primatas. Primeiro, vamos usar a variável binária, que possui dois estados: Solitary vs.~Social. Mais de 4 ind = solitary Menos de 4 ind = social

\begin{Shaded}
\begin{Highlighting}[]
\NormalTok{binsocial}\OtherTok{\textless{}{-}}\NormalTok{dados}\SpecialCharTok{$}\NormalTok{BinSocial}
\FunctionTok{names}\NormalTok{(binsocial)}\OtherTok{=}\FunctionTok{rownames}\NormalTok{(dados)}
\NormalTok{binsocial}\OtherTok{\textless{}{-}}\FunctionTok{as.factor}\NormalTok{(binsocial)}
\NormalTok{binsocial}
\end{Highlighting}
\end{Shaded}

\section{Modelos evolutivos discretos}\label{modelos-evolutivos-discretos}

Vamos usar a função fit.Discrete do pacote \emph{geiger} para ajustar modelos que têm diferentes expectativas sobre as taxas de transição entre os estados fenotípicos.

\subsection{Modelo de taxas iguais}\label{modelo-de-taxas-iguais}

Equal rates (ER) (ou all rates equal). Esse modelo tem um único parâmetro, a taxa de transição, que é igual entre qualquer par de estados.

\begin{Shaded}
\begin{Highlighting}[]
\NormalTok{ER}\OtherTok{\textless{}{-}}\FunctionTok{fitDiscrete}\NormalTok{(tree,binsocial,}\AttributeTok{model=}\StringTok{"ER"}\NormalTok{)}
\NormalTok{ER}
\FunctionTok{plot}\NormalTok{(ER,}\AttributeTok{signif=}\DecValTok{5}\NormalTok{)}
\end{Highlighting}
\end{Shaded}

\includegraphics{04.4-pcmdisc_files/figure-latex/unnamed-chunk-3-1.pdf}

\subsection{Modelo simétrico}\label{modelo-simuxe9trico}

No modelo simétrico (model=``SYM''), a taxa de mudança entre qualquer par de estados é a mesma independente da direção (estado 1 pra 2 é o mesmo que 2 pra 1; q12 = q21), mas a taxa de mudança difere entre estados (estado 1 pra 2 não é o mesmo que 1 pra 3; q12/q21 ≠ q13/q31 ≠ q23/q32). Note que, para uma variável binária, o modelo simétrico é exatamente igual ao modelo equal rates.

\subsection{Modelos com todas as taxas diferentes}\label{modelos-com-todas-as-taxas-diferentes}

All rates different (ARD). Todas as taxas são diferentes. Para S estados, o número de parâmetros (S²-S) pode aumentar rapidamente dependendo de quantos estados houver.

\begin{Shaded}
\begin{Highlighting}[]
\NormalTok{ARD}\OtherTok{\textless{}{-}}\FunctionTok{fitDiscrete}\NormalTok{(tree,binsocial,}\AttributeTok{model=}\StringTok{"ARD"}\NormalTok{)}
\NormalTok{ARD}
\FunctionTok{plot}\NormalTok{(ARD,}\AttributeTok{signif=}\DecValTok{5}\NormalTok{)}
\end{Highlighting}
\end{Shaded}

\includegraphics{04.4-pcmdisc_files/figure-latex/unnamed-chunk-4-1.pdf}

Parece que a taxa de mudança de solitário para social (0.010) é maior do que a taxa de mudança de social para solitário (0.005).

\subsection{Modelo irreversível}\label{modelo-irreversuxedvel}

Nós também podemos ajustar um modelo em que as transições de solitário para social são possíveis, mas o reverso não é possível. Para isso, podemos criar uma matriz Q representando as transições entre estados.

\begin{Shaded}
\begin{Highlighting}[]
\NormalTok{matriz}\OtherTok{\textless{}{-}}\FunctionTok{matrix}\NormalTok{(}\FunctionTok{c}\NormalTok{(}
 \DecValTok{0}\NormalTok{,}\DecValTok{0}\NormalTok{,}
 \DecValTok{1}\NormalTok{,}\DecValTok{0}\NormalTok{),}\DecValTok{2}\NormalTok{,}\DecValTok{2}\NormalTok{,}\AttributeTok{byrow=}\ConstantTok{TRUE}\NormalTok{)}
\FunctionTok{colnames}\NormalTok{(matriz)}\OtherTok{=}\FunctionTok{rownames}\NormalTok{(matriz)}\OtherTok{=}\FunctionTok{c}\NormalTok{(}\StringTok{"Social"}\NormalTok{,}\StringTok{"Solitary"}\NormalTok{)}
\NormalTok{matriz}
\NormalTok{IR}\OtherTok{\textless{}{-}}\FunctionTok{fitDiscrete}\NormalTok{(tree,binsocial,}\AttributeTok{model=}\NormalTok{matriz)}
\NormalTok{IR}
\CommentTok{\# Um warning é emitido quando qualquer parâmetro estiver nos limites (bounds) estabelecidos ou default.}
\FunctionTok{plot}\NormalTok{(IR,}\AttributeTok{signif=}\DecValTok{5}\NormalTok{)}
\end{Highlighting}
\end{Shaded}

\includegraphics{04.4-pcmdisc_files/figure-latex/unnamed-chunk-5-1.pdf}

\subsection{Comparação de modelos}\label{comparauxe7uxe3o-de-modelos-1}

Assim como nos modelos contínuos, podemos comparar os modelos usando AIC.

\begin{Shaded}
\begin{Highlighting}[]
\NormalTok{ER}\SpecialCharTok{$}\NormalTok{opt}\SpecialCharTok{$}\NormalTok{aicc}
\NormalTok{ARD}\SpecialCharTok{$}\NormalTok{opt}\SpecialCharTok{$}\NormalTok{aicc}
\NormalTok{IR}\SpecialCharTok{$}\NormalTok{opt}\SpecialCharTok{$}\NormalTok{aicc}
\NormalTok{aic.scores}\OtherTok{\textless{}{-}}\FunctionTok{setNames}\NormalTok{(}\FunctionTok{c}\NormalTok{(ER}\SpecialCharTok{$}\NormalTok{opt}\SpecialCharTok{$}\NormalTok{aicc,ARD}\SpecialCharTok{$}\NormalTok{opt}\SpecialCharTok{$}\NormalTok{aicc,IR}\SpecialCharTok{$}\NormalTok{opt}\SpecialCharTok{$}\NormalTok{aicc),}
 \FunctionTok{c}\NormalTok{(}\StringTok{"ER"}\NormalTok{,}\StringTok{"ARD"}\NormalTok{,}\StringTok{"IR"}\NormalTok{))}
\NormalTok{aic.scores}
\FunctionTok{aicw}\NormalTok{(aic.scores) }\CommentTok{\# Delta AICc e AICc weights}
\end{Highlighting}
\end{Shaded}

O modelo com taxas iguais de transição recebeu maior suporte, seguido do modelo ARD.

\subsection{Exemplo com múltiplos estados}\label{exemplo-com-muxfaltiplos-estados}

Modelos com várias categorias com mais de dois estados são implementados com exatamente a mesma lógica. Considere a variável grupo social construída com várias categorias. Menos de 4 ind = solitary Entre 4 e 10 ind = pequeno Entre 10 e 30 ind = medio Mais de 31 ind = grande

\begin{Shaded}
\begin{Highlighting}[]
\NormalTok{catsocial}\OtherTok{\textless{}{-}}\NormalTok{dados}\SpecialCharTok{$}\NormalTok{CatSocial}
\FunctionTok{names}\NormalTok{(catsocial)}\OtherTok{=}\FunctionTok{rownames}\NormalTok{(dados)}
\NormalTok{catsocial}
\CommentTok{\# Modelos}
\CommentTok{\# ER}
\NormalTok{ER.grupo}\OtherTok{\textless{}{-}}\FunctionTok{fitDiscrete}\NormalTok{(tree,catsocial,}\AttributeTok{model=}\StringTok{"ER"}\NormalTok{)}
\NormalTok{ER.grupo}
\FunctionTok{plot}\NormalTok{(ER.grupo)}
\end{Highlighting}
\end{Shaded}

\includegraphics{04.4-pcmdisc_files/figure-latex/unnamed-chunk-7-1.pdf}

\begin{Shaded}
\begin{Highlighting}[]
\CommentTok{\# SYM}
\NormalTok{SYM.grupo}\OtherTok{\textless{}{-}}\FunctionTok{fitDiscrete}\NormalTok{(tree,catsocial,}\AttributeTok{model=}\StringTok{"SYM"}\NormalTok{)}
\NormalTok{SYM.grupo}
\FunctionTok{plot}\NormalTok{(SYM.grupo)}
\end{Highlighting}
\end{Shaded}

\includegraphics{04.4-pcmdisc_files/figure-latex/unnamed-chunk-7-2.pdf}

\begin{Shaded}
\begin{Highlighting}[]
\CommentTok{\# ARD}
\NormalTok{ARD.grupo}\OtherTok{\textless{}{-}}\FunctionTok{fitDiscrete}\NormalTok{(tree,catsocial,}\AttributeTok{model=}\StringTok{"ARD"}\NormalTok{) }\CommentTok{\# Pode demorar}
\NormalTok{ARD.grupo}
\FunctionTok{plot}\NormalTok{(ARD.grupo)}
\end{Highlighting}
\end{Shaded}

\includegraphics{04.4-pcmdisc_files/figure-latex/unnamed-chunk-7-3.pdf}

Modelo Ordenado. Biologicamente, faz sentido que o tamanho do grupo social evolua de forma ordenada (e.g.~solitário\textgreater pequeno\textgreater medio\textgreater grande). Podemos modificar a matriz Q para permitir transições apenas entre estes estados.

\begin{Shaded}
\begin{Highlighting}[]
\NormalTok{matriz}\OtherTok{\textless{}{-}}\FunctionTok{matrix}\NormalTok{(}\FunctionTok{c}\NormalTok{(}
 \DecValTok{0}\NormalTok{,}\DecValTok{1}\NormalTok{,}\DecValTok{0}\NormalTok{,}\DecValTok{0}\NormalTok{,}
 \DecValTok{2}\NormalTok{,}\DecValTok{0}\NormalTok{,}\DecValTok{3}\NormalTok{,}\DecValTok{0}\NormalTok{,}
 \DecValTok{0}\NormalTok{,}\DecValTok{4}\NormalTok{,}\DecValTok{0}\NormalTok{,}\DecValTok{5}\NormalTok{,}
 \DecValTok{0}\NormalTok{,}\DecValTok{0}\NormalTok{,}\DecValTok{6}\NormalTok{,}\DecValTok{0}\NormalTok{),}\DecValTok{4}\NormalTok{,}\DecValTok{4}\NormalTok{,}\AttributeTok{byrow=}\ConstantTok{TRUE}\NormalTok{)}
\FunctionTok{colnames}\NormalTok{(matriz)}\OtherTok{=}\FunctionTok{rownames}\NormalTok{(matriz)}\OtherTok{=}\FunctionTok{c}\NormalTok{(}\StringTok{"Grande"}\NormalTok{,}\StringTok{"Medio"}\NormalTok{,}\StringTok{"Pequeno"}\NormalTok{,}\StringTok{"Solitary"}\NormalTok{)}
\NormalTok{matriz}
\NormalTok{OR.grupo}\OtherTok{\textless{}{-}}\FunctionTok{fitDiscrete}\NormalTok{(tree,catsocial,}\AttributeTok{model=}\NormalTok{matriz)}
\NormalTok{OR.grupo}
\FunctionTok{plot}\NormalTok{(OR.grupo,}\AttributeTok{show.zeros=}\ConstantTok{FALSE}\NormalTok{)}
\end{Highlighting}
\end{Shaded}

\includegraphics{04.4-pcmdisc_files/figure-latex/unnamed-chunk-8-1.pdf}

Comparando os modelos.

\begin{Shaded}
\begin{Highlighting}[]
\NormalTok{aic.scores}\OtherTok{\textless{}{-}}
\FunctionTok{setNames}\NormalTok{(}\FunctionTok{c}\NormalTok{(ER.grupo}\SpecialCharTok{$}\NormalTok{opt}\SpecialCharTok{$}\NormalTok{aicc,SYM.grupo}\SpecialCharTok{$}\NormalTok{opt}\SpecialCharTok{$}\NormalTok{aicc,ARD.grupo}\SpecialCharTok{$}\NormalTok{opt}\SpecialCharTok{$}\NormalTok{aicc,OR.grupo}\SpecialCharTok{$}\NormalTok{opt}\SpecialCharTok{$}\NormalTok{aicc),}
 \FunctionTok{c}\NormalTok{(}\StringTok{"ER"}\NormalTok{,}\StringTok{"SYM"}\NormalTok{,}\StringTok{"ARD"}\NormalTok{,}\StringTok{"OR"}\NormalTok{))}
\NormalTok{aic.scores}
\FunctionTok{aicw}\NormalTok{(aic.scores) }\CommentTok{\# Delta AICc e AICc weights}
\FunctionTok{aic.w}\NormalTok{(aic.scores)}
\end{Highlighting}
\end{Shaded}

Neste caso, o modelo em que tamanho do grupo social pode mudar de maneira ordenada (solitário\textgreater pequeno\textgreater medio\textgreater grande \& volta) teve muito mais suporte que os outros.

\section{Reconstrução de estados ancestrais}\label{reconstruuxe7uxe3o-de-estados-ancestrais}

Para reconstrução ancestral de estados discretos, o modelo subjacente será o Markoviano (Mk model). A função ace do pacote \emph{ape} realizada estimativa para caracteres discretos.

\begin{Shaded}
\begin{Highlighting}[]
\CommentTok{\# Visualização dos atributos}
\NormalTok{\{}\FunctionTok{plotTree}\NormalTok{(tree,}\AttributeTok{fsize=}\FloatTok{0.5}\NormalTok{,}\AttributeTok{ftype=}\StringTok{"i"}\NormalTok{)}
\NormalTok{cores}\OtherTok{\textless{}{-}}\FunctionTok{setNames}\NormalTok{(}\FunctionTok{c}\NormalTok{(}\StringTok{"blue"}\NormalTok{,}\StringTok{"red"}\NormalTok{),}\FunctionTok{levels}\NormalTok{(binsocial))}
\NormalTok{cores}
\FunctionTok{tiplabels}\NormalTok{(}\AttributeTok{pie=}\FunctionTok{to.matrix}\NormalTok{(binsocial[tree}\SpecialCharTok{$}\NormalTok{tip.label],}\FunctionTok{levels}\NormalTok{(binsocial)),}\AttributeTok{piecol=}\NormalTok{cores,}\AttributeTok{cex=}\FloatTok{0.3}\NormalTok{)}
\FunctionTok{legend}\NormalTok{(}\StringTok{"bottomleft"}\NormalTok{,}\AttributeTok{legend=}\FunctionTok{c}\NormalTok{(}\StringTok{"Social"}\NormalTok{,}\StringTok{"Solitary"}\NormalTok{),}\AttributeTok{pch=}\DecValTok{15}\NormalTok{,}\AttributeTok{col=}\NormalTok{cores,}\AttributeTok{bty=}\StringTok{"n"}\NormalTok{)\}}
\end{Highlighting}
\end{Shaded}

\includegraphics{04.4-pcmdisc_files/figure-latex/unnamed-chunk-10-1.pdf}

\begin{Shaded}
\begin{Highlighting}[]
\CommentTok{\# Estimativa de estados ancestrais com modelo ER}
\NormalTok{ERanc}\OtherTok{\textless{}{-}}\FunctionTok{ace}\NormalTok{(binsocial,tree,}\AttributeTok{type=}\StringTok{"discrete"}\NormalTok{,}\AttributeTok{model=}\StringTok{"ER"}\NormalTok{)}
\NormalTok{ERanc}
\CommentTok{\# Estados ancestrais}
\NormalTok{ERanc}\SpecialCharTok{$}\NormalTok{lik.anc}
\end{Highlighting}
\end{Shaded}

As linhas correspondem aos nós da árvore, e as colunas são as probabilidades de que o nó esteja no estado 1 ou 2 (social ou solitária). Note que os nós retornados pelo ace estão formatados em sequência numérica começando do 1, enquanto o formato phylo começa a dar nome aos nós seguindo a sequência n+1 (onde n é o número de espécies). É possível formatar os números das linhas para adequar ao formato phylo:

\begin{Shaded}
\begin{Highlighting}[]
\NormalTok{node.states}\OtherTok{\textless{}{-}}\NormalTok{ERanc}\SpecialCharTok{$}\NormalTok{lik.anc}
\FunctionTok{rownames}\NormalTok{(node.states) }\OtherTok{\textless{}{-}} \FunctionTok{seq}\NormalTok{(}\DecValTok{1}\SpecialCharTok{:}\FunctionTok{nrow}\NormalTok{(node.states)) }\SpecialCharTok{+} \FunctionTok{length}\NormalTok{(tree}\SpecialCharTok{$}\NormalTok{tip.label)}
\NormalTok{node.states}
\end{Highlighting}
\end{Shaded}

Visualização de estados com gráficos de ``pizza''.

\begin{Shaded}
\begin{Highlighting}[]
\NormalTok{\{}\FunctionTok{plotTree}\NormalTok{(tree,}\AttributeTok{fsize=}\FloatTok{0.5}\NormalTok{,}\AttributeTok{ftype=}\StringTok{"i"}\NormalTok{)}
\FunctionTok{nodelabels}\NormalTok{(}\AttributeTok{pie =}\NormalTok{ ERanc}\SpecialCharTok{$}\NormalTok{lik.anc, }\AttributeTok{piecol =}\NormalTok{ cores, }\AttributeTok{cex =} \FloatTok{0.5}\NormalTok{)}
\FunctionTok{tiplabels}\NormalTok{(}\AttributeTok{pie=}\FunctionTok{to.matrix}\NormalTok{(binsocial[tree}\SpecialCharTok{$}\NormalTok{tip.label],}\FunctionTok{levels}\NormalTok{(binsocial)),}\AttributeTok{piecol=}\NormalTok{cores,}\AttributeTok{cex=}\FloatTok{0.3}\NormalTok{)}
\FunctionTok{legend}\NormalTok{(}\StringTok{"bottomleft"}\NormalTok{,}\AttributeTok{legend=}\FunctionTok{c}\NormalTok{(}\StringTok{"Social"}\NormalTok{,}\StringTok{"Solitary"}\NormalTok{),}\AttributeTok{pch=}\DecValTok{15}\NormalTok{,}\AttributeTok{col=}\NormalTok{cores,}\AttributeTok{bty=}\StringTok{"n"}\NormalTok{)\}}
\end{Highlighting}
\end{Shaded}

\includegraphics{04.4-pcmdisc_files/figure-latex/unnamed-chunk-12-1.pdf}

Estimativa de estados ancestrais com o modelo ARD.

\begin{Shaded}
\begin{Highlighting}[]
\NormalTok{ARDanc}\OtherTok{\textless{}{-}}\FunctionTok{ace}\NormalTok{(binsocial,tree,}\AttributeTok{type=}\StringTok{"discrete"}\NormalTok{,}\AttributeTok{model=}\StringTok{"ARD"}\NormalTok{)}
\NormalTok{ARDanc}
\NormalTok{\{}\FunctionTok{plotTree}\NormalTok{(tree,}\AttributeTok{fsize=}\FloatTok{0.5}\NormalTok{,}\AttributeTok{ftype=}\StringTok{"i"}\NormalTok{)}
\FunctionTok{nodelabels}\NormalTok{(}\AttributeTok{pie =}\NormalTok{ ARDanc}\SpecialCharTok{$}\NormalTok{lik.anc, }\AttributeTok{piecol =}\NormalTok{ cores, }\AttributeTok{cex =} \FloatTok{0.5}\NormalTok{)}
\FunctionTok{tiplabels}\NormalTok{(}\AttributeTok{pie=}\FunctionTok{to.matrix}\NormalTok{(binsocial[tree}\SpecialCharTok{$}\NormalTok{tip.label],}\FunctionTok{levels}\NormalTok{(binsocial)),}\AttributeTok{piecol=}\NormalTok{cores,}\AttributeTok{cex=}\FloatTok{0.3}\NormalTok{)}
\FunctionTok{legend}\NormalTok{(}\StringTok{"bottomleft"}\NormalTok{,}\AttributeTok{legend=}\FunctionTok{c}\NormalTok{(}\StringTok{"Social"}\NormalTok{,}\StringTok{"Solitary"}\NormalTok{),}\AttributeTok{pch=}\DecValTok{15}\NormalTok{,}\AttributeTok{col=}\NormalTok{cores,}\AttributeTok{bty=}\StringTok{"n"}\NormalTok{)\}}
\end{Highlighting}
\end{Shaded}

\includegraphics{04.4-pcmdisc_files/figure-latex/unnamed-chunk-13-1.pdf}

\subsection{Mapeamento estocástico de caracteres}\label{mapeamento-estocuxe1stico-de-caracteres}

Existe outra alternativa às estimativas acima, que usa MCMC (Bayesiano) para amostrar histórias a partir de uma distribuição contendo muitos (milhares de) resultados (o que chamamos de distribuição posterior de probabilidades). Assim, os modelos são os mesmos usados anteriormente (ER, ARD), mas no final é obtida uma amostra de histórias que são únicas e não ambíguas, ou seja, cada etapa na árvore pode conter apenas um estado, ao invés da probabilidade de todos os estados para cada nó. Pra se aprofundar, leia Huelsenbeck, Nielsen, and Bollback (2003), Revell and Harmon (2022).

\begin{Shaded}
\begin{Highlighting}[]
\CommentTok{\# Função make.simmap do phytools}
\NormalTok{anc.tree}\OtherTok{\textless{}{-}}\FunctionTok{make.simmap}\NormalTok{(tree,binsocial,}\AttributeTok{model=}\StringTok{"ER"}\NormalTok{)}
\CommentTok{\#\textgreater{} Done.}
\NormalTok{\{}\FunctionTok{plot}\NormalTok{(anc.tree,cores,}\AttributeTok{fsize=}\FloatTok{0.5}\NormalTok{,}\AttributeTok{ftype=}\StringTok{"i"}\NormalTok{)}
\FunctionTok{add.simmap.legend}\NormalTok{(}\AttributeTok{colors=}\NormalTok{cores,}\AttributeTok{prompt=}\ConstantTok{FALSE}\NormalTok{,}\AttributeTok{x=}\SpecialCharTok{{-}}\DecValTok{2}\SpecialCharTok{*}\FunctionTok{par}\NormalTok{()}\SpecialCharTok{$}\NormalTok{usr[}\DecValTok{1}\NormalTok{],}\AttributeTok{y=}\SpecialCharTok{{-}}\DecValTok{5}\SpecialCharTok{*}\FunctionTok{par}\NormalTok{()}\SpecialCharTok{$}\NormalTok{usr[}\DecValTok{3}\NormalTok{],}\AttributeTok{fsize=}\FloatTok{0.8}\NormalTok{)\}}
\end{Highlighting}
\end{Shaded}

\includegraphics{04.4-pcmdisc_files/figure-latex/unnamed-chunk-14-1.pdf}
Esta é uma única simulação da história, é preciso analisar múltiplos mapas estocásticos.

\begin{Shaded}
\begin{Highlighting}[]
\NormalTok{anc.tree}\OtherTok{\textless{}{-}}\FunctionTok{make.simmap}\NormalTok{(tree,binsocial,}\AttributeTok{model=}\StringTok{"ER"}\NormalTok{,}\AttributeTok{nsim=}\DecValTok{100}\NormalTok{) }\CommentTok{\# ex: 100}
\CommentTok{\#\textgreater{} Done.}
\NormalTok{anc.tree}
\CommentTok{\# Resumindo múltiplas árvores}
\NormalTok{pd}\OtherTok{\textless{}{-}}\FunctionTok{summary}\NormalTok{(anc.tree)}
\NormalTok{pd}
\NormalTok{\{}\FunctionTok{plot}\NormalTok{(pd,}\AttributeTok{fsize=}\FloatTok{0.5}\NormalTok{,}\AttributeTok{ftype=}\StringTok{"i"}\NormalTok{,}\AttributeTok{colors=}\NormalTok{cores)}
\FunctionTok{add.simmap.legend}\NormalTok{(}\AttributeTok{colors=}\NormalTok{cores,}\AttributeTok{prompt=}\ConstantTok{FALSE}\NormalTok{,}\AttributeTok{x=}\DecValTok{1}\NormalTok{,}\AttributeTok{y=}\DecValTok{10}\NormalTok{)\}}
\end{Highlighting}
\end{Shaded}

\includegraphics{04.4-pcmdisc_files/figure-latex/unnamed-chunk-15-1.pdf}

\begin{Shaded}
\begin{Highlighting}[]
\CommentTok{\# Plotando pizza sobre um mapa ao acaso}
\CommentTok{\#\{plot(sample(anc.tree,1),fsize=0.5,ftype="i",colors=cores)}
\CommentTok{\#nodelabels(pie=pd$ace,piecol=cores,cex=0.5)}
\CommentTok{\#add.simmap.legend(colors=cores,prompt=FALSE,x=1,y=10)\}}
\end{Highlighting}
\end{Shaded}

Poderíamos comparar as probabilidades posteriores do stochastic mapping com as probabilidades de estados diferentes pelo ace. As probabilidades devem se tornar as mesmas conforme o número de mapas estocásticos aumenta ao infinito.

\section{Exercício - Modelos discretos}\label{exercuxedcio---modelos-discretos}

Carregue os mesmos dados:
primate-data.txt
primate-tree.txt
Usando a variável categórica com múltiplos grupos sociais (catsocial), ajuste um modelo ordenado direcional onde o tamanho do grupo social pode aumentar mas nunca diminuir (i.e.~Solitary\textgreater Pequeno\textgreater Medio\textgreater Grande mas sem retorno). Compare este modelo com os quatro anteriores e descubra qual recebe maior suporte.

\chapter{Regressão filogenética}\label{regressuxe3o-filogenuxe9tica}

Esse capítulo implementa contrastes filogenéticos independentes (PIC) (Felsenstein 1985), regressão filogenética generalizada por quadrados mínimos (PGLS) (Grafen 1989) e regressão de autovetores filogenéticos (PVR) (Diniz-Filho, de Sant'Ana, and Bini 1998).

\section{Contrastes filogenéticos independentes}\label{contrastes-filogenuxe9ticos-independentes-1}

Vamos investigar a relação evolutiva entre força da mordida e tamanho do crânio em um gênero de roedores. Pra isso, vamos precisar dos dados fenotípicos e pelo menos uma hipótese filogenética.

Carregar dados fenotípicos.

\begin{Shaded}
\begin{Highlighting}[]
\NormalTok{dados}\OtherTok{\textless{}{-}}\FunctionTok{read.table}\NormalTok{(}\StringTok{"dadospcm/akodon.txt"}\NormalTok{,}\AttributeTok{h=}\NormalTok{T)}
\NormalTok{dados}
\FunctionTok{attach}\NormalTok{(dados)}
\end{Highlighting}
\end{Shaded}

Antes, vamos ajustar um modelo ordinário (OLS), sem levar a filogenia em consideração, apenas como exemplo didático. Na prática, você não deve/precisa fazer isso.

\begin{Shaded}
\begin{Highlighting}[]
\CommentTok{\# OLS}
\NormalTok{fit.ols}\OtherTok{\textless{}{-}}\FunctionTok{lm}\NormalTok{(bite\_force}\SpecialCharTok{\textasciitilde{}}\NormalTok{skull\_length,}\AttributeTok{data=}\NormalTok{dados)}
\FunctionTok{summary}\NormalTok{(fit.ols)}
\NormalTok{\{}\FunctionTok{plot}\NormalTok{(skull\_length,bite\_force,}
 \AttributeTok{xlab=}\StringTok{"skull length (log)"}\NormalTok{,}
 \AttributeTok{ylab=}\StringTok{"bite force (log)"}\NormalTok{,}\AttributeTok{pch=}\DecValTok{21}\NormalTok{,}\AttributeTok{bg=}\StringTok{"grey"}\NormalTok{,}\AttributeTok{cex=}\FloatTok{1.5}\NormalTok{)}
\FunctionTok{abline}\NormalTok{(fit.ols,}\AttributeTok{col=}\StringTok{"red"}\NormalTok{)\}}
\end{Highlighting}
\end{Shaded}

\includegraphics{04.5-pcmreg_files/figure-latex/unnamed-chunk-2-1.pdf}

Como os dados são de espécies diferentes, o pressuposto de independência dos resíduos do modelo é violado. Uma maneira de levar em consideração a filogenia é através dos Contrastes Filogenéticos Independentes.

\begin{Shaded}
\begin{Highlighting}[]
\CommentTok{\# PIC}
\FunctionTok{require}\NormalTok{(ape)}
\CommentTok{\#\textgreater{} Carregando pacotes exigidos: ape}
\NormalTok{akodon.tree}\OtherTok{\textless{}{-}}\FunctionTok{read.tree}\NormalTok{(}\AttributeTok{file=}\StringTok{"dadospcm/akodon.tree"}\NormalTok{)}
\FunctionTok{plot}\NormalTok{(akodon.tree)}
\end{Highlighting}
\end{Shaded}

\includegraphics{04.5-pcmreg_files/figure-latex/unnamed-chunk-3-1.pdf}

\begin{Shaded}
\begin{Highlighting}[]
\NormalTok{akodon.tree}
\CommentTok{\# Atribuindo nomes para os dados}
\NormalTok{skull\_length}\OtherTok{\textless{}{-}}\FunctionTok{setNames}\NormalTok{(skull\_length,species)}
\NormalTok{skull\_length}
\NormalTok{bite\_force}\OtherTok{\textless{}{-}}\FunctionTok{setNames}\NormalTok{(bite\_force,species)}
\NormalTok{bite\_force}
\CommentTok{\# Calculando os contrastes}
\NormalTok{pic.sl}\OtherTok{\textless{}{-}}\FunctionTok{pic}\NormalTok{(skull\_length,akodon.tree)}
\NormalTok{pic.bf}\OtherTok{\textless{}{-}}\FunctionTok{pic}\NormalTok{(bite\_force,akodon.tree)}
\CommentTok{\# Modelo PIC}
\NormalTok{fit.pic}\OtherTok{\textless{}{-}}\FunctionTok{lm}\NormalTok{(pic.bf}\SpecialCharTok{\textasciitilde{}}\NormalTok{pic.sl}\SpecialCharTok{+}\DecValTok{0}\NormalTok{)}
\NormalTok{fit.pic}
\FunctionTok{summary}\NormalTok{(fit.pic)}
\FunctionTok{plot}\NormalTok{(pic.sl,pic.bf,}
 \AttributeTok{xlab=}\StringTok{"PICs for skull length"}\NormalTok{,}
 \AttributeTok{ylab=}\StringTok{"PICs for bite force"}\NormalTok{,}\AttributeTok{pch=}\DecValTok{21}\NormalTok{,}\AttributeTok{bg=}\StringTok{"grey"}\NormalTok{,}\AttributeTok{cex=}\FloatTok{1.5}\NormalTok{)}
\FunctionTok{abline}\NormalTok{(fit.pic,}\AttributeTok{col=}\StringTok{"red"}\NormalTok{)}
\end{Highlighting}
\end{Shaded}

\includegraphics{04.5-pcmreg_files/figure-latex/unnamed-chunk-3-2.pdf}

Como podemos perceber, para estes dados não existe muita diferença nos resultados levando em consideração a filogenia. Isso nem sempre vai ser o caso. Nós podemos simular condições para visualizar como levar em conta a filogenia pode ter uma consequência grande.

\begin{Shaded}
\begin{Highlighting}[]
\CommentTok{\# PIC dados simulados}
\FunctionTok{require}\NormalTok{(phytools)}
\CommentTok{\#\textgreater{} Carregando pacotes exigidos: phytools}
\CommentTok{\#\textgreater{} Carregando pacotes exigidos: maps}
\CommentTok{\# set seed (starting point na geração de sequência aleatória, assim todos teremos os mesmos resultados)}
\FunctionTok{set.seed}\NormalTok{(}\DecValTok{4}\NormalTok{)}
\CommentTok{\# simular uma árvore filogenética}
\NormalTok{tree}\OtherTok{\textless{}{-}}\FunctionTok{rcoal}\NormalTok{(}\AttributeTok{n=}\DecValTok{80}\NormalTok{)}
\FunctionTok{plotTree}\NormalTok{(tree,}\AttributeTok{ftype=}\StringTok{"off"}\NormalTok{)}
\end{Highlighting}
\end{Shaded}

\includegraphics{04.5-pcmreg_files/figure-latex/unnamed-chunk-4-1.pdf}

\begin{Shaded}
\begin{Highlighting}[]
\CommentTok{\# simular 2 atributos (independentemente) evoluindo por movimento Browniano na filogenia}
\NormalTok{x}\OtherTok{\textless{}{-}}\FunctionTok{fastBM}\NormalTok{(tree)}
\NormalTok{y}\OtherTok{\textless{}{-}}\FunctionTok{fastBM}\NormalTok{(tree)}
\FunctionTok{par}\NormalTok{(}\AttributeTok{mar=}\FunctionTok{c}\NormalTok{(}\DecValTok{5}\NormalTok{,}\DecValTok{4}\NormalTok{,}\DecValTok{2}\NormalTok{,}\DecValTok{2}\NormalTok{))}
\CommentTok{\# Ajustar OLS}
\FunctionTok{plot}\NormalTok{(x,y,}\AttributeTok{pch=}\DecValTok{21}\NormalTok{,}\AttributeTok{bg=}\StringTok{"grey"}\NormalTok{,}\AttributeTok{cex=}\FloatTok{1.5}\NormalTok{)}
\NormalTok{fit.OLS}\OtherTok{\textless{}{-}}\FunctionTok{lm}\NormalTok{(y}\SpecialCharTok{\textasciitilde{}}\NormalTok{x)}
\NormalTok{\{}\FunctionTok{plot}\NormalTok{(x,y,}\AttributeTok{pch=}\DecValTok{21}\NormalTok{,}\AttributeTok{bg=}\StringTok{"grey"}\NormalTok{,}\AttributeTok{cex=}\FloatTok{1.5}\NormalTok{)}
\FunctionTok{abline}\NormalTok{(fit.OLS,}\AttributeTok{col=}\StringTok{"red"}\NormalTok{)\}}
\end{Highlighting}
\end{Shaded}

\includegraphics{04.5-pcmreg_files/figure-latex/unnamed-chunk-4-2.pdf}

\begin{Shaded}
\begin{Highlighting}[]
\FunctionTok{summary}\NormalTok{(fit.OLS)}
\end{Highlighting}
\end{Shaded}

Mesmo os dados tendo sidos simulados independentemente (na ausência de correlação entre x e y) nós podemos ver que existe uma correlação entre os dados induzida pela filogenia (erro tipo I).
Podemos plotar a filogenia no espaço de forma para descobrir o motivo:

\begin{Shaded}
\begin{Highlighting}[]
\NormalTok{\{}\FunctionTok{phylomorphospace}\NormalTok{(tree,}\FunctionTok{cbind}\NormalTok{(x,y),}\AttributeTok{label=}\StringTok{"off"}\NormalTok{,}\AttributeTok{node.size=}\FunctionTok{c}\NormalTok{(}\DecValTok{0}\NormalTok{,}\DecValTok{0}\NormalTok{))}
\FunctionTok{points}\NormalTok{(x,y,}\AttributeTok{pch=}\DecValTok{21}\NormalTok{,}\AttributeTok{bg=}\StringTok{"grey"}\NormalTok{,}\AttributeTok{cex=}\FloatTok{1.5}\NormalTok{)}
\FunctionTok{abline}\NormalTok{(fit.OLS,}\AttributeTok{col=}\StringTok{"red"}\NormalTok{,}\AttributeTok{lwd=}\DecValTok{2}\NormalTok{)\}}
\end{Highlighting}
\end{Shaded}

\includegraphics{04.5-pcmreg_files/figure-latex/unnamed-chunk-5-1.pdf}
Neste caso, espécies próximas filogeneticamente têm valores de atributo muito similares. Isso gera dois grupos de espécies próximas com fenótipo similar. Podemos usar os contrastes independentes neste caso.

\begin{Shaded}
\begin{Highlighting}[]
\CommentTok{\# Ajustando PIC}
\NormalTok{pic.x}\OtherTok{\textless{}{-}}\FunctionTok{pic}\NormalTok{(x,tree)}
\NormalTok{pic.y}\OtherTok{\textless{}{-}}\FunctionTok{pic}\NormalTok{(y,tree)}
\NormalTok{fit.pic}\OtherTok{\textless{}{-}}\FunctionTok{lm}\NormalTok{(pic.y}\SpecialCharTok{\textasciitilde{}}\NormalTok{pic.x}\SpecialCharTok{+}\DecValTok{0}\NormalTok{)}
\FunctionTok{summary}\NormalTok{(fit.pic)}
\FunctionTok{summary.aov}\NormalTok{(fit.pic)}
\end{Highlighting}
\end{Shaded}

\section{Regressão filogenética generalisada por quadrados mínimos}\label{regressuxe3o-filogenuxe9tica-generalisada-por-quadrados-muxednimos}

Vários pacotes realizam PGLS no R. Vamos usar os mais comuns - \emph{ape}/\emph{nlme} + \emph{caper}. A PGLS usa uma matriz de covariância entre espécies de acordo com algum processo evolutivo (BM, Pagel, OU, outros).

Vamos carregar alguns dados de borboletas.

\begin{Shaded}
\begin{Highlighting}[]
\FunctionTok{require}\NormalTok{(ape)}
\NormalTok{dados}\OtherTok{\textless{}{-}}\FunctionTok{read.table}\NormalTok{(}\StringTok{"dadospcm/butterfly{-}data.txt"}\NormalTok{,}\AttributeTok{h=}\NormalTok{T,}\AttributeTok{row.names =} \DecValTok{1}\NormalTok{)}
\NormalTok{dados}
\FunctionTok{attach}\NormalTok{(dados)}
\NormalTok{wing\_length}\OtherTok{\textless{}{-}}\FunctionTok{setNames}\NormalTok{(wing\_length,}\FunctionTok{rownames}\NormalTok{(dados))}
\NormalTok{temp}\OtherTok{\textless{}{-}}\FunctionTok{setNames}\NormalTok{(temp,}\FunctionTok{rownames}\NormalTok{(dados))}
\NormalTok{eye\_width}\OtherTok{\textless{}{-}}\FunctionTok{setNames}\NormalTok{(eye\_width,}\FunctionTok{rownames}\NormalTok{(dados))}
\CommentTok{\# Carregar árvore}
\NormalTok{tree}\OtherTok{\textless{}{-}}\FunctionTok{read.tree}\NormalTok{(}\StringTok{"dadospcm/butterfly{-}tree.txt"}\NormalTok{)}
\NormalTok{\{}\FunctionTok{plot}\NormalTok{(tree)}
\FunctionTok{axisPhylo}\NormalTok{()\}}
\end{Highlighting}
\end{Shaded}

\includegraphics{04.5-pcmreg_files/figure-latex/unnamed-chunk-7-1.pdf}

Conferir correspondência entre espécies nos dados fenotópicos e na árvore filogenética.

\begin{Shaded}
\begin{Highlighting}[]
\FunctionTok{require}\NormalTok{(geiger)}
\CommentTok{\#\textgreater{} Carregando pacotes exigidos: geiger}
\NormalTok{obj}\OtherTok{\textless{}{-}}\FunctionTok{name.check}\NormalTok{(tree,dados)}
\NormalTok{obj}
\CommentTok{\# 59 espécies na árvore não contêm dados}
\CommentTok{\# remover espécies da árvore}
\NormalTok{tree.pruned}\OtherTok{\textless{}{-}}\FunctionTok{drop.tip}\NormalTok{(tree,obj}\SpecialCharTok{$}\NormalTok{tree\_not\_data)}
\FunctionTok{name.check}\NormalTok{(tree.pruned,dados)}
\FunctionTok{plot}\NormalTok{(tree.pruned)}
\end{Highlighting}
\end{Shaded}

\includegraphics{04.5-pcmreg_files/figure-latex/unnamed-chunk-8-1.pdf}

PGLS com os pacotes \emph{ape} e \emph{nlme}.

\begin{Shaded}
\begin{Highlighting}[]
\FunctionTok{require}\NormalTok{(ape)}
\FunctionTok{require}\NormalTok{(nlme)}
\CommentTok{\#\textgreater{} Carregando pacotes exigidos: nlme}
\CommentTok{\# Estrutura de covariação Browniana}
\NormalTok{BM}\OtherTok{\textless{}{-}}\FunctionTok{corBrownian}\NormalTok{(}\DecValTok{1}\NormalTok{,tree.pruned)}
\NormalTok{BM}
\CommentTok{\# Podemos ajustar um modelo PGLS para investigar a relação entre temperatura e }
\CommentTok{\#comprimento da asa}
\NormalTok{modelo1}\OtherTok{\textless{}{-}}\FunctionTok{gls}\NormalTok{(wing\_length}\SpecialCharTok{\textasciitilde{}}\NormalTok{temp,}\AttributeTok{data=}\NormalTok{dados,}\AttributeTok{correlation=}\NormalTok{BM)}
\FunctionTok{summary}\NormalTok{(modelo1)}
\CommentTok{\# Plots de diagnóstico}
\FunctionTok{qqnorm}\NormalTok{(modelo1}\SpecialCharTok{$}\NormalTok{residuals)}
\end{Highlighting}
\end{Shaded}

\includegraphics{04.5-pcmreg_files/figure-latex/unnamed-chunk-9-1.pdf}

\begin{Shaded}
\begin{Highlighting}[]
\CommentTok{\# Comparando com OLS, os resultados são similares}
\NormalTok{modelo.ols}\OtherTok{\textless{}{-}}\FunctionTok{lm}\NormalTok{(wing\_length}\SpecialCharTok{\textasciitilde{}}\NormalTok{temp)}
\FunctionTok{summary}\NormalTok{(modelo.ols)}
\CommentTok{\# Plotando os resultados}
\NormalTok{\{}\FunctionTok{plot}\NormalTok{(temp,wing\_length,}\AttributeTok{pch=}\DecValTok{21}\NormalTok{,}\AttributeTok{bg=}\StringTok{"grey"}\NormalTok{,}\AttributeTok{cex=}\FloatTok{1.5}\NormalTok{)}
\FunctionTok{abline}\NormalTok{(modelo.ols,}\AttributeTok{col=}\StringTok{"red"}\NormalTok{,}\AttributeTok{lty=}\StringTok{"dashed"}\NormalTok{,}\AttributeTok{lwd=}\DecValTok{2}\NormalTok{)}
\FunctionTok{abline}\NormalTok{(modelo1,}\AttributeTok{col=}\StringTok{"blue"}\NormalTok{,}\AttributeTok{lty=}\StringTok{"dashed"}\NormalTok{,}\AttributeTok{lwd=}\DecValTok{2}\NormalTok{)\}}
\end{Highlighting}
\end{Shaded}

\includegraphics{04.5-pcmreg_files/figure-latex/unnamed-chunk-9-2.pdf}

É possível usar outras estruturas de covariância. Um dos modelos mais comuns usa o λ de Pagel. Esse modelo tem um parâmetro adicional (λ) que estima o sinal filogenético nos resíduos; os elementos off-diagonal são multiplicados pelo valor de lambda estimado, que varia entre 0 e 1. Quando λ=0 a covariância entre espécies é zero, e a regressão filogenética = regressão comum. Quando λ=1 a covariação é igual ao esperado pelo modelo Browniano (= PIC).

\begin{Shaded}
\begin{Highlighting}[]
\NormalTok{modelo2}\OtherTok{\textless{}{-}}\FunctionTok{gls}\NormalTok{(wing\_length}\SpecialCharTok{\textasciitilde{}}\NormalTok{temp,}\AttributeTok{data=}\NormalTok{dados,}\AttributeTok{correlation=}\FunctionTok{corPagel}\NormalTok{(}\DecValTok{1}\NormalTok{,tree.pruned))}
\FunctionTok{summary}\NormalTok{(modelo2)}
\NormalTok{\{}\FunctionTok{plot}\NormalTok{(temp,wing\_length,}\AttributeTok{pch=}\DecValTok{21}\NormalTok{,}\AttributeTok{bg=}\StringTok{"grey"}\NormalTok{,}\AttributeTok{cex=}\FloatTok{1.5}\NormalTok{)}
\FunctionTok{abline}\NormalTok{(modelo.ols,}\AttributeTok{col=}\StringTok{"red"}\NormalTok{,}\AttributeTok{lty=}\StringTok{"dashed"}\NormalTok{,}\AttributeTok{lwd=}\DecValTok{2}\NormalTok{)}
\FunctionTok{abline}\NormalTok{(modelo1,}\AttributeTok{col=}\StringTok{"blue"}\NormalTok{,}\AttributeTok{lty=}\StringTok{"dashed"}\NormalTok{,}\AttributeTok{lwd=}\DecValTok{2}\NormalTok{)}
\FunctionTok{abline}\NormalTok{(modelo2,}\AttributeTok{col=}\StringTok{"green"}\NormalTok{,}\AttributeTok{lty=}\StringTok{"dashed"}\NormalTok{,}\AttributeTok{lwd=}\DecValTok{2}\NormalTok{)\}}
\end{Highlighting}
\end{Shaded}

\includegraphics{04.5-pcmreg_files/figure-latex/unnamed-chunk-10-1.pdf}

PGLS pode ser aplicada em um contexto de regressão múltipla.

\begin{Shaded}
\begin{Highlighting}[]
\NormalTok{modelo3}\OtherTok{\textless{}{-}}
\FunctionTok{gls}\NormalTok{(wing\_length}\SpecialCharTok{\textasciitilde{}}\NormalTok{eye\_width}\SpecialCharTok{+}\NormalTok{temp,}\AttributeTok{data=}\NormalTok{dados,}\AttributeTok{correlation=}\FunctionTok{corPagel}\NormalTok{(}\DecValTok{1}\NormalTok{,tree.pruned))}
\FunctionTok{summary}\NormalTok{(modelo3)}
\end{Highlighting}
\end{Shaded}

PGLS com o pacote \emph{caper}.
O caper tem um objeto especial `comparative data' para combinar todos os dados.

\begin{Shaded}
\begin{Highlighting}[]
\FunctionTok{require}\NormalTok{(caper)}
\CommentTok{\#\textgreater{} Carregando pacotes exigidos: caper}
\CommentTok{\#\textgreater{} Carregando pacotes exigidos: MASS}
\CommentTok{\#\textgreater{} Carregando pacotes exigidos: mvtnorm}
\NormalTok{dados}\OtherTok{\textless{}{-}}\FunctionTok{read.table}\NormalTok{(}\StringTok{"dadospcm/butterfly{-}data.txt"}\NormalTok{,}\AttributeTok{header=}\NormalTok{T)}
\NormalTok{comp.data}\OtherTok{\textless{}{-}}\FunctionTok{comparative.data}\NormalTok{(tree.pruned,dados,}\AttributeTok{names.col=}\StringTok{\textquotesingle{}species\textquotesingle{}}\NormalTok{)}
\CommentTok{\# pgls usando lambda}
\NormalTok{modelo4}\OtherTok{\textless{}{-}}\FunctionTok{pgls}\NormalTok{(wing\_length}\SpecialCharTok{\textasciitilde{}}\NormalTok{temp,}\AttributeTok{data =}\NormalTok{ comp.data,}\AttributeTok{lambda=}\StringTok{"ML"}\NormalTok{)}
\FunctionTok{summary}\NormalTok{(modelo4)}
\end{Highlighting}
\end{Shaded}

Com o caper é possível plotar gráficos de diagnostico facilmente.

\begin{Shaded}
\begin{Highlighting}[]
\FunctionTok{par}\NormalTok{(}\AttributeTok{mfrow =} \FunctionTok{c}\NormalTok{(}\DecValTok{2}\NormalTok{, }\DecValTok{2}\NormalTok{))}
\FunctionTok{plot}\NormalTok{(modelo4)}
\end{Highlighting}
\end{Shaded}

\includegraphics{04.5-pcmreg_files/figure-latex/unnamed-chunk-13-1.pdf}

\begin{Shaded}
\begin{Highlighting}[]
\FunctionTok{par}\NormalTok{(}\AttributeTok{mfrow =} \FunctionTok{c}\NormalTok{(}\DecValTok{1}\NormalTok{, }\DecValTok{1}\NormalTok{))}
\end{Highlighting}
\end{Shaded}

O que olhar nesses plots:
No gráfico 1, dados com resíduos maiores que +- 3 podem ser outliers. Os pontos no QQ-plot (gráfico 2) devem cair aproximadamente sobre a linha. Os gráficos 3 e 4 devem mostrar uma distribuição mais ou menos randômica, sem padrões aparentes.

Como qualquer outro modelo linear, também é preciso prestar atenção ao número de preditores e número de termos estimados no modelo em relação ao número deobservações (espécies). Uma ``regra'' simples é usar +- 10 vezes mais espécies do que termos estimados (incluindo intercepto, interações, e o parâmetro lambda, por exemplo).

\section{Regressão de autovetores filogenéticos}\label{regressuxe3o-de-autovetores-filogenuxe9ticos}

Regressão de autovetores filogenéticos e Curva PSR com o pacote \emph{PVR}. Veja o capítulo sobre Sinal Filogenético pra mais detalhes sobre as análises.

\begin{Shaded}
\begin{Highlighting}[]
\FunctionTok{require}\NormalTok{(PVR)}
\CommentTok{\#\textgreater{} Carregando pacotes exigidos: PVR}
\CommentTok{\#\textgreater{} Warning: package \textquotesingle{}PVR\textquotesingle{} was built under R version 4.3.2}
\CommentTok{\#\textgreater{} PVR 0.3 loaded}
\CommentTok{\# Decomposição de matriz de distância filogenética}
\NormalTok{pvr\_obj}\OtherTok{\textless{}{-}}\FunctionTok{PVRdecomp}\NormalTok{(tree.pruned)}
\CommentTok{\# Regressão de autovetores filogenéticos}
\NormalTok{pvr\_reg}\OtherTok{\textless{}{-}}\FunctionTok{PVR}\NormalTok{(pvr\_obj,tree.pruned,wing\_length,temp)}
\CommentTok{\# Porcentagem de cada componente}
\NormalTok{pvr\_reg}\SpecialCharTok{@}\NormalTok{VarPart}
\CommentTok{\#\textgreater{} $a}
\CommentTok{\#\textgreater{} [1] 0.3326349}
\CommentTok{\#\textgreater{} }
\CommentTok{\#\textgreater{} $b}
\CommentTok{\#\textgreater{} [1] 0.1093753}
\CommentTok{\#\textgreater{} }
\CommentTok{\#\textgreater{} $c}
\CommentTok{\#\textgreater{} [1] 0.08198286}
\CommentTok{\#\textgreater{} }
\CommentTok{\#\textgreater{} $d}
\CommentTok{\#\textgreater{} [1] 0.4760069}
\FunctionTok{VarPartplot}\NormalTok{(pvr\_reg)}
\end{Highlighting}
\end{Shaded}

\includegraphics{04.5-pcmreg_files/figure-latex/unnamed-chunk-14-1.pdf}

\begin{verbatim}
#> [a] - Variation explained by environmental variables
#> [b] - Shared variation between environmental variables and phylogeny(PVR)
#> [c] - Variation explained by phylogeny(PVR)
#> [d]     - Unexplained variation
# Sinal Filogenético com PSR Curve
psr_obj<-PSR(pvr_obj,trait=wing_length,
null.model=TRUE,Brownian.model=TRUE,times=10)
PSRplot(psr_obj,info="both")
\end{verbatim}

\includegraphics{04.5-pcmreg_files/figure-latex/unnamed-chunk-14-2.pdf}

\section{Exercício - Regressão filogenética}\label{exercuxedcio---regressuxe3o-filogenuxe9tica}

Carregue os dois conjuntos de dados:
pic-dados.txt
pic-arvore.tre
Ajuste uma regressão não-filogenética e uma regressão filogenética aos dados usando contrastes independentes. Quais foram os resultados? Qual a interpretação?

\chapter{Correlação de fenótipos discretos}\label{correlauxe7uxe3o-de-fenuxf3tipos-discretos}

Esse capítulo implementa algumas análises que são úteis quando a variável resposta é de interesse é categórica/discreta.

\section{Correlação de fenótipos binários}\label{correlauxe7uxe3o-de-fenuxf3tipos-binuxe1rios}

Vamos carregar filogenia e fenótipos de espécies de primatas e investigar a associação evolutiva entre o tamanho do grupo social (solitario vs.~social) e o tamanho da área de vida (Pequena \textless{} 300m² \textgreater{} Grande).

\begin{Shaded}
\begin{Highlighting}[]
\FunctionTok{require}\NormalTok{(ape)}
\CommentTok{\#\textgreater{} Carregando pacotes exigidos: ape}
\FunctionTok{require}\NormalTok{(phytools)}
\CommentTok{\#\textgreater{} Carregando pacotes exigidos: phytools}
\CommentTok{\#\textgreater{} Carregando pacotes exigidos: maps}
\FunctionTok{require}\NormalTok{(geiger)}
\CommentTok{\#\textgreater{} Carregando pacotes exigidos: geiger}
\CommentTok{\# Filogenia}
\NormalTok{tree}\OtherTok{\textless{}{-}}\FunctionTok{read.tree}\NormalTok{(}\StringTok{"dadospcm/primate{-}tree.txt"}\NormalTok{)}
\FunctionTok{plotTree}\NormalTok{(tree,}\AttributeTok{fsize=}\FloatTok{0.4}\NormalTok{,}\AttributeTok{ftype=}\StringTok{"i"}\NormalTok{,}\AttributeTok{type=}\StringTok{"fan"}\NormalTok{,}\AttributeTok{lwd=}\DecValTok{1}\NormalTok{)}
\end{Highlighting}
\end{Shaded}

\includegraphics{04.6-pcmcordisc_files/figure-latex/unnamed-chunk-1-1.pdf}

\begin{Shaded}
\begin{Highlighting}[]
\CommentTok{\# Atributos}
\NormalTok{dados}\OtherTok{\textless{}{-}}\FunctionTok{read.table}\NormalTok{(}\StringTok{"dadospcm/primate{-}data.txt"}\NormalTok{,}\AttributeTok{h=}\NormalTok{T,}\AttributeTok{row.names=}\DecValTok{1}\NormalTok{)}
\FunctionTok{head}\NormalTok{(dados)}
\CommentTok{\# Recortar filogenia}
\NormalTok{match.species}\OtherTok{\textless{}{-}}\FunctionTok{treedata}\NormalTok{(tree,dados)}
\NormalTok{tree}\OtherTok{\textless{}{-}}\NormalTok{match.species}\SpecialCharTok{$}\NormalTok{phy}
\FunctionTok{plotTree}\NormalTok{(tree,}\AttributeTok{fsize=}\FloatTok{0.5}\NormalTok{,}\AttributeTok{ftype=}\StringTok{"i"}\NormalTok{,}\AttributeTok{type=}\StringTok{"fan"}\NormalTok{,}\AttributeTok{lwd=}\DecValTok{1}\NormalTok{)}
\end{Highlighting}
\end{Shaded}

\includegraphics{04.6-pcmcordisc_files/figure-latex/unnamed-chunk-1-2.pdf}

\begin{Shaded}
\begin{Highlighting}[]
\CommentTok{\# Grupo Social}
\NormalTok{binsocial}\OtherTok{\textless{}{-}}\NormalTok{dados}\SpecialCharTok{$}\NormalTok{BinSocial}
\FunctionTok{names}\NormalTok{(binsocial)}\OtherTok{=}\FunctionTok{rownames}\NormalTok{(dados)}
\NormalTok{binsocial}\OtherTok{\textless{}{-}}\FunctionTok{as.factor}\NormalTok{(binsocial)}
\NormalTok{binsocial}
\CommentTok{\# Area de Vida}
\NormalTok{area}\OtherTok{\textless{}{-}}\NormalTok{dados}\SpecialCharTok{$}\NormalTok{AreaVida}
\FunctionTok{names}\NormalTok{(area)}\OtherTok{=}\FunctionTok{rownames}\NormalTok{(dados)}
\NormalTok{area}\OtherTok{\textless{}{-}}\FunctionTok{as.factor}\NormalTok{(area)}
\NormalTok{area}
\end{Highlighting}
\end{Shaded}

O primeiro passo é visualizar os dados junto com a árvore.

\begin{Shaded}
\begin{Highlighting}[]
\CommentTok{\# Visualizando os dados binários}
\FunctionTok{par}\NormalTok{(}\AttributeTok{mfrow=}\FunctionTok{c}\NormalTok{(}\DecValTok{1}\NormalTok{,}\DecValTok{2}\NormalTok{))}
\NormalTok{\{}\FunctionTok{plot}\NormalTok{(tree,}\AttributeTok{show.tip.label=}\ConstantTok{FALSE}\NormalTok{,}\AttributeTok{no.margin=}\ConstantTok{TRUE}\NormalTok{)}
\NormalTok{cores}\OtherTok{\textless{}{-}}\FunctionTok{setNames}\NormalTok{(}\FunctionTok{c}\NormalTok{(}\StringTok{"blue"}\NormalTok{,}\StringTok{"red"}\NormalTok{),}\FunctionTok{levels}\NormalTok{(binsocial))}
\FunctionTok{tiplabels}\NormalTok{(}\AttributeTok{pie=}\FunctionTok{to.matrix}\NormalTok{(binsocial[tree}\SpecialCharTok{$}\NormalTok{tip.label],}\FunctionTok{levels}\NormalTok{(binsocial)),}
 \AttributeTok{piecol=}\NormalTok{cores,}\AttributeTok{cex=}\FloatTok{0.5}\NormalTok{)}
\FunctionTok{legend}\NormalTok{(}\StringTok{"bottomleft"}\NormalTok{,}\AttributeTok{legend=}\FunctionTok{c}\NormalTok{(}\StringTok{"Social"}\NormalTok{,}\StringTok{"Solitary"}\NormalTok{),}\AttributeTok{pch=}\DecValTok{15}\NormalTok{,}\AttributeTok{col=}\NormalTok{cores,}\AttributeTok{bty=}\StringTok{"n"}\NormalTok{)}
\FunctionTok{plot}\NormalTok{(tree,}\AttributeTok{show.tip.label=}\ConstantTok{FALSE}\NormalTok{,}\AttributeTok{no.margin=}\ConstantTok{TRUE}\NormalTok{,}\AttributeTok{direction=}\StringTok{"leftwards"}\NormalTok{)}
\NormalTok{cores}\OtherTok{\textless{}{-}}\FunctionTok{setNames}\NormalTok{(}\FunctionTok{c}\NormalTok{(}\StringTok{"blue"}\NormalTok{,}\StringTok{"red"}\NormalTok{),}\FunctionTok{levels}\NormalTok{(area))}
\FunctionTok{tiplabels}\NormalTok{(}\AttributeTok{pie=}\FunctionTok{to.matrix}\NormalTok{(area[tree}\SpecialCharTok{$}\NormalTok{tip.label],}\FunctionTok{levels}\NormalTok{(area)),}
 \AttributeTok{piecol=}\NormalTok{cores,}\AttributeTok{cex=}\FloatTok{0.5}\NormalTok{)}
\FunctionTok{legend}\NormalTok{(}\StringTok{"bottomright"}\NormalTok{,}\AttributeTok{legend=}\FunctionTok{c}\NormalTok{(}\StringTok{"Grande"}\NormalTok{,}\StringTok{"Pequena"}\NormalTok{),}\AttributeTok{pch=}\DecValTok{15}\NormalTok{,}\AttributeTok{col=}\NormalTok{cores,}\AttributeTok{bty=}\StringTok{"n"}\NormalTok{)\}}
\end{Highlighting}
\end{Shaded}

\includegraphics{04.6-pcmcordisc_files/figure-latex/unnamed-chunk-2-1.pdf}

\begin{Shaded}
\begin{Highlighting}[]
\FunctionTok{par}\NormalTok{(}\AttributeTok{mfrow=}\FunctionTok{c}\NormalTok{(}\DecValTok{1}\NormalTok{,}\DecValTok{1}\NormalTok{))}
\end{Highlighting}
\end{Shaded}

Agora, vamos implementar o teste de Pagel (1994), que compara um modelo de evolução independente entre os caracteres (ausência de correlação) com um modelo de evolução dependente (presença de correlação).

\begin{Shaded}
\begin{Highlighting}[]
\NormalTok{fit.AS}\OtherTok{\textless{}{-}}\FunctionTok{fitPagel}\NormalTok{(tree,binsocial,area,}\AttributeTok{model=}\StringTok{"ARD"}\NormalTok{)}
\NormalTok{fit.AS}
\FunctionTok{plot}\NormalTok{(fit.AS)}
\end{Highlighting}
\end{Shaded}

\includegraphics{04.6-pcmcordisc_files/figure-latex/unnamed-chunk-3-1.pdf}

\begin{Shaded}
\begin{Highlighting}[]
\FunctionTok{plot}\NormalTok{(fit.AS,}\AttributeTok{lwd.by.rate=}\ConstantTok{TRUE}\NormalTok{)}
\end{Highlighting}
\end{Shaded}

\includegraphics{04.6-pcmcordisc_files/figure-latex/unnamed-chunk-3-2.pdf}
O \emph{P} valor é um indício de que o modelo dependente é melhor (podemos rejeitar o modelo independente com \emph{P}\textless alpha).

Podemos comparar os modelos dependente e independente com modelos onde a taxa de mudança em Grupo Social depende de Área (dep.var=``binsocial'') mas não o contrário, ou onde a taxa de mudança em Área depende do Grupo Social (dep.var=``area'') mas não o contrário.

\begin{Shaded}
\begin{Highlighting}[]
\CommentTok{\# Dependente X {-} Mudança em Grupo Social depende da Area}
\NormalTok{fit.GS}\OtherTok{\textless{}{-}}\FunctionTok{fitPagel}\NormalTok{(tree,binsocial,area,}\AttributeTok{dep.var=}\StringTok{"x"}\NormalTok{)}
\NormalTok{fit.GS}
\FunctionTok{plot}\NormalTok{(fit.GS,}\AttributeTok{lwd.by.rate=}\ConstantTok{TRUE}\NormalTok{)}
\end{Highlighting}
\end{Shaded}

\includegraphics{04.6-pcmcordisc_files/figure-latex/unnamed-chunk-4-1.pdf}

\begin{Shaded}
\begin{Highlighting}[]
\CommentTok{\# Dependente Y {-} Mudança em Area depende de Grupo Social}
\NormalTok{fit.AR}\OtherTok{\textless{}{-}}\FunctionTok{fitPagel}\NormalTok{(tree,binsocial,area,}\AttributeTok{dep.var=}\StringTok{"y"}\NormalTok{)}
\NormalTok{fit.AR}
\FunctionTok{plot}\NormalTok{(fit.AR,}\AttributeTok{lwd.by.rate=}\ConstantTok{TRUE}\NormalTok{)}
\end{Highlighting}
\end{Shaded}

\includegraphics{04.6-pcmcordisc_files/figure-latex/unnamed-chunk-4-2.pdf}

Comparação de modelos.

\begin{Shaded}
\begin{Highlighting}[]
\NormalTok{aic.scores}\OtherTok{\textless{}{-}}\FunctionTok{setNames}\NormalTok{(}\FunctionTok{c}\NormalTok{(fit.AS}\SpecialCharTok{$}\NormalTok{independent.AIC,}
\NormalTok{ fit.GS}\SpecialCharTok{$}\NormalTok{dependent.AIC,}
\NormalTok{ fit.AR}\SpecialCharTok{$}\NormalTok{dependent.AIC,}
\NormalTok{ fit.AS}\SpecialCharTok{$}\NormalTok{dependent.AIC),}
 \FunctionTok{c}\NormalTok{(}\StringTok{"independente"}\NormalTok{,}\StringTok{"dependente Grupo"}\NormalTok{,}
 \StringTok{"dependente Area"}\NormalTok{,}\StringTok{"dependente Grupo\&Area"}\NormalTok{))}
\NormalTok{aic.scores}
\FunctionTok{aicw}\NormalTok{(aic.scores)}
\end{Highlighting}
\end{Shaded}

O maior peso está para o modelo de dependência de Área para Grupo, mas o modelo de dependência mútua também recebeu forte suporte.

\section{Correlação de fenótipos categóricos multiestado}\label{correlauxe7uxe3o-de-fenuxf3tipos-categuxf3ricos-multiestado}

A correlação entre atributos categóricos multiestado pode ser feita com a estatística D de Huelsenbeck, Nielsen, and Bollback (2003), implementada no pacote \emph{phytools} com a função Dtest.

\section{Exercício - Correlação de fenótipos discretos}\label{exercuxedcio---correlauxe7uxe3o-de-fenuxf3tipos-discretos}

Carregue os dados: Árvore: Pagel.tre Dados: Pagel.txt Teste a correlação entre os caracteres x1 e y1 usando fitPagel. Quais foram os resultados? Visualize a distribuição dos caracteres na árvore. O que isso te informa sobre os resultados?

\chapter{Evolução multi-taxa e multi-regime}\label{evoluuxe7uxe3o-multi-taxa-e-multi-regime}

Modelos mais complexos que permitem que os parâmetros σ2, α, e θ (taxa evolutiva, força de seleção e valor do ótimo, respectivamente) variem entre regimes evolutivos/ótimos ambientais podem ser implementados no R. É possível criar hipóteses onde os ótimos, as taxas ou os alpha variem, em qualquer combinação. No entanto, é preciso estar atento ao fato de que modelos mais complexos vão estimar muitos parâmetros e podem gerar \emph{over-fitting} se o número de parâmetros for parecido com o
número de observações (espécies).

Neste capítulo, vamos implementar modelos onde a variável resposta é univariada com o pacote \emph{OUwie}, que tem a seguinte estrutura para os argumentos/modelos:
1. BM1 - Modelo Browniano com um único valor de taxa evolutiva (rate) para toda a árvore;
2. BMS - Modelo Browniano com diferentes taxas evolutivas de acordo com os estados na árvore;
3. OU1 - Modelo Ornstein-Uhlenbeck com um único ótimo para todas as espécies;
4. OUM - Modelo OU com diferentes ótimos de acordo com os estados, mas um único alpha e sigma2 para todos os regimes seletivos;
4.1. OUMV - Modelo OU com diferentes médias de estados (ótimos), diferentes sigma2, mesmo alpha;
4.2. OUMA - Modelo OU com diferentes médias de estados (ótimos), diferentes alpha, mesmo sigma2;
4.3. OUMVA - Modelo OU com múltiplos theta, alpha e sigma2.

O OUwie precisa de um dataframe com três colunas: nomes das espécies, o código do regime, e o atributo contínuo de interesse.

\begin{Shaded}
\begin{Highlighting}[]
\FunctionTok{require}\NormalTok{(ape)}
\CommentTok{\#\textgreater{} Carregando pacotes exigidos: ape}
\FunctionTok{require}\NormalTok{(geiger)}
\CommentTok{\#\textgreater{} Carregando pacotes exigidos: geiger}
\CommentTok{\#\textgreater{} Carregando pacotes exigidos: phytools}
\CommentTok{\#\textgreater{} Carregando pacotes exigidos: maps}
\FunctionTok{require}\NormalTok{(phytools)}
\FunctionTok{require}\NormalTok{(OUwie)}
\CommentTok{\#\textgreater{} Carregando pacotes exigidos: OUwie}
\CommentTok{\#\textgreater{} Carregando pacotes exigidos: corpcor}
\CommentTok{\#\textgreater{} Carregando pacotes exigidos: nloptr}
\CommentTok{\#\textgreater{} Carregando pacotes exigidos: RColorBrewer}
\CommentTok{\# Carregar dados}
\NormalTok{dados}\OtherTok{\textless{}{-}}\FunctionTok{read.table}\NormalTok{(}\StringTok{"dadospcm/sigmodontinae{-}atributos.txt"}\NormalTok{,}\AttributeTok{h=}\NormalTok{T,}\AttributeTok{row.names=}\DecValTok{1}\NormalTok{)}
\FunctionTok{head}\NormalTok{(dados)}
\NormalTok{diet}\OtherTok{\textless{}{-}}\FunctionTok{as.factor}\NormalTok{(dados}\SpecialCharTok{$}\NormalTok{Diet) }\CommentTok{\# A dieta é um código (Onivoro{-}1; Herbivoro{-}2; Insetivoro{-}3)}
\FunctionTok{names}\NormalTok{(diet)}\OtherTok{=}\FunctionTok{rownames}\NormalTok{(dados)}
\NormalTok{size}\OtherTok{\textless{}{-}}\NormalTok{dados}\SpecialCharTok{$}\NormalTok{Size}
\FunctionTok{names}\NormalTok{(size)}\OtherTok{=}\FunctionTok{rownames}\NormalTok{(dados)}
\end{Highlighting}
\end{Shaded}

O OUwie requer uma filogenia onde os nós internos estejam rotulados com o regime seletivo ancestral, que foi o que fizemos na reconstrução ancestral de atributos discretos.

\begin{Shaded}
\begin{Highlighting}[]
\CommentTok{\# Carregar filogenia}
\NormalTok{tree}\OtherTok{\textless{}{-}}\FunctionTok{read.nexus}\NormalTok{(}\StringTok{"dadospcm/sigmodontinae{-}tree.tre"}\NormalTok{)}
\FunctionTok{plotTree}\NormalTok{(tree,}\AttributeTok{fsize=}\FloatTok{0.4}\NormalTok{,}\AttributeTok{ftype=}\StringTok{"i"}\NormalTok{,}\AttributeTok{type=}\StringTok{"fan"}\NormalTok{,}\AttributeTok{lwd=}\DecValTok{1}\NormalTok{)}
\end{Highlighting}
\end{Shaded}

\includegraphics{04.7-pcmmulti_files/figure-latex/unnamed-chunk-2-1.pdf}

\begin{Shaded}
\begin{Highlighting}[]
\FunctionTok{str}\NormalTok{(tree)}
\CommentTok{\# Match Species}
\NormalTok{match.species}\OtherTok{\textless{}{-}}\FunctionTok{treedata}\NormalTok{(tree,dados)}
\NormalTok{tree}\OtherTok{\textless{}{-}}\NormalTok{match.species}\SpecialCharTok{$}\NormalTok{phy}
\FunctionTok{str}\NormalTok{(tree)}
\FunctionTok{plotTree}\NormalTok{(tree,}\AttributeTok{fsize=}\FloatTok{0.5}\NormalTok{,}\AttributeTok{ftype=}\StringTok{"i"}\NormalTok{,}\AttributeTok{type=}\StringTok{"fan"}\NormalTok{,}\AttributeTok{lwd=}\DecValTok{1}\NormalTok{)}
\end{Highlighting}
\end{Shaded}

\includegraphics{04.7-pcmmulti_files/figure-latex/unnamed-chunk-2-2.pdf}

\begin{Shaded}
\begin{Highlighting}[]
\CommentTok{\# Reconstruir evolução da dieta}
\NormalTok{diet.anc}\OtherTok{\textless{}{-}}\FunctionTok{ace}\NormalTok{(diet,tree,}\AttributeTok{type=}\StringTok{"discrete"}\NormalTok{,}\AttributeTok{model=}\StringTok{"ER"}\NormalTok{)}
\NormalTok{diet.anc}
\end{Highlighting}
\end{Shaded}

O ace retorna a probabilidade de cada estado em cada nó. Precisamos de uma árvore única com o ``melhor'' estado para cada nó (alternativamente, podemos usar o make.simmap como fizemos anteriormente, e ajustar um modelo para cada N histórias). Nós podemos encontrar o estado com maior probabilidade em cada nó, e atribuir esse valor aos nós de uma árvore em formato phylo.

\begin{Shaded}
\begin{Highlighting}[]
\CommentTok{\# Primeiro, corrigimos a ordem númerica no nome dos nós para n+1}
\NormalTok{node.states}\OtherTok{\textless{}{-}}\NormalTok{diet.anc}\SpecialCharTok{$}\NormalTok{lik.anc}
\FunctionTok{rownames}\NormalTok{(node.states) }\OtherTok{\textless{}{-}} \FunctionTok{seq}\NormalTok{(}\DecValTok{1}\SpecialCharTok{:}\FunctionTok{nrow}\NormalTok{(node.states)) }\SpecialCharTok{+} \FunctionTok{length}\NormalTok{(tree}\SpecialCharTok{$}\NormalTok{tip.label)}
\NormalTok{node.states}
\CommentTok{\# Extraímos o estado mais provável de cada nó}
\NormalTok{best}\OtherTok{\textless{}{-}}\FunctionTok{apply}\NormalTok{(node.states,}\DecValTok{1}\NormalTok{,which.max)}
\NormalTok{best}
\CommentTok{\# Incluimos esses valores na árvore}
\NormalTok{tree}\SpecialCharTok{$}\NormalTok{node.label}\OtherTok{\textless{}{-}}\NormalTok{best}
\NormalTok{tree}
\end{Highlighting}
\end{Shaded}

\section{Ajuste de modelos evolutivos}\label{ajuste-de-modelos-evolutivos}

\subsection{Browniano com taxa única}\label{browniano-com-taxa-uxfanica}

\begin{Shaded}
\begin{Highlighting}[]
\NormalTok{BM1}\OtherTok{\textless{}{-}}\FunctionTok{OUwie}\NormalTok{(tree,dados,}\AttributeTok{model=}\StringTok{"BM1"}\NormalTok{)}
\NormalTok{BM1}
\end{Highlighting}
\end{Shaded}

Apha=NA pois este é um modelo Browniano. Mesmo valor de sigma.sq = taxa única. Optima é o valor estimado na raiz.

\subsection{Browniano com múltiplas taxas}\label{browniano-com-muxfaltiplas-taxas}

\begin{Shaded}
\begin{Highlighting}[]
\NormalTok{BMS}\OtherTok{\textless{}{-}}\FunctionTok{OUwie}\NormalTok{(tree,dados,}\AttributeTok{model=}\StringTok{"BMS"}\NormalTok{)}
\NormalTok{BMS}
\end{Highlighting}
\end{Shaded}

Taxa de evolução de tamanho para insetívoros é mais rápida que para outros grupos. Mas ainda não sabemos se esse é o modelo evolutivo com maior suporte.

\subsection{OU com um único ótimo}\label{ou-com-um-uxfanico-uxf3timo}

\begin{Shaded}
\begin{Highlighting}[]
\NormalTok{OU1}\OtherTok{\textless{}{-}}\FunctionTok{OUwie}\NormalTok{(tree,dados,}\AttributeTok{model=}\StringTok{"OU1"}\NormalTok{)}
\NormalTok{OU1}
\end{Highlighting}
\end{Shaded}

\subsection{OU com múltiplos ótimos}\label{ou-com-muxfaltiplos-uxf3timos}

\begin{Shaded}
\begin{Highlighting}[]
\NormalTok{OUM}\OtherTok{\textless{}{-}}\FunctionTok{OUwie}\NormalTok{(tree,dados,}\AttributeTok{model=}\StringTok{"OUM"}\NormalTok{)}
\NormalTok{OUM}
\end{Highlighting}
\end{Shaded}

Tamanho `ótimo' para cada dieta é estimado, porém o alpha neste caso é muito próximo de zero.

\subsection{OU com múltiplos ótimos, taxas, e forças de seleção}\label{ou-com-muxfaltiplos-uxf3timos-taxas-e-foruxe7as-de-seleuxe7uxe3o}

\begin{Shaded}
\begin{Highlighting}[]
\CommentTok{\# OU múltiplos otimos, multiplos sigma2}
\NormalTok{OUMV}\OtherTok{\textless{}{-}}\FunctionTok{OUwie}\NormalTok{(tree,dados,}\AttributeTok{model=}\StringTok{"OUMV"}\NormalTok{)}
\NormalTok{OUMV}
\CommentTok{\# OU múltiplos otimos, multiplos alpha}
\NormalTok{OUMA}\OtherTok{\textless{}{-}}\FunctionTok{OUwie}\NormalTok{(tree,dados,}\AttributeTok{model=}\StringTok{"OUMA"}\NormalTok{)}
\NormalTok{OUMA}
\CommentTok{\# OU multiplos otimos, alpha e sigma2}
\NormalTok{OUMVA}\OtherTok{\textless{}{-}}\FunctionTok{OUwie}\NormalTok{(tree,dados,}\AttributeTok{model=}\StringTok{"OUMVA"}\NormalTok{)}
\NormalTok{OUMVA}
\end{Highlighting}
\end{Shaded}

\subsection{Comparação de modelos}\label{comparauxe7uxe3o-de-modelos-2}

\begin{Shaded}
\begin{Highlighting}[]
\NormalTok{aic.scores}\OtherTok{\textless{}{-}}
\FunctionTok{setNames}\NormalTok{(}\FunctionTok{c}\NormalTok{(BM1}\SpecialCharTok{$}\NormalTok{AICc,BMS}\SpecialCharTok{$}\NormalTok{AICc,OU1}\SpecialCharTok{$}\NormalTok{AICc,OUM}\SpecialCharTok{$}\NormalTok{AICc,OUMV}\SpecialCharTok{$}\NormalTok{AICc,OUMA}\SpecialCharTok{$}\NormalTok{AICc,OUMVA}\SpecialCharTok{$}\NormalTok{AICc),}
 \FunctionTok{c}\NormalTok{(}\StringTok{"BM1"}\NormalTok{,}\StringTok{"BMS"}\NormalTok{,}\StringTok{"OU1"}\NormalTok{,}\StringTok{"OUM"}\NormalTok{,}\StringTok{"OUMV"}\NormalTok{,}\StringTok{"OUMA"}\NormalTok{,}\StringTok{"OUMVA"}\NormalTok{))}
\NormalTok{aic.scores}
\FunctionTok{aicw}\NormalTok{(aic.scores) }\CommentTok{\# Delta AICc e AICc weights}
\end{Highlighting}
\end{Shaded}

\chapter{Métodos filogenéticos multivariados}\label{muxe9todos-filogenuxe9ticos-multivariados}

A evolução morfológica é inerentemente multivariada \hyperref[cap3]{Capítulo 3}. Contudo, variáveis resposta multivariadas representam um desafio analítico, especialmente para dados altamente multidimensionais. Neste capítulo, vamos explorar alguns métodos que permitem responder questões biológicas interessantes quando o nosso fenótipo de interesse é multivariado.

Primeiro, vamos carregar uma matriz resposta multivariada e uma filogenia. No exemplo, a matriz contém 26 colunas com coordenadas de Procrustes que descrevem a forma da mandíbula de 176 espécies de roedores sigmodontineos.

\begin{Shaded}
\begin{Highlighting}[]
\FunctionTok{require}\NormalTok{(ape)}
\CommentTok{\#\textgreater{} Carregando pacotes exigidos: ape}
\FunctionTok{require}\NormalTok{(geiger)}
\CommentTok{\#\textgreater{} Carregando pacotes exigidos: geiger}
\CommentTok{\#\textgreater{} Carregando pacotes exigidos: phytools}
\CommentTok{\#\textgreater{} Carregando pacotes exigidos: maps}
\FunctionTok{require}\NormalTok{(phytools)}
\FunctionTok{require}\NormalTok{(geomorph)}
\CommentTok{\#\textgreater{} Carregando pacotes exigidos: geomorph}
\CommentTok{\#\textgreater{} Carregando pacotes exigidos: RRPP}
\CommentTok{\#\textgreater{} Carregando pacotes exigidos: rgl}
\CommentTok{\#\textgreater{} Carregando pacotes exigidos: Matrix}
\CommentTok{\# Carregar dados}
\NormalTok{dados}\OtherTok{\textless{}{-}}\FunctionTok{read.table}\NormalTok{(}\StringTok{"dadospcm/sigmodontinae{-}shape.txt"}\NormalTok{,}\AttributeTok{h=}\NormalTok{T,}\AttributeTok{row.names=}\DecValTok{1}\NormalTok{)}
\FunctionTok{head}\NormalTok{(dados)}
\NormalTok{diet}\OtherTok{\textless{}{-}}\FunctionTok{as.factor}\NormalTok{(dados}\SpecialCharTok{$}\NormalTok{Diet) }\CommentTok{\# A dieta é um código (Onivoro{-}1; Herbivoro{-}2; Insetivoro{-}3)}
\FunctionTok{names}\NormalTok{(diet)}\OtherTok{=}\FunctionTok{rownames}\NormalTok{(dados)}
\NormalTok{forma}\OtherTok{\textless{}{-}}\NormalTok{dados[,}\DecValTok{3}\SpecialCharTok{:}\DecValTok{28}\NormalTok{]}
\NormalTok{forma}
\CommentTok{\# Carregar filogenia}
\NormalTok{tree}\OtherTok{\textless{}{-}}\FunctionTok{read.nexus}\NormalTok{(}\StringTok{"dadospcm/sigmodontinae{-}tree.tre"}\NormalTok{)}
\FunctionTok{plotTree}\NormalTok{(tree,}\AttributeTok{fsize=}\FloatTok{0.4}\NormalTok{,}\AttributeTok{ftype=}\StringTok{"i"}\NormalTok{,}\AttributeTok{type=}\StringTok{"fan"}\NormalTok{,}\AttributeTok{lwd=}\DecValTok{1}\NormalTok{)}
\end{Highlighting}
\end{Shaded}

\includegraphics{04.8-pcmmvar_files/figure-latex/unnamed-chunk-1-1.pdf}

\begin{Shaded}
\begin{Highlighting}[]
\CommentTok{\# Match Species}
\NormalTok{match.species}\OtherTok{\textless{}{-}}\FunctionTok{treedata}\NormalTok{(tree,dados)}
\NormalTok{tree}\OtherTok{\textless{}{-}}\NormalTok{match.species}\SpecialCharTok{$}\NormalTok{phy}
\FunctionTok{str}\NormalTok{(tree)}
\FunctionTok{plotTree}\NormalTok{(tree,}\AttributeTok{fsize=}\FloatTok{0.5}\NormalTok{,}\AttributeTok{ftype=}\StringTok{"i"}\NormalTok{,}\AttributeTok{type=}\StringTok{"fan"}\NormalTok{,}\AttributeTok{lwd=}\DecValTok{1}\NormalTok{)}
\end{Highlighting}
\end{Shaded}

\includegraphics{04.8-pcmmvar_files/figure-latex/unnamed-chunk-1-2.pdf}

\section{Sinal filogenético}\label{sinal-filogenuxe9tico-1}

\begin{Shaded}
\begin{Highlighting}[]
\FunctionTok{physignal}\NormalTok{(}\FunctionTok{as.matrix}\NormalTok{(forma),tree)}
\end{Highlighting}
\end{Shaded}

\section{Filomorfoespaço}\label{filomorfoespauxe7o-1}

\begin{Shaded}
\begin{Highlighting}[]
\CommentTok{\# Phylomorphospace}
\NormalTok{PCA}\OtherTok{\textless{}{-}}\FunctionTok{princomp}\NormalTok{(forma)}
\FunctionTok{par}\NormalTok{(}\AttributeTok{mfrow=}\FunctionTok{c}\NormalTok{(}\DecValTok{1}\NormalTok{,}\DecValTok{1}\NormalTok{),}\AttributeTok{mar=}\FunctionTok{c}\NormalTok{(}\DecValTok{5}\NormalTok{,}\DecValTok{4}\NormalTok{,}\DecValTok{2}\NormalTok{,}\DecValTok{2}\NormalTok{))}
\FunctionTok{phylomorphospace}\NormalTok{(tree,PCA}\SpecialCharTok{$}\NormalTok{scores[,}\DecValTok{1}\SpecialCharTok{:}\DecValTok{2}\NormalTok{],}\AttributeTok{label=}\StringTok{"off"}\NormalTok{,}\AttributeTok{control=}\FunctionTok{list}\NormalTok{())}
\end{Highlighting}
\end{Shaded}

\includegraphics{04.8-pcmmvar_files/figure-latex/unnamed-chunk-3-1.pdf}

\begin{Shaded}
\begin{Highlighting}[]
\CommentTok{\# terminais com cores}
\NormalTok{\{}\FunctionTok{phylomorphospace}\NormalTok{(tree,PCA}\SpecialCharTok{$}\NormalTok{scores[,}\DecValTok{1}\SpecialCharTok{:}\DecValTok{2}\NormalTok{],}\AttributeTok{label=}\StringTok{"off"}\NormalTok{,}\AttributeTok{control=}\FunctionTok{list}\NormalTok{(}\AttributeTok{col.node=}\ConstantTok{NULL}\NormalTok{))}
\NormalTok{cores}\OtherTok{\textless{}{-}}\FunctionTok{setNames}\NormalTok{(}\FunctionTok{c}\NormalTok{(}\StringTok{"blue"}\NormalTok{,}\StringTok{"red"}\NormalTok{,}\StringTok{"green"}\NormalTok{),}\FunctionTok{levels}\NormalTok{(diet))}
\NormalTok{col.group}\OtherTok{\textless{}{-}}\NormalTok{cores[}\FunctionTok{match}\NormalTok{(diet,}\FunctionTok{names}\NormalTok{(cores))]}
\FunctionTok{points}\NormalTok{(PCA}\SpecialCharTok{$}\NormalTok{scores[,}\DecValTok{1}\NormalTok{],PCA}\SpecialCharTok{$}\NormalTok{scores[,}\DecValTok{2}\NormalTok{],}\AttributeTok{col=}\NormalTok{col.group,}\AttributeTok{pch=}\DecValTok{16}\NormalTok{,}\AttributeTok{cex=}\FloatTok{1.7}\NormalTok{)}
\FunctionTok{legend}\NormalTok{(}\SpecialCharTok{{-}}\FloatTok{0.10}\NormalTok{,}\SpecialCharTok{{-}}\FloatTok{0.06}\NormalTok{,}\AttributeTok{legend=}\FunctionTok{c}\NormalTok{(}\StringTok{"On"}\NormalTok{,}\StringTok{"Hb"}\NormalTok{,}\StringTok{"In"}\NormalTok{),}\AttributeTok{pch=}\DecValTok{19}\NormalTok{,}\AttributeTok{col=}\FunctionTok{c}\NormalTok{(}\StringTok{"blue"}\NormalTok{,}\StringTok{"red"}\NormalTok{,}\StringTok{"green"}\NormalTok{))\}}
\end{Highlighting}
\end{Shaded}

\includegraphics{04.8-pcmmvar_files/figure-latex/unnamed-chunk-3-2.pdf}

\begin{Shaded}
\begin{Highlighting}[]
\CommentTok{\# 3D}
\CommentTok{\#phylomorphospace3d(tree,PCA$scores[,1:3],control=list(ftype="off"))}
\end{Highlighting}
\end{Shaded}

\section{PCA filogenética}\label{pca-filogenuxe9tica-1}

\begin{Shaded}
\begin{Highlighting}[]
\NormalTok{pPCA}\OtherTok{\textless{}{-}}\FunctionTok{phyl.pca}\NormalTok{(tree,forma,}\AttributeTok{method=}\StringTok{"lambda"}\NormalTok{,}\AttributeTok{opt=}\StringTok{"REML"}\NormalTok{)}
\NormalTok{\{}\FunctionTok{phylomorphospace}\NormalTok{(tree,pPCA}\SpecialCharTok{$}\NormalTok{S[,}\DecValTok{1}\SpecialCharTok{:}\DecValTok{2}\NormalTok{],}\AttributeTok{label=}\StringTok{"off"}\NormalTok{)}
\FunctionTok{points}\NormalTok{(pPCA}\SpecialCharTok{$}\NormalTok{S[,}\DecValTok{1}\NormalTok{],pPCA}\SpecialCharTok{$}\NormalTok{S[,}\DecValTok{2}\NormalTok{],}\AttributeTok{col=}\NormalTok{col.group,}\AttributeTok{pch=}\DecValTok{16}\NormalTok{,}\AttributeTok{cex=}\FloatTok{1.7}\NormalTok{)}
\FunctionTok{legend}\NormalTok{(}\SpecialCharTok{{-}}\FloatTok{0.10}\NormalTok{,}\SpecialCharTok{{-}}\FloatTok{0.05}\NormalTok{,}\AttributeTok{legend=}\FunctionTok{c}\NormalTok{(}\StringTok{"On"}\NormalTok{,}\StringTok{"Hb"}\NormalTok{,}\StringTok{"In"}\NormalTok{),}\AttributeTok{pch=}\DecValTok{19}\NormalTok{,}\AttributeTok{col=}\FunctionTok{c}\NormalTok{(}\StringTok{"blue"}\NormalTok{,}\StringTok{"red"}\NormalTok{,}\StringTok{"green"}\NormalTok{))\}}
\end{Highlighting}
\end{Shaded}

\includegraphics{04.8-pcmmvar_files/figure-latex/unnamed-chunk-4-1.pdf}

\section{PGLS}\label{pgls}

\begin{Shaded}
\begin{Highlighting}[]
\CommentTok{\# PGLS Multivariada}
\CommentTok{\# Ajuste de formato da matriz para tornar compatível com o geomorph}
\NormalTok{forma.array}\OtherTok{\textless{}{-}}\FunctionTok{arrayspecs}\NormalTok{(forma,}\AttributeTok{p=}\DecValTok{13}\NormalTok{,}\AttributeTok{k=}\DecValTok{2}\NormalTok{) }\CommentTok{\# transformação em array}
\NormalTok{forma}\OtherTok{\textless{}{-}}\FunctionTok{two.d.array}\NormalTok{(forma.array) }\CommentTok{\# transformar de volta em matriz 2D}
\CommentTok{\# PGLS}
\NormalTok{fit.pgls}\OtherTok{\textless{}{-}}\FunctionTok{procD.pgls}\NormalTok{(forma}\SpecialCharTok{\textasciitilde{}}\NormalTok{diet,tree)}
\FunctionTok{summary}\NormalTok{(fit.pgls)}
\end{Highlighting}
\end{Shaded}

\section{Comparação de taxas evolutivas}\label{comparauxe7uxe3o-de-taxas-evolutivas}

\begin{Shaded}
\begin{Highlighting}[]
\NormalTok{forma.array}\OtherTok{\textless{}{-}}\FunctionTok{arrayspecs}\NormalTok{(forma,}\AttributeTok{p=}\DecValTok{13}\NormalTok{,}\AttributeTok{k=}\DecValTok{2}\NormalTok{) }\CommentTok{\# transformação em array}
\FunctionTok{compare.evol.rates}\NormalTok{(forma.array,tree,diet)}
\end{Highlighting}
\end{Shaded}

\section{Modelos contínuos multivariados}\label{modelos-contuxednuos-multivariados}

O pacote mvMORPH (\textbf{clavel2015?}) permite ajustar modelos evolutivos (e.g.~BM,OU,EB) onde a variável de interesse é multivariada, e comparar os modelos usando verosimilhança (\emph{likelihood}) e AIC. De forma mais acentuada do que no caso univariado, o número de parâmetros estimados pelo modelo pode ser um problema quando a variável de interesse é multi-dimensional (\emph{highly-dimensional datasets}), como no caso de dados obtidos por morfometria geométrica. A confiabilidade do ajuste de modelos para Componentes Principais (PCs) individuais é baixa e pode ser enganosa, levando ao suporte incorreto para EB e OU (\textbf{uyeda2015?}). A confiabilidade do ajuste de modelos para um \emph{conjunto} de PCs selecionados sob algum critério ainda é debatida. Atualmente, não existem métodos totalmente adequados (extensões diretas de métodos univariados) para ajustar modelos complexos para dados multidimensionais (\textbf{adams\&collyer2018?}), porém veja os modelos com verossimilhança penalizada (\textbf{clavel2019?}), (\textbf{clavel-morlon2020?}).

Como exemplo, e para reduzir o tempo computacional, vamos considerar os dois primeiros componentes principais da forma PCA\$scores{[},1:2{]}. Vamos trabalhar com modelos que são análogos ao univariados:
1. BM1 - Browniano taxa única
2. BMM - Browniano múltiplas taxas
3. EB - Early Burst (taxas evolutivas desaceleram com o tempo) 4. OU1 - OU com único ótimo adaptativo
5. OUM - OU com múltiplos ótimos

\begin{Shaded}
\begin{Highlighting}[]
\FunctionTok{require}\NormalTok{(mvMORPH)}
\CommentTok{\#\textgreater{} Carregando pacotes exigidos: mvMORPH}
\CommentTok{\#\textgreater{} Carregando pacotes exigidos: corpcor}
\CommentTok{\#\textgreater{} Carregando pacotes exigidos: subplex}
\CommentTok{\#\textgreater{} \#\#}
\CommentTok{\#\textgreater{} \#\# mvMORPH package (1.1.7)}
\CommentTok{\#\textgreater{} \#\# Multivariate evolutionary models}
\CommentTok{\#\textgreater{} \#\#}
\CommentTok{\#\textgreater{} \#\# See the tutorials: browseVignettes("mvMORPH")}
\CommentTok{\#\textgreater{} \#\#}
\CommentTok{\#\textgreater{} \#\# To cite package \textquotesingle{}mvMORPH\textquotesingle{}: citation("mvMORPH")}
\CommentTok{\#\textgreater{} \#\#}
\CommentTok{\#\textgreater{} }
\CommentTok{\#\textgreater{} Attaching package: \textquotesingle{}mvMORPH\textquotesingle{}}
\CommentTok{\#\textgreater{} The following object is masked from \textquotesingle{}package:geiger\textquotesingle{}:}
\CommentTok{\#\textgreater{} }
\CommentTok{\#\textgreater{}     aicw}
\CommentTok{\# Árvore filogenética com reconstrução em formato SIMMAP}
\NormalTok{anc.tree}\OtherTok{\textless{}{-}}\FunctionTok{make.simmap}\NormalTok{(tree,diet,}\AttributeTok{model=}\StringTok{"ER"}\NormalTok{,}\AttributeTok{nsim=}\DecValTok{1}\NormalTok{)}
\CommentTok{\#\textgreater{} Done.}
\NormalTok{anc.tree}
\FunctionTok{plot}\NormalTok{(anc.tree,cores,}\AttributeTok{fsize=}\FloatTok{0.5}\NormalTok{,}\AttributeTok{ftype=}\StringTok{"i"}\NormalTok{)}
\end{Highlighting}
\end{Shaded}

\includegraphics{04.8-pcmmvar_files/figure-latex/unnamed-chunk-7-1.pdf}

\begin{Shaded}
\begin{Highlighting}[]
\CommentTok{\# Ajuste modelos evolutivos multivariados}
\CommentTok{\# BM taxa única}
\NormalTok{multi.BM1}\OtherTok{\textless{}{-}}\FunctionTok{mvBM}\NormalTok{(anc.tree,PCA}\SpecialCharTok{$}\NormalTok{scores[,}\DecValTok{1}\SpecialCharTok{:}\DecValTok{2}\NormalTok{],}\AttributeTok{model=}\StringTok{"BM1"}\NormalTok{)}
\NormalTok{multi.BM1}
\CommentTok{\# BM múltiplas taxas}
\NormalTok{multi.BMM}\OtherTok{\textless{}{-}}\FunctionTok{mvBM}\NormalTok{(anc.tree,PCA}\SpecialCharTok{$}\NormalTok{scores[,}\DecValTok{1}\SpecialCharTok{:}\DecValTok{2}\NormalTok{],}\AttributeTok{model=}\StringTok{"BMM"}\NormalTok{)}
\NormalTok{multi.BMM}
\NormalTok{multi.BMM}\SpecialCharTok{$}\NormalTok{param}\SpecialCharTok{$}\NormalTok{nparam}
\CommentTok{\# Early{-}Burst}
\NormalTok{multi.EB}\OtherTok{\textless{}{-}}\FunctionTok{mvEB}\NormalTok{(anc.tree,PCA}\SpecialCharTok{$}\NormalTok{scores[,}\DecValTok{1}\SpecialCharTok{:}\DecValTok{2}\NormalTok{])}
\NormalTok{multi.EB}
\CommentTok{\# OU otimo unico}
\NormalTok{multi.OU1}\OtherTok{\textless{}{-}}\FunctionTok{mvOU}\NormalTok{(anc.tree,PCA}\SpecialCharTok{$}\NormalTok{scores[,}\DecValTok{1}\SpecialCharTok{:}\DecValTok{2}\NormalTok{],}\AttributeTok{model=}\StringTok{"OU1"}\NormalTok{)}
\NormalTok{multi.OU1}
\CommentTok{\# OU otimos multiplos}
\NormalTok{multi.OUM}\OtherTok{\textless{}{-}}\FunctionTok{mvOU}\NormalTok{(anc.tree,PCA}\SpecialCharTok{$}\NormalTok{scores[,}\DecValTok{1}\SpecialCharTok{:}\DecValTok{2}\NormalTok{],}\AttributeTok{model=}\StringTok{"OUM"}\NormalTok{)}
\NormalTok{multi.OUM}
\end{Highlighting}
\end{Shaded}

Comparando modelos.

\begin{Shaded}
\begin{Highlighting}[]
\CommentTok{\# Comparando modelos}
\NormalTok{aic.scores}\OtherTok{\textless{}{-}}
\FunctionTok{setNames}\NormalTok{(}\FunctionTok{c}\NormalTok{(multi.BM1}\SpecialCharTok{$}\NormalTok{AICc,multi.BMM}\SpecialCharTok{$}\NormalTok{AICc,multi.EB}\SpecialCharTok{$}\NormalTok{AICc,multi.OU1}\SpecialCharTok{$}\NormalTok{AICc,multi.OUM}\SpecialCharTok{$}\NormalTok{AICc),}
\FunctionTok{c}\NormalTok{(}\StringTok{"BM1"}\NormalTok{,}\StringTok{"BMM"}\NormalTok{,}\StringTok{"EB"}\NormalTok{,}\StringTok{"OU1"}\NormalTok{,}\StringTok{"OUM"}\NormalTok{))}
\NormalTok{aic.scores}
\FunctionTok{aicw}\NormalTok{(aic.scores) }\CommentTok{\# Delta AICc e AICc weights}
\FunctionTok{aic.w}\NormalTok{(aic.scores)}
\end{Highlighting}
\end{Shaded}

\section{Exercício - Modelos multivariados}\label{exercuxedcio---modelos-multivariados}

Com os mesmos dados de forma e dieta para roedores sigmodontíneos:
sigmodontinae-shape.txt
sigmodontinae-tree.tre
Simule 10 estimativas de estado ancestral para o caráter dieta. Usando 2 componentes principais da forma da mandíbula, ajuste modelos (com mvMORPH) brownianos com múltiplas taxas (BMM) e OU com múltiplos ótimos (OUM) usando os regimes reconstruídos em cada uma das 10 simulações. Calcule a média de AICc para cada simulação e compare o ajuste entre todos os modelos (BM1,BMM,EB,OU1,OUM).
Você pode usar, por exemplo:

list.AIC\textless-matrix(NA,length(trees),2)
for(i in 1:length(trees))\{
fit.BMM\textless-mvBM(trees{[}{[}i{]}{]},PCA\(scores[,1:2],model="BMM")
fit.OUM<-mvOU(trees[[i]],PCA\)scores{[},1:2{]},model=``OUM'')
list.AIC{[}i,1{]}\textless-fit.BMM\(AICc
list.AIC[i,2]<-fit.OUM\)AICc
\}

E a função apply para calcular a média do AICc.

\chapter{Disparidade fenotípica ao longo do tempo}\label{disparidade-fenotuxedpica-ao-longo-do-tempo}

Esse exemplo implementa a análise de disparidade morfológica ao longo do tempo (\emph{Disparity through time plots}) de Harmon et al. (2003).

Carregar dados.

\begin{Shaded}
\begin{Highlighting}[]
\FunctionTok{require}\NormalTok{(ape)}
\CommentTok{\#\textgreater{} Carregando pacotes exigidos: ape}
\FunctionTok{require}\NormalTok{(geiger)}
\CommentTok{\#\textgreater{} Carregando pacotes exigidos: geiger}
\CommentTok{\#\textgreater{} Carregando pacotes exigidos: phytools}
\CommentTok{\#\textgreater{} Carregando pacotes exigidos: maps}
\FunctionTok{require}\NormalTok{(phytools)}
\CommentTok{\# Filogenia}
\NormalTok{tree}\OtherTok{\textless{}{-}}\FunctionTok{read.nexus}\NormalTok{(}\StringTok{"dadospcm/sigmodontinae{-}tree.tre"}\NormalTok{)}
\FunctionTok{plotTree}\NormalTok{(tree,}\AttributeTok{fsize=}\FloatTok{0.4}\NormalTok{,}\AttributeTok{ftype=}\StringTok{"i"}\NormalTok{,}\AttributeTok{type=}\StringTok{"fan"}\NormalTok{,}\AttributeTok{lwd=}\DecValTok{1}\NormalTok{)}
\end{Highlighting}
\end{Shaded}

\includegraphics{04.9-pcmdtt_files/figure-latex/unnamed-chunk-1-1.pdf}

\begin{Shaded}
\begin{Highlighting}[]
\CommentTok{\# Dados morfológicos}
\NormalTok{dados}\OtherTok{\textless{}{-}}\FunctionTok{read.table}\NormalTok{(}\StringTok{"dadospcm/sigmodontinae{-}atributos.txt"}\NormalTok{,}\AttributeTok{h=}\NormalTok{T,}\AttributeTok{row.names=}\DecValTok{1}\NormalTok{)}
\FunctionTok{head}\NormalTok{(dados)}
\NormalTok{size}\OtherTok{\textless{}{-}}\NormalTok{dados}\SpecialCharTok{$}\NormalTok{Size}
\FunctionTok{names}\NormalTok{(size)}\OtherTok{=}\FunctionTok{rownames}\NormalTok{(dados)}
\NormalTok{size}
\end{Highlighting}
\end{Shaded}

\emph{Disparity through time plots}.

\begin{Shaded}
\begin{Highlighting}[]
\CommentTok{\# Tamanho}
\NormalTok{dtt.size}\OtherTok{\textless{}{-}}\FunctionTok{dtt}\NormalTok{(tree,size,}\AttributeTok{nsim=}\DecValTok{999}\NormalTok{,}\AttributeTok{calculateMDIp =} \ConstantTok{TRUE}\NormalTok{)}
\end{Highlighting}
\end{Shaded}

\includegraphics{04.9-pcmdtt_files/figure-latex/unnamed-chunk-2-1.pdf}

\begin{Shaded}
\begin{Highlighting}[]
\NormalTok{dtt.size}\SpecialCharTok{$}\NormalTok{MDI}
\NormalTok{dtt.size}\SpecialCharTok{$}\NormalTok{MDIpVal}
\CommentTok{\# DTT multivariado}
\NormalTok{dados}\OtherTok{\textless{}{-}}\FunctionTok{read.table}\NormalTok{(}\StringTok{"dadospcm/sigmodontinae{-}shape.txt"}\NormalTok{,}\AttributeTok{h=}\NormalTok{T,}\AttributeTok{row.names=}\DecValTok{1}\NormalTok{)}
\FunctionTok{head}\NormalTok{(dados)}
\NormalTok{diet}\OtherTok{\textless{}{-}}\FunctionTok{as.factor}\NormalTok{(dados}\SpecialCharTok{$}\NormalTok{Diet) }\CommentTok{\# A dieta é um código (Onivoro{-}1; Herbivoro{-}2; Insetivoro{-}3)}
\FunctionTok{names}\NormalTok{(diet)}\OtherTok{=}\FunctionTok{rownames}\NormalTok{(dados)}
\NormalTok{forma}\OtherTok{\textless{}{-}}\NormalTok{dados[,}\DecValTok{3}\SpecialCharTok{:}\DecValTok{28}\NormalTok{]}
\NormalTok{forma}
\NormalTok{dtt.shape}\OtherTok{\textless{}{-}}\FunctionTok{dtt}\NormalTok{(tree,}\FunctionTok{as.matrix}\NormalTok{(forma),}\AttributeTok{nsim=}\DecValTok{999}\NormalTok{,}\AttributeTok{calculateMDIp =} \ConstantTok{TRUE}\NormalTok{)}
\end{Highlighting}
\end{Shaded}

\includegraphics{04.9-pcmdtt_files/figure-latex/unnamed-chunk-2-2.pdf}

\begin{Shaded}
\begin{Highlighting}[]
\NormalTok{dtt.shape}\SpecialCharTok{$}\NormalTok{MDI}
\NormalTok{dtt.shape}\SpecialCharTok{$}\NormalTok{MDIpVal}
\end{Highlighting}
\end{Shaded}

\chapter{Diversificação de linhagens}\label{diversificauxe7uxe3o-de-linhagens}

Esse exemplo implementa os gráficos de acúmulo de linhagens ao longo do tempo (\emph{lineage through time plots}) e a estatística γ. Aqui, estamos assumindo especiação constante ao longo do tempo e entre linhagens (\emph{pure-birth}), e nenhuma extinção.

Carregar hipótese filogenética.

\begin{Shaded}
\begin{Highlighting}[]
\FunctionTok{require}\NormalTok{(ape)}
\CommentTok{\#\textgreater{} Carregando pacotes exigidos: ape}
\FunctionTok{require}\NormalTok{(geiger)}
\CommentTok{\#\textgreater{} Carregando pacotes exigidos: geiger}
\CommentTok{\#\textgreater{} Carregando pacotes exigidos: phytools}
\CommentTok{\#\textgreater{} Carregando pacotes exigidos: maps}
\FunctionTok{require}\NormalTok{(phytools)}
\CommentTok{\# Filogenia}
\NormalTok{tree}\OtherTok{\textless{}{-}}\FunctionTok{read.nexus}\NormalTok{(}\StringTok{"dadospcm/sigmodontinae{-}tree.tre"}\NormalTok{)}
\FunctionTok{plotTree}\NormalTok{(tree,}\AttributeTok{fsize=}\FloatTok{0.4}\NormalTok{,}\AttributeTok{ftype=}\StringTok{"i"}\NormalTok{,}\AttributeTok{type=}\StringTok{"fan"}\NormalTok{,}\AttributeTok{lwd=}\DecValTok{1}\NormalTok{)}
\end{Highlighting}
\end{Shaded}

\includegraphics{04.10-pcmltt_files/figure-latex/unnamed-chunk-1-1.pdf}
A árvore tem 279 espécies, de um total de 400 (70\%). Idealmente deveria haver 100\% de amostragem.

\emph{Lineage through time plots}.

\begin{Shaded}
\begin{Highlighting}[]
\CommentTok{\# LTT plot}
\FunctionTok{par}\NormalTok{(}\AttributeTok{mfrow=}\FunctionTok{c}\NormalTok{(}\DecValTok{1}\NormalTok{,}\DecValTok{1}\NormalTok{),}\AttributeTok{mar=}\FunctionTok{c}\NormalTok{(}\DecValTok{5}\NormalTok{,}\DecValTok{4}\NormalTok{,}\DecValTok{2}\NormalTok{,}\DecValTok{2}\NormalTok{))}
\NormalTok{obj.ltt}\OtherTok{\textless{}{-}}\FunctionTok{ltt}\NormalTok{(tree,}\AttributeTok{log.lineages=}\ConstantTok{FALSE}\NormalTok{)}
\NormalTok{obj.ltt}
\CommentTok{\# Tree background}
\FunctionTok{plot}\NormalTok{(obj.ltt,}\AttributeTok{show.tree=}\ConstantTok{TRUE}\NormalTok{,}\AttributeTok{log.lineages=}\ConstantTok{FALSE}\NormalTok{)}
\end{Highlighting}
\end{Shaded}

\includegraphics{04.10-pcmltt_files/figure-latex/unnamed-chunk-2-1.pdf}

Sob \emph{pure-birth} (só especiação) o número de linhagens deve crescer exponencialmente com o tempo. Isso se traduz em uma expectativa linear em escala semi-logarítmica: o eixo vertical em log e o horizontal em escala linear.

\begin{Shaded}
\begin{Highlighting}[]
\NormalTok{\{}\FunctionTok{plot}\NormalTok{(obj.ltt,}\AttributeTok{log.lineages=}\ConstantTok{FALSE}\NormalTok{,}\AttributeTok{log=}\StringTok{"y"}\NormalTok{,}\AttributeTok{ylim=}\FunctionTok{c}\NormalTok{(}\DecValTok{2}\NormalTok{,}\FunctionTok{Ntip}\NormalTok{(tree)))}
\CommentTok{\# comparando com o que seria esperado sob pure{-}birth (linha reta)}
\NormalTok{h}\OtherTok{\textless{}{-}}\FunctionTok{max}\NormalTok{(}\FunctionTok{nodeHeights}\NormalTok{(tree)) }\CommentTok{\# idade máxima da arvore}
\NormalTok{x}\OtherTok{\textless{}{-}}\FunctionTok{seq}\NormalTok{(}\DecValTok{0}\NormalTok{,h,}\AttributeTok{by=}\NormalTok{h}\SpecialCharTok{/}\DecValTok{100}\NormalTok{)}
\NormalTok{b}\OtherTok{\textless{}{-}}\NormalTok{(}\FunctionTok{log}\NormalTok{(}\FunctionTok{Ntip}\NormalTok{(tree))}\SpecialCharTok{{-}}\FunctionTok{log}\NormalTok{(}\DecValTok{2}\NormalTok{))}\SpecialCharTok{/}\NormalTok{h }\CommentTok{\# birth{-}rate}
\FunctionTok{lines}\NormalTok{(x,}\DecValTok{2}\SpecialCharTok{*}\FunctionTok{exp}\NormalTok{(b}\SpecialCharTok{*}\NormalTok{x),}\AttributeTok{col=}\StringTok{"red"}\NormalTok{,}\AttributeTok{lty=}\StringTok{"dashed"}\NormalTok{,}\AttributeTok{lwd=}\DecValTok{2}\NormalTok{)\}}
\end{Highlighting}
\end{Shaded}

\includegraphics{04.10-pcmltt_files/figure-latex/unnamed-chunk-3-1.pdf}

É possível comparar o LTT observado com múltiplas simulações de LTT sob um modelo \emph{pure-birth} assumindo o mesmo período de tempo e resultando no mesmo número de espécies que a nossa árvore de interesse.

\begin{Shaded}
\begin{Highlighting}[]
\NormalTok{trees}\OtherTok{\textless{}{-}}\FunctionTok{pbtree}\NormalTok{(}\AttributeTok{n=}\FunctionTok{Ntip}\NormalTok{(tree),}\AttributeTok{scale=}\NormalTok{h,}\AttributeTok{nsim=}\DecValTok{100}\NormalTok{,}\AttributeTok{quiet=}\ConstantTok{TRUE}\NormalTok{) }\CommentTok{\# nsim=999}
\NormalTok{obj.trees}\OtherTok{\textless{}{-}}\FunctionTok{ltt}\NormalTok{(trees,}\AttributeTok{plot=}\ConstantTok{FALSE}\NormalTok{)}
\NormalTok{\{}\FunctionTok{plot}\NormalTok{(obj.trees,}\AttributeTok{col=}\StringTok{"grey"}\NormalTok{)}
\FunctionTok{lines}\NormalTok{(}\FunctionTok{c}\NormalTok{(}\DecValTok{0}\NormalTok{,h),}\FunctionTok{log}\NormalTok{(}\FunctionTok{c}\NormalTok{(}\DecValTok{2}\NormalTok{,}\FunctionTok{Ntip}\NormalTok{(tree))),}\AttributeTok{col=}\StringTok{"blue"}\NormalTok{)}
\CommentTok{\# arvore original}
\FunctionTok{ltt}\NormalTok{(tree,}\AttributeTok{add=}\ConstantTok{TRUE}\NormalTok{,}\AttributeTok{col=}\StringTok{"red"}\NormalTok{,}\AttributeTok{lty=}\StringTok{"dashed"}\NormalTok{,}\AttributeTok{lwd=}\DecValTok{2}\NormalTok{)\}}
\end{Highlighting}
\end{Shaded}

\includegraphics{04.10-pcmltt_files/figure-latex/unnamed-chunk-4-1.pdf}

\begin{Shaded}
\begin{Highlighting}[]
\CommentTok{\# Intervalo de confiança 95\%}
\NormalTok{\{}\FunctionTok{ltt95}\NormalTok{(trees,}\AttributeTok{log=}\ConstantTok{TRUE}\NormalTok{)}
\FunctionTok{plot}\NormalTok{(obj.ltt,}\AttributeTok{add=}\ConstantTok{TRUE}\NormalTok{,}\AttributeTok{col=}\StringTok{"red"}\NormalTok{,}\AttributeTok{lty=}\StringTok{"dashed"}\NormalTok{,}\AttributeTok{lwd=}\DecValTok{2}\NormalTok{,}\AttributeTok{log.lineages=}\ConstantTok{FALSE}\NormalTok{)\}}
\end{Highlighting}
\end{Shaded}

\includegraphics{04.10-pcmltt_files/figure-latex/unnamed-chunk-4-2.pdf}

Estatística γ.

\begin{Shaded}
\begin{Highlighting}[]
\NormalTok{obj.ltt}\OtherTok{\textless{}{-}}\FunctionTok{ltt}\NormalTok{(tree,}\AttributeTok{log.lineages=}\ConstantTok{FALSE}\NormalTok{,}\AttributeTok{plot=}\ConstantTok{FALSE}\NormalTok{)}
\NormalTok{obj.ltt}
\FunctionTok{gammatest}\NormalTok{(obj.ltt)}
\end{Highlighting}
\end{Shaded}

Como exemplo, vamos simular uma árvore que retorne um gamma positivo.

\begin{Shaded}
\begin{Highlighting}[]
\CommentTok{\# Exemplo de gamma positivo}
\NormalTok{coal.tree}\OtherTok{\textless{}{-}}\FunctionTok{rcoal}\NormalTok{(}\AttributeTok{n=}\DecValTok{100}\NormalTok{)}
\FunctionTok{plotTree}\NormalTok{(coal.tree,}\AttributeTok{ftype=}\StringTok{"off"}\NormalTok{)}
\end{Highlighting}
\end{Shaded}

\includegraphics{04.10-pcmltt_files/figure-latex/unnamed-chunk-6-1.pdf}

\begin{Shaded}
\begin{Highlighting}[]
\FunctionTok{par}\NormalTok{(}\AttributeTok{mfrow=}\FunctionTok{c}\NormalTok{(}\DecValTok{1}\NormalTok{,}\DecValTok{1}\NormalTok{),}\AttributeTok{mar=}\FunctionTok{c}\NormalTok{(}\DecValTok{5}\NormalTok{,}\DecValTok{4}\NormalTok{,}\DecValTok{2}\NormalTok{,}\DecValTok{2}\NormalTok{))}
\NormalTok{obj}\OtherTok{\textless{}{-}}\FunctionTok{ltt}\NormalTok{(coal.tree,}\AttributeTok{log.lineages=}\ConstantTok{FALSE}\NormalTok{,}\AttributeTok{log=}\StringTok{"y"}\NormalTok{,}\AttributeTok{col=}\StringTok{"red"}\NormalTok{,}\AttributeTok{lty=}\StringTok{"dashed"}\NormalTok{)}
\end{Highlighting}
\end{Shaded}

\includegraphics{04.10-pcmltt_files/figure-latex/unnamed-chunk-6-2.pdf}

\begin{Shaded}
\begin{Highlighting}[]
\NormalTok{obj}
\CommentTok{\# comparação com pure{-}birth}
\NormalTok{trees}\OtherTok{\textless{}{-}}\FunctionTok{pbtree}\NormalTok{(}\AttributeTok{n=}\DecValTok{100}\NormalTok{,}\AttributeTok{nsim=}\DecValTok{200}\NormalTok{,}\AttributeTok{scale=}\FunctionTok{max}\NormalTok{(}\FunctionTok{nodeHeights}\NormalTok{(coal.tree)))}
\NormalTok{\{}\FunctionTok{ltt95}\NormalTok{(trees,}\AttributeTok{log=}\ConstantTok{TRUE}\NormalTok{)}
\FunctionTok{ltt}\NormalTok{(coal.tree,}\AttributeTok{add=}\ConstantTok{TRUE}\NormalTok{,}\AttributeTok{log.lineages=}\ConstantTok{FALSE}\NormalTok{,}\AttributeTok{col=}\StringTok{"red"}\NormalTok{,}\AttributeTok{lwd=}\DecValTok{2}\NormalTok{,}\AttributeTok{lty=}\StringTok{"dashed"}\NormalTok{)\}}
\end{Highlighting}
\end{Shaded}

\includegraphics{04.10-pcmltt_files/figure-latex/unnamed-chunk-6-3.pdf}

\chapter{Taxas de especiação e extinção constantes e variáveis}\label{taxas-de-especiauxe7uxe3o-e-extinuxe7uxe3o-constantes-e-variuxe1veis}

\section{Diversificação constante}\label{diversificauxe7uxe3o-constante}

Primeiro, vamos carregar alguns dados e calcular uma diversificação
líquida simples, esperada sobr um modelo \emph{pure-birth} (Magallon and Sanderson 2001).

\begin{Shaded}
\begin{Highlighting}[]
\FunctionTok{require}\NormalTok{(ape)}
\CommentTok{\#\textgreater{} Carregando pacotes exigidos: ape}
\FunctionTok{require}\NormalTok{(phytools)}
\CommentTok{\#\textgreater{} Carregando pacotes exigidos: phytools}
\CommentTok{\#\textgreater{} Carregando pacotes exigidos: maps}
\CommentTok{\# Filogenia}
\NormalTok{tree}\OtherTok{\textless{}{-}}\FunctionTok{read.tree}\NormalTok{(}\StringTok{"dadospcm/primate{-}tree.txt"}\NormalTok{)}
\FunctionTok{plotTree}\NormalTok{(tree,}\AttributeTok{fsize=}\FloatTok{0.4}\NormalTok{,}\AttributeTok{ftype=}\StringTok{"i"}\NormalTok{,}\AttributeTok{type=}\StringTok{"fan"}\NormalTok{,}\AttributeTok{lwd=}\DecValTok{1}\NormalTok{)}
\end{Highlighting}
\end{Shaded}

\includegraphics{04.11-pcmdiv_files/figure-latex/unnamed-chunk-1-1.pdf}

\begin{Shaded}
\begin{Highlighting}[]
\CommentTok{\# Número de espécies}
\NormalTok{n\_sp}\OtherTok{\textless{}{-}}\FunctionTok{length}\NormalTok{(tree}\SpecialCharTok{$}\NormalTok{tip.label)}
\CommentTok{\# Stem age}
\NormalTok{stem\_age}\OtherTok{\textless{}{-}}\FunctionTok{max}\NormalTok{(}\FunctionTok{nodeHeights}\NormalTok{(tree))}
\CommentTok{\# Diversificação (Magallon \& Sanderson 2001)}
\NormalTok{r1}\OtherTok{\textless{}{-}}\FunctionTok{log}\NormalTok{(n\_sp)}\SpecialCharTok{/}\NormalTok{stem\_age}
\CommentTok{\# Com essa diversificação obtemos o número de espécies visto na árvore}
\NormalTok{r1}\SpecialCharTok{*}\NormalTok{stem\_age}
\FunctionTok{log}\NormalTok{(n\_sp)}
\end{Highlighting}
\end{Shaded}

Agora, vamos estimar extinção constante com um modelo \emph{birth-death}. No pacote \emph{ape}, a função birthdeath estima taxas de extinção e especiação em uma filogenia (método de Nee, May, and Harvey (1994)). As estimativas são confiáveis se as premissas do modelo forem reais (especiação e extinção constantes no tempo). O método também assume uma filogenia completa. \emph{b} é a taxa de especiação e \emph{d} é a taxa de extinção.

\begin{Shaded}
\begin{Highlighting}[]
\CommentTok{\#nee\textless{}{-}birthdeath(tree)}
\end{Highlighting}
\end{Shaded}

A árvore contem algumas politomias. Vamos resolver as politomias e
tentar novamente.

\begin{Shaded}
\begin{Highlighting}[]
\NormalTok{tree}\OtherTok{\textless{}{-}}\FunctionTok{multi2di}\NormalTok{(tree)}
\FunctionTok{is.binary}\NormalTok{(tree)}
\NormalTok{nee}\OtherTok{\textless{}{-}}\FunctionTok{birthdeath}\NormalTok{(tree)}
\NormalTok{nee}
\end{Highlighting}
\end{Shaded}

Por motivos computacionais, a função birthdeath estima os parâmetros
derivados b-d e d/b (a diferença entre especiação e extinção (\emph{net
diversification}) e a razão entre extinção e especiação (\emph{relative
extinction rate}), respectivamente). Para encontrar b e d podemos usar a função bd do pacote \emph{phytools}:

\begin{Shaded}
\begin{Highlighting}[]
\FunctionTok{bd}\NormalTok{(nee)}
\end{Highlighting}
\end{Shaded}

As unidades dos valores dependem da unidade do comprimento dos ramos (ln para número de espécies).Só é possível comparar árvores diferentes quando o comprimento dos ramos está na mesma unidade.

Pacote \emph{diversitree}. O pacote diversitree possui múltiplas funções para realizar vários tipos de análises de diversificação com filogenias. Uma das vantagens do diversitree é que ele permite levar em consideração a incompletude no número de espécies amostradas. Na nossa filogenia dos primatas, existem 226 espécies amostradas de um total de aprox. 300 (o número variava de 190 até 350 da última vez que conferi; vamos considerar 300 para exemplo). Isso é cerca de 75\% de espécies presentes. O diversitree também permite ajustar modelos pure-birth (Yule) e birth-death (BD) e compará-los.

\begin{Shaded}
\begin{Highlighting}[]
\FunctionTok{require}\NormalTok{(diversitree)}
\CommentTok{\#\textgreater{} Carregando pacotes exigidos: diversitree}
\CommentTok{\# Yule (pure{-}birth {-}{-} somente especiação)}
\NormalTok{modelo.pb}\OtherTok{\textless{}{-}}\FunctionTok{make.yule}\NormalTok{(tree,}\AttributeTok{sampling.f =} \DecValTok{226}\SpecialCharTok{/}\DecValTok{300}\NormalTok{)}
\NormalTok{param.pb}\OtherTok{\textless{}{-}}\FunctionTok{find.mle}\NormalTok{(modelo.pb,}\AttributeTok{x.init=}\FloatTok{0.1}\NormalTok{) }\CommentTok{\# Estimar os valores dos parâmetros por ML \# x.init é um valor arbitrário para iniciar o otimização}
\NormalTok{param.pb}
\end{Highlighting}
\end{Shaded}

lambda = taxa de especiação; mu = taxa de extinção.

\begin{Shaded}
\begin{Highlighting}[]
\CommentTok{\# Modelo birth{-}death (especiação + extinção)}
\NormalTok{modelo.bd}\OtherTok{\textless{}{-}}\FunctionTok{make.bd}\NormalTok{(tree,}\AttributeTok{sampling.f =} \DecValTok{226}\SpecialCharTok{/}\DecValTok{300}\NormalTok{)}
\NormalTok{param.bd}\OtherTok{\textless{}{-}}\FunctionTok{find.mle}\NormalTok{(modelo.bd,}\AttributeTok{x.init=}\FunctionTok{c}\NormalTok{(}\FloatTok{0.1}\NormalTok{,}\FloatTok{0.05}\NormalTok{))}
\NormalTok{param.bd}
\end{Highlighting}
\end{Shaded}

Comparar modelos.

\begin{Shaded}
\begin{Highlighting}[]
\FunctionTok{anova}\NormalTok{(param.bd,}\AttributeTok{pure.birth=}\NormalTok{param.pb)}
\end{Highlighting}
\end{Shaded}

O modelo que incluí extinção (birth-death) aumentou significativamente o seu ajuste em comparação com o modelo \emph{pure-birth}. Esse é essencialmente um sinal de que a extinção é importante (o parâmetro adicional \emph{d} se justifica).

MCMC.O pacote \emph{diversitree} também permite rodar analises bayesianas
para estimar os parâmetros de diversificação usando a distribuição
posterior probabilidades.

\begin{Shaded}
\begin{Highlighting}[]
\NormalTok{bd.amostras }\OtherTok{\textless{}{-}} \FunctionTok{mcmc}\NormalTok{(modelo.bd, param.bd}\SpecialCharTok{$}\NormalTok{par, }\AttributeTok{nsteps =} \FloatTok{1e+05}\NormalTok{, }\AttributeTok{lower =} \FunctionTok{c}\NormalTok{(}\DecValTok{0}\NormalTok{, }\DecValTok{0}\NormalTok{),}
 \AttributeTok{upper =} \FunctionTok{c}\NormalTok{(}\ConstantTok{Inf}\NormalTok{,}\ConstantTok{Inf}\NormalTok{), }\AttributeTok{w =} \FunctionTok{c}\NormalTok{(}\FloatTok{0.1}\NormalTok{, }\FloatTok{0.1}\NormalTok{), }\AttributeTok{fail.value =} \SpecialCharTok{{-}}\ConstantTok{Inf}\NormalTok{,}
 \AttributeTok{print.every =} \DecValTok{10000}\NormalTok{)}
\CommentTok{\# visualizar distribuição}
\NormalTok{postSamples}\OtherTok{\textless{}{-}}\NormalTok{bd.amostras[}\FunctionTok{c}\NormalTok{(}\StringTok{"lambda"}\NormalTok{, }\StringTok{"mu"}\NormalTok{)]}
\FunctionTok{par}\NormalTok{(}\AttributeTok{mfrow=}\FunctionTok{c}\NormalTok{(}\DecValTok{1}\NormalTok{,}\DecValTok{1}\NormalTok{),}\AttributeTok{mar=}\FunctionTok{c}\NormalTok{(}\DecValTok{5}\NormalTok{,}\DecValTok{4}\NormalTok{,}\DecValTok{2}\NormalTok{,}\DecValTok{2}\NormalTok{))}
\FunctionTok{profiles.plot}\NormalTok{(postSamples,}\AttributeTok{col.line=}\FunctionTok{c}\NormalTok{(}\StringTok{"red"}\NormalTok{, }\StringTok{"blue"}\NormalTok{),}\AttributeTok{las=}\DecValTok{1}\NormalTok{,}\AttributeTok{legend =}\StringTok{"topright"}\NormalTok{)}
\end{Highlighting}
\end{Shaded}

\includegraphics{04.11-pcmdiv_files/figure-latex/unnamed-chunk-8-1.pdf}

\section{\texorpdfstring{Diversificação variando no tempo (\emph{time-dependent})}{Diversificação variando no tempo (time-dependent)}}\label{diversificauxe7uxe3o-variando-no-tempo-time-dependent}

Carregar dados.

\begin{Shaded}
\begin{Highlighting}[]
\FunctionTok{require}\NormalTok{(ape)}
\FunctionTok{require}\NormalTok{(phytools)}
\FunctionTok{require}\NormalTok{(geiger)}
\CommentTok{\#\textgreater{} Carregando pacotes exigidos: geiger}
\CommentTok{\# Arvore marsupiais}
\NormalTok{tree}\OtherTok{\textless{}{-}}\FunctionTok{read.tree}\NormalTok{(}\StringTok{"dadospcm/marsupial{-}tree.tre"}\NormalTok{)}
\NormalTok{\{}\FunctionTok{plot}\NormalTok{(tree,}\AttributeTok{cex=}\FloatTok{0.4}\NormalTok{)}
\FunctionTok{axisPhylo}\NormalTok{()\}}
\end{Highlighting}
\end{Shaded}

\includegraphics{04.11-pcmdiv_files/figure-latex/unnamed-chunk-9-1.pdf}

O \emph{RPANDA} (Morlon et al. 2016) é um pacote que permite ajustar modelos com
taxas de especiação e/ou extinção variáveis, e comparar os modelos
usando verossimilhança.

\begin{Shaded}
\begin{Highlighting}[]
\FunctionTok{require}\NormalTok{(RPANDA)}
\CommentTok{\#\textgreater{} Carregando pacotes exigidos: RPANDA}
\CommentTok{\#\textgreater{} Carregando pacotes exigidos: picante}
\CommentTok{\#\textgreater{} Carregando pacotes exigidos: vegan}
\CommentTok{\#\textgreater{} Carregando pacotes exigidos: permute}
\CommentTok{\#\textgreater{} Carregando pacotes exigidos: lattice}
\CommentTok{\#\textgreater{} This is vegan 2.6{-}4}
\CommentTok{\#\textgreater{} }
\CommentTok{\#\textgreater{} Attaching package: \textquotesingle{}vegan\textquotesingle{}}
\CommentTok{\#\textgreater{} The following object is masked from \textquotesingle{}package:phytools\textquotesingle{}:}
\CommentTok{\#\textgreater{} }
\CommentTok{\#\textgreater{}     scores}
\CommentTok{\#\textgreater{} Carregando pacotes exigidos: nlme}
\end{Highlighting}
\end{Shaded}

Nós vamos ajustar modelos com combinações de lambda e mu constantes e variáveis. Os modelos podem ser consultados no help do \emph{RPANDA}.

\begin{Shaded}
\begin{Highlighting}[]
\NormalTok{lambda.cst }\OtherTok{\textless{}{-}} \ControlFlowTok{function}\NormalTok{(t,y)\{y[}\DecValTok{1}\NormalTok{]\} }\CommentTok{\# lambda constante}
\NormalTok{lambda.var }\OtherTok{\textless{}{-}} \ControlFlowTok{function}\NormalTok{(t,y)\{y[}\DecValTok{1}\NormalTok{]}\SpecialCharTok{*}\FunctionTok{exp}\NormalTok{(y[}\DecValTok{2}\NormalTok{]}\SpecialCharTok{*}\NormalTok{t)\} }\CommentTok{\# lambda com variação exponencial (em função do tempo t)}
\NormalTok{mu.cst }\OtherTok{\textless{}{-}} \ControlFlowTok{function}\NormalTok{(t,y)\{y[}\DecValTok{1}\NormalTok{]\}}
\NormalTok{mu.var }\OtherTok{\textless{}{-}} \ControlFlowTok{function}\NormalTok{(t,y)\{y[}\DecValTok{1}\NormalTok{]}\SpecialCharTok{*}\FunctionTok{exp}\NormalTok{(y[}\DecValTok{2}\NormalTok{]}\SpecialCharTok{*}\NormalTok{t)\}}
\end{Highlighting}
\end{Shaded}

Ajustando os modelos para a árvore completa.

\begin{Shaded}
\begin{Highlighting}[]
\CommentTok{\# birth{-}death especiação e extinção constante}
\NormalTok{bcdc}\OtherTok{\textless{}{-}}\FunctionTok{fit\_bd}\NormalTok{(tree,}\AttributeTok{tot\_time=}\FunctionTok{max}\NormalTok{(}\FunctionTok{branching.times}\NormalTok{(tree)),}\AttributeTok{f.lamb=}\NormalTok{lambda.cst,}\AttributeTok{f.mu=}\NormalTok{mu.cst,}\AttributeTok{lamb\_par=}\FloatTok{0.01}\NormalTok{,}\AttributeTok{mu\_par=}\FloatTok{0.001}\NormalTok{,}\AttributeTok{cst.lamb=}\ConstantTok{TRUE}\NormalTok{,}\AttributeTok{cst.mu=}\ConstantTok{TRUE}\NormalTok{)}
\NormalTok{bcdc}\SpecialCharTok{$}\NormalTok{model}\OtherTok{\textless{}{-}}\StringTok{"birth{-}death taxa de especiação e extinção constantes"}
\NormalTok{bcdc}
\CommentTok{\# birth{-}death especiação exponencial e extinção constante}
\NormalTok{bvdc}\OtherTok{\textless{}{-}}\FunctionTok{fit\_bd}\NormalTok{(tree,}\AttributeTok{tot\_time=}\FunctionTok{max}\NormalTok{(}\FunctionTok{branching.times}\NormalTok{(tree)),}\AttributeTok{f.lamb=}\NormalTok{lambda.var,}\AttributeTok{f.mu=}\NormalTok{mu.cst,}\AttributeTok{lamb\_par=}\FunctionTok{c}\NormalTok{(}\FloatTok{0.05}\NormalTok{,}\FloatTok{0.01}\NormalTok{),}\AttributeTok{mu\_par=}\FloatTok{0.001}\NormalTok{,}\AttributeTok{expo.lamb=}\ConstantTok{TRUE}\NormalTok{,}\AttributeTok{cst.mu=}\ConstantTok{TRUE}\NormalTok{)}
\NormalTok{bvdc}\SpecialCharTok{$}\NormalTok{model}\OtherTok{\textless{}{-}}\StringTok{"birth{-}death taxa de especiação exponencial e extinção constante"}
\CommentTok{\# birth{-}death especiação constante e extinção exponencial}
\NormalTok{bcdv}\OtherTok{\textless{}{-}}\FunctionTok{fit\_bd}\NormalTok{(tree,}\AttributeTok{tot\_time=}\FunctionTok{max}\NormalTok{(}\FunctionTok{branching.times}\NormalTok{(tree)),}\AttributeTok{f.lamb=}\NormalTok{lambda.cst,}\AttributeTok{f.mu=}\NormalTok{mu.var,}\AttributeTok{lamb\_par=}\FloatTok{0.01}\NormalTok{,}\AttributeTok{mu\_par=}\FunctionTok{c}\NormalTok{(}\FloatTok{0.005}\NormalTok{,}\FloatTok{0.001}\NormalTok{),}\AttributeTok{cst.lamb=}\ConstantTok{TRUE}\NormalTok{,}\AttributeTok{expo.mu=}\ConstantTok{TRUE}\NormalTok{)}
\NormalTok{bcdv}\SpecialCharTok{$}\NormalTok{model}\OtherTok{\textless{}{-}}\StringTok{"birth{-}death taxa de especiação constante e extinção exponencial"}
\CommentTok{\# birth{-}death especiação e extinção exponencial}
\NormalTok{bvdv}\OtherTok{\textless{}{-}}\FunctionTok{fit\_bd}\NormalTok{(tree,}\AttributeTok{tot\_time=}\FunctionTok{max}\NormalTok{(}\FunctionTok{branching.times}\NormalTok{(tree)),}\AttributeTok{f.lamb=}\NormalTok{lambda.var,}\AttributeTok{f.mu=}\NormalTok{mu.var,}\AttributeTok{lamb\_par=}\FunctionTok{c}\NormalTok{(}\FloatTok{0.05}\NormalTok{,}\FloatTok{0.01}\NormalTok{),}\AttributeTok{mu\_par=}\FunctionTok{c}\NormalTok{(}\FloatTok{0.005}\NormalTok{,}\FloatTok{0.001}\NormalTok{),}\AttributeTok{expo.lamb=}\ConstantTok{TRUE}\NormalTok{,}\AttributeTok{expo.mu=}\ConstantTok{TRUE}\NormalTok{)}
\NormalTok{bvdv}\SpecialCharTok{$}\NormalTok{model}\OtherTok{\textless{}{-}}\StringTok{"birth{-}death taxa de especiação e extinção exponencial"}
\end{Highlighting}
\end{Shaded}

Comparando os modelos.

\begin{Shaded}
\begin{Highlighting}[]
\NormalTok{aic.scores}\OtherTok{\textless{}{-}}\FunctionTok{setNames}\NormalTok{(}\FunctionTok{c}\NormalTok{(bcdc}\SpecialCharTok{$}\NormalTok{aicc,bvdc}\SpecialCharTok{$}\NormalTok{aicc,bcdv}\SpecialCharTok{$}\NormalTok{aicc,bvdv}\SpecialCharTok{$}\NormalTok{aicc),}
 \FunctionTok{c}\NormalTok{(}\StringTok{"bcdc"}\NormalTok{,}\StringTok{"bvdc"}\NormalTok{,}\StringTok{"bcdv"}\NormalTok{,}\StringTok{"bvdv"}\NormalTok{))}
\NormalTok{aic.scores}
\FunctionTok{aic.w}\NormalTok{(aic.scores)}
\end{Highlighting}
\end{Shaded}

Resultados um pouco ambíguos, mas sugerem que taxas constantes são
improváveis.

É possível ajustar os modelos para toda a árvore, como acima, ou definir
partes da árvore a priori (clados) e testar diferenças nas taxas entre
clados. Se não temos uma hipótese a priori de quais clados tem
diversificação mais rápida, podemos identificar mudanças rápidas olhando
pra árvore, ou usando métodos como o MEDUSA ou o BAMM para detectar
mudanças na diversificação em hipóteses a priori.

\section{\texorpdfstring{Diversificação variando entre clados (\emph{clade-dependent})}{Diversificação variando entre clados (clade-dependent)}}\label{diversificauxe7uxe3o-variando-entre-clados-clade-dependent}

O exemplo implementa o MEDUSA (Alfaro et al. 2009).

\begin{Shaded}
\begin{Highlighting}[]
\NormalTok{mar}\OtherTok{\textless{}{-}}\FunctionTok{medusa}\NormalTok{(tree)}
\NormalTok{mar}
\FunctionTok{plot}\NormalTok{(mar)}
\end{Highlighting}
\end{Shaded}

\includegraphics{04.11-pcmdiv_files/figure-latex/unnamed-chunk-14-1.pdf}

Pelo menos dois grandes `shifts' de diversificação podem ser percebidos.
Um no clado dos Macropodidae (cangurus e afins - nó 301) e outro no
clado dos Dasyuromorpha (marsupiais carnívoros - nó 355). r= net
diversification rate. Agora podemos hipotetizar, como exemplo, que o
clado dos cangurus e aliados (Macropodidae) parece ter se diversificado
mais rapidamente.

\begin{Shaded}
\begin{Highlighting}[]
\FunctionTok{par}\NormalTok{(}\AttributeTok{mfrow=}\FunctionTok{c}\NormalTok{(}\DecValTok{1}\NormalTok{,}\DecValTok{1}\NormalTok{),}\AttributeTok{mar=}\FunctionTok{c}\NormalTok{(}\DecValTok{5}\NormalTok{,}\DecValTok{4}\NormalTok{,}\DecValTok{2}\NormalTok{,}\DecValTok{2}\NormalTok{))}
\NormalTok{\{}\FunctionTok{plot}\NormalTok{(mar)}
\FunctionTok{nodelabels}\NormalTok{()\}}
\end{Highlighting}
\end{Shaded}

\includegraphics{04.11-pcmdiv_files/figure-latex/unnamed-chunk-15-1.pdf}

\begin{Shaded}
\begin{Highlighting}[]
\NormalTok{macropodidae.tree }\OtherTok{\textless{}{-}} \FunctionTok{extract.clade}\NormalTok{(tree, }\DecValTok{301}\NormalTok{)}
\FunctionTok{plot}\NormalTok{(macropodidae.tree)}
\end{Highlighting}
\end{Shaded}

\includegraphics{04.11-pcmdiv_files/figure-latex/unnamed-chunk-15-2.pdf}

\begin{Shaded}
\begin{Highlighting}[]
\NormalTok{resto.tree}\OtherTok{\textless{}{-}}\FunctionTok{drop.tip}\NormalTok{(tree,macropodidae.tree}\SpecialCharTok{$}\NormalTok{tip.label)}
\FunctionTok{plot}\NormalTok{(resto.tree)}
\end{Highlighting}
\end{Shaded}

\includegraphics{04.11-pcmdiv_files/figure-latex/unnamed-chunk-15-3.pdf}

\section{Diversificação variando no tempo e entre clados}\label{diversificauxe7uxe3o-variando-no-tempo-e-entre-clados}

Vamos usar RPANDA para ajustar modelos com taxas variáveis entre clados.
Para facilitar, vamos unir os modelos em uma função.

\begin{Shaded}
\begin{Highlighting}[]
\NormalTok{fit.multi.rpanda }\OtherTok{\textless{}{-}} \ControlFlowTok{function}\NormalTok{(tree)\{}
\NormalTok{ bcdc}\OtherTok{\textless{}{-}}\FunctionTok{fit\_bd}\NormalTok{(tree,}\AttributeTok{tot\_time=}\FunctionTok{max}\NormalTok{(}\FunctionTok{branching.times}\NormalTok{(tree)),}\AttributeTok{f.lamb=}\NormalTok{lambda.cst,}\AttributeTok{f.mu=}\NormalTok{mu.cst,}\AttributeTok{lamb\_par=}\FloatTok{0.01}\NormalTok{,}\AttributeTok{mu\_par=}\FloatTok{0.001}\NormalTok{,}\AttributeTok{cst.lamb=}\ConstantTok{TRUE}\NormalTok{,}\AttributeTok{cst.mu=}\ConstantTok{TRUE}\NormalTok{)}
\NormalTok{ bvdc}\OtherTok{\textless{}{-}}\FunctionTok{fit\_bd}\NormalTok{(tree,}\AttributeTok{tot\_time=}\FunctionTok{max}\NormalTok{(}\FunctionTok{branching.times}\NormalTok{(tree)),}\AttributeTok{f.lamb=}\NormalTok{lambda.var,}\AttributeTok{f.mu=}\NormalTok{mu.cst,}\AttributeTok{lamb\_par=}\FunctionTok{c}\NormalTok{(}\FloatTok{0.05}\NormalTok{,}\FloatTok{0.01}\NormalTok{),}\AttributeTok{mu\_par=}\FloatTok{0.001}\NormalTok{,}\AttributeTok{expo.lamb=}\ConstantTok{TRUE}\NormalTok{,}\AttributeTok{cst.mu=}\ConstantTok{TRUE}\NormalTok{)}
\NormalTok{ bcdv}\OtherTok{\textless{}{-}}\FunctionTok{fit\_bd}\NormalTok{(tree,}\AttributeTok{tot\_time=}\FunctionTok{max}\NormalTok{(}\FunctionTok{branching.times}\NormalTok{(tree)),}\AttributeTok{f.lamb=}\NormalTok{lambda.cst,}\AttributeTok{f.mu=}\NormalTok{mu.var,}\AttributeTok{lamb\_par=}\FloatTok{0.01}\NormalTok{,}\AttributeTok{mu\_par=}\FunctionTok{c}\NormalTok{(}\FloatTok{0.005}\NormalTok{,}\FloatTok{0.001}\NormalTok{),}\AttributeTok{cst.lamb=}\ConstantTok{TRUE}\NormalTok{,}\AttributeTok{expo.mu=}\ConstantTok{TRUE}\NormalTok{)}
\NormalTok{ bvdv}\OtherTok{\textless{}{-}}\FunctionTok{fit\_bd}\NormalTok{(tree,}\AttributeTok{tot\_time=}\FunctionTok{max}\NormalTok{(}\FunctionTok{branching.times}\NormalTok{(tree)),}\AttributeTok{f.lamb=}\NormalTok{lambda.var,}\AttributeTok{f.mu=}\NormalTok{mu.var,}\AttributeTok{lamb\_par=}\FunctionTok{c}\NormalTok{(}\FloatTok{0.05}\NormalTok{,}\FloatTok{0.01}\NormalTok{),}\AttributeTok{mu\_par=}\FunctionTok{c}\NormalTok{(}\FloatTok{0.005}\NormalTok{,}\FloatTok{0.001}\NormalTok{),}\AttributeTok{expo.lamb=}\ConstantTok{TRUE}\NormalTok{,}\AttributeTok{expo.mu=}\ConstantTok{TRUE}\NormalTok{)}
 \FunctionTok{return}\NormalTok{(}\FunctionTok{list}\NormalTok{(}\StringTok{"bcdc"}\OtherTok{=}\NormalTok{bcdc,}\StringTok{"bvdc"}\OtherTok{=}\NormalTok{bvdc,}\StringTok{"bcdv"}\OtherTok{=}\NormalTok{bcdv,}\StringTok{"bvdv"}\OtherTok{=}\NormalTok{bvdv))}
\NormalTok{ \}}
\NormalTok{macropodidae\_panda}\OtherTok{\textless{}{-}}\FunctionTok{fit.multi.rpanda}\NormalTok{(macropodidae.tree)}
\NormalTok{resto\_panda}\OtherTok{\textless{}{-}}\FunctionTok{fit.multi.rpanda}\NormalTok{(resto.tree)}
\end{Highlighting}
\end{Shaded}

Comparando modelos.

\begin{Shaded}
\begin{Highlighting}[]
\CommentTok{\# Extraindo aicc}
\NormalTok{macro\_aic}\OtherTok{\textless{}{-}}\FunctionTok{numeric}\NormalTok{(}\DecValTok{4}\NormalTok{) }\CommentTok{\#vetor vazio}
\NormalTok{resto\_aic}\OtherTok{\textless{}{-}}\FunctionTok{numeric}\NormalTok{(}\DecValTok{4}\NormalTok{)}
\FunctionTok{names}\NormalTok{(macro\_aic)}\OtherTok{\textless{}{-}}\FunctionTok{c}\NormalTok{(}\StringTok{"bcdc"}\NormalTok{, }\StringTok{"bvdc"}\NormalTok{, }\StringTok{"bcdv"}\NormalTok{, }\StringTok{"bvdv"}\NormalTok{)}
\FunctionTok{names}\NormalTok{(resto\_aic)}\OtherTok{\textless{}{-}}\FunctionTok{names}\NormalTok{(macro\_aic)}
\ControlFlowTok{for}\NormalTok{(i }\ControlFlowTok{in} \DecValTok{1}\SpecialCharTok{:}\DecValTok{4}\NormalTok{) \{}
\NormalTok{ macro\_aic[i]}\OtherTok{\textless{}{-}}\NormalTok{macropodidae\_panda[[i]]}\SpecialCharTok{$}\NormalTok{aicc}
\NormalTok{ resto\_aic[i]}\OtherTok{\textless{}{-}}\NormalTok{resto\_panda[[i]]}\SpecialCharTok{$}\NormalTok{aicc}
\NormalTok{\}}
\FunctionTok{aic.w}\NormalTok{(macro\_aic)}
\FunctionTok{aic.w}\NormalTok{(resto\_aic)}
\end{Highlighting}
\end{Shaded}

Novamente os resultados são um pouco inconclusivos. Parece haver maior suporte para modelos com taxas variáveis do que constantes. Para o clado Macropodidae uma variação exponencial nas taxas de especiação parece bem provável, enquanto que para o resto da árvore, uma variação em ambas as taxas (especiação e extinção) é mais provável.

Podemos plotar as taxas de especiação, extinção, e diversificação (b-d) ao longo do tempo para os modelos mais prováveis de cada clado.

\begin{Shaded}
\begin{Highlighting}[]
\FunctionTok{plot\_fit\_bd}\NormalTok{(macropodidae\_panda}\SpecialCharTok{$}\NormalTok{bvdc,}\FunctionTok{max}\NormalTok{(}\FunctionTok{branching.times}\NormalTok{(macropodidae.tree)))}
\FunctionTok{plot\_fit\_bd}\NormalTok{(resto\_panda}\SpecialCharTok{$}\NormalTok{bvdv,}\FunctionTok{max}\NormalTok{(}\FunctionTok{branching.times}\NormalTok{(resto.tree)))}
\FunctionTok{plot}\NormalTok{(macropodidae.tree)}
\FunctionTok{plot}\NormalTok{(resto.tree)}
\end{Highlighting}
\end{Shaded}

\section{Exercício - Diversificação dependente de variável ambiental}\label{exercuxedcio---diversificauxe7uxe3o-dependente-de-variuxe1vel-ambiental}

RPANDA pode ser usado para ajustar modelos de diversificação dependente de variável ambiental (fit\_env). Como variável adicional, esse modelo requer uma estimativa da variação em uma variável ambiental de interesse ao longo do tempo. Zachos et al (2008) fornecem uma estimativa de variação de temperatura ao longo do Cenozóico baseada em isótopos de oxigênio.

Baseado nos dados:
Árvore para roedores sigmodontíneos:
sigmodontinae-tree.tre
Dados de temperatura, duas colunas com tempo e
temperatura:
Cenozoic\_Temp.txt

Ajuste um modelo com variação nas taxas de especiação dependentes da
temperatura, sem dependência temporal e sem extinção:

lambda.env\textless-function(t,x,y)\{y{[}1{]}\emph{exp(y{[}2{]}}x)\}
e
mu.zero\textless-function(t,x,y)\{0\}

Com os parâmetros iniciais:

lamb\_par\_init\textless-c(0.10,0.01)

mu\_par\_init\textless-c()

O modelo deve ser ajustado com a função fit\_env do pacote RPANDA. Faça um gráfico dos resultados com plot\_fit\_env e interprete os resultados.

\chapter{Diversificação dependente de estado}\label{diversificauxe7uxe3o-dependente-de-estado}

Esse capítulo exemplifica a lógica dos modelos de diversificação dependentes de estado (SSE) usando o modelo mais simples e primeiro a ser desenvolvido, o BiSSE. Há uma grande literatura sobre modelos dependentes de estado, em constante atualização, que precisa ser consultada antes de utilizar um modelo SSE. Veja os artigos nas referências recomendadas na disciplina.

Carregar dados.

\begin{Shaded}
\begin{Highlighting}[]
\FunctionTok{require}\NormalTok{(ape)}
\CommentTok{\#\textgreater{} Carregando pacotes exigidos: ape}
\FunctionTok{require}\NormalTok{(phytools)}
\CommentTok{\#\textgreater{} Carregando pacotes exigidos: phytools}
\CommentTok{\#\textgreater{} Carregando pacotes exigidos: maps}
\FunctionTok{require}\NormalTok{(geiger)}
\CommentTok{\#\textgreater{} Carregando pacotes exigidos: geiger}
\FunctionTok{require}\NormalTok{(diversitree)}
\CommentTok{\#\textgreater{} Carregando pacotes exigidos: diversitree}
\CommentTok{\# Filogenia}
\NormalTok{tree}\OtherTok{\textless{}{-}}\FunctionTok{read.tree}\NormalTok{(}\StringTok{"dadospcm/primate{-}tree.txt"}\NormalTok{)}
\FunctionTok{plotTree}\NormalTok{(tree,}\AttributeTok{fsize=}\FloatTok{0.4}\NormalTok{,}\AttributeTok{ftype=}\StringTok{"i"}\NormalTok{,}\AttributeTok{type=}\StringTok{"fan"}\NormalTok{,}\AttributeTok{lwd=}\DecValTok{1}\NormalTok{)}
\end{Highlighting}
\end{Shaded}

\includegraphics{04.12-pcmsse_files/figure-latex/unnamed-chunk-1-1.pdf}

\begin{Shaded}
\begin{Highlighting}[]
\CommentTok{\# Atributos}
\NormalTok{dados}\OtherTok{\textless{}{-}}\FunctionTok{read.table}\NormalTok{(}\StringTok{"dadospcm/primate{-}data.txt"}\NormalTok{,}\AttributeTok{h=}\NormalTok{T,}\AttributeTok{row.names=}\DecValTok{1}\NormalTok{)}
\FunctionTok{head}\NormalTok{(dados)}
\CommentTok{\# Recortar filogenia}
\NormalTok{match.species}\OtherTok{\textless{}{-}}\FunctionTok{treedata}\NormalTok{(tree,dados)}
\NormalTok{tree}\OtherTok{\textless{}{-}}\NormalTok{match.species}\SpecialCharTok{$}\NormalTok{phy}
\FunctionTok{plotTree}\NormalTok{(tree,}\AttributeTok{fsize=}\FloatTok{0.5}\NormalTok{,}\AttributeTok{ftype=}\StringTok{"i"}\NormalTok{,}\AttributeTok{type=}\StringTok{"fan"}\NormalTok{,}\AttributeTok{lwd=}\DecValTok{1}\NormalTok{)}
\end{Highlighting}
\end{Shaded}

\includegraphics{04.12-pcmsse_files/figure-latex/unnamed-chunk-1-2.pdf}

Vamos utilizar o BiSSE (\emph{binary-state speciation and extinction}) para testar hipóteses sobre a diversificação de primatas e o tamanho do grupo social.
Considere a variável binária, que possui dois estados: Solitary vs.~Social.
Menos de 4 ind = solitary
Mais de 4 ind = social

\begin{Shaded}
\begin{Highlighting}[]
\NormalTok{binsocial}\OtherTok{\textless{}{-}}\NormalTok{dados}\SpecialCharTok{$}\NormalTok{BinSocial}
\FunctionTok{names}\NormalTok{(binsocial)}\OtherTok{=}\FunctionTok{rownames}\NormalTok{(dados)}
\NormalTok{binsocial}\OtherTok{\textless{}{-}}\FunctionTok{as.factor}\NormalTok{(binsocial)}
\NormalTok{binsocial}
\end{Highlighting}
\end{Shaded}

Uma dessas categorias promove maior diversificação em primatas?

\begin{Shaded}
\begin{Highlighting}[]
\CommentTok{\# O diversitree requer que os códigos sejam numéricos.}
\FunctionTok{levels}\NormalTok{(binsocial) }\CommentTok{\# 1{-}Social, 2{-}Solitário}
\NormalTok{binsocial}\OtherTok{\textless{}{-}}\FunctionTok{as.numeric}\NormalTok{(binsocial)}\SpecialCharTok{{-}}\DecValTok{1} \CommentTok{\# valores tem que ser 0 e 1}
\NormalTok{binsocial }
\FunctionTok{names}\NormalTok{(binsocial)}\OtherTok{=}\FunctionTok{rownames}\NormalTok{(dados)}
\NormalTok{binsocial }\CommentTok{\# 0{-}Social, 1{-}Solitário}
\end{Highlighting}
\end{Shaded}

O BiSSE tem uma função para estimar valores iniciais `razoáveis' dos parâmetros de interesse.

\begin{Shaded}
\begin{Highlighting}[]
\NormalTok{p}\OtherTok{\textless{}{-}}\FunctionTok{starting.point.bisse}\NormalTok{(tree)}
\NormalTok{p}
\end{Highlighting}
\end{Shaded}

Para rodar o BiSSE temos que criar a função de verossimilhança (\emph{likelihood}) para posteriormente otimizar os parâmetros.

\begin{Shaded}
\begin{Highlighting}[]
\NormalTok{lik.bisse}\OtherTok{\textless{}{-}}\FunctionTok{make.bisse}\NormalTok{(tree,binsocial)}
\NormalTok{lik.bisse}
\end{Highlighting}
\end{Shaded}

Agora podemos restringir a função geral para criar diferentes modelos.
Por exemplo, podemos criar um modelo ``nulo'' onde as taxas de especiação e extinção sejam iguais entre estados (b0\textsubscript{b1;d0}d1) e que existe uma única taxa de transição (q01=q10).

\begin{Shaded}
\begin{Highlighting}[]
\NormalTok{lik.null}\OtherTok{\textless{}{-}}\FunctionTok{constrain}\NormalTok{(lik.bisse,lambda0}\SpecialCharTok{\textasciitilde{}}\NormalTok{lambda1,mu0}\SpecialCharTok{\textasciitilde{}}\NormalTok{mu1,q10}\SpecialCharTok{\textasciitilde{}}\NormalTok{q01)}
\NormalTok{lik.null}
\end{Highlighting}
\end{Shaded}

Ajustar os modelos.

\begin{Shaded}
\begin{Highlighting}[]
\CommentTok{\# Ajustar modelo nulo}
\NormalTok{p[}\FunctionTok{argnames}\NormalTok{(lik.null)] }\CommentTok{\# ajustar número de parâmetros inicial}
\NormalTok{fit.null}\OtherTok{\textless{}{-}}\FunctionTok{find.mle}\NormalTok{(lik.null,}\AttributeTok{x.init=}\NormalTok{p[}\FunctionTok{argnames}\NormalTok{(lik.null)])}
\NormalTok{fit.null}
\CommentTok{\# Ajustar o modelo mais complexo, onde as taxas de especiação e extinção dependem do estado do fenótipo}
\NormalTok{fit.full}\OtherTok{\textless{}{-}}\FunctionTok{find.mle}\NormalTok{(lik.bisse,}\AttributeTok{x.init=}\NormalTok{p)}
\NormalTok{fit.full}
\end{Highlighting}
\end{Shaded}

Outros modelos podem ser ajustados entre o mais simples e o mais complexo, dependendo das hipóteses biológicas a priori. Em um estudo empírico, especialmente com múltiplos estados (MuSSE e similares), os modelos devem ser identificados a priori ao invés de ajustar todos os modelos possíveis para identificar aqueles com melhor ajuste.
Como exemplo, vamos ajustar um modelo onde somente a taxa de especiação varie ente estados, um modelo onde somente a taxa de extinção varie, e um modelo com b e d constantes, mas com taxas de transição diferentes entre estados.

\begin{Shaded}
\begin{Highlighting}[]
\CommentTok{\# especiação variável}
\NormalTok{lik.lambda}\OtherTok{\textless{}{-}}\FunctionTok{constrain}\NormalTok{(lik.bisse,mu0}\SpecialCharTok{\textasciitilde{}}\NormalTok{mu1,q10}\SpecialCharTok{\textasciitilde{}}\NormalTok{q01)}
\NormalTok{lik.lambda}
\CommentTok{\# ajuste}
\NormalTok{fit.lambda}\OtherTok{\textless{}{-}}\FunctionTok{find.mle}\NormalTok{(lik.lambda,}\AttributeTok{x.init=}\NormalTok{p[}\FunctionTok{argnames}\NormalTok{(lik.lambda)])}
\NormalTok{fit.lambda}
\CommentTok{\# extinção variável}
\NormalTok{lik.mu}\OtherTok{\textless{}{-}}\FunctionTok{constrain}\NormalTok{(lik.bisse,lambda0}\SpecialCharTok{\textasciitilde{}}\NormalTok{lambda1,q10}\SpecialCharTok{\textasciitilde{}}\NormalTok{q01)}
\NormalTok{lik.mu}
\CommentTok{\# ajuste}
\NormalTok{fit.mu}\OtherTok{\textless{}{-}}\FunctionTok{find.mle}\NormalTok{(lik.mu,}\AttributeTok{x.init=}\NormalTok{p[}\FunctionTok{argnames}\NormalTok{(lik.mu)])}
\NormalTok{fit.mu}
\CommentTok{\# especiação e extinção variável, única taxa de transição entre estados}
\NormalTok{lik.lambda.mu}\OtherTok{\textless{}{-}}\FunctionTok{constrain}\NormalTok{(lik.bisse,q10}\SpecialCharTok{\textasciitilde{}}\NormalTok{q01)}
\NormalTok{lik.lambda.mu}
\CommentTok{\# ajuste}
\NormalTok{fit.lambda.mu}\OtherTok{\textless{}{-}}\FunctionTok{find.mle}\NormalTok{(lik.lambda.mu,}\AttributeTok{x.init=}\NormalTok{p[}\FunctionTok{argnames}\NormalTok{(lik.lambda.mu)])}
\NormalTok{fit.lambda.mu}
\CommentTok{\# b e d constantes, taxas de transição flexíveis}
\NormalTok{lik.q}\OtherTok{\textless{}{-}}\FunctionTok{constrain}\NormalTok{(lik.bisse,lambda0}\SpecialCharTok{\textasciitilde{}}\NormalTok{lambda1,mu0}\SpecialCharTok{\textasciitilde{}}\NormalTok{mu1)}
\NormalTok{lik.q}
\CommentTok{\# ajuste}
\NormalTok{fit.q}\OtherTok{\textless{}{-}}\FunctionTok{find.mle}\NormalTok{(lik.q,}\AttributeTok{x.init=}\NormalTok{p[}\FunctionTok{argnames}\NormalTok{(lik.q)])}
\NormalTok{fit.q}
\end{Highlighting}
\end{Shaded}

Agora podemos comparar os modelos ajustados. A função genérica anova permite
obter os valores de AIC para uma lista.

\begin{Shaded}
\begin{Highlighting}[]
\CommentTok{\# Comparando os modelos}
\NormalTok{resultados}\OtherTok{\textless{}{-}}\FunctionTok{anova}\NormalTok{(fit.null,}
 \AttributeTok{completo=}\NormalTok{fit.full,}
 \AttributeTok{lambda.variavel=}\NormalTok{fit.lambda,}
 \AttributeTok{mu.variavel=}\NormalTok{fit.mu,}
 \AttributeTok{lambda.mu.variavel=}\NormalTok{fit.lambda.mu,}
 \AttributeTok{q.variavel=}\NormalTok{fit.q)}
\NormalTok{resultados}
\FunctionTok{aicw}\NormalTok{(}\FunctionTok{setNames}\NormalTok{(resultados}\SpecialCharTok{$}\NormalTok{AIC,}\FunctionTok{rownames}\NormalTok{(resultados)))}
\FunctionTok{aic.w}\NormalTok{(}\FunctionTok{setNames}\NormalTok{(resultados}\SpecialCharTok{$}\NormalTok{AIC,}\FunctionTok{rownames}\NormalTok{(resultados)))}
\end{Highlighting}
\end{Shaded}

Os modelos com especiação variável e com especiação e extinção variável (e taxa de transição única) receberam maior suporte.

\begin{Shaded}
\begin{Highlighting}[]
\FunctionTok{coef}\NormalTok{(fit.lambda.mu)}
\end{Highlighting}
\end{Shaded}

Parece que tanto as taxas de especiação como de extinção são maiores para o estado 0 (social) do que para o estado 1 (solitário).
E a taxa de diversificação?

\begin{Shaded}
\begin{Highlighting}[]
\NormalTok{lambda0}\OtherTok{\textless{}{-}}\FunctionTok{coef}\NormalTok{(fit.lambda.mu)[[}\DecValTok{1}\NormalTok{]]}
\NormalTok{mu0}\OtherTok{\textless{}{-}}\FunctionTok{coef}\NormalTok{(fit.lambda.mu)[[}\DecValTok{3}\NormalTok{]]}
\NormalTok{r}\FloatTok{.0}\OtherTok{\textless{}{-}}\NormalTok{lambda0}\SpecialCharTok{{-}}\NormalTok{mu0}
\NormalTok{r}\FloatTok{.0}
\NormalTok{lambda1}\OtherTok{\textless{}{-}}\FunctionTok{coef}\NormalTok{(fit.lambda.mu)[[}\DecValTok{2}\NormalTok{]]}
\NormalTok{mu1}\OtherTok{\textless{}{-}}\FunctionTok{coef}\NormalTok{(fit.lambda.mu)[[}\DecValTok{4}\NormalTok{]]}
\NormalTok{r}\FloatTok{.1}\OtherTok{\textless{}{-}}\NormalTok{lambda1}\SpecialCharTok{{-}}\NormalTok{mu1}
\NormalTok{r}\FloatTok{.1}
\end{Highlighting}
\end{Shaded}

Como as taxas de extinção estimadas são muito baixas, as taxas de diversificação
são próximas das de especiação.

Assim como vimos no modelo bd simples, com o pacote diversitree podemos rodar análises bayesianas com os nossos modelos. Vamos fazer isso para o nosso melhor modelo para estimar os valores dos parâmetros.

\begin{Shaded}
\begin{Highlighting}[]
\CommentTok{\# MCMC}
\NormalTok{mcmc.lambda.mu}\OtherTok{\textless{}{-}}\FunctionTok{mcmc}\NormalTok{(lik.lambda.mu,fit.lambda.mu}\SpecialCharTok{$}\NormalTok{par, }\AttributeTok{nsteps =} \DecValTok{1000}\NormalTok{, }\AttributeTok{w =} \DecValTok{1}\NormalTok{,}\AttributeTok{print.every=}\DecValTok{100}\NormalTok{)}
\CommentTok{\# visualizar probabilidades}
\NormalTok{colors}\OtherTok{\textless{}{-}}\FunctionTok{setNames}\NormalTok{(}\FunctionTok{c}\NormalTok{(}\StringTok{"red"}\NormalTok{,}\StringTok{"blue"}\NormalTok{),}\DecValTok{1}\SpecialCharTok{:}\DecValTok{2}\NormalTok{)}
\FunctionTok{par}\NormalTok{(}\AttributeTok{mfrow=}\FunctionTok{c}\NormalTok{(}\DecValTok{1}\NormalTok{,}\DecValTok{1}\NormalTok{),}\AttributeTok{mar=}\FunctionTok{c}\NormalTok{(}\DecValTok{5}\NormalTok{,}\DecValTok{4}\NormalTok{,}\DecValTok{2}\NormalTok{,}\DecValTok{2}\NormalTok{))}
\CommentTok{\# especiação}
\FunctionTok{profiles.plot}\NormalTok{(mcmc.lambda.mu[,}\FunctionTok{grep}\NormalTok{(}\StringTok{"lambda"}\NormalTok{,}\FunctionTok{colnames}\NormalTok{(mcmc.lambda.mu))], }
 \AttributeTok{col.line=}\NormalTok{colors, }\AttributeTok{las=}\DecValTok{1}\NormalTok{, }\AttributeTok{legend.pos=}\StringTok{"topright"}\NormalTok{)}
\end{Highlighting}
\end{Shaded}

\includegraphics{04.12-pcmsse_files/figure-latex/unnamed-chunk-12-1.pdf}

\begin{Shaded}
\begin{Highlighting}[]
\CommentTok{\# extinção}
\FunctionTok{profiles.plot}\NormalTok{(mcmc.lambda.mu[,}\FunctionTok{grep}\NormalTok{(}\StringTok{"mu"}\NormalTok{,}\FunctionTok{colnames}\NormalTok{(mcmc.lambda.mu))],}
 \AttributeTok{col.line=}\NormalTok{colors, }\AttributeTok{las=}\DecValTok{1}\NormalTok{, }\AttributeTok{legend.pos=}\StringTok{"topright"}\NormalTok{)}
\end{Highlighting}
\end{Shaded}

\includegraphics{04.12-pcmsse_files/figure-latex/unnamed-chunk-12-2.pdf}

\begin{Shaded}
\begin{Highlighting}[]
\CommentTok{\# diversificação líquida}
\NormalTok{div}\OtherTok{\textless{}{-}}\NormalTok{mcmc.lambda.mu[,}\FunctionTok{grep}\NormalTok{(}\StringTok{"lambda"}\NormalTok{,}\FunctionTok{colnames}\NormalTok{(mcmc.lambda.mu))]}\SpecialCharTok{{-}}
\NormalTok{ mcmc.lambda.mu[,}\FunctionTok{grep}\NormalTok{(}\StringTok{"mu"}\NormalTok{,}\FunctionTok{colnames}\NormalTok{(mcmc.lambda.mu))]}
\FunctionTok{colnames}\NormalTok{(div)}\OtherTok{\textless{}{-}}\FunctionTok{paste}\NormalTok{(}\StringTok{"lambda{-}mu("}\NormalTok{,}\DecValTok{0}\SpecialCharTok{:}\DecValTok{1}\NormalTok{,}\StringTok{")"}\NormalTok{,}\AttributeTok{sep=}\StringTok{""}\NormalTok{)}
\FunctionTok{profiles.plot}\NormalTok{(div,}
 \AttributeTok{xlab=}\StringTok{"Net diversification rate"}\NormalTok{, }\AttributeTok{ylab=}\StringTok{"Probability density"}\NormalTok{,}
 \AttributeTok{legend.pos=}\StringTok{"topleft"}\NormalTok{,}\AttributeTok{col.line=}\FunctionTok{setNames}\NormalTok{(colors,}\FunctionTok{colnames}\NormalTok{(div)),}
 \AttributeTok{lty=}\DecValTok{1}\NormalTok{)}
\end{Highlighting}
\end{Shaded}

\includegraphics{04.12-pcmsse_files/figure-latex/unnamed-chunk-12-3.pdf}

\chapter{\texorpdfstring{Diversificação baseada em terminais (\emph{tip-speciation})}{Diversificação baseada em terminais (tip-speciation)}}\label{diversificauxe7uxe3o-baseada-em-terminais-tip-speciation}

Uma das medidas mais comuns de taxa de especiação/diversificação nos terminais (tipDR) é a de Jetz et al. (2012), que é o inverso da medida equal splits de Redding and Mooers (2006). Veja também uma outra medida de diversificação em terminais no ClaDS Maliet, Hartig, and Morlon (2019).

Carregar dados.

\begin{Shaded}
\begin{Highlighting}[]
\FunctionTok{require}\NormalTok{(ape)}
\CommentTok{\#\textgreater{} Carregando pacotes exigidos: ape}
\FunctionTok{require}\NormalTok{(phytools)}
\CommentTok{\#\textgreater{} Carregando pacotes exigidos: phytools}
\CommentTok{\#\textgreater{} Carregando pacotes exigidos: maps}
\FunctionTok{require}\NormalTok{(picante)}
\CommentTok{\#\textgreater{} Carregando pacotes exigidos: picante}
\CommentTok{\#\textgreater{} Carregando pacotes exigidos: vegan}
\CommentTok{\#\textgreater{} Carregando pacotes exigidos: permute}
\CommentTok{\#\textgreater{} Carregando pacotes exigidos: lattice}
\CommentTok{\#\textgreater{} This is vegan 2.6{-}4}
\CommentTok{\#\textgreater{} }
\CommentTok{\#\textgreater{} Attaching package: \textquotesingle{}vegan\textquotesingle{}}
\CommentTok{\#\textgreater{} The following object is masked from \textquotesingle{}package:phytools\textquotesingle{}:}
\CommentTok{\#\textgreater{} }
\CommentTok{\#\textgreater{}     scores}
\CommentTok{\#\textgreater{} Carregando pacotes exigidos: nlme}
\FunctionTok{require}\NormalTok{(geiger)}
\CommentTok{\#\textgreater{} Carregando pacotes exigidos: geiger}
\CommentTok{\# Filogenia}
\NormalTok{tree}\OtherTok{\textless{}{-}}\FunctionTok{read.tree}\NormalTok{(}\StringTok{"dadospcm/primate{-}tree.txt"}\NormalTok{)}
\FunctionTok{plotTree}\NormalTok{(tree,}\AttributeTok{fsize=}\FloatTok{0.4}\NormalTok{,}\AttributeTok{ftype=}\StringTok{"i"}\NormalTok{,}\AttributeTok{type=}\StringTok{"fan"}\NormalTok{,}\AttributeTok{lwd=}\DecValTok{1}\NormalTok{)}
\end{Highlighting}
\end{Shaded}

\includegraphics{04.13-pcmtip_files/figure-latex/unnamed-chunk-1-1.pdf}

\begin{Shaded}
\begin{Highlighting}[]
\CommentTok{\# Atributos}
\NormalTok{dados}\OtherTok{\textless{}{-}}\FunctionTok{read.table}\NormalTok{(}\StringTok{"dadospcm/primate{-}data.txt"}\NormalTok{,}\AttributeTok{h=}\NormalTok{T,}\AttributeTok{row.names=}\DecValTok{1}\NormalTok{)}
\CommentTok{\# Recortar filogenia}
\NormalTok{match.species}\OtherTok{\textless{}{-}}\FunctionTok{treedata}\NormalTok{(tree,dados)}
\NormalTok{tree}\OtherTok{\textless{}{-}}\NormalTok{match.species}\SpecialCharTok{$}\NormalTok{phy}
\FunctionTok{plotTree}\NormalTok{(tree,}\AttributeTok{fsize=}\FloatTok{0.5}\NormalTok{,}\AttributeTok{ftype=}\StringTok{"i"}\NormalTok{,}\AttributeTok{type=}\StringTok{"fan"}\NormalTok{,}\AttributeTok{lwd=}\DecValTok{1}\NormalTok{)}
\end{Highlighting}
\end{Shaded}

\includegraphics{04.13-pcmtip_files/figure-latex/unnamed-chunk-1-2.pdf}

DR.

\begin{Shaded}
\begin{Highlighting}[]
\NormalTok{ES }\OtherTok{=} \FunctionTok{evol.distinct}\NormalTok{(tree, }\AttributeTok{type =} \StringTok{"equal.splits"}\NormalTok{)}\SpecialCharTok{$}\NormalTok{w}
\NormalTok{DR }\OtherTok{=} \DecValTok{1}\SpecialCharTok{/}\NormalTok{ES}
\CommentTok{\# Plot}
\FunctionTok{names}\NormalTok{(DR)}\OtherTok{=}\NormalTok{tree}\SpecialCharTok{$}\NormalTok{tip.label}
\FunctionTok{contMap}\NormalTok{(tree,DR)}
\end{Highlighting}
\end{Shaded}

\includegraphics{04.13-pcmtip_files/figure-latex/unnamed-chunk-2-1.pdf}

O tamanho corporal influencia a diversificação?

\begin{Shaded}
\begin{Highlighting}[]
\CommentTok{\# Massa corporal}
\NormalTok{massa}\OtherTok{\textless{}{-}}\NormalTok{dados}\SpecialCharTok{$}\NormalTok{AdultBodyMass\_g}
\FunctionTok{names}\NormalTok{(massa)}\OtherTok{=}\FunctionTok{rownames}\NormalTok{(dados)}
\FunctionTok{contMap}\NormalTok{(tree,}\FunctionTok{log}\NormalTok{(massa))}
\end{Highlighting}
\end{Shaded}

\includegraphics{04.13-pcmtip_files/figure-latex/unnamed-chunk-3-1.pdf}

\begin{Shaded}
\begin{Highlighting}[]
\CommentTok{\# Correlação}
\FunctionTok{cor}\NormalTok{(DR,}\FunctionTok{log}\NormalTok{(massa))}
\FunctionTok{plot}\NormalTok{(}\FunctionTok{log}\NormalTok{(massa),DR)}
\end{Highlighting}
\end{Shaded}

\includegraphics{04.13-pcmtip_files/figure-latex/unnamed-chunk-3-2.pdf}

\chapter{Co-diversificação}\label{co-diversificauxe7uxe3o}

Análises cofilogenéticas investigam a associação/dependência entre a diversificação de dois clados, normalmente um clado de parasitas e um clado de hospedeiros; para uma revisão veja Perez-Lamarque and Morlon (2022). Aqui, usaremos um teste baseado em Mantel para testar o sinal filogenético na associação entre dois clados (Perez-Lamarque and Morlon 2022). Vamos carregar dados de associação entre coronavirus (parasitas) e mamíferos (hospedeiros) de Maestri et al. (n.d.).

Carregar dados.

\begin{Shaded}
\begin{Highlighting}[]
\CommentTok{\# Carregar árvore de Coronaviridae}
\FunctionTok{require}\NormalTok{(ape)}
\CommentTok{\#\textgreater{} Carregando pacotes exigidos: ape}
\NormalTok{cov\_tree}\OtherTok{\textless{}{-}}\FunctionTok{read.nexus}\NormalTok{(}\StringTok{"dadospcm/combined\_alignment\_cov\_trimal\_tree\_35.tre"}\NormalTok{)}
\CommentTok{\# Carregar matriz de associação}
\NormalTok{hosts\_matrix}\OtherTok{\textless{}{-}}\FunctionTok{read.table}\NormalTok{(}\StringTok{"dadospcm/association\_matrix.txt"}\NormalTok{)}
\CommentTok{\# Carregar árvore de mamíferos hospedeiros}
\NormalTok{hostmammaltree}\OtherTok{\textless{}{-}}\FunctionTok{read.tree}\NormalTok{(}\StringTok{"dadospcm/hostsMammalTree.tre"}\NormalTok{)}
\end{Highlighting}
\end{Shaded}

Sinal filogenético na associação, com permutações mantendo constante o número de parceiros.

\begin{Shaded}
\begin{Highlighting}[]
\FunctionTok{require}\NormalTok{(RPANDA)}
\CommentTok{\#\textgreater{} Carregando pacotes exigidos: RPANDA}
\CommentTok{\#\textgreater{} Carregando pacotes exigidos: picante}
\CommentTok{\#\textgreater{} Carregando pacotes exigidos: vegan}
\CommentTok{\#\textgreater{} Carregando pacotes exigidos: permute}
\CommentTok{\#\textgreater{} Carregando pacotes exigidos: lattice}
\CommentTok{\#\textgreater{} This is vegan 2.6{-}4}
\CommentTok{\#\textgreater{} Carregando pacotes exigidos: nlme}
\NormalTok{psuni\_n}\OtherTok{\textless{}{-}}\FunctionTok{phylosignal\_network}\NormalTok{(hosts\_matrix, }\AttributeTok{tree\_A =}\NormalTok{ cov\_tree, }
 \AttributeTok{tree\_B =}\NormalTok{ hostmammaltree, }
 \AttributeTok{method =} \StringTok{"GUniFrac"}\NormalTok{, }
 \AttributeTok{correlation =} \StringTok{"Pearson"}\NormalTok{, }
 \AttributeTok{nperm=}\DecValTok{1000}\NormalTok{, }\AttributeTok{permutation =} \StringTok{"nbpartners"}\NormalTok{)}
\NormalTok{psuni\_n}
\end{Highlighting}
\end{Shaded}

Coronavirus próximos filogenéticamente tendem a se associar (infectar) mamíferos proximos filogenéticamente (r=0.377) e vice-versa (r=0.295).

Também podemos testar se parasitas proximamente relacionados tendem a infectar um \emph{número} de hospeiros similar.

\begin{Shaded}
\begin{Highlighting}[]
\CommentTok{\# Sinal filogenético no número de parceiros}
\NormalTok{npart\_cov}\OtherTok{\textless{}{-}}\FunctionTok{phylosignal\_network}\NormalTok{(hosts\_matrix, }
 \AttributeTok{tree\_A =}\NormalTok{ cov\_tree, }
 \AttributeTok{method =} \StringTok{"degree"}\NormalTok{, }
 \AttributeTok{correlation =} \StringTok{"Pearson"}\NormalTok{, }\AttributeTok{nperm =} \DecValTok{1000}\NormalTok{)}
\NormalTok{npart\_mammals}\OtherTok{\textless{}{-}}\FunctionTok{phylosignal\_network}\NormalTok{(}\FunctionTok{t}\NormalTok{(hosts\_matrix), }
 \AttributeTok{tree\_A =}\NormalTok{ hostmammaltree, }
 \AttributeTok{method =} \StringTok{"degree"}\NormalTok{, }
 \AttributeTok{correlation =} \StringTok{"Pearson"}\NormalTok{, }\AttributeTok{nperm =} \DecValTok{1000}\NormalTok{)}
\NormalTok{npart\_cov}
\NormalTok{npart\_mammals}
\end{Highlighting}
\end{Shaded}

Coronavirus próximos filogenéticamente tendem a infectar um número de hospeiros similar, porém, mamíferos proximos filogenéticamente não tendem a hospedar um número similar de coronavirus. Isso sugere que a especificidade de coronavírus em relação aos seus hospedeiros é conservada evolutivamente, enquanto que a suscetibilidade dos hospedeiros aos coronavírus não é.

Por fim, podemos testar em que escala filogenética o sinal é mais forte.

\begin{Shaded}
\begin{Highlighting}[]
\CommentTok{\# Calcular sinal filogenético por subclados}
\CommentTok{\# Para parasitas}
\NormalTok{results\_clade\_A }\OtherTok{\textless{}{-}} \FunctionTok{phylosignal\_sub\_network}\NormalTok{(hosts\_matrix, }
 \AttributeTok{tree\_A =}\NormalTok{ cov\_tree, }
 \AttributeTok{tree\_B =}\NormalTok{ hostmammaltree, }
 \AttributeTok{method =} \StringTok{"GUniFrac"}\NormalTok{, }
 \AttributeTok{correlation =} \StringTok{"Pearson"}\NormalTok{, }
 \AttributeTok{nperm=}\DecValTok{1000}\NormalTok{,}\AttributeTok{minimum=}\DecValTok{10}\NormalTok{, }\AttributeTok{degree=}\NormalTok{F,}
 \AttributeTok{permutation =} \StringTok{"shuffle"}\NormalTok{)}
\FunctionTok{plot\_phylosignal\_sub\_network}\NormalTok{(}\AttributeTok{tree\_A =}\NormalTok{ cov\_tree, results\_clade\_A,}
 \AttributeTok{network =}\NormalTok{ hosts\_matrix)}
\end{Highlighting}
\end{Shaded}

\includegraphics{04.14-pcmcodiv_files/figure-latex/unnamed-chunk-4-1.pdf}

\begin{Shaded}
\begin{Highlighting}[]
\CommentTok{\# Para hospedeiros}
\NormalTok{results\_clade\_B }\OtherTok{\textless{}{-}} \FunctionTok{phylosignal\_sub\_network}\NormalTok{(}\FunctionTok{t}\NormalTok{(hosts\_matrix), }
 \AttributeTok{tree\_A =}\NormalTok{ hostmammaltree, }
 \AttributeTok{tree\_B =}\NormalTok{ cov\_tree, }
 \AttributeTok{method =} \StringTok{"GUniFrac"}\NormalTok{, }
 \AttributeTok{correlation =} \StringTok{"Pearson"}\NormalTok{, }
 \AttributeTok{nperm=}\DecValTok{1000}\NormalTok{, }
 \AttributeTok{minimum=}\DecValTok{10}\NormalTok{, }\AttributeTok{degree=}\NormalTok{F,}
 \AttributeTok{permutation =} \StringTok{"shuffle"}\NormalTok{)}
\FunctionTok{plot\_phylosignal\_sub\_network}\NormalTok{(}\AttributeTok{tree\_A =}\NormalTok{ hostmammaltree, results\_clade\_B,}
 \AttributeTok{network =} \FunctionTok{t}\NormalTok{(hosts\_matrix),}
 \AttributeTok{legend=}\ConstantTok{FALSE}\NormalTok{,}\AttributeTok{where =} \StringTok{"topleft"}\NormalTok{)}
\end{Highlighting}
\end{Shaded}

\includegraphics{04.14-pcmcodiv_files/figure-latex/unnamed-chunk-4-2.pdf}

\begin{Shaded}
\begin{Highlighting}[]
\CommentTok{\# Plot conjunto }
\FunctionTok{layout}\NormalTok{(}\FunctionTok{matrix}\NormalTok{(}\FunctionTok{c}\NormalTok{(}\DecValTok{1}\NormalTok{,}\DecValTok{2}\NormalTok{),}\DecValTok{1}\NormalTok{,}\DecValTok{2}\NormalTok{),}\AttributeTok{widths=}\FunctionTok{c}\NormalTok{(}\FloatTok{0.5}\NormalTok{,}\FloatTok{0.5}\NormalTok{),}\AttributeTok{heights=}\FunctionTok{c}\NormalTok{(}\FloatTok{0.1}\NormalTok{,}\FloatTok{0.1}\NormalTok{))}
\FunctionTok{plot\_phylosignal\_sub\_network}\NormalTok{(}\AttributeTok{tree\_A =}\NormalTok{ cov\_tree, results\_clade\_A,}
 \AttributeTok{network =}\NormalTok{ hosts\_matrix)}
\FunctionTok{plot\_phylosignal\_sub\_network}\NormalTok{(}\AttributeTok{tree\_A =}\NormalTok{ hostmammaltree, results\_clade\_B,}
 \AttributeTok{network =} \FunctionTok{t}\NormalTok{(hosts\_matrix),}
 \AttributeTok{legend=}\ConstantTok{FALSE}\NormalTok{,}\AttributeTok{where =} \StringTok{"topleft"}\NormalTok{)}
\end{Highlighting}
\end{Shaded}

\includegraphics{04.14-pcmcodiv_files/figure-latex/unnamed-chunk-4-3.pdf}

Em geral, o sinal filogenético é mais forte nos nós mais antigos das duas filogenias.

\part{PARTE IV - Biogeografia e Macroecologia}\label{part-parte-iv---biogeografia-e-macroecologia}

\chapter{Biogeografia e Macroecologia}\label{biogeografia-e-macroecologia}

Em construção.

\part{Conclusão}\label{part-conclusuxe3o}

\chapter{Considerações finais}\label{considerauxe7uxf5es-finais}

Em construção.

\phantomsection\label{refs}
\begin{CSLReferences}{1}{0}
\bibitem[\citeproctext]{ref-alfaro2009}
Alfaro, Michael E., Francesco Santini, Chad Brock, Hugo Alamillo, Alex Dornburg, Daniel L. Rabosky, Giorgio Carnevale, and Luke J. Harmon. 2009. {``Nine exceptional radiations plus high turnover explain species diversity in jawed vertebrates.''} \emph{Proceedings of the National Academy of Sciences of the United States of America} 106 (32): 13410--14. \url{https://doi.org/10.1073/pnas.0811087106}.

\bibitem[\citeproctext]{ref-blomberg2003}
Blomberg, Simon P., Theodore Garland JR., and Anthony R. Ives. 2003. {``Testing for Phylogenetic Signal in Comparative Data: Behavioral Traits Are More Labile.''} \emph{Evolution} 57 (4): 717--45. \url{https://doi.org/10.1111/j.0014-3820.2003.tb00285.x}.

\bibitem[\citeproctext]{ref-butler-king2004}
Butler, Marguerite~A., and Aaron~A. King. 2004. {``Phylogenetic Comparative Analysis: A Modeling Approach for Adaptive Evolution.''} \emph{The American Naturalist} 164 (6): 683--95. \url{https://doi.org/10.1086/426002}.

\bibitem[\citeproctext]{ref-diniz-filho1998}
Diniz-Filho, José Alexandre Felizola, Carlos Eduardo Ramos de Sant'Ana, and Luis Mauricio Bini. 1998. {``An Eigenvector Method for Estimating Phylogenetic Inertia.''} \emph{Evolution} 52 (5): 1247--62. \url{https://doi.org/10.1111/j.1558-5646.1998.tb02006.x}.

\bibitem[\citeproctext]{ref-diniz-filho2012}
Diniz-Filho, José Alexandre Felizola, Thiago F. Rangel, Thiago Santos, and Luis Mauricio Bini. 2012a. {``EXPLORING PATTERNS OF INTERSPECIFIC VARIATION IN QUANTITATIVE TRAITS USING SEQUENTIAL PHYLOGENETIC EIGENVECTOR REGRESSIONS.''} \emph{Evolution} 66 (4): 1079--90. \url{https://doi.org/10.1111/j.1558-5646.2011.01499.x}.

\bibitem[\citeproctext]{ref-diniz-filho2012a}
---------. 2012b. {``EXPLORING PATTERNS OF INTERSPECIFIC VARIATION IN QUANTITATIVE TRAITS USING SEQUENTIAL PHYLOGENETIC EIGENVECTOR REGRESSIONS.''} \emph{Evolution} 66 (4): 1079--90. \url{https://doi.org/10.1111/j.1558-5646.2011.01499.x}.

\bibitem[\citeproctext]{ref-felsenstein1985}
Felsenstein, Joseph. 1985. {``Phylogenies and the Comparative Method.''} \emph{The American Naturalist} 125 (1): 1--15. \url{https://www.jstor.org/stable/2461605}.

\bibitem[\citeproctext]{ref-grafen1989}
Grafen, A. 1989. {``The phylogenetic regression.''} \emph{Philosophical Transactions of the Royal Society of London. Series B, Biological Sciences} 326 (1233): 119--57. \url{https://doi.org/10.1098/rstb.1989.0106}.

\bibitem[\citeproctext]{ref-hansen1997}
Hansen, Thomas F. 1997. {``Stabilizing Selection and the Comparative Analysis of Adaptation.''} \emph{Evolution} 51 (5): 1341--51. \url{https://doi.org/10.1111/j.1558-5646.1997.tb01457.x}.

\bibitem[\citeproctext]{ref-hansen-martins1996}
Hansen, Thomas F., and Emília P. Martins. 1996. {``Translating Between Microevolutionary Process and Macroevolutionary Patterns: The Correlation Structure of Interspecific Data.''} \emph{Evolution} 50 (4): 1404--17. \url{https://doi.org/10.1111/j.1558-5646.1996.tb03914.x}.

\bibitem[\citeproctext]{ref-harmon2003}
Harmon, Luke J., James A. Schulte, Allan Larson, and Jonathan B. Losos. 2003. {``Tempo and mode of evolutionary radiation in iguanian lizards.''} \emph{Science (New York, N.Y.)} 301 (5635): 961--64. \url{https://doi.org/10.1126/science.1084786}.

\bibitem[\citeproctext]{ref-huelsenbeck2003}
Huelsenbeck, John P., Rasmus Nielsen, and Jonathan P. Bollback. 2003. {``Stochastic Mapping of Morphological Characters.''} \emph{Systematic Biology} 52 (2): 131--58. \url{https://doi.org/10.1080/10635150390192780}.

\bibitem[\citeproctext]{ref-jetz2012}
Jetz, W., G. H. Thomas, J. B. Joy, K. Hartmann, and A. O. Mooers. 2012. {``The Global Diversity of Birds in Space and Time.''} \emph{Nature} 491 (7424): 444--48. \url{https://doi.org/10.1038/nature11631}.

\bibitem[\citeproctext]{ref-maestri2017}
Maestri, Renan, Leandro Rabello Monteiro, Rodrigo Fornel, Nathan S. Upham, Bruce D. Patterson, and Thales Renato Ochotorena Freitas. 2017. {``The Ecology of a Continental Evolutionary Radiation: Is the Radiation of Sigmodontine Rodents Adaptive?''} \emph{Evolution} 71 (3): 610--32. \url{https://doi.org/10.1111/evo.13155}.

\bibitem[\citeproctext]{ref-maestricovpreprint}
Maestri, Renan, Benoît Perez-Lamarque, Anna Zhukova, and Hélène Morlon. n.d. {``Recent Evolutionary Origin and Localized Diversity Hotspots of Mammalian Coronaviruses.''} \url{https://doi.org/10.1101/2023.03.09.531875}.

\bibitem[\citeproctext]{ref-magallon2001}
Magallon, Susana, and Michael J. Sanderson. 2001. {``Absolute Diversification Rates in Angiosperm Clades.''} \emph{Evolution} 55 (9): 1762--80. \url{https://doi.org/10.1111/j.0014-3820.2001.tb00826.x}.

\bibitem[\citeproctext]{ref-maliet2019}
Maliet, Odile, Florian Hartig, and Hélène Morlon. 2019. {``A Model with Many Small Shifts for Estimating Species-Specific Diversification Rates.''} \emph{Nature Ecology \& Evolution} 3 (7): 1086--92. \url{https://doi.org/10.1038/s41559-019-0908-0}.

\bibitem[\citeproctext]{ref-morlon2016}
Morlon, Hélène, Eric Lewitus, Fabien L. Condamine, Marc Manceau, Julien Clavel, and Jonathan Drury. 2016. {``RPANDA: An R Package for Macroevolutionary Analyses on Phylogenetic Trees.''} \emph{Methods in Ecology and Evolution} 7 (5): 589--97. \url{https://doi.org/10.1111/2041-210X.12526}.

\bibitem[\citeproctext]{ref-nee1994}
Nee, Sean, Robert M. May, and Paul H. Harvey. 1994. {``The Reconstructed Evolutionary Process.''} \emph{Philosophical Transactions: Biological Sciences} 344 (1309): 305--11. \url{https://www.jstor.org/stable/55922}.

\bibitem[\citeproctext]{ref-pagel1994}
Pagel, Mark. 1994. {``Detecting Correlated Evolution on Phylogenies: A General Method for the Comparative Analysis of Discrete Characters.''} \emph{Proceedings of the Royal Society of London. Series B: Biological Sciences} 255 (1342): 37--45. \url{https://doi.org/10.1098/rspb.1994.0006}.

\bibitem[\citeproctext]{ref-pagel1999}
---------. 1999. {``Inferring the Historical Patterns of Biological Evolution.''} \emph{Nature} 401 (6756): 877--84. \url{https://doi.org/10.1038/44766}.

\bibitem[\citeproctext]{ref-perez-lamarque-morlon2022}
Perez-Lamarque, Benoît, and Hélène Morlon. 2022. {``Comparing Different Computational Approaches for Detecting Long-Term Vertical Transmission in Host-Associated Microbiota.''} \emph{Molecular Ecology} n/a (n/a). \url{https://doi.org/10.1111/mec.16681}.

\bibitem[\citeproctext]{ref-rangel2015}
Rangel, Thiago F., Robert K. Colwell, Gary R. Graves, Karolina Fučíková, Carsten Rahbek, and José Alexandre F. Diniz-Filho. 2015. {``Phylogenetic Uncertainty Revisited: Implications for Ecological Analyses.''} \emph{Evolution} 69 (5): 1301--12. \url{https://doi.org/10.1111/evo.12644}.

\bibitem[\citeproctext]{ref-redding2006}
Redding, David W., and Arne Ø. Mooers. 2006. {``Incorporating Evolutionary Measures into Conservation Prioritization.''} \emph{Conservation Biology} 20 (6): 1670--78. \url{https://www.jstor.org/stable/4124694}.

\bibitem[\citeproctext]{ref-revell2010}
Revell, Liam J. 2010. {``Phylogenetic Signal and Linear Regression on Species Data.''} \emph{Methods in Ecology and Evolution} 1 (4): 319--29. \url{https://doi.org/10.1111/j.2041-210X.2010.00044.x}.

\bibitem[\citeproctext]{ref-revell-harmon2022}
Revell, Liam J., and Luke J. Harmon. 2022. \emph{Phylogenetic Comparative Methods in R}. Princeton University Press.

\bibitem[\citeproctext]{ref-rohlf2001}
Rohlf, F. James. 2001. {``Comparative Methods for the Analysis of Continuous Variables: Geometric Interpretations.''} \emph{Evolution} 55 (11): 2143--60. \url{https://doi.org/10.1111/j.0014-3820.2001.tb00731.x}.

\bibitem[\citeproctext]{ref-yang2006}
Yang, Ziheng. 2006. \emph{Computational Molecular Evolution}. Oxford: Oxford University Pres.

\end{CSLReferences}

\end{document}
